% =====================================================================
%  Fondamenti di Unreal Engine, lezione 4 — Vettori, rotazioni e la
%  torretta: il mulino che gira e la torretta che segue il player.
%  Quarta lezione del corso breve di introduzione a Unreal Engine (focus
%  sui Blueprint), online su Zoom, dalle 16 circa alle 19. Dopo la
%  sostituzione di Davide nella terza lezione, torna il docente
%  titolare (Igor, che però nella trascrizione non viene mai chiamato
%  per nome).
%  Contenuti, nell'ordine della lezione: l'attesa (Quixel Megascans, la
%  lava di Davide); il trucco del CCD per la pallina della seconda
%  lezione; il programma (il mulino e la torretta che segue lo
%  sguardo); i vettori (direzione, magnitudine, verso), addizione,
%  sottrazione, moltiplicazione, prodotto scalare, vettori
%  normalizzati; i rotator, il giroscopio, il gimbal lock e i
%  quaternioni; il Rotating Movement Component e la mobilità Movable;
%  l'esercizio dei tre cubi; il Blueprint del mulino (BP_Windmill, il
%  corpo e le due pale), la viewport persa e l'editor in inglese; far
%  girare una pala con l'Event Tick, i delta seconds, Make Rotator,
%  Combine Rotators, Get e Set Relative Rotation; la gerarchia dei
%  componenti (la seconda pala figlia della prima, il cubetto
%  sull'estremità); Find Look at Rotation, Get Player Character,
%  RInterp To e l'offset di 90 gradi: il mulino che guarda il player;
%  il Clamp e il debug con Print String; la pausa; la torretta in tre
%  pezzi (BP_Turret: Base, Pivot, Cannon) e Browse to Asset; la Sphere
%  Collision, lo Sphere Radius e Hidden in Game; gli eventi dei
%  componenti e On Component Begin Overlap; Other Actor; Is Valid;
%  Actor Has Tag e i tag; la variabile Player Target; i Custom Event;
%  Set Timer by Event, Get World Delta Seconds, Looping e Max Once Per
%  Frame; l'evento Move Turret; il tuning dell'interpolazione; i
%  problemi degli studenti (il pivot reimportato a 113, la gerarchia
%  sbagliata, il Blueprint non in scena, Other Actor scollegato); la
%  chiusura e la prossima lezione (proiettili, spawn e line trace).
%  Fonte: trascrizione automatica della quarta lezione (UE4), circa tre
%  ore. Integrazioni dai libri "Blueprints Visual Scripting for Unreal
%  Engine 5" (Romero, Sewell; capitoli 1, 4, 5 e 14) e "Blueprints
%  Visual Scripting for Unreal Engine", seconda edizione (stessi autori,
%  Unreal Engine 4; capitoli 2 e 8), e nozioni tecniche generali su
%  Unreal Engine 5, sempre segnalate come approfondimenti.
%
%  Pacchetti richiesti nel preambolo del libro (già presenti nel
%  main.tex): \usepackage{booktabs}, \usepackage{amsmath},
%  \usepackage{tikz} con
%    \usetikzlibrary{shapes.geometric,shapes.symbols,shapes.misc,
%      arrows.meta,positioning,calc,fit}
%  (le librerie vanno caricate nel preambolo, non nel capitolo).
%  Consigliati, come per gli altri capitoli: \usepackage[italian]{babel}
%  e \usepackage[T1]{fontenc} (il capitolo compila anche senza).
%
%  Collocazione nel libro: parte "Fondamenti di Unreal Engine", quarto
%  capitolo (numerazione propria 1.4). Nel file principale del libro,
%  subito dopo la riga della terza lezione:
%
%    % =====================================================================
%  Fondamenti di Unreal Engine, lezione 4 — Vettori, rotazioni e la
%  torretta: il mulino che gira e la torretta che segue il player.
%  Quarta lezione del corso breve di introduzione a Unreal Engine (focus
%  sui Blueprint), online su Zoom, dalle 16 circa alle 19. Dopo la
%  sostituzione di Davide nella terza lezione, torna il docente
%  titolare (Igor, che però nella trascrizione non viene mai chiamato
%  per nome).
%  Contenuti, nell'ordine della lezione: l'attesa (Quixel Megascans, la
%  lava di Davide); il trucco del CCD per la pallina della seconda
%  lezione; il programma (il mulino e la torretta che segue lo
%  sguardo); i vettori (direzione, magnitudine, verso), addizione,
%  sottrazione, moltiplicazione, prodotto scalare, vettori
%  normalizzati; i rotator, il giroscopio, il gimbal lock e i
%  quaternioni; il Rotating Movement Component e la mobilità Movable;
%  l'esercizio dei tre cubi; il Blueprint del mulino (BP_Windmill, il
%  corpo e le due pale), la viewport persa e l'editor in inglese; far
%  girare una pala con l'Event Tick, i delta seconds, Make Rotator,
%  Combine Rotators, Get e Set Relative Rotation; la gerarchia dei
%  componenti (la seconda pala figlia della prima, il cubetto
%  sull'estremità); Find Look at Rotation, Get Player Character,
%  RInterp To e l'offset di 90 gradi: il mulino che guarda il player;
%  il Clamp e il debug con Print String; la pausa; la torretta in tre
%  pezzi (BP_Turret: Base, Pivot, Cannon) e Browse to Asset; la Sphere
%  Collision, lo Sphere Radius e Hidden in Game; gli eventi dei
%  componenti e On Component Begin Overlap; Other Actor; Is Valid;
%  Actor Has Tag e i tag; la variabile Player Target; i Custom Event;
%  Set Timer by Event, Get World Delta Seconds, Looping e Max Once Per
%  Frame; l'evento Move Turret; il tuning dell'interpolazione; i
%  problemi degli studenti (il pivot reimportato a 113, la gerarchia
%  sbagliata, il Blueprint non in scena, Other Actor scollegato); la
%  chiusura e la prossima lezione (proiettili, spawn e line trace).
%  Fonte: trascrizione automatica della quarta lezione (UE4), circa tre
%  ore. Integrazioni dai libri "Blueprints Visual Scripting for Unreal
%  Engine 5" (Romero, Sewell; capitoli 1, 4, 5 e 14) e "Blueprints
%  Visual Scripting for Unreal Engine", seconda edizione (stessi autori,
%  Unreal Engine 4; capitoli 2 e 8), e nozioni tecniche generali su
%  Unreal Engine 5, sempre segnalate come approfondimenti.
%
%  Pacchetti richiesti nel preambolo del libro (già presenti nel
%  main.tex): \usepackage{booktabs}, \usepackage{amsmath},
%  \usepackage{tikz} con
%    \usetikzlibrary{shapes.geometric,shapes.symbols,shapes.misc,
%      arrows.meta,positioning,calc,fit}
%  (le librerie vanno caricate nel preambolo, non nel capitolo).
%  Consigliati, come per gli altri capitoli: \usepackage[italian]{babel}
%  e \usepackage[T1]{fontenc} (il capitolo compila anche senza).
%
%  Collocazione nel libro: parte "Fondamenti di Unreal Engine", quarto
%  capitolo (numerazione propria 1.4). Nel file principale del libro,
%  subito dopo la riga della terza lezione:
%
%    % =====================================================================
%  Fondamenti di Unreal Engine, lezione 4 — Vettori, rotazioni e la
%  torretta: il mulino che gira e la torretta che segue il player.
%  Quarta lezione del corso breve di introduzione a Unreal Engine (focus
%  sui Blueprint), online su Zoom, dalle 16 circa alle 19. Dopo la
%  sostituzione di Davide nella terza lezione, torna il docente
%  titolare (Igor, che però nella trascrizione non viene mai chiamato
%  per nome).
%  Contenuti, nell'ordine della lezione: l'attesa (Quixel Megascans, la
%  lava di Davide); il trucco del CCD per la pallina della seconda
%  lezione; il programma (il mulino e la torretta che segue lo
%  sguardo); i vettori (direzione, magnitudine, verso), addizione,
%  sottrazione, moltiplicazione, prodotto scalare, vettori
%  normalizzati; i rotator, il giroscopio, il gimbal lock e i
%  quaternioni; il Rotating Movement Component e la mobilità Movable;
%  l'esercizio dei tre cubi; il Blueprint del mulino (BP_Windmill, il
%  corpo e le due pale), la viewport persa e l'editor in inglese; far
%  girare una pala con l'Event Tick, i delta seconds, Make Rotator,
%  Combine Rotators, Get e Set Relative Rotation; la gerarchia dei
%  componenti (la seconda pala figlia della prima, il cubetto
%  sull'estremità); Find Look at Rotation, Get Player Character,
%  RInterp To e l'offset di 90 gradi: il mulino che guarda il player;
%  il Clamp e il debug con Print String; la pausa; la torretta in tre
%  pezzi (BP_Turret: Base, Pivot, Cannon) e Browse to Asset; la Sphere
%  Collision, lo Sphere Radius e Hidden in Game; gli eventi dei
%  componenti e On Component Begin Overlap; Other Actor; Is Valid;
%  Actor Has Tag e i tag; la variabile Player Target; i Custom Event;
%  Set Timer by Event, Get World Delta Seconds, Looping e Max Once Per
%  Frame; l'evento Move Turret; il tuning dell'interpolazione; i
%  problemi degli studenti (il pivot reimportato a 113, la gerarchia
%  sbagliata, il Blueprint non in scena, Other Actor scollegato); la
%  chiusura e la prossima lezione (proiettili, spawn e line trace).
%  Fonte: trascrizione automatica della quarta lezione (UE4), circa tre
%  ore. Integrazioni dai libri "Blueprints Visual Scripting for Unreal
%  Engine 5" (Romero, Sewell; capitoli 1, 4, 5 e 14) e "Blueprints
%  Visual Scripting for Unreal Engine", seconda edizione (stessi autori,
%  Unreal Engine 4; capitoli 2 e 8), e nozioni tecniche generali su
%  Unreal Engine 5, sempre segnalate come approfondimenti.
%
%  Pacchetti richiesti nel preambolo del libro (già presenti nel
%  main.tex): \usepackage{booktabs}, \usepackage{amsmath},
%  \usepackage{tikz} con
%    \usetikzlibrary{shapes.geometric,shapes.symbols,shapes.misc,
%      arrows.meta,positioning,calc,fit}
%  (le librerie vanno caricate nel preambolo, non nel capitolo).
%  Consigliati, come per gli altri capitoli: \usepackage[italian]{babel}
%  e \usepackage[T1]{fontenc} (il capitolo compila anche senza).
%
%  Collocazione nel libro: parte "Fondamenti di Unreal Engine", quarto
%  capitolo (numerazione propria 1.4). Nel file principale del libro,
%  subito dopo la riga della terza lezione:
%
%    % =====================================================================
%  Fondamenti di Unreal Engine, lezione 4 — Vettori, rotazioni e la
%  torretta: il mulino che gira e la torretta che segue il player.
%  Quarta lezione del corso breve di introduzione a Unreal Engine (focus
%  sui Blueprint), online su Zoom, dalle 16 circa alle 19. Dopo la
%  sostituzione di Davide nella terza lezione, torna il docente
%  titolare (Igor, che però nella trascrizione non viene mai chiamato
%  per nome).
%  Contenuti, nell'ordine della lezione: l'attesa (Quixel Megascans, la
%  lava di Davide); il trucco del CCD per la pallina della seconda
%  lezione; il programma (il mulino e la torretta che segue lo
%  sguardo); i vettori (direzione, magnitudine, verso), addizione,
%  sottrazione, moltiplicazione, prodotto scalare, vettori
%  normalizzati; i rotator, il giroscopio, il gimbal lock e i
%  quaternioni; il Rotating Movement Component e la mobilità Movable;
%  l'esercizio dei tre cubi; il Blueprint del mulino (BP_Windmill, il
%  corpo e le due pale), la viewport persa e l'editor in inglese; far
%  girare una pala con l'Event Tick, i delta seconds, Make Rotator,
%  Combine Rotators, Get e Set Relative Rotation; la gerarchia dei
%  componenti (la seconda pala figlia della prima, il cubetto
%  sull'estremità); Find Look at Rotation, Get Player Character,
%  RInterp To e l'offset di 90 gradi: il mulino che guarda il player;
%  il Clamp e il debug con Print String; la pausa; la torretta in tre
%  pezzi (BP_Turret: Base, Pivot, Cannon) e Browse to Asset; la Sphere
%  Collision, lo Sphere Radius e Hidden in Game; gli eventi dei
%  componenti e On Component Begin Overlap; Other Actor; Is Valid;
%  Actor Has Tag e i tag; la variabile Player Target; i Custom Event;
%  Set Timer by Event, Get World Delta Seconds, Looping e Max Once Per
%  Frame; l'evento Move Turret; il tuning dell'interpolazione; i
%  problemi degli studenti (il pivot reimportato a 113, la gerarchia
%  sbagliata, il Blueprint non in scena, Other Actor scollegato); la
%  chiusura e la prossima lezione (proiettili, spawn e line trace).
%  Fonte: trascrizione automatica della quarta lezione (UE4), circa tre
%  ore. Integrazioni dai libri "Blueprints Visual Scripting for Unreal
%  Engine 5" (Romero, Sewell; capitoli 1, 4, 5 e 14) e "Blueprints
%  Visual Scripting for Unreal Engine", seconda edizione (stessi autori,
%  Unreal Engine 4; capitoli 2 e 8), e nozioni tecniche generali su
%  Unreal Engine 5, sempre segnalate come approfondimenti.
%
%  Pacchetti richiesti nel preambolo del libro (già presenti nel
%  main.tex): \usepackage{booktabs}, \usepackage{amsmath},
%  \usepackage{tikz} con
%    \usetikzlibrary{shapes.geometric,shapes.symbols,shapes.misc,
%      arrows.meta,positioning,calc,fit}
%  (le librerie vanno caricate nel preambolo, non nel capitolo).
%  Consigliati, come per gli altri capitoli: \usepackage[italian]{babel}
%  e \usepackage[T1]{fontenc} (il capitolo compila anche senza).
%
%  Collocazione nel libro: parte "Fondamenti di Unreal Engine", quarto
%  capitolo (numerazione propria 1.4). Nel file principale del libro,
%  subito dopo la riga della terza lezione:
%
%    \input{UE/UE4_vettori_rotazioni_e_torretta}
%
%  Tutte le etichette di questo capitolo hanno il prefisso "ue4:"
%  (ue4:cap:..., ue4:sec:..., ue4:fig:..., ue4:tab:...), e tutti gli
%  stili TikZ hanno il prefisso "ue4" (ue4node, ue4exec, ue4data, ...)
%  per non collidere con chiavi native di TikZ (come "step") né con
%  quelli delle lezioni precedenti. Il capitolo è autonomo: non contiene
%  \ref verso altri capitoli (i rimandi alle lezioni precedenti sono a
%  parole).
%  Le figure sono ricostruzioni schematiche di quanto mostrato e detto a
%  lezione, non copie delle slide o delle schermate del docente. Dove la
%  trascrizione è troppo rovinata per ricostruire un dettaglio, il
%  capitolo lo dice invece di inventarlo.
% =====================================================================

\chapter[Vettori, rotazioni e la torretta]{Vettori, rotazioni e la torretta: il
mulino che gira e la torretta che ti segue}
\label{ue4:cap:torretta}

Questa è la \textbf{quarta lezione} del corso breve di \textbf{introduzione a Unreal
Engine}, e segna il ritorno del docente titolare dopo la lezione di Davide sulla
grafica. Se la terza lezione guardava a ciò che ``fa il grafico all'interno di Unreal''
(texture, materiali, l'import della torretta), questa torna alla \textbf{logica}, cioè
ai \textbf{Blueprint}, e lo fa con un obiettivo preciso, mostrato in apertura: un
\textbf{mulino} le cui pale girano, e la \textbf{torretta} importata con Davide che
``adesso ci segue con lo sguardo''. Per arrivarci serve un po' di matematica dello
spazio 3D: i \textbf{vettori}, le loro operazioni, i \textbf{rotator}.

La lezione ha due metà separate da una pausa di venti minuti, presa alle 17:28. La
prima metà parte dalla \textbf{teoria}: che cos'è un \textbf{vettore} (``una freccia
con una direzione, una magnitudine e un verso''), le quattro operazioni
(\textbf{addizione}, \textbf{sottrazione}, \textbf{moltiplicazione} e \textbf{prodotto
scalare}, il \emph{dot product}), i \textbf{vettori normalizzati}, i
\textbf{rotator} e il problema del \textbf{gimbal lock}. Poi si passa all'engine con un
mulino: prima il modo più semplice di far ruotare un oggetto, il \textbf{Rotating
Movement Component}, poi un Blueprint vero, \textbf{BP\_Windmill}, in cui una pala
gira grazie all'\textbf{Event Tick}, ai \emph{delta seconds} e ai nodi \textbf{Make
Rotator}, \textbf{Combine Rotators} e \textbf{Set Relative Rotation}; la seconda pala
gira ``gratis'' grazie alla \textbf{gerarchia} dei componenti. Infine il mulino
diventa ``magico'' e si gira verso il giocatore con \textbf{Find Look at Rotation} e
\textbf{RInterp To}, con un \textbf{Clamp} per limitarne la rotazione. La seconda metà
costruisce la \textbf{torretta}: un Blueprint in tre pezzi (\textbf{Base},
\textbf{Pivot}, \textbf{Cannon}), una \textbf{Sphere Collision} come raggio d'azione,
l'evento \textbf{On Component Begin Overlap}, i controlli \textbf{Is Valid} e
\textbf{Actor Has Tag}, una variabile per ricordare il bersaglio, un \textbf{Custom
Event} e un \textbf{timer} che, a ogni \emph{frame}, fa ruotare il pivot verso il
player.

La lezione è \textbf{online}, su Zoom, come le precedenti. Il docente alterna
spiegazione e piccole esercitazioni ``botta e risposta'', da cinque minuti l'una, e
passa gran parte della seconda metà a risolvere i problemi degli studenti in
condivisione dello schermo. Finisce verso le 19, ``andiamo a cena tutti'', con
l'appuntamento a \textbf{martedì}.

Tre avvertenze di lettura, le stesse delle lezioni precedenti. La prima: la
trascrizione è automatica e storpia moltissimi nomi. Unreal è ``a Real'', ``Arreal'',
``Arry''; l'\textbf{Event Tick} è ``Leventic'', ``l'eventic'', ``eventHIC'' e
``l'event kick''; i \textbf{vettori} sono una volta ``vittori''; i \emph{rotator} sono
``rotatempi'' e ``rotatori''; \textbf{Find Look at Rotation} è ``findLookAtRotation'',
``find all the rotation'' e ``Find, look, rotation''; \textbf{RInterp To} è ``Er inter
tu'', ``R Inter2'' e ``r interp tu''; \textbf{Get Player Character} è
``getLayerCharacter''; \textbf{Get Actor Location} è ``Get. Hector. Location'' e ``get
active location''; il \textbf{Rotating Movement} è a un certo punto ``vetro in
movement''; \textbf{Is Valid} è ``Svalid'' e ``isvalid''; \textbf{Other Actor} è
``OderActor'' e ``Oder Actor''; la \textbf{gun} della torretta è ``Ganna'' e ``GAN'';
il \textbf{Pivot} è spesso ``pivo''; il \emph{Third Person Character} è ``ter person
character'' e ``Peter per personacarte''; \textbf{Quixel Megascans} è ``QuickSell'' e
``Quixel megaschanico''. Dove la ricostruzione è sicura il capitolo usa la forma
corretta e, la prima volta, lo segnala in nota; dove non lo è riporta la forma della
trascrizione e lo dice. La seconda: all'inizio della registrazione, in un punto della
prima esercitazione e durante tutta la pausa, il riconoscimento vocale produce lunghi
blocchi di testo senza senso e ripetuto (``qual è la risposta che ci può essere?'',
``grazie a tutti per l'attenzione'', discorsi sul ``mercato del lavoro'', sul ``governo
di Dalai Lama'' e su ``la migliore dal 2005''): sono artefatti, non parole di nessuno,
e il capitolo li ignora. La terza: gran parte della lezione è una dimostrazione dal
vivo nell'editor dei Blueprint, spesso costruita a voce (``tirate fuori'',
``collegate'', ``splittate''). Il capitolo ricostruisce i grafi da ciò che il docente
dice mentre li costruisce, e li completa, dove serve, con i nomi esatti dei nodi e dei
pin di Unreal Engine 5.

\paragraph{Come leggere questo capitolo.} Il capitolo segue l'ordine della lezione.
Le figure ricostruiscono gli schemi spiegati a voce o alla lavagna (``perdonate le mie
doti molto brutte di disegno''): le operazioni tra vettori, i casi del prodotto
scalare, gli assi di rotazione, la gerarchia del mulino, i grafi della pala, del
mulino che guarda il player e della torretta, i casi dell'\emph{overlap}. Le tabelle
riassumono le operazioni sui vettori, i nodi della lezione e gli errori più comuni
incontrati dagli studenti. Qua e là, dei paragrafi di \textbf{approfondimento}
collegano la lezione ai libri di Marcos Romero e Brenden Sewell (\emph{Blueprints
Visual Scripting for Unreal Engine 5} e la seconda edizione dello stesso testo, scritta
per Unreal Engine 4), i testi di riferimento per questa parte del libro, e a qualche
nozione tecnica generale su Unreal Engine 5; sono segnalati come tali e non fanno parte
della lezione.

Il percorso del capitolo:
\begin{enumerate}
  \item l'attesa: \textbf{Quixel Megascans} e il materiale della lava; il trucco del
        \textbf{CCD}; il programma della giornata (sezioni da~\ref{ue4:sec:attesa}
        a~\ref{ue4:sec:programma});
  \item la teoria: i \textbf{vettori} e le loro operazioni, il \textbf{prodotto
        scalare}, la \textbf{normalizzazione}, i \textbf{rotator} e il \textbf{gimbal
        lock} (sezioni da~\ref{ue4:sec:vettori} a~\ref{ue4:sec:rotator});
  \item il \textbf{Rotating Movement Component} e l'esercizio dei tre cubi (sezione
        \ref{ue4:sec:rotating});
  \item il Blueprint del \textbf{mulino}: i componenti, la pala che gira con
        l'Event Tick, la \textbf{gerarchia} (sezioni da~\ref{ue4:sec:mulino}
        a~\ref{ue4:sec:gerarchia});
  \item il mulino che guarda il player: \textbf{Find Look at Rotation},
        \textbf{RInterp To}, il \textbf{Clamp} e il debug (sezioni
        \ref{ue4:sec:lookat} e~\ref{ue4:sec:clamp});
  \item dopo la pausa, la \textbf{torretta}: i tre pezzi, la \textbf{Sphere
        Collision}, l'\textbf{overlap} (sezioni da~\ref{ue4:sec:torretta}
        a~\ref{ue4:sec:overlap});
  \item \textbf{Is Valid}, i \textbf{tag}, la variabile del bersaglio e i
        \textbf{Custom Event} (sezioni da~\ref{ue4:sec:isvalid}
        a~\ref{ue4:sec:custom});
  \item il \textbf{timer} e l'evento \textbf{Move Turret}: la torretta che segue il
        player (sezioni \ref{ue4:sec:timer} e~\ref{ue4:sec:move});
  \item il \emph{tuning}, i problemi degli studenti e la chiusura (sezioni
        \ref{ue4:sec:problemi} e~\ref{ue4:sec:chiusura}).
\end{enumerate}

% ---------------------------------------------------------------------
\section{L'attesa: Quixel Megascans e la lava}
\label{ue4:sec:attesa}

La registrazione parte mentre il docente fa entrare gli studenti nella stanza di Zoom
(``ho aperto io la prima a far entrare un po' tutti'') e chiacchiera con il primo
arrivato. ``Come va con Unreal? Hai sperimentato un po', l'hai toccato un pochetto?''
Lo studente ci ha giocato con un collega, che gli ha insegnato le texture, ma non ha
trovato online ``alcuna fonte per prendere cose utili, file utili'': per l'esercizio
della lava assegnato da Davide nella lezione precedente, un materiale di lava animata,
ha usato quello che c'era nella cartella del Drive.

\paragraph{Quixel Megascans.} La risposta del docente è \textbf{Quixel
Megascans}\footnote{La trascrizione dice ``QuickSell'', ``Quixel megaschanico'' e
``Materials pixel mega scan''. Il nome corretto è Quixel Megascans: la libreria di
asset scansionati di Quixel, società acquisita da Epic Games, gratuita per l'uso con
Unreal Engine.}, di cui passa il link in chat. Quixel ``ti permette di avere proprio
dei materiali'': ``una marea di materiali di qualsiasi tipo'', ma anche rocce,
\emph{mesh} e texture di alta qualità. Il docente fatica un po' a ritrovare la pagina
giusta, ``perché hanno cambiato il sito'', e ricorda che serve accedere con l'account
di Unreal (cioè di Epic Games). ``Ormai fanno la cooperazione tra di loro'': c'è un
collegamento diretto tra Quixel e Unreal, ``proprio tipo \emph{face to face}''. Cercando
\emph{lava} escono diverse risorse; una, \emph{lava story} o qualcosa del genere,
``sembra carina''.

\paragraph{Come sono fatte.} Le texture di Megascans ``sono proprio \emph{scan}'': c'è
``gente che si mette là con queste macchine che costano una vita umana'' a
scansionare le rocce. ``Quixel è un grande, grande calderone di risorse.'' Da lì si
può partire anche per la lava senza trovare la texture perfetta: si prende una
\textbf{roccia} e, ``come avete visto con Davide, con la matematica'', cioè con un
\textbf{Multiply} e con il \textbf{Lerp} pilotato da un'\textbf{alfa}, si inserisce
``il rosso della lava all'interno di queste rocce''. Un'altra strada sono i video su
YouTube: alcuni tutorial lasciano in descrizione le texture usate (altri le fanno
pagare, ``il maledetto''), e ``ti interessano le texture, il materiale poi lo fai tu''.

\paragraph{La lava degli studenti.} Arrivano gli altri (``buongiorno, Filippo, Marco,
Patrick''), e il docente chiede chi ha provato il materiale della lava ``che vi è stato
dato da Davide''\footnote{La trascrizione dice ``da Dario'', ma subito dopo e più
avanti il nome è Davide, il docente della terza lezione.}. Uno studente l'ha mandato sul
Discord, e il docente lo commenta: ``bello il movimento che hai fatto''. L'hai messo
su un piano, in engine, per vedere come reagiva? No. ``Quello vedilo un po' tu'',
perché a volte il risultato in \emph{viewport} cambia a seconda di come è costruito il
materiale (``che so, hai messo il Multiply sbagliato''): il consiglio è di provarlo
sempre in \emph{viewport}, su un oggetto vero. E ``sicuramente una texture differente
potrebbe essere più carina'': torna utile Megascans.

\paragraph{L'appello.} Il docente aspetta ``2--3 minuti'' che entrino gli altri. Uno
studente avvisa che anche oggi dovrà staccarsi venti minuti prima per andare a
lavorare: ``assolutamente nessun problema''.

% ---------------------------------------------------------------------
\section{Le domande e il trucco del CCD}
\label{ue4:sec:ccd}

Prima di cominciare, il docente apre lo spazio delle domande: ``se qualcuno di voi ha
avuto qualche particolare cosa di cui vuole parlare, domandare già prima che magari
iniziamo, così facciamo il Q\&A già da subito''. Nessuno: ``mi sembrate tutti molto
tranquilli''.

\paragraph{La pallina che entra nel piano.} Allora è il docente a recuperare ``un
piccolo trucchetto che mi sono dimenticato di dirvi'' per il \textbf{gioco della
pallina} della seconda lezione, quello della piattaforma inclinabile: a volte la
pallina, con la fisica attiva, \textbf{penetra} nel piano, ``come abbiamo visto''. Nei
dettagli del componente, nella sezione della fisica (vicino a \textbf{Simulate
Physics}), c'è un'opzione da spuntare: il \textbf{CCD}, \emph{Continuous Collision
Detection}\footnote{La trascrizione dice ``Trucco Collision Collision Detection''.}.
``Fondamentalmente, è un \emph{tick} che, se lo mettiamo su questo particolare
\emph{asset}, lui in modo continuo cerca di vedere le collisioni'': le rende ``più
precise'', ed evita che la pallina attraversi il terreno. ``Se fate questo, vi risolve
tranquillamente il problema.'' Chi vuole, può sistemare il gioco per la prossima volta.

\paragraph{Approfondimento: perché la pallina attraversa il piano.} La fisica di un
gioco non è continua: il motore fisico controlla le collisioni a passi discreti, una
volta per \emph{frame} (o per \emph{substep}). Se in un passo la pallina si sposta più
dello spessore del piano, può trovarsi ``prima'' sopra e ``dopo'' già dentro o sotto,
senza che nessun controllo l'abbia mai vista toccare la superficie: è il fenomeno del
\emph{tunneling}, tanto più probabile quanto più l'oggetto è piccolo e veloce e il piano
sottile. Il CCD (in Unreal la casella \emph{Use CCD} nella sezione \emph{Collision} del
componente, nelle impostazioni avanzate) fa sì che il motore fisico consideri anche il
percorso tra le due posizioni, la ``scia'' dell'oggetto, e non solo le posizioni
finali. Costa un po' di più, ed è per questo che non è attivo di default: lo si accende
solo sugli oggetti che ne hanno bisogno, come proiettili e palline veloci. Un'alternativa
complementare è dare al piano uno spessore maggiore.

% ---------------------------------------------------------------------
\section{Il programma: il mulino e la torretta}
\label{ue4:sec:programma}

``Oggi vi dico già da subito cosa faremo.'' Gli studenti hanno scaricato il
\textbf{materiale didattico}: dentro il file estratto c'è il progetto
\emph{UE5 Intro Start}\footnote{La trascrizione dice ``Hive UE5 intro start'', la stessa
parola iniziale incerta già incontrata nella terza lezione.}, con il suo file di
progetto. Lo si apre, ``così almeno gli \emph{shader} si cominciano a buildare'', si fa
``una piccola spiegazione teorica'' e poi ``partiamo direttamente con la pratica''.

\paragraph{I due obiettivi.} Il docente sposta il \emph{Player Start} e mostra il
livello. A sinistra c'è ``questo bellissimo \textbf{mulino}'', fatto ``facendo ruotare
le pale''. Poi c'è la \textbf{torretta} che gli studenti hanno costruito la volta prima
con Davide (``avete fatto il \emph{merge}, giusto?''), che ``adesso ci segue con lo
sguardo''. O meglio, dovrebbe: nella dimostrazione non lo fa. ``Abbiamo scoperto un
piccolo bug. Poi adesso lo aggiustiamo. E giustamente i bug ogni volta appaiono
praticamente subito.'' Il programma della giornata è questo: le pale del mulino, e la
torretta che segue il giocatore.

% ---------------------------------------------------------------------
\section{I vettori}
\label{ue4:sec:vettori}

``Avete sentito mai parlare di vettori? Matematica, i grafici cartesiani?'' Sì.

\paragraph{Che cos'è un vettore.} ``In generale, un vettore può essere semplificato, in
uno spazio, come letteralmente una \textbf{freccia}'' che ha tre caratteristiche: una
\textbf{direzione}, una \textbf{magnitudine} (``la sua ampiezza'', cioè la lunghezza) e
un \textbf{verso}. In Unreal i vettori si sono già visti senza nominarli: nella
\textbf{location} e nella \textbf{rotation} di ogni oggetto, ``in particolare molto di
più sulla location''. La location è fatta di tre coordinate, \textbf{X}, \textbf{Y} e
\textbf{Z}, e ``queste particolari coordinate possono formare tra di loro, nello spazio
3D, un vettore''. Il docente parte dalla parola ``velocità di movimento'' e si corregge:
``più che una velocità, è più una direzione'', con la sua ampiezza e il suo verso.
Con i vettori si possono ``fare tante cose in Unreal'', per esempio muovere degli
oggetti; le pale del mulino girano in un verso, e anche se in quel caso si parla di
rotazione, e quindi di \emph{rotator}, ``sono sempre formate da delle coordinate''.

\paragraph{Un punto non è ancora un vettore.} Più avanti, alla lavagna, il docente
aggiunge una precisazione importante: ``ricordatevi che quello che vediamo è un
\textbf{punto} nello spazio, non è una freccia, quindi non è già un vettore''. Un
punto, per esempio $(5, 4, 1)$, è solo una posizione; il vettore è la freccia che
collega due punti, per esempio dall'origine $(0,0,0)$ a $(5,4,1)$\footnote{A lezione il
docente dice che ``un vettore è formato dalla somma di due punti''. La relazione
corretta è la \emph{differenza}: il vettore che va da A a B è B $-$ A (si veda la
sottrazione, poco più sotto). Unreal, peraltro, usa la stessa struttura
\emph{Vector} sia per i punti sia per gli spostamenti.}.

\paragraph{Le operazioni.} Con i vettori ``possiamo fare varie operazioni matematiche,
che sono praticamente quelle classiche'': \textbf{addizione}, \textbf{sottrazione},
\textbf{moltiplicazione} e \textbf{prodotto scalare} (\emph{dot product}), ``che è
un'altra cosa, che poi spiegheremo''. La slide le mostra con due vettori, A e B
(figura~\ref{ue4:fig:vettori}).
\begin{itemize}
  \item \textbf{Addizione.} ``Se andiamo a prendere due vettori e li addizioniamo,
        avremo un vettore unico.'' Il vettore A più il vettore B ``dà come risultato
        il vettore C, che è $(5,4)$'': si mette la coda di B sulla punta di A, e C va
        dall'inizio di A alla fine di B. ``Come fare l'ipotenusa, se vi ricordate i
        triangoli.''
  \item \textbf{Sottrazione.} Dà ``la \textbf{differenza} nello spazio tra due
        punti''. Il risultato dipende dall'ordine: A $-$ B è il vettore che va da B
        ad A, B $-$ A quello che va da A a B. La sua lunghezza è la distanza tra i due
        punti.
  \item \textbf{Moltiplicazione.} ``La useremo tantissimo'': si moltiplica tutto ciò
        che fa parte del vettore, X, Y e Z, e quindi il vettore ``viene allungato nella
        sua direzione''.
  \item \textbf{Prodotto scalare.} ``Statemi un attimo a sentire: è l'unica operazione
        che ci dà un valore \textbf{scalare}'', cioè ``non ci ridà un vettore indietro,
        ma un valore in \emph{float}''. (``Chi si ricorda che è il float?'' ``I numeri
        decimali.'' ``Perfetto.'') Lo si vede nella prossima sezione.
\end{itemize}

\begin{figure}[htbp]
\centering
\begin{tikzpicture}[font=\small, scale=0.8,
  ue4vec/.style={-{Stealth[length=2.5mm]}, very thick},
  ue4grid/.style={gray!25, very thin},
  ue4ax/.style={-{Stealth[length=2mm]}, gray!70}]
  % --- addizione
  \begin{scope}
    \draw[ue4grid] (0,0) grid (6,5);
    \draw[ue4ax] (0,0) -- (6.3,0) node[right, font=\scriptsize]{X};
    \draw[ue4ax] (0,0) -- (0,5.3) node[above, font=\scriptsize]{Y};
    \draw[ue4vec, blue!70!black] (0,0) -- node[left, pos=0.55]{A} (2,3);
    \draw[ue4vec, red!70!black] (2,3) -- node[above, pos=0.5]{B} (5,4);
    \draw[ue4vec, green!45!black] (0,0) -- node[below right, pos=0.6]{C $=$ A $+$ B}
      (5,4);
    \fill (5,4) circle (1.5pt) node[above right, font=\scriptsize]{$(5,4)$};
    \node[font=\bfseries] at (3,-0.9) {addizione};
  \end{scope}
  % --- sottrazione
  \begin{scope}[shift={(8.5,0)}]
    \draw[ue4grid] (0,0) grid (6,5);
    \draw[ue4ax] (0,0) -- (6.3,0) node[right, font=\scriptsize]{X};
    \draw[ue4ax] (0,0) -- (0,5.3) node[above, font=\scriptsize]{Y};
    \fill[blue!70!black] (1,1) circle (2.5pt) node[below left]{A};
    \fill[red!70!black] (5,4) circle (2.5pt) node[above right]{B};
    \draw[ue4vec, green!45!black] (1.15,1.11) -- node[above left, pos=0.5]
      {B $-$ A: da A a B} (4.85,3.89);
    \draw[ue4vec, orange!80!black, dashed] (4.9,3.6) -- node[below right, pos=0.5]
      {A $-$ B: da B ad A} (1.2,0.85);
    \node[font=\bfseries] at (3,-0.9) {sottrazione};
  \end{scope}
  % --- moltiplicazione
  \begin{scope}[shift={(0,-7.2)}]
    \draw[ue4grid] (0,0) grid (6,5);
    \draw[ue4ax] (0,0) -- (6.3,0) node[right, font=\scriptsize]{X};
    \draw[ue4ax] (0,0) -- (0,5.3) node[above, font=\scriptsize]{Y};
    \draw[ue4vec, green!45!black, line width=3pt, opacity=0.35] (0,0) -- (5,4);
    \draw[ue4vec, blue!70!black] (0,0) -- (2.5,2) node[pos=0.5, above left]{V};
    \node[green!45!black, font=\small] at (4.6,4.55) {$2 \times$ V};
    \node[font=\bfseries] at (3,-0.9) {moltiplicazione per uno scalare};
  \end{scope}
  % --- normalizzazione
  \begin{scope}[shift={(8.5,-7.2)}]
    \draw[ue4grid] (0,0) grid (6,5);
    \draw[ue4ax] (0,0) -- (6.3,0) node[right, font=\scriptsize]{X};
    \draw[ue4ax] (0,0) -- (0,5.3) node[above, font=\scriptsize]{Y};
    \draw[gray!60, dashed] (0,0) circle (1);
    \node[gray, font=\scriptsize] at (1.35,-0.3) {lunghezza 1};
    \draw[ue4vec, blue!70!black, opacity=0.45] (0,0) -- (4,3)
      node[above right, font=\scriptsize]{V, lunghezza 5};
    \draw[ue4vec, red!70!black] (0,0) -- (0.8,0.6);
    \node[red!70!black, font=\scriptsize, align=left] at (2.7,0.45)
      {Normalize(V)\\$= (0{,}8;\ 0{,}6)$};
    \node[font=\bfseries] at (3,-0.9) {normalizzazione};
  \end{scope}
\end{tikzpicture}
\caption{Le operazioni sui vettori viste a lezione, nel solo piano XY per semplicità.
In alto a sinistra la somma della slide: A seguito da B porta al punto $(5,4)$. In
alto a destra la sottrazione, il cui verso dipende dall'ordine dei fattori. In basso,
la moltiplicazione per un numero, che allunga il vettore senza cambiarne la direzione,
e la normalizzazione, che lo riporta a lunghezza 1 (sezione~\ref{ue4:sec:dot}). I
valori di A e B della slide non si leggono nella trascrizione e sono scelti per
l'esempio; il risultato $(5,4)$ è quello detto a lezione.}
\label{ue4:fig:vettori}
\end{figure}

\paragraph{Approfondimento: i vettori nel libro di Romero e Sewell.} Il
quattordicesimo capitolo di \emph{Blueprints Visual Scripting for Unreal Engine 5},
dedicato ai nodi matematici e alle \emph{trace}, presenta i vettori con un esempio
molto vicino a quello della lezione. In Unreal la struttura \textbf{Vector} contiene
tre \emph{float}, X, Y e Z, e può rappresentare sia un \textbf{punto} nello spazio
(una \emph{location}) sia uno \textbf{spostamento} o una velocità. Un personaggio sta
in $(50, 0, 20)$, un divano in $(450, 0, 20)$, e un'unità di Unreal vale un
centimetro. Per arrivare al divano il personaggio deve spostarsi di $(400, 0, 0)$: 400
centimetri lungo X. Le due regole che ne derivano sono quelle della lezione, scritte
in formula:
\[
  \text{destinazione} = \text{partenza} + \text{spostamento},
  \qquad
  \text{spostamento} = \text{destinazione} - \text{partenza}.
\]
Ogni operazione si fa componente per componente: $(3,5,0) + (5,2,0) = (8,7,0)$ e
$(6,8,0) - (1,4,0) = (5,4,0)$. Il libro aggiunge due esempi pratici: se V1 è la
posizione nel mondo di un personaggio e V2 la posizione \emph{relativa} della sua arma,
V1 $+$ V2 è la posizione nel mondo dell'arma; e la \textbf{lunghezza} della differenza
tra due posizioni (il nodo \emph{Vector Length}) è la loro distanza. Moltiplicare per un
numero negativo inverte il verso: moltiplicando per $-1$ il vettore ``in avanti'' di un
attore si ottiene quello ``all'indietro'', ed è lo stesso trucco usato nel quinto
capitolo per far cambiare direzione a un bersaglio mobile. Il nodo \textbf{Multiply},
avvertono gli autori, è un nodo \emph{wildcard}: se gli si collega un vettore diventa
un \emph{Vector $\times$ Vector}, e per moltiplicare un vettore per un \emph{float}
bisogna fare tasto destro sul secondo pin e scegliere \emph{Convert Pin\ldots} e poi
\emph{Float}.

\begin{table}[htbp]
\centering
\small
\begin{tabular}{@{}p{3.1cm}p{3.4cm}p{2.0cm}p{5.6cm}@{}}
\toprule
\textbf{Operazione} & \textbf{In formula} & \textbf{Risultato} &
\textbf{A che cosa serve in un gioco} \\
\midrule
Addizione & $(a_x+b_x,\ a_y+b_y,\ a_z+b_z)$ & vettore & spostare un punto: posizione
$+$ spostamento $=$ nuova posizione \\
Sottrazione & $(b_x-a_x,\ b_y-a_y,\ b_z-a_z)$ & vettore & la direzione e la distanza
da A a B (``dove sta il player rispetto a me?'') \\
Moltiplicazione per uno scalare & $(k\,a_x,\ k\,a_y,\ k\,a_z)$ & vettore & allungare o
accorciare; con $-1$ invertire il verso; velocità $\times$ \emph{delta seconds} \\
Lunghezza (\emph{Vector Length}) & $\sqrt{a_x^2+a_y^2+a_z^2}$ & \emph{float} & la
magnitudine; della differenza tra due punti, la distanza \\
Normalizzazione (\emph{Normalize}) & $\vec a \,/\, |\vec a|$ & vettore lungo 1 & tenere
solo la direzione, da moltiplicare poi per una velocità \\
Prodotto scalare (\emph{Dot Product}) & $a_x b_x + a_y b_y + a_z b_z$ & \emph{float} &
capire se due direzioni sono concordi, perpendicolari od opposte \\
\bottomrule
\end{tabular}
\caption{Le operazioni sui vettori della lezione, con il nome del nodo Blueprint
corrispondente dove serve e un uso tipico in un gioco.}
\label{ue4:tab:vettori}
\end{table}

% ---------------------------------------------------------------------
\section{Il prodotto scalare e i vettori normalizzati}
\label{ue4:sec:dot}

Il \textbf{prodotto scalare} è l'operazione che restituisce un \emph{float} invece di un
vettore. ``Che cosa ci dà?'' Ci dice come sono orientati due vettori l'uno rispetto
all'altro. Con i vettori A e B della slide (figura~\ref{ue4:fig:dot}):
\begin{itemize}
  \item se il prodotto scalare è \textbf{maggiore di 0}, i due vettori sono
        ``direzionati nella stessa direzione'', grosso modo concordi; se è
        \textbf{1}, ``sono proprio uguali'', cioè puntano esattamente nello stesso
        verso;
  \item se è \textbf{uguale a 0}, sono ``perfettamente \textbf{perpendicolari}'',
        formano un angolo retto;
  \item se è \textbf{minore di 0}, formano un angolo ottuso, ``al verso opposto'': per
        due vettori come quelli della slide ``magari $-0{,}4$, più o meno''; ``se
        proprio sono così'', uno contro l'altro, è $-1$.
\end{itemize}
Questi valori ``puliti'', tra $-1$ e 1, valgono per i \textbf{vettori normalizzati}, ed
è il motivo per cui la slide successiva li introduce.

\paragraph{I vettori normalizzati.} ``Qualcuno di voi sa cos'è un vettore
normalizzato?'' Uno studente: ``un vettore tra 0 e 1''. ``Esatto'': normalizzando un
vettore, ``indipendentemente dalla sua lunghezza, lo ridurremo a 1''. Un vettore ``che
è lungo 5'', normalizzato, ``lo portiamo a 1''\footnote{Più precisamente: un vettore
normalizzato ha lunghezza \emph{esattamente} 1; sono le sue componenti, prese una a
una, a stare tra $-1$ e 1. Il vettore $(4,3,0)$, lungo 5, normalizzato diventa
$(0{,}8;\ 0{,}6;\ 0)$, come nella figura~\ref{ue4:fig:vettori}.}. La direzione resta la
stessa, cambia solo la lunghezza.

\begin{figure}[htbp]
\centering
\begin{tikzpicture}[font=\small,
  ue4vec/.style={-{Stealth[length=2.5mm]}, very thick}]
  \foreach \xs/\ang/\val/\txt in {0/0/{$= 1$}/{stessa direzione},
                                  3.4/40/{$\approx 0{,}77$}/{concordi},
                                  6.8/90/{$= 0$}/{perpendicolari},
                                  10.2/140/{$\approx -0{,}77$}/{angolo ottuso},
                                  13.6/180/{$= -1$}/{opposti}} {
    \begin{scope}[shift={(\xs,0)}]
      \draw[ue4vec, blue!70!black] (0,0) -- (1.5,0) node[below, pos=0.8,
        font=\scriptsize]{A};
      \draw[ue4vec, red!70!black] (0,0) -- (\ang:1.35) node[above,
        font=\scriptsize]{B};
      \fill (0,0) circle (1.5pt);
      \node[font=\small] at (0.3,-1.0) {A $\cdot$ B \val};
      \node[font=\scriptsize, gray!50!black] at (0.3,-1.5) {\txt};
    \end{scope}
  }
\end{tikzpicture}
\caption{Il prodotto scalare tra due vettori normalizzati A e B, al variare
dell'angolo tra di loro: 1 se puntano nello stesso verso, positivo se sono concordi,
0 se sono perpendicolari, negativo se formano un angolo ottuso, $-1$ se sono opposti.
I valori intermedi sono il coseno dell'angolo ($\cos 40^\circ \approx 0{,}77$).}
\label{ue4:fig:dot}
\end{figure}

\paragraph{Approfondimento: il prodotto scalare è un coseno.} Per due vettori qualsiasi
vale
\[
  \vec a \cdot \vec b \;=\; a_x b_x + a_y b_y + a_z b_z \;=\; |\vec a|\,|\vec b|\,
  \cos\theta ,
\]
dove $\theta$ è l'angolo tra i due. Se i vettori sono normalizzati, le lunghezze valgono
1 e il prodotto scalare \emph{è} il coseno dell'angolo: per questo varia tra $-1$ e 1,
come dicono Romero e Sewell nel quattordicesimo capitolo (``il prodotto scalare di due
vettori normalizzati è uguale al coseno dell'angolo che formano''), e per questo i casi
della slide sono 1, 0 e $-1$. Con vettori non normalizzati il \emph{segno} resta
affidabile (positivo, zero, negativo) ma il valore non è più confrontabile. Nei giochi
il prodotto scalare risponde a domande come ``il player è davanti o dietro al
nemico?'': si normalizza la differenza tra la posizione del player e quella del nemico
e la si confronta con il \emph{Forward Vector} del nemico, che Unreal restituisce già
normalizzato (nodi \emph{Get Actor Forward Vector}, \emph{Right Vector}, \emph{Up
Vector}). Un risultato maggiore di $\cos 45^\circ \approx 0{,}71$, per esempio,
significa ``il player è dentro un cono di visione di 90 gradi''. Il quinto capitolo del
libro usa invece la normalizzazione per il movimento: si normalizza una direzione, la si
moltiplica per una velocità e per il \emph{delta time} (con il nodo \textbf{Get World
Delta Seconds}, che tornerà nella torretta) e la si somma alla posizione a ogni
\emph{tick} con \emph{Add Actor World Offset}. ``Normalizzare'', scrivono, ``è una
procedura comune che assicura che il vettore abbia lunghezza 1, il che lo rende
compatibile con il resto dei calcoli''.

% ---------------------------------------------------------------------
\section{I rotator, il gimbal lock e i quaternioni}
\label{ue4:sec:rotator}

``Andiamo avanti: poi abbiamo i \textbf{rotator}.'' Per spiegarli il docente usa
l'immagine del \textbf{giroscopio}: tre ``anelli che girano tra di loro e formano
proprio quella che è la rotazione''. Ogni anello corrisponde a un asse, con i colori di
sempre: la X in rosso, la Y in verde, la Z in blu (figura~\ref{ue4:fig:rotator}).

\paragraph{Il gimbal lock.} Con le ``regole classiche di Eulero'', cioè descrivendo la
rotazione come tre angoli applicati uno dopo l'altro, c'è ``un possibile problema'': il
\textbf{gimbal lock}. ``Non vi sto a spiegare troppo, perché fondamentalmente andremo a
parlare di quattro dimensioni invece di tre, ed è tutta una cosa molto, molto
complessa.'' In breve, ``potrebbe succedere che due assi si blocchino tra di loro'' e
comincino a ruotare ``sempre dipendenti l'uno dall'altro''. Si risolve con i
\textbf{quaternioni}, che usano quattro valori invece di tre, ma ``parliamo di
matematica molto complessa, e non andiamo molto avanti su questa cosa''. Quello che
conta, per ora: per descrivere una rotazione nello spazio servono sempre la X, la Y e
la Z, ``perché abbiamo tre giroscopi che girano di per sé''. Nell'editor, su un
qualsiasi \emph{asset}, si vedono proprio questi tre valori.

\paragraph{L'obiettivo.} ``Cosa vogliamo fare oggi? Vogliamo creare qualcosa che giri'':
qualcosa che usi il \emph{rotator} per far girare le pale del mulino.

\begin{figure}[htbp]
\centering
\begin{tikzpicture}[font=\small,
  ue4ax/.style={-{Stealth[length=2.5mm]}, very thick},
  ue4rot/.style={-{Stealth[length=2mm]}, thick}]
  % assi in proiezione
  \coordinate (O) at (0,0);
  \draw[ue4ax, red!75!black] (O) -- (-1.9,-1.1) node[below left]{X};
  \draw[ue4ax, green!50!black] (O) -- (2.6,0) node[right]{Y};
  \draw[ue4ax, blue!75!black] (O) -- (0,2.6) node[above]{Z};
  % anelli di rotazione
  \draw[ue4rot, red!75!black] (-1.25,-0.25) arc[x radius=0.35, y radius=0.6,
    start angle=90, end angle=400];
  \node[red!75!black, font=\scriptsize, align=center] at (-2.9,-0.1)
    {\textbf{Roll}\\attorno a X};
  \draw[ue4rot, green!50!black] (1.7,0.55) arc[x radius=0.3, y radius=0.55,
    start angle=90, end angle=400];
  \node[green!50!black, font=\scriptsize, align=center] at (2.4,1.05)
    {\textbf{Pitch}\\attorno a Y};
  \draw[ue4rot, blue!75!black] (0.55,1.75) arc[x radius=0.55, y radius=0.2,
    start angle=0, end angle=310];
  \node[blue!75!black, font=\scriptsize, align=center] at (1.65,2.15)
    {\textbf{Yaw}\\attorno a Z};
  % giroscopio stilizzato
  \begin{scope}[shift={(7.0,0.6)}]
    \draw[very thick, red!75!black] (0,0) ellipse (1.9 and 1.9);
    \draw[very thick, green!50!black] (0,0) ellipse (1.45 and 0.55);
    \draw[very thick, blue!75!black] (0,0) ellipse (0.45 and 1.05);
    \fill (0,0) circle (2pt);
    \node[font=\scriptsize, align=center] at (0,-2.4) {il giroscopio: tre anelli
      annidati,\\ognuno ruota attorno al proprio asse};
  \end{scope}
\end{tikzpicture}
\caption{I tre assi di Unreal con i loro colori e i nomi delle tre rotazioni di un
\emph{rotator}: \emph{roll} attorno a X, \emph{pitch} attorno a Y, \emph{yaw} attorno a
Z. A destra, l'immagine del giroscopio usata a lezione. Schema ricostruito.}
\label{ue4:fig:rotator}
\end{figure}

\paragraph{Approfondimento: roll, pitch, yaw e il gimbal lock in pratica.} Nella seconda
edizione di Romero e Sewell, l'elenco dei tipi di variabile descrive il
\textbf{Rotator} così: contiene i \emph{float} X (\textbf{Roll}), Y (\textbf{Pitch}) e
Z (\textbf{Yaw}), espressi in gradi. I tre nomi vengono dall'aviazione: il
\emph{roll} è il rollio (l'aereo che inclina le ali), il \emph{pitch} il beccheggio (il
muso che sale e scende), lo \emph{yaw} l'imbardata (il muso che gira a destra e a
sinistra). Una torretta che insegue un bersaglio usa lo \emph{yaw} per girarsi e il
\emph{pitch} per alzare la canna; le pale del mulino, come si vedrà, girano in
\emph{pitch}. Il \textbf{gimbal lock} succede quando l'anello di mezzo ruota di 90
gradi: il primo e il terzo anello si ritrovano allineati, e ruotarli produce lo stesso
movimento, così un grado di libertà va perso. Con gli angoli di Unreal il caso tipico
è il \emph{pitch} a $\pm 90$: con il muso dritto verso l'alto, \emph{yaw} e \emph{roll}
fanno la stessa cosa. I \textbf{quaternioni} (in Unreal la struttura \emph{Quat})
rappresentano invece ogni rotazione come una rotazione unica attorno a un asse
qualsiasi, con quattro numeri, e non hanno questo problema. Unreal li usa
internamente: un \emph{rotator} è il formato comodo per le persone, il quaternione
quello sicuro per i calcoli, e nodi come \emph{Combine Rotators} convertono i
\emph{rotator} in quaternioni, li combinano e li riconvertono. Un effetto collaterale,
che conviene conoscere: dato che il \emph{pitch} di un \emph{rotator} viene sempre
riportato tra $-90$ e $90$ gradi, una pala che ruota in \emph{pitch} oltre i 90 gradi
viene descritta ``capovolta'' (\emph{yaw} e \emph{roll} a 180), e chi stampasse il suo
\emph{rotator} vedrebbe numeri che saltano. La rotazione sullo schermo, però, resta
corretta.

% ---------------------------------------------------------------------
\section{Il Rotating Movement Component}
\label{ue4:sec:rotating}

``Allora, avete aperto l'editor di Unreal? Sì? Perfetto.'' Prima di costruire il
mulino, il docente mostra ``quello più semplice'' tra i ``vari modi per far ruotare un
actor in generale''; un modo che però, avverte, ``è difficile usarlo su vari
componenti''.

\paragraph{Un componente aggiunto dai Details.} Si trascina in scena un \textbf{cubo}
(con il tasto \textbf{Fine}, \emph{End}, l'oggetto ``va diretto sul
terreno''\footnote{La trascrizione dice ``il tasto AND'': è il tasto \emph{End} (Fine),
che in Unreal appoggia l'oggetto selezionato sulla superficie sottostante.}). Ora,
``qui sui \textbf{Details} sulla destra, potete aggiungere dei componenti anche non
avendo una schermata di Blueprint'': con il pulsante \textbf{Add} dei Details si può
mettere un componente direttamente sull'actor in scena. Si scrive \emph{rotating
movement} e si sceglie il \textbf{Rotating Movement Component}. ``Trovato? Sì?
Perfetto.''

\paragraph{I parametri.} Il componente ha una serie di dati, tra cui il più importante
è il \textbf{Rotation Rate}, con i tre assi X, Y e Z: la velocità di rotazione su
ciascun asse. ``Di default il Rotating Movement, come potete vedere, è messo a
\textbf{180}'' sulla Z: mezzo giro al secondo attorno all'asse verticale.

\paragraph{La regola della mobilità.} ``Intanto vi dico, ragazzi: come abbiamo visto,
se un \emph{asset} si deve muovere all'interno di uno spazio, la sua mobilità dovrà
essere \textbf{Movable}. Ricordatevelo sempre.'' Se lo si lascia \textbf{Static} e si
va in Play, ``prima di tutto non si muove'', e poi l'engine restituisce un errore,
che il docente legge: \emph{Static Mesh Component has to be movable if you'd like to
move it with Rotating Movement Component}. Con la mobilità corretta, in Play il cubo
gira su sé stesso: ``abbiamo fatto le munizioni di \emph{Metal Gear Solid}'', scherza
il docente, pensando agli oggetti da raccogliere che ruotano a mezz'aria. ``Questo è un
modo per far ruotare'': Unreal ``ci dà la possibilità di poter ruotare direttamente un
nostro bellissimo actor tramite un componente che loro hanno già creato''. ``Un oggetto
ruotante che ruota. Perché `ruotante' è brutto, sembra.''

\paragraph{Esercitazione: i tre cubi.} ``Piccola esercitazione velocissima, faccio
queste botte e risposte con le esercitazioni.'' Si copia il cubo (``tanto vi ricordate
come si copiano'') e si fa in modo che \textbf{un cubo ruoti sull'asse X} e
\textbf{un altro sull'asse Y}, alla velocità che si preferisce. ``Ricordatevi di
metterli Movable. Vi do cinque minuti.'' La correzione: sul secondo cubo si azzera la Z
e si mette 180 sulla X; sul terzo, 180 sulla Y. Si controlla che siano tutti e tre
Movable, si va in Play, ``ed ecco le nostre tre rotazioni''. ``Possiamo farlo con
qualsiasi \emph{asset}.''

\paragraph{Approfondimento: la sedia che gira.} Il Rotating Movement Component è il
primo componente che Romero e Sewell fanno usare nel libro, già nel primo capitolo: un
Blueprint \emph{BP\_RotatingChair} con una \emph{Static Mesh} (la sedia dello
\emph{Starter Content}) e un Rotating Movement. ``Di default'', scrivono, ``questo
componente ruoterà il Blueprint attorno all'asse Z, quindi non serve cambiarne le
proprietà'': si compila, si trascina la sedia nel livello e la si guarda girare in Play.
Il componente fa parte della famiglia dei \emph{movement component}, componenti che non
hanno una forma propria ma muovono un altro componente, quello indicato come
\emph{Updated Component}, che di default è la \textbf{radice} dell'actor. È qui la
ragione del limite che il docente scoprirà tra poco con il mulino: il Rotating Movement
fa girare la radice, e con essa tutto ciò che vi è attaccato. Il componente ha anche
un'opzione \emph{Rotation in Local Space} e un \emph{Pivot Translation}, che sposta il
centro di rotazione; ma per far girare una parte sola di un actor la strada più chiara
resta quella che segue, il Blueprint.

% ---------------------------------------------------------------------
\section{Il Blueprint del mulino}
\label{ue4:sec:mulino}

``Proviamo a fare intanto il nostro \emph{windmill}.'' Il mulino del livello di esempio
è formato da più cose, si vede nei Details. ``Chi si ricorda come si chiamano queste
cose qua, che sono tutte insieme in un singolo \emph{asset}?'' ``Componenti.''
``Sono i \textbf{componenti}'': come si è visto nella seconda lezione, un actor ``può
essere, scusate la ridondanza, composto da più componenti''. In questo caso un
\textbf{body}, il corpo, ``la nostra \emph{mesh} principale, il cubone'', e
\textbf{due pale}.

\paragraph{La classe.} Si crea un nuovo \textbf{Blueprint Class} di tipo
\textbf{Actor}, ``perché non sarà disponibile al giocatore, ma sarà solamente là come
sfondo'': non va controllato, non è un \emph{Pawn}. Il docente lo chiama
\textbf{BP\_Windmill}\footnote{La trascrizione dice ``Windmill, io in questo caso lo
chiamerò 1'': probabilmente \emph{BP\_Windmill1}, per non confonderlo con il mulino già
presente nel progetto.}.

\paragraph{Il corpo.} Nella \textbf{Viewport} del Blueprint si aggiunge una
\textbf{Static Mesh} (Add, in alto a sinistra, \emph{Static Mesh}) e la si chiama
\textbf{Body}. Nei dettagli, alla voce \emph{Static Mesh}, si sceglie un cubo, per
esempio \emph{SM\_Cube}. Il primo cubo scelto dal docente ``ha un \emph{pivot}
particolarmente strano''; ne sceglie un altro, ``molto meglio, perché c'ha il
\emph{pivot} al centro''. Lo solleva di circa 40 sulla Z e gli dà un materiale: il
docente mette la \emph{prototype grid}, ``voi mettete quello che volete''. Poi lo scala
per farne il corpo del mulino: ``3 metri per 5'' di altezza, ``di larghezza facciamo 2
metri'', cioè una scala di circa 3, 2 e 5 sui tre assi\footnote{La trascrizione dice
``x32 sulla y5 sull'altezza'' e ``3:25'': la lettura più probabile è X 3, Y 2, Z 5. Il
cubo di base di Unreal misura un metro di lato, quindi una scala di 3 corrisponde a 3
metri.}.

\paragraph{La viewport che non si trova.} Uno studente non ricorda come si apre la
viewport del Blueprint. ``O ce l'hai qua, o ce l'hai qua, oppure dovresti avere
\textbf{Window} e dovresti vedere \textbf{Viewport}.'' La voce è già spuntata, ma la
finestra non c'è: forse è ``nascosta da qualche schermata di Windows'' (``prova a fare
Alt+Tab''). Si passa alla condivisione dello schermo, e si scopre che il problema è un
altro: lo studente ha creato il Blueprint ma non lo ritrova. La soluzione è il
\textbf{filtro} del Content Browser: le ``tre stanghettine'' vicino alla ricerca,
categoria \emph{Blueprint}, ed ecco il file; doppio clic e, se si apre la versione
ridotta, \emph{Open Full Blueprint Editor}. Poi Add, ma nell'editor in italiano la voce
non si trova subito (è ``Componente Mesh Statica''). ``Effettivamente, ragazzi, vi
consiglio di mettervi l'editor in \textbf{inglese}, perché è molto più semplice. Anzi,
ve lo consiglio proprio'': i nomi dei nodi e dei componenti nei tutorial, nei libri e
nella documentazione sono quasi sempre in inglese.

\paragraph{Le pale.} ``Adesso dobbiamo creare le nostre palette.'' Si copia e incolla il
corpo (Ctrl+C, Ctrl+V): le copie finiscono dentro il corpo, ``quindi cliccatele e
tiratele fuori, perché sono uguali''. Si scalano a forma di pala, una orizzontale e una
verticale, si dà loro un materiale diverso (il docente usa \emph{M\_Grid}, ``questa
qua violacea'') e le si rinomina \textbf{Pala1} e \textbf{Pala2}.

\paragraph{Attaccare le pale al corpo.} ``Importante, ragazzi: queste pale devono essere
attaccate al corpo. Come facciamo? Abbiamo visto l'ereditarietà.'' È la
\textbf{gerarchia} dei componenti: ``il fatto che io posso mettere degli \emph{asset}
attaccati l'uno con l'altro''. Si prende una pala e la si \textbf{trascina sopra} Body,
nell'elenco dei componenti: ora è sua figlia. ``Scusami, non ho capito come si fa.''
``Se hai questa pala, la prendi e la trascini sopra Body.''

\paragraph{Il Rotating Movement non basta.} ``Adesso dobbiamo dare una rotazione a
queste pale. Potremmo provare con il Rotating Movement.'' Il docente lo aggiunge al
Blueprint, mette il mulino in scena (``molto più grosso'') e va in Play. ``Guardate che
succede'': il Rotating Movement è impostato sulla Z, e \textbf{ruota tutto il corpo}.
``Una particolarità del Rotating Movement è che non può essere attaccato a
componenti: o ruota tutto, o non ruota niente.'' Quindi bisogna programmare il
movimento delle pale. Il Rotating Movement va tolto: ``se l'avete messo e avete adesso
il vostro mulino che ruota su sé stesso, lo dovete eliminare''.

% ---------------------------------------------------------------------
\section{Far girare una pala: Event Tick e Combine Rotators}
\label{ue4:sec:pala}

``Domanda. Questo movimento deve essere fatto per ogni \emph{frame}. Quindi deve essere
un evento che riprende ogni \emph{frame}. Chi si ricorda qual è l'evento che lo fa? L'abbiamo
usato per la piattaforma dell'altra volta.'' ``L'\textbf{Event Tick}.'' ``Ottimo.''
L'Event Tick, ``che sarà il nostro protagonista'', è già presente nell'Event Graph di un
nuovo Blueprint.

\paragraph{I delta seconds, di nuovo.} ``Chi si ricorda che sono i \textbf{Delta
Seconds}?'' Uno studente: ``determina quanti tick al minuto\ldots''. Un altro: ``i
secondi effettivi, e non quelli basati su \emph{frame}''. Il docente riprende: è ``lo
spazio che intercorre tra un \emph{frame} e l'altro'', la differenza di tempo tra due
\emph{frame}, ``in questo caso si parla di millisecondi''. Si moltiplica un valore per i
\emph{delta seconds} perché la durata di un \emph{frame} cambia (a 30 \emph{frame} al
secondo è molto più lunga che a 120), e così gli spostamenti ``non diventano più lenti
se abbiamo meno \emph{frame}, o più veloci se ne abbiamo di più''. (Il docente prova a
disegnarlo sullo schermo, ma lo strumento ``a quanto pare non mi funziona benissimo''.)

\paragraph{La velocità.} Si tira fuori un cavo da Delta Seconds e si crea un
\textbf{Multiply}: o scrivendo \emph{multiply}, o digitando semplicemente
l'\textbf{asterisco}, ``e vi dà subito il Multiply''. Il secondo valore è la
\textbf{velocità}: ``più è alto, più è veloce''. Per ora, \textbf{90}.

\paragraph{Make Rotator.} ``Adesso vi dico una cosa.'' Il nodo \textbf{Make Rotator}
costruisce un \emph{rotator}, ``quello che muove a livello di rotazione un
\emph{asset}'', a partire da tre \emph{float}: X, Y e Z. Quale asse serve? ``Prendiamo
la nostra bellissima pala e guardiamo subito in che asse la vogliamo ruotare. La X non
mi sembra, la Z neanche: è la nostra carissima \textbf{Y}'', cioè il \textbf{pitch}.
Si collega il risultato del Multiply al pin Y (Pitch): il \emph{rotator} creato ha X e
Z a 0, ``perché non si devono muovere'', e la Y uguale a velocità per \emph{delta
seconds}.

\paragraph{Combine Rotators.} ``Possiamo farlo in tanti modi, ma oggi ne facciamo tipo
tre.'' Il primo è \textbf{Combine Rotators}, che ``combina due \emph{rotator} per
darne uno in uscita''. Il primo dei due è l'\textbf{origine}: ``a livello concettuale,
qual è l'origine che dobbiamo prendere?'' ``Le coordinate del componente.'' ``Sì:
quale componente?'' Per ora una sola pala, la Pala1. Siccome serve un componente e non
l'intero actor, si trascina la pala dall'elenco dei componenti nel grafo: si ottiene
una \textbf{referenza} a quel componente, ``un \emph{get}'' che ``si riferisce a quel
componente preciso''.

\paragraph{Relativa o del mondo?} Da quella referenza si chiede la rotazione. Ma quale?
``Abbiamo visto la differenza tra \textbf{rotazione relativa} e \textbf{rotazione del
mondo}.'' La pala è attaccata al corpo, ma deve muoversi da sola, indipendentemente dal
corpo. ``Relativa al corpo, secondo me'', propone uno studente. ``Relativa a sé
stessa, quindi la sua asse singola. Se prendiamo quella del mondo, girerà in modo
strano.'' Si usa dunque \textbf{Get Relative Rotation}, la si collega al primo
ingresso di Combine Rotators, e il \emph{rotator} di Make Rotator al secondo.

\paragraph{Set Relative Rotation.} ``Chi si ricorda qual è il blocco che serve per
settare la rotazione? Tanto è un \emph{set}.'' È \textbf{Set Relative Rotation},
``perché, ovviamente, abbiamo preso la relativa''. Lo si crea tirando fuori un cavo da
Pala1 (così il \emph{target} è già la pala), lo si collega all'Event Tick e gli si dà
in ingresso la nuova rotazione, l'uscita di Combine Rotators
(figura~\ref{ue4:fig:pala}). ``E mi raccomando, deve essere \textbf{Movable}.'' Si
mette il mulino in scena, si va in Play: ``ed ecco che una pala ruota''.

\begin{figure}[htbp]
\centering
\begin{tikzpicture}[font=\small,
  ue4node/.style={draw, thick, rounded corners=4pt, align=center, fill=gray!8,
                  minimum height=0.9cm},
  ue4event/.style={ue4node, fill=red!12},
  ue4exec/.style={-{Stealth[length=2.5mm]}, very thick, black!80},
  ue4flt/.style={-{Stealth[length=2mm]}, thick, green!50!black},
  ue4rot/.style={-{Stealth[length=2mm]}, thick, violet!70!black},
  ue4obj/.style={-{Stealth[length=2mm]}, thick, blue!60!black}]
  \node[ue4event, text width=2.1cm] (tick) at (0,0)
    {\textbf{Event Tick}\\\footnotesize Delta Seconds};
  \node[ue4node, text width=1.5cm] (mul) at (2.9,-1.6)
    {\textbf{Multiply}\\\footnotesize $\times$ 90};
  \node[ue4node, text width=1.9cm] (make) at (5.6,-1.6)
    {\textbf{Make Rotator}\\\footnotesize X 0, \textbf{Y}, Z 0};
  \node[ue4node, fill=blue!8, text width=1.3cm] (pala) at (2.9,-3.5)
    {\textbf{Pala1}};
  \node[ue4node, text width=2.1cm] (get) at (5.6,-3.5)
    {\textbf{Get Relative}\\\textbf{Rotation}};
  \node[ue4node, text width=2.0cm, minimum height=1.5cm] (comb) at (8.6,-2.5)
    {\textbf{Combine}\\\textbf{Rotators}\\\footnotesize A, B};
  \node[ue4node, text width=2.2cm, minimum height=1.6cm] (set) at (11.6,0)
    {\textbf{Set Relative}\\\textbf{Rotation}\\\footnotesize Target: Pala1};
  \draw[ue4exec] (tick.east) -- (set.west |- tick.east);
  \draw[ue4flt] ([yshift=-3mm]tick.east) -- ++(0.4,0) |- (mul.west);
  \draw[ue4flt] (mul.east) -- node[above, font=\scriptsize]{Pitch} (make.west);
  \draw[ue4obj] (pala.east) -- (get.west);
  \draw[ue4rot] (get.east) -- ++(0.4,0) |- node[pos=0.75, below,
    font=\scriptsize]{A} ([yshift=-3mm]comb.west);
  \draw[ue4rot] (make.east) -- ++(0.4,0) |- node[pos=0.75, above,
    font=\scriptsize]{B} ([yshift=3mm]comb.west);
  \draw[ue4rot] (comb.east) -- node[above, font=\scriptsize]{New Rotation}
    ++(1.6,0) |- ([yshift=-4mm]set.west);
  \draw[ue4obj, dashed] (pala.south) -- ++(0,-0.5) -| node[pos=0.25, below,
    font=\scriptsize]{Target} (set.south);
\end{tikzpicture}
\caption{Il grafo della prima pala. A ogni \emph{frame}, l'Event Tick esegue Set
Relative Rotation su Pala1; la nuova rotazione è la rotazione relativa attuale della
pala combinata con un piccolo incremento in \emph{pitch}, pari a 90 gradi al secondo
moltiplicati per i \emph{delta seconds}. Il cavo nero è l'esecuzione, gli altri sono
dati (verde i \emph{float}, viola i \emph{rotator}, blu le referenze). Schema
ricostruito.}
\label{ue4:fig:pala}
\end{figure}

\paragraph{I problemi degli studenti.} ``Avete la vostra paletta che ruota?'' Non
tutti. Una studentessa chiede di rivedere il codice e non trova il blocco Get Relative
Rotation: ``tira fuori da Pala1 e scrivi \emph{get relative rotation}''. Un altro
studente vede la pala ``staccarsi da tutto l'asset'' quando preme Play, e in
condivisione si scoprono due errori. Il primo: ha calcolato la rotazione con Combine
Rotators ma non l'ha collegata al Set (``hai calcolato la rotazione, ma non
gliel'hai assegnata''). Il secondo: anche dopo, la pala gira attorno a un punto
sbagliato, perché la \emph{mesh} che ha scelto ha il \textbf{pivot} spostato. ``Vedi che
il pivot sta quasi lì? Cambia la \emph{mesh}'': con un cubo che ha il pivot al centro, la
pala gira correttamente. (Nel frattempo il docente accenna anche a un \textbf{Set World
Rotation}, per mostrare la differenza con la rotazione relativa.)

\paragraph{Approfondimento: relativo e mondo nel libro.} Il quattordicesimo capitolo di
Romero e Sewell spiega la differenza con un esempio molto simile al mulino. Un
Blueprint ha come radice un \emph{DefaultSceneRoot}; sotto c'è un componente
\emph{Table}, un tavolo, e sotto il tavolo un componente \emph{Statue}, una statua. Il
\emph{transform} di Table è \textbf{relativo} a DefaultSceneRoot, e quello di Statue è
relativo a Table: ``se muovete il componente Table nella Viewport, anche il componente
Statue si muoverà; ma se cambiate il \emph{transform} relativo di Statue, Table resterà
dov'è''. Lo stesso vale per il mulino: la rotazione relativa di Pala1 è misurata
rispetto a Body. Se il mulino viene ruotato nel livello, la rotazione relativa della
pala non cambia, e il grafo continua a farla girare attorno al proprio asse; usando la
rotazione del mondo, invece, l'incremento in \emph{pitch} verrebbe applicato attorno
all'asse Y del \emph{mondo}, che coincide con quello della pala solo se il mulino non è
ruotato: è il ``girerà in modo strano'' del docente. Nei vettori lo si vede ancora più
chiaramente: il punto iniziale di ogni posizione nel mondo è sempre $(0,0,0)$, mentre
quello di una posizione relativa è il punto in cui sta il genitore.

\paragraph{Approfondimento: una scorciatoia, Add Local Rotation.} Il trio Get Relative
Rotation, Combine Rotators e Set Relative Rotation è istruttivo perché mostra i tre
passaggi di ogni animazione fatta a mano: \emph{leggi} lo stato attuale, \emph{calcola}
il nuovo, \emph{scrivi} il risultato. Unreal offre anche nodi che fanno tutto in uno:
\textbf{Add Local Rotation} (o \emph{Add Relative Rotation}) su un componente aggiunge
un \emph{Delta Rotation} alla rotazione attuale, così come \emph{Add Actor World
Rotation}, usato nella seconda lezione per la piattaforma, lo fa per un intero actor.
Collegando l'Event Tick a un Add Local Rotation su Pala1, con un Delta Rotation fatto
con Make Rotator come sopra, si ottiene lo stesso mulino con un nodo in meno. Il
risultato è lo stesso; cambia la chiarezza di ciò che succede, ed è per questo che a
lezione conviene vedere prima la versione lunga.

% ---------------------------------------------------------------------
\section{La gerarchia: la seconda pala e il cubetto}
\label{ue4:sec:gerarchia}

``Possiamo fare una cosa bella: invece di dover fare la copia esatta per quanto
riguarda la nostra prima pala, possiamo vedere quello che la \textbf{gerarchia} ci
permette di fare in modo molto facile.'' La Pala2 deve seguire in ogni caso il movimento
della Pala1. Invece di programmare anche lei, la si mette ``in eredità'' (``non quella
con Carlo Conti, ragazzi'') alla Pala1: la si trascina sopra Pala1, e diventa sua
\textbf{figlia}. Play: ``non abbiamo programmato niente'', ma la Pala2 segue la
rotazione della Pala1. ``È la Pala1 che gira, non la Pala2; la Pala2 segue solamente.''

\paragraph{Esercitazione: il cubetto sulla pala.} ``Se noi volessimo attaccare un
cubettino su un estremo della pala\ldots anzi, sapete che c'è? Lo fate voi.'' La
consegna: aggiungere una \textbf{Static Mesh} qualsiasi (``un cubo, una \emph{mesh},
quello che vi pare''), piccola, su un'estremità della pala, e attaccarla alla
\textbf{Pala2}, in modo che segua anche lei il movimento. ``Vi do cinque minuti
classici. Ricordatevi l'ereditarietà: per spostarlo dovete prenderlo proprio e metterlo
sopra.''

\paragraph{La domanda sul pivot.} Uno studente l'ha fatto, ma è sorpreso: ``mi
aspettavo di dover spostare il pivot al centro della rotazione. Invece,
automaticamente, è come se si agganciasse alla pala.'' ``Esatto. Perché comunque
\textbf{mantiene la sua posizione relativa}'': il cubetto ha il suo pivot al proprio
centro, ma la sua posizione è espressa rispetto al genitore, e quando il genitore
ruota, il figlio ruota con lui attorno al pivot del genitore. ``Se lo dovessimo
spostare, andrebbe insieme: manterrebbe la sua origine, il fatto che sta sulla parte
destra, sopra, giù, quello che è, ma si muoverebbe relativo all'altro.'' Play: ``eccolo
che gira insieme''. ``Così abbiamo capito come si ereditano piano piano gli oggetti tra
di loro'' (figura~\ref{ue4:fig:gerarchia}).

\begin{figure}[htbp]
\centering
\begin{tikzpicture}[font=\small,
  ue4comp/.style={draw, thick, rounded corners=3pt, minimum width=2.6cm,
                  minimum height=0.7cm, align=center},
  ue4tree/.style={thick, gray!60!black},
  ue4arr/.style={-{Stealth[length=2.5mm]}, thick}]
  % albero dei componenti
  \node[ue4comp, fill=gray!10] (root) at (0,0) {BP\_Windmill (actor)};
  \node[ue4comp, fill=orange!15] (body) at (0.6,-1.0) {Body};
  \node[ue4comp, fill=violet!12] (p1) at (1.2,-2.0) {Pala1 \footnotesize(ruota)};
  \node[ue4comp, fill=violet!6] (p2) at (1.8,-3.0) {Pala2 \footnotesize(segue)};
  \node[ue4comp, fill=green!10] (cub) at (2.4,-4.0) {Cubetto \footnotesize(segue)};
  \foreach \a/\b in {root/body, body/p1, p1/p2, p2/cub}
    \draw[ue4tree] ([xshift=-1.0cm]\a.south) |- (\b.west);
  \node[font=\footnotesize, align=center] at (1.2,-5.0) {l'elenco dei componenti};
  % mulino stilizzato
  \begin{scope}[shift={(8.2,-2.2)}]
    \draw[thick, fill=orange!15] (-0.75,-2.5) rectangle (0.75,1.0);
    \node[font=\scriptsize] at (0,-2.2) {Body};
    \begin{scope}[rotate=25]
      \draw[thick, fill=violet!25] (-2.2,-0.18) rectangle (2.2,0.18);
      \draw[thick, fill=violet!12] (-0.18,-2.2) rectangle (0.18,2.2);
      \draw[thick, fill=green!35] (-0.25,1.85) rectangle (0.25,2.35);
    \end{scope}
    \fill (0,0) circle (2.2pt);
    \draw[ue4arr, red!70!black] (2.5,0.4) arc[start angle=10, end angle=70,
      radius=2.5];
    \node[font=\scriptsize, align=left] at (3.6,-0.4) {Pala1 orizzontale\\Pala2
      verticale\\cubetto in cima};
  \end{scope}
\end{tikzpicture}
\caption{La gerarchia del mulino. Solo Pala1 è programmata per ruotare; Pala2, sua
figlia, e il cubetto, figlio di Pala2, ne seguono la rotazione senza alcun nodo,
mantenendo la loro posizione relativa. Schema ricostruito (nel mulino della lezione le
pale sono di fronte al corpo).}
\label{ue4:fig:gerarchia}
\end{figure}

% ---------------------------------------------------------------------
\section{Il mulino che guarda il player: Find Look at Rotation}
\label{ue4:sec:lookat}

``Adesso vi farò vedere un altro metodino'' per creare una rotazione, sempre
sull'Event Tick, ``così vi faccio vedere anche un altro blocchettino interessante,
perché sono tanti''. Nei Blueprint, come nella programmazione in generale, ``non esiste
un qualcosa che sia fisso'': ``essendo che è un mondo di spazi, punti, vettori, ce lo
possiamo giocare come ci pare''. Il docente ricorda una persona che una volta ha fatto
``un materiale con la matematica senza usare texture. Non so come abbia fatto, ma si
può fare''.

\paragraph{Staccare i cavi.} Si stacca l'Event Tick dal grafo precedente. ``Ricordatevi:
per staccare i cavettini, tenete premuto \textbf{Alt} e cliccate sui puntini.''

\paragraph{Find Look at Rotation.} ``Un blocco chiamato \textbf{Find Look at Rotation}.
Interessante, molto interessante. Questo è un sacco utile, ve lo dico.'' Il docente
legge la descrizione del nodo: \emph{find a rotation for an object at Start location to
point at Target location}. Alla lavagna (``scrittura incredibile, mi dispiace'') disegna
due punti nello spazio: lo \textbf{Start}, ``da dove parto'', e il \textbf{Target}, il
punto da raggiungere. Il nodo restituisce ``la rotazione verso quel punto'': la
rotazione che un oggetto posto in Start deve avere per guardare Target.

\paragraph{Un falso inizio.} Il docente parte con un esempio sbagliato (lo Start e un
Target ottenuto sommando 360, poi 359, alla posizione), e si ferma: ``scusate, ragazzi,
perdonatemi: il Find Look at Rotation lo facciamo per un esercizio che era diverso.
Ho sbagliato io, perdonate, perché non ho gli appunti.'' L'esercizio vero: far ruotare
il mulino \textbf{in direzione del character}, ``un mulino molto magico'' che si gira
verso il giocatore.

\paragraph{Il grafo.} I due punti sono:
\begin{itemize}
  \item lo \textbf{Start}: la posizione del mulino, con \textbf{Get Actor Location}
        (senza \emph{target}, cioè il Blueprint stesso);
  \item il \textbf{Target}: la posizione del giocatore, con \textbf{Get Player
        Character} e, da questo, Get Actor Location.
\end{itemize}
Il Find Look at Rotation ``ci darà il vettore verso la posizione del nostro player'',
o meglio la rotazione che lo punta.

\paragraph{RInterp To: l'interpolazione.} ``Possiamo usare anche un altro blocco, che ci
darà un'\textbf{interpolazione}.'' Si chiama \textbf{RInterp To} (R come
\emph{rotator}). ``Qualcuno sa cosa vuol dire interpolazione? No.'' ``È letteralmente
il passaggio tra due punti'': da un valore A a un valore B. I nodi di questo tipo
(\emph{Lerp}, \emph{VLerp}, \emph{Interp}) sono ovunque: ``avete presente le animazioni?
Un'animazione è un passaggio da un punto A a un punto B'', e di base è fatta a livello
di valori, perché ruotare un braccio è cambiare i valori di mano, braccio, gamba. Il
nodo ha quattro ingressi:
\begin{itemize}
  \item \textbf{Current}, la rotazione attuale del mulino (Get Actor Rotation, o la
        rotazione relativa del corpo);
  \item \textbf{Target}, la rotazione appena trovata da Find Look at Rotation;
  \item \textbf{Delta Time}, che ``ci dice in quanto tempo'' si fa l'interpolazione;
  \item \textbf{Interp Speed}, ``la velocità in cui lo fa''.
\end{itemize}

\paragraph{Solo lo yaw.} ``Però una cosa importante, ragazzi: non vogliamo che si muova
tutto sulla stessa asse.'' Il mulino deve girare ``al massimo così'', in orizzontale,
non inclinarsi verso l'alto o verso il basso. ``Quale asse dovrò influenzare?''
``La Z?'' ``Esatto, perfetto.'' Si fa \textbf{Split Struct Pin} (tasto destro sul pin)
sull'uscita di Find Look at Rotation e sull'ingresso Target, e si collega solo la Z,
lo \emph{yaw}. Il risultato dell'interpolazione si applica con un nodo di rotazione
(nella trascrizione, \emph{Set Relative Rotation})\footnote{La trascrizione dice ``set
relative rotation di pala 1'', ma a girare verso il player è tutto il mulino. Con ogni
probabilità il nodo è un Set Relative Rotation sul corpo (Body), a cui sono attaccate le
pale, o un Set Actor Rotation sull'intero Blueprint: il risultato visibile è lo
stesso.}, collegato all'Event Tick (figura~\ref{ue4:fig:lookat}).

\paragraph{L'offset di 90 gradi.} In Play il mulino si gira verso il player, ma di
lato. ``Possiamo mettere più 90'': si somma \textbf{90} alla Z prima di collegarla, ``in
modo tale che sia\ldots ecco''. Ora il mulino guarda il giocatore. ``Quello che sta
facendo è continuare a calcolare, visto che lo fa sull'Event Tick, il vettore che
indirizza in un punto, che in questo caso è il player, in modo continuo.'' ``Questo è un
blocco che serve moltissimo, ragazzi'': serve ogni volta che bisogna fare ``delle robe a
\emph{target}'', e servirà per la torretta. ``Vedete che segue continuamente la nostra
torretta? Questo è fatto con il Find Look at Rotation.''

\begin{figure}[htbp]
\centering
\begin{tikzpicture}[font=\small,
  ue4node/.style={draw, thick, rounded corners=4pt, align=center, fill=gray!8,
                  minimum height=0.9cm},
  ue4event/.style={ue4node, fill=red!12},
  ue4exec/.style={-{Stealth[length=2.5mm]}, very thick, black!80},
  ue4flt/.style={-{Stealth[length=2mm]}, thick, green!50!black},
  ue4rot/.style={-{Stealth[length=2mm]}, thick, violet!70!black},
  ue4vecw/.style={-{Stealth[length=2mm]}, thick, yellow!50!black},
  ue4obj/.style={-{Stealth[length=2mm]}, thick, blue!60!black}]
  \node[ue4event, text width=2.0cm] (tick) at (0,0)
    {\textbf{Event Tick}\\\footnotesize Delta Seconds};
  \node[ue4node, text width=2.0cm] (self) at (0,-1.9)
    {\textbf{Get Actor}\\\textbf{Location}\\\footnotesize (self)};
  \node[ue4node, fill=blue!8, text width=2.0cm] (pc) at (0,-3.9)
    {\textbf{Get Player}\\\textbf{Character}};
  \node[ue4node, text width=2.0cm] (pcl) at (2.8,-3.9)
    {\textbf{Get Actor}\\\textbf{Location}};
  \node[ue4node, text width=2.1cm, minimum height=1.4cm] (look) at (6.0,-2.6)
    {\textbf{Find Look at}\\\textbf{Rotation}};
  \node[ue4node, text width=1.2cm] (add) at (8.7,-2.6) {\textbf{$+$ 90}\\
    \footnotesize su Z};
  \node[ue4node, text width=2.0cm] (cur) at (8.7,-4.5)
    {\textbf{Get Actor}\\\textbf{Rotation}};
  \node[ue4node, text width=2.6cm, minimum height=2.2cm] (interp) at (11.5,-2.9)
    {\textbf{RInterp To}\\[2pt]\footnotesize Current\\\footnotesize Target (Z)\\
    \footnotesize Delta Time (Tick)\\\footnotesize Interp Speed};
  \node[ue4node, text width=1.5cm, minimum height=1.3cm] (set) at (14.3,0)
    {\textbf{Set}\\\textbf{Rotation}\\\footnotesize (Z)};
  \draw[ue4exec] (tick.east) -- (set.west |- tick.east);
  \draw[ue4vecw] (self.east) -- node[above, font=\scriptsize]{Start} (4.3,-1.9)
    |- ([yshift=3mm]look.west);
  \draw[ue4obj] (pc.east) -- (pcl.west);
  \draw[ue4vecw] (pcl.east) -- node[below, font=\scriptsize, pos=0.6]{Target}
    (4.5,-3.9) |- ([yshift=-3mm]look.west);
  \draw[ue4flt] (look.east) -- node[above, font=\scriptsize]{Z} (add.west);
  \draw[ue4flt] (add.east) -- ++(0.5,0) |- ([yshift=2mm]interp.west);
  \draw[ue4rot] (cur.east) -- ++(0.5,0) |- ([yshift=-5mm]interp.west);
  \draw[ue4rot] (interp.east) -- ++(0.4,0) |- ([yshift=-4mm]set.west);
\end{tikzpicture}
\caption{Il mulino che guarda il player. A ogni \emph{frame}, Find Look at Rotation
calcola la rotazione che dal mulino (Start) punta al giocatore (Target); se ne tiene
solo lo \emph{yaw} (Z), gli si sommano 90 gradi per compensare l'orientamento della
\emph{mesh}, e RInterp To fa avvicinare gradualmente la rotazione attuale a quella
cercata. Il Delta Time può venire dall'Event Tick, come qui (un cavo non disegnato per
chiarezza), o essere un numero fisso (a lezione, 1); si veda l'approfondimento. Schema ricostruito.}
\label{ue4:fig:lookat}
\end{figure}

\paragraph{``A me non va.''} Seguono tre studenti in condivisione.
\begin{itemize}
  \item Al primo il mulino resta fermo: ``ti manca il Delta Time e
        l'Interp Speed''. Li riempie (``metti uno e due''), compila (``non ho
        compilato, vabbè'') e aggiunge l'offset: ``tira fuori dalla Z, fai più, metti
        90''. ``Perfetto, adesso ti segue.''
  \item Al secondo manca anche a lui il Delta Time: ``metti uno''.
  \item Il terzo, collegando il nodo a Find Look at Rotation, si ritrova un pin
        diviso che non voleva: ha preso \textbf{Get Actor Rotation} invece di
        \textbf{Get Actor Location}. Poi non riesce a collegare l'uscita alla Z:
        ``lo devi splittare: tasto destro sul pinnetto, quello violaceo, quello
        lilla, Split Struct Pin''. E infine l'offset: da Return Value di Find Look at
        Rotation, la Z, un ``più'', 90, e al Target.
\end{itemize}

\paragraph{Approfondimento: che cosa calcola Find Look at Rotation.} Il nodo fa una cosa
semplice: sottrae la posizione di partenza da quella di arrivo, ottenendo il vettore
direzione $\vec d = T - S$ (la sottrazione della sezione~\ref{ue4:sec:vettori}), e
trasforma quella direzione in un \emph{rotator} con \emph{roll} nullo. Per la sola
rotazione orizzontale, quella che interessa al mulino e alla torretta, il conto è
\[
  \text{yaw} \;=\; \operatorname{atan2}(d_y,\ d_x) \quad \text{(in gradi)},
\]
cioè l'angolo, misurato nel piano XY a partire dall'asse X, sotto cui il bersaglio è
visto. Il risultato presuppone che la ``faccia'' dell'oggetto guardi lungo il suo asse
\textbf{X} locale, che in Unreal è la direzione ``in avanti'' (il \emph{Forward
Vector}). Se la \emph{mesh} è stata modellata guardando lungo la Y, come le pale del
mulino della lezione, la rotazione calcolata è sfasata di 90 gradi, e il ``più 90'' del
docente non fa che correggere questo sfasamento. Lo stesso problema, come si vedrà, si
ripresenta con la torretta.

\paragraph{Approfondimento: come funziona RInterp To.} A ogni chiamata, RInterp To
calcola quanto manca per arrivare al Target e ne percorre una frazione, pari
(semplificando) a
\[
  \alpha \;=\; \min\bigl(1,\ \text{Delta Time} \times \text{Interp Speed}\bigr),
  \qquad
  \text{nuova rotazione} = \text{Current} + \alpha \cdot (\text{Target} - \text{Current}),
\]
con la differenza presa sempre ``dal lato più corto'' (da 350 a 10 gradi si passa per
0, non per 180). Ne seguono tre cose utili. Primo: il movimento è veloce quando il
bersaglio è lontano e rallenta avvicinandosi, perché a ogni \emph{frame} si percorre
una frazione di una distanza sempre più piccola; è l'effetto ``morbido'' che si vede in
Play. Secondo: il nodo è pensato per ricevere come \textbf{Delta Time} proprio i
\emph{delta seconds} (dall'Event Tick o da \emph{Get World Delta Seconds}) e come
\textbf{Interp Speed} un numero intorno a 1--10: così la velocità non dipende dal
\emph{frame rate}. Con un Delta Time fisso, come l'1 messo a lezione, la frazione
percorsa a ogni \emph{frame} è sempre la stessa, e il mulino gira più veloce su un PC
che fa più \emph{frame}. Terzo: se il prodotto supera 1, la frazione è 1 e la rotazione
\textbf{scatta} subito al Target; e se l'Interp Speed è 0 o negativo, Unreal restituisce
direttamente il Target. Con Delta Time 1 e Interp Speed 2 il prodotto è 2: il mulino,
in pratica, ``snappa''. Il fenomeno tornerà con la torretta
(sezione~\ref{ue4:sec:problemi}).

% ---------------------------------------------------------------------
\section{Il Clamp, il debug con Print String e la pausa}
\label{ue4:sec:clamp}

``Abbiamo fatto il nostro bellissimo Find Look at Rotation. Potremmo anche
\textbf{clamparlo}.'' ``Si ricorda il Clamp qualcuno di voi?'' ``È quello per dare un
valore minimo'', e massimo. ``Esatto, bravissimo.'' Lo scopo: limitare la rotazione del
mulino, in modo che non possa girarsi del tutto.

\paragraph{Il grafo.} Si splitta il pin con Split Struct Pin, si inserisce un
\textbf{Clamp (float)}, ``perché è un \emph{float}'', sulla Z, e la si ricompone con un
\textbf{Make Rotator} prima di collegarla. Il primo tentativo è un intervallo stretto,
da $-25$ a $25$ gradi, ma qualcosa non torna: i valori in gioco sono molto più grandi
(``mi dice: Pala1 è un botto di valori'').

\paragraph{Il debug con Print String.} ``Questo come ho fatto, ragazzi? Questo è un modo
per \textbf{debuggare} le cose.'' Non sapendo di quanto ruotasse il mulino e quali
limiti dargli, il docente collega un \textbf{Print String} al valore e lo stampa a
schermo mentre si muove attorno al mulino: ``abbiamo stampato letteralmente il
valore''. I numeri che compaiono vanno più o meno da $-160$ a $160$. Il docente prova
quei limiti, poi li allarga: ``snappa un pochino. Facciamo duecento massimo. Molto
meglio. Non dovrebbe andare oltre''\footnote{La trascrizione di questo passaggio è
molto confusa (``facciamo zero, no, mettere meno duecento\ldots{} facciamo duecento
massimo''), e non permette di ricostruire con certezza i valori finali del Clamp.
Quello che conta è il metodo: stampare i valori, osservarli, scegliere i limiti.}.

\paragraph{Approfondimento: Clamp e gli angoli.} Il Clamp restituisce il valore stesso
se è compreso tra Min e Max, altrimenti il limite più vicino:
$\operatorname{clamp}(v, m, M) = \min(\max(v, m), M)$. Con gli angoli c'è una
sottigliezza che spiega i valori ``strani'' della lezione: lo \emph{yaw} restituito da
Find Look at Rotation sta tra $-180$ e $180$, e passando ``dietro'' al mulino salta da
$180$ a $-180$; sommandoci 90 gradi di offset, l'intervallo si sposta (da $-90$ a
$270$). Un limite come $-25 \ldots 25$, applicato a questi numeri, non corrisponde a
``25 gradi a destra e a sinistra della direzione di riposo'', ma a una fetta del cerchio
che dipende da dove sta il player e da come è orientato il mulino nel livello. Il
modo pulito è ragionare sulla differenza tra l'angolo voluto e quello di riposo,
riportarla tra $-180$ e $180$ (Unreal ha il nodo \emph{Normalize Axis}, o \emph{Delta
(Rotator)} per due \emph{rotator}), e solo allora fare il Clamp. Il metodo del docente,
stampare i numeri prima di decidere, resta comunque la prima cosa da fare.

\paragraph{La pausa.} ``Ci siete? C'è qualche cosa che non avete capito? C'è qualche
blocco che vi risulta ancora un pochino nebuloso?'' Tutto chiaro. ``Sono le 17 e 28.
Facciamo pausa, così poi riprendiamo e cominciamo a fare la nostra bellissima
torrettina. Ci vediamo tra 20 minuti.''

% ---------------------------------------------------------------------
\section{Dopo la pausa: la torretta in tre pezzi}
\label{ue4:sec:torretta}

``Eccoci qua. Allora, ragazzi, adesso andremo a fare una bellissima torrettina.'' Il
docente mostra la versione finita. Attorno alla torretta c'è una \textbf{sfera} disegnata
con linee rosse: è il suo raggio d'azione. Quando il giocatore ci entra, la torretta
``fa lo stesso processo che abbiamo visto'', cioè ruota verso il \emph{character}, e
``in questo caso spara dei proiettili''. Non ci sono effetti speciali (``non abbiamo
Niagara''\footnote{La trascrizione dice ``BevoFX'': con ogni probabilità
\emph{Niagara VFX}, il sistema di effetti visivi di Unreal Engine, di cui Davide aveva
parlato nella lezione precedente.}), ma si costruisce ``la logica che ci servirà con la
torretta che avete \emph{mergiato}''. Si vedono anche delle linee che partono dal
cannone: sono un \textbf{line trace}, ``che permette a noi di avere una corrispondenza
dipendente da cosa incontra''. Oggi si fa la rotazione; proiettili e \emph{line trace}
saranno la prossima volta.

\paragraph{Il Blueprint.} ``La prima cosa che ovviamente dobbiamo fare è creare un
nuovo Blueprint, sempre \textbf{Actor}'', chiamato \textbf{BP\_Turret}. Ma come hanno
composto la torretta gli studenti, con Davide? ``Abbiamo preso i tre pezzi forniti e li
abbiamo uniti'' con \textbf{Merge Actors}: una \emph{mesh} unica. ``Allora facciamo una
cosa: la facciamo a \textbf{tre pezzi}.'' Per far girare il supporto e il cannone
indipendentemente dalla base servono tre componenti distinti, quindi si torna ai tre
pezzi importati con Davide: tre \textbf{Static Mesh} nel Blueprint, una per la
\textbf{base}, una per il supporto (``l'\emph{holder}, chiamiamolo'', poi rinominato
\textbf{Pivot}) e una per il \textbf{cannone} (``la \emph{gun}'', poi
\textbf{Cannon}).

\paragraph{Chi segue chi.} I tre pezzi vanno ``uno dentro l'altro''. Il docente parte
con l'idea che il supporto segua il cannone, poi si corregge: ``fate il contrario, che
è meglio: fate che la \emph{gun} segua l'\emph{holder}''. La gerarchia giusta è quindi:
la \textbf{base} come radice, il \textbf{Pivot} figlio della base, il \textbf{Cannon}
figlio del Pivot (figura~\ref{ue4:fig:turretbp}). Così, quando il Pivot ruota
orizzontalmente, il cannone gira con lui; e un giorno il cannone potrà alzarsi e
abbassarsi da solo, senza trascinarsi dietro il supporto.

\paragraph{Un trucco: Browse to Asset.} Per assegnare le \emph{mesh} ai componenti non
serve cercarle a mano nel Content Browser ``con l'elenchino e tutto quanto''. ``Andate
qua sull'\emph{asset} che avete in scena, cliccate tasto destro e fate \textbf{Browse to
Asset}'': il Content Browser salta direttamente alla posizione del file, e da lì lo si
può \textbf{trascinare} (\emph{drag and drop}) nello slot \emph{Static Mesh} del
componente. In alternativa, come sempre, si clicca sullo slot e si scrive il nome:
\emph{turret base}, \emph{turret pivot}, \emph{cannon}. Poi si aggiustano le posizioni
(``un buchettino sta un po' più giù'') finché i pezzi combaciano.

\paragraph{Un problema di Add.} Una studentessa, premendo Add, vede solo la \emph{mesh}
intera della torretta (\emph{SM\_\allowbreak Torretta}) e nessun'altra opzione. In condivisione si
prova a compilare e a riaprire il Blueprint, e alla fine il componente si riesce ad
aggiungere: ``va benissimo così, abbiamo risolto. A volte Unreal può fare dei giochi
simpatici.''

% ---------------------------------------------------------------------
\section{Il raggio d'azione: la Sphere Collision}
\label{ue4:sec:sfera}

``Adesso dobbiamo creare quello che abbiamo visto qua: vedi questo cerchio rosso intorno
alla torretta? Questo è una \textbf{Sphere Collision}.'' È il \textbf{range}: entrandoci
dentro, la torretta si attiva.

\paragraph{Aggiungerla.} In BP\_Turret si aggiunge, tra i componenti, una
\textbf{Sphere Collision}, ``che tranquillamente può stare anche esterna alla base'',
cioè non deve essere figlia di un pezzo che ruota. Appena aggiunta non si vede quasi,
``semplicemente perché è molto piccola''. Cliccando su \emph{Sphere}, nei dettagli c'è
lo \textbf{Sphere Radius}: di default è 32; il docente lo porta a \textbf{800}. C'è anche
la \textbf{Line Thickness}, lo spessore delle linee con cui la sfera è disegnata
nell'editor, ``che serve magari solo per vederla un po' più spessa''.

\paragraph{Hidden in Game.} In Play, però, la sfera non si vede. ``Di base, questa è una
\textbf{collisione}'', una cosa interna agli \emph{asset}, come le collisioni della
\emph{mesh} viste con Davide: esiste per la logica, non per essere guardata. Per
vederla mentre si prova il gioco (``per avere facilità nel vedere quando entriamo e
usciamo'') si cerca nei dettagli \textbf{Hidden in Game} e lo si toglie. ``Parlando di
\emph{rendering}, Hidden in Game dice: so che logicamente esiste all'interno del mondo,
ma se non è \emph{hidden in game}, la renderizzo pure.'' Tolta la spunta, il valore è
\emph{false}, ``quindi non è nascosto in gioco'', e in Play ecco il range. È la stessa
opzione usata nella seconda lezione, al contrario, per rendere invisibile il soffitto
della vaschetta.

\paragraph{Ripetere i passaggi.} Uno studente ha perso questa parte mentre cercava i
pezzi della torretta, e chiede di rivedere come si mettono le \emph{mesh} nei tre
componenti. Il docente lo rifà: si clicca sul componente, si scrive il nome della
\emph{mesh} nello slot Static Mesh (\emph{turret base}, \emph{turret pivot}), si
rinominano i componenti in modo che siano coerenti (\textbf{Pivot}, \textbf{Cannon},
``come a me è chiamata Cannon nella nostra \emph{mesh}'').

% ---------------------------------------------------------------------
\section{Gli eventi dei componenti e l'overlap}
\label{ue4:sec:overlap}

``Adesso andiamo a vedere un evento in particolare.'' Oltre agli eventi già visti, che
appartengono all'intero actor (l'Event Begin Play, l'Event Actor Begin Overlap,
l'Event Tick), ``possiamo richiamare degli eventi anche all'interno dei
\textbf{componenti}''. Selezionando \emph{Sphere} e scorrendo in fondo ai dettagli, c'è
una sezione \textbf{Events} piena di pulsanti: \emph{On Component Begin Overlap},
\emph{On Component End Overlap}, \emph{On Component Wake}, \emph{On Component Sleep},
gli eventi del cursore e altri. ``Tutta una serie di eventi'' che si possono usare
``per far sì che avvenga qualcosa in un preciso momento''. Qui la torretta deve partire
quando il player entra nella sfera: si clicca sul \textbf{+} di \textbf{On Component
Begin Overlap}, e nel grafo compare l'evento, già legato alla Sphere.

\paragraph{Che cos'è un overlap.} ``Sapete cosa si intende per \emph{overlap}?''
Uno studente: ``tipo appena si passa un confine, qualcosa del genere''. ``In un certo
senso.'' Alla lavagna (``perdonate le mie doti molto brutte di disegno'') il docente
disegna la sfera con la torretta al centro: un \emph{overlap} avviene quando il player,
``o un qualsiasi actor'', \textbf{tocca o entra} in quel volume. Un actor fuori non sta
sovrapponendosi; uno dentro sì. E uno a cavallo del bordo? ``Sta overlappando o non sta
overlappando?'' ``Sì.'' ``Qualcun altro ha un'opinione differente?'' Sta
sovrapponendosi, ``perché anche un qualsiasi poligono da cui è formato è parte
integrante di esso'': basta un pezzo (figura~\ref{ue4:fig:overlap}). Ci sono quindi tre
situazioni: non si sovrappone, \textbf{comincia} a sovrapporsi (Begin Overlap),
\textbf{smette} di sovrapporsi (End Overlap, ``quando ci esce''). Per ora serve solo la
prima: ``vogliamo che quando lui entra in questo volume la torretta faccia tipo: ti ho
visto, mo' comincio a ruotare verso di te''.

\begin{figure}[htbp]
\centering
\begin{tikzpicture}[font=\small,
  ue4arr/.style={-{Stealth[length=2.5mm]}, thick}]
  \draw[red!70!black, thick, fill=red!4] (0,0) circle (2.4);
  \node[red!70!black, font=\scriptsize] at (0,2.65) {Sphere Collision (raggio 800)};
  \draw[thick, fill=orange!35] (-0.35,-0.3) rectangle (0.35,0.3);
  \draw[thick, fill=orange!20] (0.35,-0.08) rectangle (0.9,0.08);
  \node[font=\scriptsize] at (0,-0.55) {torretta};
  % fuori
  \draw[thick, fill=blue!15] (-4.3,0.6) circle (0.32);
  \node[font=\scriptsize, align=center] at (-4.3,-0.2) {fuori:\\nessun overlap};
  % a cavallo
  \draw[thick, fill=blue!15] (1.75,1.65) circle (0.32);
  \node[font=\scriptsize, align=left, anchor=west] at (2.3,1.9) {a cavallo del bordo:\\
    \textbf{overlap} (basta un pezzo)};
  % dentro
  \draw[thick, fill=blue!15] (-1.0,-1.25) circle (0.32);
  \node[font=\scriptsize, align=center] at (-1.0,-1.95) {dentro: overlap};
  % frecce degli eventi
  \draw[ue4arr, green!45!black] (-3.9,0.45) -- (-2.55,0.25);
  \node[green!45!black, font=\scriptsize, align=center] at (-3.4,1.45)
    {entra:\\\textbf{Begin Overlap}};
  \draw[ue4arr, orange!80!black] (1.95,-1.55) -- (3.2,-2.3);
  \node[orange!80!black, font=\scriptsize, align=left, anchor=west] at (3.3,-2.3)
    {esce: \textbf{End Overlap}};
\end{tikzpicture}
\caption{Le situazioni dell'\emph{overlap} disegnate a lezione. Un actor si sovrappone
alla sfera appena una parte qualsiasi della sua forma di collisione ne entra; Begin
Overlap scatta una volta, nel momento in cui entra, End Overlap una volta, quando ne
esce del tutto. Schema ricostruito.}
\label{ue4:fig:overlap}
\end{figure}

\paragraph{Other Actor.} ``Come facciamo a capire quale actor ci dà indietro?'' L'evento
ha una serie di uscite: \emph{Overlapped Component}, \textbf{Other Actor}, \emph{Other
Comp}, \emph{Other Body Index}, \emph{From Sweep}, \emph{Sweep Result}. ``Sono tutti
input che sono risultanti dall'overlap'', ``in parole povere'' le informazioni su chi si
è sovrapposto a che cosa. Per mostrarlo, il docente collega Other Actor a un
\textbf{Print String} (Unreal inserisce da solo un \emph{Get Display Name}, per
trasformare l'actor in testo), va in Play ed entra nella sfera: in alto a sinistra
compare \emph{BP\_ThirdPersonCharacter\_0}\footnote{La trascrizione dice ``BP ha perso
un character 0''. Il suffisso \_0 è il numero che Unreal assegna a ogni istanza di una
classe presente nel livello.}. ``Mi ridà cosa sta overlappando, che in questo caso è il
\emph{character}.'' Esistono anche gli \emph{overlapped component} e gli \emph{other
component}, ma quello che interessa è l'Other Actor, ``perché il nostro
\emph{character} è un actor''.

\paragraph{Approfondimento: perché l'overlap funziona ``da solo''.} Perché Unreal generi
un evento di \emph{overlap} servono tre condizioni: entrambi i componenti devono avere
attivo \emph{Generate Overlap Events}; le loro risposte di collisione reciproche devono
essere \emph{Overlap} (e non \emph{Ignore}); e almeno uno dei due deve muoversi. La
Sphere Collision nasce con il profilo di collisione \emph{OverlapAllDynamic}, che si
sovrappone a tutto ciò che si muove, e la capsula del \emph{character} è un oggetto
dinamico di tipo \emph{Pawn}: per questo, nella lezione, non c'è stato bisogno di
toccare nulla. Se un giorno un evento di \emph{overlap} non scatta, è la prima cosa da
controllare, nei dettagli di entrambi i componenti, sezione \emph{Collision}. La seconda
edizione di Romero e Sewell elenca proprio gli eventi di collisione (``eseguiti quando
due actor collidono o si sovrappongono'') tra le famiglie principali di eventi, accanto
a quelli di input e all'Event Begin Play.

% ---------------------------------------------------------------------
\section{Is Valid: esiste davvero?}
\label{ue4:sec:isvalid}

``Vogliamo essere sicuri che il nostro actor in questione sia proprio il player. Perché
un actor può essere anche un cubo, volendo.'' Prima, però, un controllo più semplice. Si
tira fuori da Other Actor un \textbf{Is Valid}: ``ce ne avete tanti, ma quello che ci
interessa è quello col \textbf{punto interrogativo}'', la \emph{macro} con due uscite di
esecuzione, \textbf{Is Valid} e \textbf{Is Not Valid}.

\paragraph{Che cosa controlla.} ``Is Valid, statemi a sentire, ci dà una risultante: è
valido o non è valido?'' Qui chiede: ``hai overlappato qualcosa? È vero, esiste questa
cosa?'' In generale, Is Valid serve a dire se esiste o no qualcosa, ``che può essere
anche una \textbf{variabile}, quindi una variabile vuota''.

\paragraph{L'esempio della variabile.} Per farlo vedere il docente costruisce un piccolo
esperimento:
\begin{enumerate}
  \item da Other Actor crea una variabile con \textbf{Promote to Variable};
  \item mette un Is Valid su quella variabile, all'inizio dell'evento di \emph{overlap};
  \item sul ramo \textbf{Is Not Valid} assegna alla variabile l'actor appena
        sovrapposto (il \emph{Third Person Character});
  \item sul ramo \textbf{Is Valid} stampa con un Print String ``valore registrato'' e il
        nome contenuto nella variabile.
\end{enumerate}
Prima entrata nella sfera: non succede nulla di visibile, ``è normale''. Il Blueprint ha
chiesto: ``questa variabile ha qualcosa dentro? No, non c'è niente. Quindi non è
valida'', e allora le ha assegnato l'actor. Si esce e si rientra: ``ah, me l'hai fatto
prima, quindi questa cosa c'ha qualcosa dentro'', e il ramo è Is Valid: compare
\emph{valore registrato: BP\_ThirdPersonCharacter\_0}. ``Qualcuno ha qualche dubbio
riguardo all'Is Valid? Abbastanza chiaro?'' L'esperimento si cancella: ``ci serve
solamente a capire cosa fa l'Is Valid''.

\paragraph{Approfondimento: la regola di Romero e Sewell.} Il quarto capitolo di
\emph{Blueprints Visual Scripting for Unreal Engine 5}, dedicato alla comunicazione tra
Blueprint, mette in evidenza una nota: ``usate sempre la macro Is Valid prima di
eseguire una funzione usando una variabile che contiene una referenza a un oggetto. La
variabile può non essere valida per diversi motivi, ed eseguire una funzione con una
variabile non valida causerà un errore durante l'esecuzione''. I motivi tipici sono
due: la variabile non è mai stata riempita (come nella prima entrata dell'esempio), o
l'oggetto a cui puntava è stato distrutto nel frattempo. In entrambi i casi un nodo che
prova a leggere la posizione di ``nessuno'' produce nel registro il classico errore
\emph{Accessed None}. Il libro usa lo stesso schema in molti punti: un Is Valid prima di
distruggere un actor, prima di applicare un danno a un bersaglio, prima di leggere un
salvataggio (se non è valido, lo si crea). La torretta lo userà due volte, nella
prossima sezione e nell'evento Move Turret.

% ---------------------------------------------------------------------
\section{I tag e Actor Has Tag}
\label{ue4:sec:tag}

Is Valid dice che \emph{qualcosa} è entrato. Per sapere \emph{chi}, si tira fuori da
Other Actor un \textbf{Actor Has Tag}, ``un blocco che va a esaminare una cosa chiamata
\textbf{tag}''. I tag ``sono dei \emph{label}, delle \textbf{etichette} che noi possiamo
mettere all'interno di vari \emph{asset} del gioco''.

\paragraph{Mettere il tag al player.} Si apre il \textbf{Third Person Character} e,
``non dobbiamo fare niente di particolare'', si scrive \emph{tag} nella ricerca dei
dettagli: compare \textbf{Actor} $\rightarrow$ \textbf{Tags}, un elenco vuoto con un
\textbf{+}. Si aggiunge un elemento e lo si chiama \textbf{Player}. ``Cosa abbiamo
fatto? Abbiamo letteralmente etichettato il player con un tag Player'': è come avere un
cartellino di riconoscimento sulla schiena, con scritto sopra ``tu sei il player''.

\paragraph{Il controllo.} Tornati nella torretta, Actor Has Tag chiede: il
\emph{target}, cioè l'actor sovrapposto (Other Actor), ha il tag \textbf{Player}? Il
nodo restituisce un valore \textbf{booleano}: ``sì, questo actor che ho overlappato ha
il tag Player'', e si va nel \textbf{True} di un \textbf{Branch}; ``se no, invece andrà
in False''. Il docente lo mostra con due Print String, ``Player trovato'' sul ramo True e
``Player non trovato'' sul ramo False: entrando nella sfera, ``Player trovato''.

\paragraph{Attenzione all'ortografia.} ``Non lo fate, ragazzi, faccio solo vedere per
esempio.'' Il docente modifica il tag, togliendo la R finale, e rientra: ``Player non
trovato''. ``Questo vi fa vedere anche che dovete essere precisi nel dire qual è il
tag.'' Il consiglio: ``fate Ctrl+C, Ctrl+V del tag che avete scritto'', copiandolo dal
\emph{character} nel nodo invece di riscriverlo a mano.

\paragraph{Approfondimento: tag, Cast o confronto?} Ci sono almeno tre modi per
rispondere alla domanda ``l'actor che è entrato è il player?''. Il primo è il
\textbf{tag}, quello della lezione: semplice, e indipendente dalla classe (lo stesso tag
\emph{Player} può stare su personaggi diversi, e un tag \emph{Enemy} su nemici di classi
diverse); esiste anche \emph{Component Has Tag}, per etichettare singoli componenti. Il
secondo è il \textbf{Cast To} (per esempio \emph{Cast To BP\_ThirdPersonCharacter}),
che verifica la classe e, se riesce, restituisce l'actor già ``convertito'', con
accesso alle sue variabili e funzioni. Il terzo è il confronto diretto tra oggetti: nel
capitolo sull'intelligenza artificiale della seconda edizione di Romero e Sewell, un
nemico che pattuglia usa un \emph{Event Actor Begin Overlap} seguito da un Branch su un
nodo \textbf{Equal (Object)}, che controlla se l'Other Actor è proprio il punto di
pattuglia salvato in una variabile. Una curiosità sui tag: sono di tipo \emph{Name}, e
in Unreal il confronto tra \emph{Name} \textbf{non distingue maiuscole e minuscole},
per cui \emph{player} e \emph{Player} risultano uguali; quello che conta, come ha
mostrato il docente, è che le lettere siano le stesse. Copiare e incollare il tag
resta comunque la via più sicura.

% ---------------------------------------------------------------------
\section{La variabile del bersaglio e i Custom Event}
\label{ue4:sec:custom}

Ora che il player è stato trovato, bisogna ricordarselo: ``andiamo a crearci una
variabile in cui mettiamo il nostro player che è stato trovato'', di tipo
\textbf{Actor}. Il modo più veloce: tirare fuori un cavo da Other Actor e scegliere
\textbf{Promote to Variable}. Si crea un nodo \emph{Set} di una variabile nuova, che si
rinomina \textbf{PlayerTarget}. Per rinominare, ``o vai qua e fai \textbf{F2}, che è più
facile, oppure fai tasto destro e Rinomina''.

\paragraph{Un Custom Event.} ``Manca mezz'ora'', e il docente si chiede se si riuscirà a
finire (``questo lo continueremo la prossima volta''), ma va avanti. ``Adesso qua dentro
andremo a fare un \textbf{Custom Event}.'' Si scrive \emph{custom event} nel grafo, si
sceglie \emph{Add Custom Event} e gli si dà un nome: \textbf{TrackTargetTimer}. ``Come
potete vedere, questo bellissimo evento è un Custom Event: è uguale agli altri'', come
il Begin Overlap, il Tick o il Begin Play, ``ma è un evento che abbiamo creato noi'',
``che possiamo richiamare quando vogliamo, con delle regole specifiche''. Gli eventi di
Unreal sono ``prefissati'': l'\emph{overlap} scatta perché l'editor di Unreal lo
controlla già per conto suo. I Custom Event, invece, partono quando li si chiama.

\paragraph{Chiamarlo.} Per mostrarlo, il docente collega al Custom Event un Print String
che stampa ``Hello'' e lo chiama dal ramo del Begin Overlap, subito dopo il Set di
PlayerTarget: si scrive il nome dell'evento (\emph{TrackTargetTimer}) e compare il nodo
che lo invoca. In Play, entrando nella sfera: ``Hello''. ``Ho creato un evento
customizzato che richiamo quando entro nella Sphere Collision. Compreso?'' ``Sappiate
che di Custom Event ne potete creare a milioni'': basta scrivere \emph{custom event} e
dare un nome.

\paragraph{Approfondimento: eventi, funzioni e macro.} La seconda edizione di Romero e
Sewell ricorda che nell'Event Graph ogni evento di Unreal può comparire una sola volta,
mentre i \textbf{Custom Event} si possono creare a volontà; e offre una tabella per
scegliere tra \emph{macro}, \emph{funzioni} ed \emph{eventi}. In sintesi: tutti e tre
accettano parametri in ingresso; solo macro e funzioni restituiscono valori; le funzioni
e gli eventi possono essere chiamati da altri Blueprint, le macro no; e solo macro ed
eventi possono contenere azioni \emph{latenti}, cioè che aspettano (come un
\emph{Delay}), mentre le \emph{Timeline} stanno solo negli eventi. Il capitolo sui
vincoli di gioco della stessa edizione usa proprio un Custom Event, \emph{Sprint Drain},
richiamato da un timer, per consumare la stamina del giocatore mentre corre: ``un Custom
Event permette di attivare i nodi collegati a lui ogni volta che un altro nodo lo
chiama; così gruppi di nodi dello stesso Blueprint possono comunicare tra loro anche
quando non sono collegati direttamente''. È lo schema esatto della torretta.

% ---------------------------------------------------------------------
\section{Il timer: Set Timer by Event}
\label{ue4:sec:timer}

``Cosa vogliamo fare con questo Custom Event?'' Durante tutto il tempo, ``come un Event
Tick'', la torretta deve cercare la posizione verso cui ruotare per seguire il player.
``Lo potremmo fare tranquillamente nell'Event Tick, volendo, ma possiamo fare una cosa
interessante: possiamo fare un \textbf{timer}. L'abbiamo chiamato timer apposta.''

\paragraph{Il nodo.} Dal Custom Event si tira fuori \textbf{Set Timer by Event}, un
blocco ``che ci permette di far partire un timer che ci chiede ogni quanto''. I suoi
parametri, così come li imposta il docente:
\begin{itemize}
  \item \textbf{Time}: ogni quanti secondi scatta. ``Vogliamo che lo faccia allo stesso
        tempo del Tick.'' Si potrebbe portare fin qui il cavo dei \emph{delta seconds}
        dall'Event Tick (``ma è un po' brutto''); c'è invece un nodo apposta,
        \textbf{Get World Delta Seconds}, ``un blocchetto semplicemente per comodità'',
        ``per non far sì di filare il Tick così, a cavolo'';
  \item \textbf{Looping}: spuntato, ``vogliamo che sia un timer che \emph{loopi}'';
  \item \textbf{Max Once Per Frame}: spuntato anche questo, ``massimo una volta per
        ogni \emph{frame}''. ``Ci permette almeno di dire: ok, per ogni \emph{frame} si
        calcola, tanto effettivamente si muoverà per ogni \emph{frame}. Questa è una
        cosa un po' più nel micro, però noi mettiamo \emph{max}, in modo che sia più
        ottimizzato.''
\end{itemize}

\paragraph{Quale evento?} ``Set Timer by Event: quale evento?'' Sul nodo c'è un
\textbf{quadratino rosso}, il pin \textbf{Event}, ``che guarda caso è uguale a quello del
Custom Event'': un pin di tipo \emph{delegato}. Vuol dire che lo si può tirare fuori e
creare da lì un \textbf{altro Custom Event}, quello che il timer chiamerà a ogni giro:
lo si chiama \textbf{MoveTurret}. È lì che si farà muovere la torretta.

\paragraph{Come ragiona.} Per farlo vedere, a MoveTurret si collega il solito Print
String, ``Hello''. La catena completa è: Begin Overlap, Is Valid, Branch su Actor Has
Tag, Set PlayerTarget, TrackTargetTimer; e TrackTargetTimer fa partire un timer, alla
velocità dei \emph{delta seconds}, in \emph{loop}, al massimo una volta per
\emph{frame}, che chiama MoveTurret, che stampa ``Hello''. ``Domanda: quando mi
stamperà Hello?'' ``Da quando entri in range.'' ``Ottimo, esatto.'' Appena il player
entra, Hello a ogni \emph{frame}. ``Esco fuori: si cancellano\ldots{} ah, no, scusate,
non si stanno cancellando, è vero. Perché non stoppiamo mai il timer.'' Dalla prima
entrata in poi, il timer continua all'infinito: ``ciò vuol dire che il timer a un certo
punto va disattivato; lo vedremo dopo''.

\paragraph{Approfondimento: timer o Event Tick?} Romero e Sewell presentano i timer, già
nel secondo capitolo, come un'alternativa all'Event Tick: tutto ciò che è collegato al
Tick viene eseguito a ogni \emph{frame}, e ciò che non ha bisogno di aggiornarsi così
spesso può girare a intervalli più lunghi. Nel quinto capitolo un Custom Event
\emph{ChangeDirection}, chiamato da un Set Timer by Event con Looping attivo, fa
invertire periodicamente la direzione di un bersaglio mobile; nella seconda edizione un
timer chiama ripetutamente \emph{Sprint Drain}, e un \emph{Clear Timer by Function Name}
lo ferma quando il giocatore smette di correre. Nella torretta della lezione l'intervallo
è pari ai \emph{delta seconds}, quindi il timer scatta (circa) una volta per
\emph{frame}, come un Tick: il vantaggio non è la frequenza, ma il fatto che il timer
\textbf{parte solo quando serve} (all'ingresso del player) e \textbf{si può fermare}.
Tre dettagli utili. Primo: il valore di Get World Delta Seconds viene letto una sola
volta, quando il timer viene impostato, e resta fisso anche se poi il \emph{frame rate}
cambia. Secondo: se l'intervallo di un timer è più breve di un \emph{frame}, Unreal può
eseguirlo più volte nello stesso \emph{frame} per ``recuperare''; \textbf{Max Once Per
Frame}, un'opzione delle versioni recenti di Unreal Engine 5, impedisce proprio questo.
Terzo: per fermarlo, Set Timer by Event restituisce un \textbf{Timer Handle}, una
``maniglia'' che conviene salvare in una variabile (Promote to Variable) e passare,
quando serve, a \textbf{Clear and Invalidate Timer by Handle}. Il posto naturale per
farlo è l'evento \textbf{On Component End Overlap} della sfera, insieme a
\emph{svuotare} la variabile PlayerTarget (un Set senza nulla collegato): altrimenti
l'Is Valid di MoveTurret continuerà a trovare un bersaglio valido anche dopo che il
player è uscito dal range.

% ---------------------------------------------------------------------
\section{Move Turret: la torretta che segue il player}
\label{ue4:sec:move}

``Cosa dobbiamo fare? L'abbiamo fatto prima, praticamente'': un \textbf{Find Look at
Rotation}, ``perché dobbiamo ruotare il nostro Pivot in direzione del player''.

\paragraph{Prima, un Is Valid.} ``Intanto facciamo un Is Valid. Perché? Perché diciamo:
ma ce l'ho un target?'' Si trascina nel grafo la variabile \textbf{PlayerTarget} e la si
collega a un Is Valid, in testa a MoveTurret. ``Ce l'ho un target effettivo, o non ce
l'ho più, l'ho perso, sono uscito fuori? Se ce l'ho, lo faccio e mi muovo. Se non ce
l'ho, non lo faccio'': senza questo controllo la torretta ``continuerebbe a ruotare a
cavolo''.

\paragraph{Start e Target.} ``Del Find Look at Rotation ci servono due punti nello
spazio. Quali sono?'' ``Lo Start è l'\emph{holder}, che tanto la \emph{gun} è collegata
all'\emph{holder}, e il Target è il player.'' ``Esatto.'' Per lo Start si usa
\textbf{Get Actor Location} della torretta stessa, ``la locazione generale
dell'actor''. Per il Target, il docente comincia con Get Player Character e Get Actor
Location, come per il mulino, ma uno studente fa notare che si può usare
\textbf{PlayerTarget}: ``bravo, mi piace che l'hai notato, effettivamente sarebbe stato
il secondo metodo''. Da PlayerTarget si tira fuori Get Actor Location, ``una sua
particolarità'', e la si collega al Target.

\paragraph{Lo yaw e RInterp To.} Il Pivot deve ruotare attorno alla \textbf{Z}:
``guardate un po' dove si modifica: sulla Z''. Si splitta l'uscita di Find Look at
Rotation e si fa ``quello che abbiamo fatto prima'', un \textbf{RInterp To}:
\begin{itemize}
  \item \textbf{Current}: la rotazione attuale del Pivot, con \textbf{Get Relative
        Rotation} sul componente Pivot;
  \item \textbf{Target}: splittato, con la sola Z collegata, ``perché non ci interessa
        la X e la Y'';
  \item \textbf{Delta Time}: 5;
  \item \textbf{Interp Speed}: \textbf{0,04}, ``perché, essendo che sta sui
        \emph{frame}, è molto veloce. Se voi mettete uno, questa \emph{snappa}
        diretto''.
\end{itemize}
Infine \textbf{Set Relative Rotation} sul Pivot (``prendete Pivot, Set Relative
Rotation''), con i pin splittati ``per sicurezza'', collegato al ramo Is Valid; e un
Print String per il debug (figura~\ref{ue4:fig:turretbp}).

\begin{figure}[htbp]
\centering
\begin{tikzpicture}[font=\footnotesize,
  ue4node/.style={draw, thick, rounded corners=4pt, align=center, fill=gray!8,
                  minimum height=0.85cm, text width=2.3cm},
  ue4event/.style={ue4node, fill=red!12},
  ue4exec/.style={-{Stealth[length=2.5mm]}, very thick, black!80},
  ue4data/.style={-{Stealth[length=2mm]}, thick, blue!60!black},
  ue4flt/.style={-{Stealth[length=2mm]}, thick, green!50!black},
  ue4dlg/.style={-{Stealth[length=2mm]}, thick, red!70!black, dashed},
  ue4comp/.style={draw, thick, rounded corners=3pt, minimum width=2.3cm,
                  minimum height=0.55cm, align=center, font=\scriptsize}]
  % --- riga 1: overlap
  \node[ue4event] (ov) at (0,0)
    {\textbf{On Component Begin Overlap}\\(Sphere)};
  \node[ue4node] (iv1) at (2.95,0) {\textbf{Is Valid}\\Other Actor};
  \node[ue4node] (br) at (5.9,0) {\textbf{Branch}\\Actor Has Tag ``Player''};
  \node[ue4node] (set) at (8.85,0) {\textbf{Set}\\\textbf{PlayerTarget}};
  \node[ue4node, text width=2.8cm] (call) at (12.05,0) {\textbf{TrackTargetTimer}\\(chiamata)};
  \draw[ue4exec] (ov) -- (iv1);
  \draw[ue4exec] (iv1) -- (br);
  \draw[ue4exec] (br) -- node[above, font=\scriptsize]{True} (set);
  \draw[ue4exec] (set) -- (call);
  \draw[ue4data] (ov.south) -- (0,-1.15) -| node[pos=0.2, below,
    font=\scriptsize]{Other Actor} (set.south);
  % --- riga 2: timer
  \node[ue4event, text width=2.8cm] (cev) at (0,-2.9)
    {\textbf{Custom Event}\\\textbf{TrackTargetTimer}};
  \node[ue4node, text width=3.0cm] (timer) at (3.65,-2.9)
    {\textbf{Set Timer by Event}\\Time: Get World Delta Seconds\\Looping \checkmark\\
    Max Once Per Frame \checkmark};
  \node[ue4event] (mv) at (7.1,-2.9) {\textbf{Custom Event}\\\textbf{MoveTurret}};
  \draw[ue4exec] (cev) -- (timer);
  \draw[ue4dlg] (timer.east) -- node[above, font=\scriptsize]{Event} (mv.west);
  % --- riga 3: move turret
  \node[ue4node] (iv2) at (0,-5.7) {\textbf{Is Valid}\\PlayerTarget};
  \node[ue4node] (look) at (3.1,-5.7) {\textbf{Find Look at Rotation}\\
    Start: self\\Target: PlayerTarget};
  \node[ue4node] (interp) at (6.3,-5.7) {\textbf{RInterp To}\\
    Current: Pivot\\Target: solo Z\\Speed 0,04--0,05};
  \node[ue4node] (srr) at (9.5,-5.7) {\textbf{Set Relative Rotation}\\Target: Pivot};
  \draw[ue4exec] (mv.south) -- (7.1,-4.35) -| (iv2.north);
  \draw[ue4exec] (iv2.south) -- (0,-7.1) -| node[pos=0.2, below,
    font=\scriptsize]{Is Valid} (srr.south);
  \draw[ue4flt] (look) -- node[above, font=\scriptsize]{Z} (interp);
  \draw[ue4flt] (interp) -- (srr);
  % --- gerarchia componenti
  \node[font=\scriptsize] at (11.6,-1.95) {i componenti};
  \node[ue4comp, fill=orange!25] (cb) at (11.4,-2.5) {Base (radice)};
  \node[ue4comp, fill=gray!20] (cp) at (11.9,-3.2) {Pivot};
  \node[ue4comp, fill=orange!12] (cc) at (12.4,-3.9) {Cannon};
  \node[ue4comp, fill=red!8] (cs) at (11.9,-4.6) {Sphere (r. 800)};
  \draw[thick, gray!60!black] ([xshift=-0.85cm]cb.south) |- (cp.west);
  \draw[thick, gray!60!black] ([xshift=-0.85cm]cp.south) |- (cc.west);
  \draw[thick, gray!60!black] ([xshift=-0.85cm]cb.south) |- (cs.west);
\end{tikzpicture}
\caption{La torretta della lezione. In alto: quando qualcosa entra nella sfera, se
esiste ed è etichettato \emph{Player}, lo si salva in PlayerTarget e si chiama
TrackTargetTimer. In mezzo: TrackTargetTimer fa partire un timer che, a ogni
\emph{frame}, chiama MoveTurret (il pin rosso tratteggiato è il delegato Event). In
basso: MoveTurret, se c'è ancora un bersaglio, calcola la rotazione verso di lui, ne
tiene lo \emph{yaw}, la interpola e la applica al Pivot (la Current dell'interpolazione
viene da Get Relative Rotation del Pivot; l'eventuale offset è discusso nel testo). A
destra la gerarchia dei componenti: il Cannon segue il Pivot, la sfera non ruota. Schema
ricostruito.}
\label{ue4:fig:turretbp}
\end{figure}

\paragraph{L'offset, e il pivot reimportato.} In Play la torretta ruota, ma ``ha
bisogno di un offset, come abbiamo visto con il nostro mulino''. Il docente prova con
$-90$, poi con $+90$: ``eccolo qua. Non è esattamente precisissimo''. Uno studente
osserva: ``noi questa cosa non la dobbiamo fare, perché l'altra volta abbiamo fatto un
\textbf{reimport} modificando i valori''. Ed è la chiave: senza alcun offset, la
torretta del docente ``guarda direttamente in faccia'', mentre quelle degli studenti no,
perché con Davide avevano reimportato il \textbf{Pivot} e il \textbf{Cannon} con una
\textbf{rotazione di 113} gradi sulla Z (la base no). Il docente ricostruisce i numeri
con la classe (si parla anche di $-67$ e di $-271$, ``proprio perfetti''), reimporta
prima il pivot, poi il cannone, prova con il solo pivot, e alla fine: ``ecco a voi.
Adesso la nostra torretta ci segue perfettamente''. Per gli studenti la regola pratica,
ripetuta più volte nel finale, è: se la torretta ha ``un \emph{offsettino}'',
\textbf{reimportare il Pivot mettendo $-113$} sulla rotazione Z\footnote{La trascrizione
di questo passaggio è tra le più confuse della lezione: i valori citati sono 113, $-113$,
$-67$, 118, $-270$ e $-271$, e non è chiaro in che ordine siano stati provati. Il senso
è comunque chiaro: la rotazione di 113 gradi aggiunta in import nella terza lezione
sposta la direzione ``in avanti'' della \emph{mesh}, e Find Look at Rotation, che
presuppone che l'oggetto guardi lungo la sua X, ne risulta sfasato. Le due correzioni
possibili sono equivalenti: annullare la rotazione in import, o sommare un offset fisso
allo \emph{yaw} nel grafo, come per il mulino.}.

\paragraph{Approfondimento: Find Look at Rotation e lo spazio relativo.} Il grafo della
lezione mescola due spazi: Find Look at Rotation lavora nello spazio del
\textbf{mondo} (le posizioni di torretta e player sono posizioni nel mondo), mentre
Set Relative Rotation imposta la rotazione del Pivot \textbf{relativa} alla base. Le
due cose coincidono finché la torretta nel livello non è ruotata: se la si piazza
girata di 30 gradi, il Pivot punterà 30 gradi a lato del player. Le soluzioni robuste
sono due: usare \textbf{Set World Rotation} sul Pivot (la rotazione calcolata è già
quella del mondo), oppure sottrarre allo \emph{yaw} calcolato quello dell'actor (Get
Actor Rotation). Anche la Current dell'interpolazione andrebbe presa nello stesso
spazio del Target: Get World Rotation con Set World Rotation, Get Relative Rotation con
Set Relative Rotation. Lo Start, infine, è la posizione dell'actor, cioè della base; per
una torretta alta conviene usare la posizione del Pivot o del cannone (\emph{Get World
Location} sul componente), specie quando si vorrà calcolare anche il \emph{pitch} per
alzare la canna.

% ---------------------------------------------------------------------
\section{Il tuning e i problemi degli studenti}
\label{ue4:sec:problemi}

\paragraph{La velocità di inseguimento.} ``Giratele intorno.'' Per aumentare o diminuire
la velocità con cui la torretta ``ci mette a cercare il target'' si cambia l'\textbf{Interp
Speed}. A 0,07 è ``molto più reattiva''; a 2 ``sarà proprio quasi uno \emph{snap}''
(``pam, \emph{snappa} proprio''). Il consiglio: ``mettetela tipo 0,05''. E, volendo, la
si può trasformare in variabile (Promote to Variable) ``e la potete settare
esternamente'', nei dettagli dell'istanza in scena, come le variabili \emph{Instance
Editable} della seconda lezione.

\paragraph{``A me snappa lo stesso.''} Uno studente ha Delta Time 5 e Interp Speed 0,05,
e la torretta scatta comunque. Si prova 0,02 (``vedo leggerissima differenza''), poi
Delta Time 7: ``stampa sempre la stessa velocità''. In condivisione il docente controlla
Get World Delta Seconds, Looping e Max Once Per Frame: tutto giusto. ``Prova ad andare di
\emph{tweak}, di \emph{tuning}: vedi un po' tu i valori.'' Abbassando ancora l'Interp
Speed, verso lo zero, la torretta rallenta: ``ora sì, già un po' più così''. (Nel
frattempo lo studente che lavora deve andare: ``ciao, ciao''.)

\paragraph{Approfondimento: perché ``snappa''.} La formula di RInterp To vista per il
mulino spiega tutto. Con Delta Time fisso a 5 e Interp Speed 0,04, la frazione di
strada percorsa a ogni chiamata è $5 \times 0{,}04 = 0{,}2$: il 20\% di quel che manca.
Dopo dieci chiamate manca $0{,}8^{10} \approx 11\%$ dell'angolo iniziale; a 100
\emph{frame} al secondo, sono un decimo di secondo. Con 0,05 la frazione è 0,25; con 7
e 0,02 è 0,14; con Interp Speed 2 il prodotto è 10, cioè più di 1, e la torretta arriva
in un solo \emph{frame}: lo \emph{snap}. In tutti questi casi, inoltre, la velocità
dipende dai \emph{frame}: con il doppio dei \emph{frame} la torretta è due volte più
rapida. L'uso previsto del nodo è il contrario di quello della lezione: \textbf{Delta
Time} collegato a \textbf{Get World Delta Seconds} (circa 0,01 a 100 \emph{fps}) e
\textbf{Interp Speed} nell'ordine delle unità (da 2 a 5 per un inseguimento morbido),
così che il prodotto resti piccolo e la velocità sia la stessa su ogni macchina. È lo
stesso principio dei \emph{delta seconds} della seconda lezione, applicato a
un'interpolazione. E un'avvertenza: un Interp Speed esattamente 0 non ferma la torretta,
ma la fa scattare subito sul bersaglio, perché Unreal lo interpreta come ``nessuna
interpolazione''.

\paragraph{La lista degli errori.} Il resto della lezione è una carrellata di
condivisioni dello schermo, molto istruttiva perché mostra gli errori più frequenti. La
tabella~\ref{ue4:tab:errori} li raccoglie. Il caso più ricco è quello di uno studente
a cui ``non va perché probabilmente ho mancato delle parti'':
\begin{itemize}
  \item per la Current di RInterp To aveva usato un \emph{Set} invece di \textbf{Get
        Relative Rotation} (``no, no, no, non \emph{set}: \emph{get}'');
  \item l'Is Valid di MoveTurret non era collegato alla variabile \textbf{PlayerTarget}
        (``a sinistra, proprio sulle variabili, non sul grafo'');
  \item MoveTurret non era un evento separato. Per riorganizzare il grafo il docente
        gli fa selezionare tutto tranne TrackTargetTimer, Set Timer by Event e Get
        World Delta Seconds (``tieni premuto Ctrl e clicca su quelli che vuoi
        escludere''), spostare in basso la selezione, tirare fuori dal pin Event del
        timer un nuovo Custom Event chiamato \textbf{MoveTurret} e collegarlo all'Is
        Valid;
  \item il timer non era in \textbf{Looping};
  \item il \textbf{Target} di Find Look at Rotation era scollegato: da PlayerTarget,
        Get Actor Location, al Target;
  \item il \emph{character} non aveva il \textbf{tag}: lo si aggiunge nel Third Person
        Character, ``\emph{Player} esattamente come l'hai scritto, quindi con la P
        maiuscola'';
  \item la sfera aveva ancora Hidden in Game attivo e raggio \textbf{32}: ``devi
        metterla a 800'';
  \item in scena c'era solo la \textbf{mesh} della torretta, non il
        \textbf{Blueprint}: ``questa è solo la \emph{mesh}, è solo il modello 3D che hai
        messo in scena''; si cancella e si trascina nel livello BP\_Turret;
  \item il player partiva già dentro il range, e quindi il Begin Overlap non scattava:
        ``il player parte da fuori, mettilo fuori'';
  \item in Play \textbf{girava la base}: nella gerarchia, la base era figlia del Pivot.
        ``Devi staccare la base'': trascinando il Pivot dentro la base, ``perfetto, ce
        l'hai perfettamente'';
  \item infine, l'\emph{offset} del pivot reimportato: ``devi reimportare il pivot a
        $-113$''.
\end{itemize}
Un altro studente ha la torretta ferma: ``Other Actor devi attaccarlo a PlayerTarget,
perché così non gli stai impostando il \emph{target} che è overlappato: non avrà un
\emph{target} da seguire''. Un terzo non vede la sfera in gioco, solo in modalità
\emph{spectator}: aveva saltato il passaggio di Hidden in Game.

\begin{table}[htbp]
\centering
\small
\begin{tabular}{@{}p{5.4cm}p{8.6cm}@{}}
\toprule
\textbf{Sintomo} & \textbf{Causa e rimedio} \\
\midrule
La pala gira attorno a un punto sbagliato & la \emph{mesh} ha il pivot fuori
centro: scegliere un cubo con il pivot al centro (o correggere il pivot) \\
La pala non gira & il risultato di Combine Rotators non è collegato a Set
Relative Rotation; oppure il componente non è Movable \\
Il mulino o la torretta restano fermi & mancano Delta Time e Interp Speed in
RInterp To; oppure non si è compilato \\
Non si riesce a collegare un valore a una sola componente & fare Split Struct Pin
(tasto destro sul pin) \\
Gira tutto l'actor, non solo una parte & il Rotating Movement fa girare la radice:
programmare la rotazione del componente \\
La torretta non si accorge del player & manca il tag \emph{Player} (o è scritto
diversamente); Other Actor non collegato al Set di PlayerTarget; il player parte già
dentro la sfera \\
La sfera non si vede in Play & Hidden in Game ancora attivo; raggio di default 32 \\
Non succede niente in Play & in scena c'è la \emph{mesh}, non il Blueprint \\
Gira la base & gerarchia sbagliata: la base deve essere la radice, il Pivot suo
figlio, il Cannon figlio del Pivot \\
La torretta guarda di lato & la \emph{mesh} non guarda lungo X (Pivot reimportato
a 113): reimportare a $-113$ o sommare un offset allo \emph{yaw} \\
La torretta scatta (\emph{snap}) & prodotto Delta Time $\times$ Interp Speed troppo
grande: usare Get World Delta Seconds e un Interp Speed di qualche unità \\
La torretta continua a seguire dopo l'uscita & il timer non viene mai fermato e
PlayerTarget resta pieno: End Overlap, Clear and Invalidate Timer, svuotare la
variabile \\
\bottomrule
\end{tabular}
\caption{I problemi incontrati durante la lezione, con la causa e il rimedio. Le ultime
due righe riguardano punti che il docente ha rimandato alla lezione successiva.}
\label{ue4:tab:errori}
\end{table}

% ---------------------------------------------------------------------
\section{La chiusura e la prossima lezione}
\label{ue4:sec:chiusura}

``Ok, perfetto, ragazzi.'' Per questa volta, ``nessun particolare compito da casa o
dimostrazione''. Chi vuole può fare una \textbf{torretta di colore diverso}, con un
materiale suo, verde o blu, ripassando quello che si è fatto con Davide.

\paragraph{La prossima volta.} ``La prossima volta faremo la nostra bellissima
\textbf{torretta che spara}'': \textbf{spawnerà} dei \textbf{proiettili}. Si vedrà
``come funziona il sistema di \emph{spawning} di Unreal'': come si basa una cosa
\emph{spawnata} ``con le sue regole'', come si fa apparire un oggetto in generale, e
come si fa il \textbf{line trace}. ``Altre domande? Se no, andiamo a cena tutti,
ragazzi.'' ``Vi auguro buona serata. E a martedì.''

\begin{table}[htbp]
\centering
\small
\begin{tabular}{@{}p{4.0cm}p{10.0cm}@{}}
\toprule
\textbf{Nodo o componente} & \textbf{Che cosa fa} \\
\midrule
Rotating Movement Component & fa ruotare la radice dell'actor a velocità costante
(Rotation Rate, default Z 180) \\
Event Tick, Delta Seconds & esegue a ogni \emph{frame}; i \emph{delta seconds} sono il
tempo tra due \emph{frame} \\
Get World Delta Seconds & gli stessi \emph{delta seconds}, ovunque nel grafo \\
Make Rotator & costruisce un \emph{rotator} da tre \emph{float} (Roll X, Pitch Y,
Yaw Z) \\
Combine Rotators & combina due rotazioni in una \\
Get / Set Relative Rotation & legge e scrive la rotazione di un componente rispetto al
genitore \\
Get Actor Location & la posizione dell'actor nel mondo \\
Get Player Character & una referenza al personaggio del giocatore \\
Find Look at Rotation & la rotazione per guardare da Start verso Target \\
RInterp To & avvicina gradualmente un \emph{rotator} a un altro \\
Clamp (float) & limita un valore tra un minimo e un massimo \\
Split Struct Pin & divide un pin di struttura (vettore, \emph{rotator}) nelle sue
componenti \\
Sphere Collision & un volume sferico per rilevare \emph{overlap} (Sphere Radius,
Hidden in Game) \\
On Component Begin Overlap & scatta quando qualcosa entra nel volume del componente;
Other Actor è chi è entrato \\
Is Valid (macro con il ``?'') & controlla che una referenza punti a qualcosa che
esiste \\
Actor Has Tag & controlla se un actor porta un'etichetta (Actor $\rightarrow$ Tags) \\
Promote to Variable & crea una variabile del tipo giusto da un pin \\
Custom Event & un evento creato da noi, che parte quando lo si chiama \\
Set Timer by Event & chiama un evento dopo un tempo, una volta o in \emph{loop} (Max
Once Per Frame) \\
Print String & stampa a schermo: lo strumento base del debug \\
\bottomrule
\end{tabular}
\caption{I nodi e i componenti della quarta lezione.}
\label{ue4:tab:nodi}
\end{table}

% ---------------------------------------------------------------------
\section*{Punti chiave del capitolo}
\addcontentsline{toc}{section}{Punti chiave del capitolo}

\begin{itemize}
  \item \textbf{Quixel Megascans} è un ``calderone'' gratuito di materiali, texture e
        \emph{mesh} scansionati, da usare con l'account Epic. Per la pallina che
        attraversa il piano: il \textbf{CCD} (\emph{Continuous Collision Detection}).
  \item Un \textbf{vettore} ha direzione, magnitudine e verso; in Unreal la struttura
        Vector (X, Y, Z) rappresenta sia punti sia spostamenti. Somma: si spostano i
        punti; \textbf{sottrazione}: B $-$ A va da A a B; \textbf{moltiplicazione}:
        allunga; \textbf{prodotto scalare}: un \emph{float}, 1 se concordi, 0 se
        perpendicolari, $-1$ se opposti (per vettori normalizzati).
  \item Un vettore \textbf{normalizzato} ha lunghezza 1: tiene solo la direzione.
  \item Un \textbf{rotator} ha tre angoli: \textbf{Roll} (X), \textbf{Pitch} (Y),
        \textbf{Yaw} (Z). Gli angoli di Eulero soffrono di \textbf{gimbal lock}; i
        \textbf{quaternioni}, usati internamente da Unreal, no.
  \item Il \textbf{Rotating Movement Component} fa girare un actor senza programmare,
        ma solo per intero. Ogni oggetto che si muove deve essere \textbf{Movable}.
  \item Per far girare un componente: \textbf{Event Tick} $\rightarrow$ velocità
        $\times$ \emph{delta seconds} $\rightarrow$ \textbf{Make Rotator}
        $\rightarrow$ \textbf{Combine Rotators} con \textbf{Get Relative Rotation}
        $\rightarrow$ \textbf{Set Relative Rotation}. Le rotazioni dei componenti si
        lavorano in spazio \textbf{relativo}.
  \item La \textbf{gerarchia} fa il resto: un componente figlio segue il genitore,
        mantenendo la propria posizione relativa (la seconda pala, il cubetto).
  \item \textbf{Find Look at Rotation} calcola la rotazione per guardare un punto;
        dello \emph{yaw} si tiene la Z (\textbf{Split Struct Pin}), con un
        \textbf{offset} se la \emph{mesh} non guarda lungo X. \textbf{RInterp To}
        rende il movimento graduale: Delta Time dai \emph{delta seconds}, Interp Speed
        di qualche unità, altrimenti ``snappa''.
  \item Il \textbf{Clamp} limita un valore; prima di sceglierne i limiti, si
        \textbf{stampano} i valori con Print String.
  \item La torretta è un Blueprint in \textbf{tre pezzi}: Base (radice), Pivot
        (figlio, ruota in \emph{yaw}), Cannon (figlio del Pivot). Browse to Asset aiuta
        a trovare le \emph{mesh}.
  \item Il range è una \textbf{Sphere Collision} (raggio 800, Hidden in Game tolto per
        vederla). \textbf{On Component Begin Overlap} scatta quando qualcosa ci entra;
        \textbf{Other Actor} dice chi.
  \item \textbf{Is Valid} prima di usare una referenza; \textbf{Actor Has Tag} con il tag
        \emph{Player} (scritto identico) per riconoscere il giocatore;
        \textbf{Promote to Variable} per ricordarlo in \emph{PlayerTarget}.
  \item I \textbf{Custom Event} sono eventi nostri, chiamati quando vogliamo;
        \textbf{Set Timer by Event} (Looping, Max Once Per Frame) chiama a ripetizione
        l'evento \textbf{MoveTurret}, che ruota il Pivot verso il bersaglio. Il timer
        va fermato (End Overlap, Clear and Invalidate Timer): sarà il prossimo passo.
  \item Prossima lezione: la torretta che \textbf{spara}, lo \textbf{spawn} dei
        proiettili e il \textbf{line trace}.
\end{itemize}

%
%  Tutte le etichette di questo capitolo hanno il prefisso "ue4:"
%  (ue4:cap:..., ue4:sec:..., ue4:fig:..., ue4:tab:...), e tutti gli
%  stili TikZ hanno il prefisso "ue4" (ue4node, ue4exec, ue4data, ...)
%  per non collidere con chiavi native di TikZ (come "step") né con
%  quelli delle lezioni precedenti. Il capitolo è autonomo: non contiene
%  \ref verso altri capitoli (i rimandi alle lezioni precedenti sono a
%  parole).
%  Le figure sono ricostruzioni schematiche di quanto mostrato e detto a
%  lezione, non copie delle slide o delle schermate del docente. Dove la
%  trascrizione è troppo rovinata per ricostruire un dettaglio, il
%  capitolo lo dice invece di inventarlo.
% =====================================================================

\chapter[Vettori, rotazioni e la torretta]{Vettori, rotazioni e la torretta: il
mulino che gira e la torretta che ti segue}
\label{ue4:cap:torretta}

Questa è la \textbf{quarta lezione} del corso breve di \textbf{introduzione a Unreal
Engine}, e segna il ritorno del docente titolare dopo la lezione di Davide sulla
grafica. Se la terza lezione guardava a ciò che ``fa il grafico all'interno di Unreal''
(texture, materiali, l'import della torretta), questa torna alla \textbf{logica}, cioè
ai \textbf{Blueprint}, e lo fa con un obiettivo preciso, mostrato in apertura: un
\textbf{mulino} le cui pale girano, e la \textbf{torretta} importata con Davide che
``adesso ci segue con lo sguardo''. Per arrivarci serve un po' di matematica dello
spazio 3D: i \textbf{vettori}, le loro operazioni, i \textbf{rotator}.

La lezione ha due metà separate da una pausa di venti minuti, presa alle 17:28. La
prima metà parte dalla \textbf{teoria}: che cos'è un \textbf{vettore} (``una freccia
con una direzione, una magnitudine e un verso''), le quattro operazioni
(\textbf{addizione}, \textbf{sottrazione}, \textbf{moltiplicazione} e \textbf{prodotto
scalare}, il \emph{dot product}), i \textbf{vettori normalizzati}, i
\textbf{rotator} e il problema del \textbf{gimbal lock}. Poi si passa all'engine con un
mulino: prima il modo più semplice di far ruotare un oggetto, il \textbf{Rotating
Movement Component}, poi un Blueprint vero, \textbf{BP\_Windmill}, in cui una pala
gira grazie all'\textbf{Event Tick}, ai \emph{delta seconds} e ai nodi \textbf{Make
Rotator}, \textbf{Combine Rotators} e \textbf{Set Relative Rotation}; la seconda pala
gira ``gratis'' grazie alla \textbf{gerarchia} dei componenti. Infine il mulino
diventa ``magico'' e si gira verso il giocatore con \textbf{Find Look at Rotation} e
\textbf{RInterp To}, con un \textbf{Clamp} per limitarne la rotazione. La seconda metà
costruisce la \textbf{torretta}: un Blueprint in tre pezzi (\textbf{Base},
\textbf{Pivot}, \textbf{Cannon}), una \textbf{Sphere Collision} come raggio d'azione,
l'evento \textbf{On Component Begin Overlap}, i controlli \textbf{Is Valid} e
\textbf{Actor Has Tag}, una variabile per ricordare il bersaglio, un \textbf{Custom
Event} e un \textbf{timer} che, a ogni \emph{frame}, fa ruotare il pivot verso il
player.

La lezione è \textbf{online}, su Zoom, come le precedenti. Il docente alterna
spiegazione e piccole esercitazioni ``botta e risposta'', da cinque minuti l'una, e
passa gran parte della seconda metà a risolvere i problemi degli studenti in
condivisione dello schermo. Finisce verso le 19, ``andiamo a cena tutti'', con
l'appuntamento a \textbf{martedì}.

Tre avvertenze di lettura, le stesse delle lezioni precedenti. La prima: la
trascrizione è automatica e storpia moltissimi nomi. Unreal è ``a Real'', ``Arreal'',
``Arry''; l'\textbf{Event Tick} è ``Leventic'', ``l'eventic'', ``eventHIC'' e
``l'event kick''; i \textbf{vettori} sono una volta ``vittori''; i \emph{rotator} sono
``rotatempi'' e ``rotatori''; \textbf{Find Look at Rotation} è ``findLookAtRotation'',
``find all the rotation'' e ``Find, look, rotation''; \textbf{RInterp To} è ``Er inter
tu'', ``R Inter2'' e ``r interp tu''; \textbf{Get Player Character} è
``getLayerCharacter''; \textbf{Get Actor Location} è ``Get. Hector. Location'' e ``get
active location''; il \textbf{Rotating Movement} è a un certo punto ``vetro in
movement''; \textbf{Is Valid} è ``Svalid'' e ``isvalid''; \textbf{Other Actor} è
``OderActor'' e ``Oder Actor''; la \textbf{gun} della torretta è ``Ganna'' e ``GAN'';
il \textbf{Pivot} è spesso ``pivo''; il \emph{Third Person Character} è ``ter person
character'' e ``Peter per personacarte''; \textbf{Quixel Megascans} è ``QuickSell'' e
``Quixel megaschanico''. Dove la ricostruzione è sicura il capitolo usa la forma
corretta e, la prima volta, lo segnala in nota; dove non lo è riporta la forma della
trascrizione e lo dice. La seconda: all'inizio della registrazione, in un punto della
prima esercitazione e durante tutta la pausa, il riconoscimento vocale produce lunghi
blocchi di testo senza senso e ripetuto (``qual è la risposta che ci può essere?'',
``grazie a tutti per l'attenzione'', discorsi sul ``mercato del lavoro'', sul ``governo
di Dalai Lama'' e su ``la migliore dal 2005''): sono artefatti, non parole di nessuno,
e il capitolo li ignora. La terza: gran parte della lezione è una dimostrazione dal
vivo nell'editor dei Blueprint, spesso costruita a voce (``tirate fuori'',
``collegate'', ``splittate''). Il capitolo ricostruisce i grafi da ciò che il docente
dice mentre li costruisce, e li completa, dove serve, con i nomi esatti dei nodi e dei
pin di Unreal Engine 5.

\paragraph{Come leggere questo capitolo.} Il capitolo segue l'ordine della lezione.
Le figure ricostruiscono gli schemi spiegati a voce o alla lavagna (``perdonate le mie
doti molto brutte di disegno''): le operazioni tra vettori, i casi del prodotto
scalare, gli assi di rotazione, la gerarchia del mulino, i grafi della pala, del
mulino che guarda il player e della torretta, i casi dell'\emph{overlap}. Le tabelle
riassumono le operazioni sui vettori, i nodi della lezione e gli errori più comuni
incontrati dagli studenti. Qua e là, dei paragrafi di \textbf{approfondimento}
collegano la lezione ai libri di Marcos Romero e Brenden Sewell (\emph{Blueprints
Visual Scripting for Unreal Engine 5} e la seconda edizione dello stesso testo, scritta
per Unreal Engine 4), i testi di riferimento per questa parte del libro, e a qualche
nozione tecnica generale su Unreal Engine 5; sono segnalati come tali e non fanno parte
della lezione.

Il percorso del capitolo:
\begin{enumerate}
  \item l'attesa: \textbf{Quixel Megascans} e il materiale della lava; il trucco del
        \textbf{CCD}; il programma della giornata (sezioni da~\ref{ue4:sec:attesa}
        a~\ref{ue4:sec:programma});
  \item la teoria: i \textbf{vettori} e le loro operazioni, il \textbf{prodotto
        scalare}, la \textbf{normalizzazione}, i \textbf{rotator} e il \textbf{gimbal
        lock} (sezioni da~\ref{ue4:sec:vettori} a~\ref{ue4:sec:rotator});
  \item il \textbf{Rotating Movement Component} e l'esercizio dei tre cubi (sezione
        \ref{ue4:sec:rotating});
  \item il Blueprint del \textbf{mulino}: i componenti, la pala che gira con
        l'Event Tick, la \textbf{gerarchia} (sezioni da~\ref{ue4:sec:mulino}
        a~\ref{ue4:sec:gerarchia});
  \item il mulino che guarda il player: \textbf{Find Look at Rotation},
        \textbf{RInterp To}, il \textbf{Clamp} e il debug (sezioni
        \ref{ue4:sec:lookat} e~\ref{ue4:sec:clamp});
  \item dopo la pausa, la \textbf{torretta}: i tre pezzi, la \textbf{Sphere
        Collision}, l'\textbf{overlap} (sezioni da~\ref{ue4:sec:torretta}
        a~\ref{ue4:sec:overlap});
  \item \textbf{Is Valid}, i \textbf{tag}, la variabile del bersaglio e i
        \textbf{Custom Event} (sezioni da~\ref{ue4:sec:isvalid}
        a~\ref{ue4:sec:custom});
  \item il \textbf{timer} e l'evento \textbf{Move Turret}: la torretta che segue il
        player (sezioni \ref{ue4:sec:timer} e~\ref{ue4:sec:move});
  \item il \emph{tuning}, i problemi degli studenti e la chiusura (sezioni
        \ref{ue4:sec:problemi} e~\ref{ue4:sec:chiusura}).
\end{enumerate}

% ---------------------------------------------------------------------
\section{L'attesa: Quixel Megascans e la lava}
\label{ue4:sec:attesa}

La registrazione parte mentre il docente fa entrare gli studenti nella stanza di Zoom
(``ho aperto io la prima a far entrare un po' tutti'') e chiacchiera con il primo
arrivato. ``Come va con Unreal? Hai sperimentato un po', l'hai toccato un pochetto?''
Lo studente ci ha giocato con un collega, che gli ha insegnato le texture, ma non ha
trovato online ``alcuna fonte per prendere cose utili, file utili'': per l'esercizio
della lava assegnato da Davide nella lezione precedente, un materiale di lava animata,
ha usato quello che c'era nella cartella del Drive.

\paragraph{Quixel Megascans.} La risposta del docente è \textbf{Quixel
Megascans}\footnote{La trascrizione dice ``QuickSell'', ``Quixel megaschanico'' e
``Materials pixel mega scan''. Il nome corretto è Quixel Megascans: la libreria di
asset scansionati di Quixel, società acquisita da Epic Games, gratuita per l'uso con
Unreal Engine.}, di cui passa il link in chat. Quixel ``ti permette di avere proprio
dei materiali'': ``una marea di materiali di qualsiasi tipo'', ma anche rocce,
\emph{mesh} e texture di alta qualità. Il docente fatica un po' a ritrovare la pagina
giusta, ``perché hanno cambiato il sito'', e ricorda che serve accedere con l'account
di Unreal (cioè di Epic Games). ``Ormai fanno la cooperazione tra di loro'': c'è un
collegamento diretto tra Quixel e Unreal, ``proprio tipo \emph{face to face}''. Cercando
\emph{lava} escono diverse risorse; una, \emph{lava story} o qualcosa del genere,
``sembra carina''.

\paragraph{Come sono fatte.} Le texture di Megascans ``sono proprio \emph{scan}'': c'è
``gente che si mette là con queste macchine che costano una vita umana'' a
scansionare le rocce. ``Quixel è un grande, grande calderone di risorse.'' Da lì si
può partire anche per la lava senza trovare la texture perfetta: si prende una
\textbf{roccia} e, ``come avete visto con Davide, con la matematica'', cioè con un
\textbf{Multiply} e con il \textbf{Lerp} pilotato da un'\textbf{alfa}, si inserisce
``il rosso della lava all'interno di queste rocce''. Un'altra strada sono i video su
YouTube: alcuni tutorial lasciano in descrizione le texture usate (altri le fanno
pagare, ``il maledetto''), e ``ti interessano le texture, il materiale poi lo fai tu''.

\paragraph{La lava degli studenti.} Arrivano gli altri (``buongiorno, Filippo, Marco,
Patrick''), e il docente chiede chi ha provato il materiale della lava ``che vi è stato
dato da Davide''\footnote{La trascrizione dice ``da Dario'', ma subito dopo e più
avanti il nome è Davide, il docente della terza lezione.}. Uno studente l'ha mandato sul
Discord, e il docente lo commenta: ``bello il movimento che hai fatto''. L'hai messo
su un piano, in engine, per vedere come reagiva? No. ``Quello vedilo un po' tu'',
perché a volte il risultato in \emph{viewport} cambia a seconda di come è costruito il
materiale (``che so, hai messo il Multiply sbagliato''): il consiglio è di provarlo
sempre in \emph{viewport}, su un oggetto vero. E ``sicuramente una texture differente
potrebbe essere più carina'': torna utile Megascans.

\paragraph{L'appello.} Il docente aspetta ``2--3 minuti'' che entrino gli altri. Uno
studente avvisa che anche oggi dovrà staccarsi venti minuti prima per andare a
lavorare: ``assolutamente nessun problema''.

% ---------------------------------------------------------------------
\section{Le domande e il trucco del CCD}
\label{ue4:sec:ccd}

Prima di cominciare, il docente apre lo spazio delle domande: ``se qualcuno di voi ha
avuto qualche particolare cosa di cui vuole parlare, domandare già prima che magari
iniziamo, così facciamo il Q\&A già da subito''. Nessuno: ``mi sembrate tutti molto
tranquilli''.

\paragraph{La pallina che entra nel piano.} Allora è il docente a recuperare ``un
piccolo trucchetto che mi sono dimenticato di dirvi'' per il \textbf{gioco della
pallina} della seconda lezione, quello della piattaforma inclinabile: a volte la
pallina, con la fisica attiva, \textbf{penetra} nel piano, ``come abbiamo visto''. Nei
dettagli del componente, nella sezione della fisica (vicino a \textbf{Simulate
Physics}), c'è un'opzione da spuntare: il \textbf{CCD}, \emph{Continuous Collision
Detection}\footnote{La trascrizione dice ``Trucco Collision Collision Detection''.}.
``Fondamentalmente, è un \emph{tick} che, se lo mettiamo su questo particolare
\emph{asset}, lui in modo continuo cerca di vedere le collisioni'': le rende ``più
precise'', ed evita che la pallina attraversi il terreno. ``Se fate questo, vi risolve
tranquillamente il problema.'' Chi vuole, può sistemare il gioco per la prossima volta.

\paragraph{Approfondimento: perché la pallina attraversa il piano.} La fisica di un
gioco non è continua: il motore fisico controlla le collisioni a passi discreti, una
volta per \emph{frame} (o per \emph{substep}). Se in un passo la pallina si sposta più
dello spessore del piano, può trovarsi ``prima'' sopra e ``dopo'' già dentro o sotto,
senza che nessun controllo l'abbia mai vista toccare la superficie: è il fenomeno del
\emph{tunneling}, tanto più probabile quanto più l'oggetto è piccolo e veloce e il piano
sottile. Il CCD (in Unreal la casella \emph{Use CCD} nella sezione \emph{Collision} del
componente, nelle impostazioni avanzate) fa sì che il motore fisico consideri anche il
percorso tra le due posizioni, la ``scia'' dell'oggetto, e non solo le posizioni
finali. Costa un po' di più, ed è per questo che non è attivo di default: lo si accende
solo sugli oggetti che ne hanno bisogno, come proiettili e palline veloci. Un'alternativa
complementare è dare al piano uno spessore maggiore.

% ---------------------------------------------------------------------
\section{Il programma: il mulino e la torretta}
\label{ue4:sec:programma}

``Oggi vi dico già da subito cosa faremo.'' Gli studenti hanno scaricato il
\textbf{materiale didattico}: dentro il file estratto c'è il progetto
\emph{UE5 Intro Start}\footnote{La trascrizione dice ``Hive UE5 intro start'', la stessa
parola iniziale incerta già incontrata nella terza lezione.}, con il suo file di
progetto. Lo si apre, ``così almeno gli \emph{shader} si cominciano a buildare'', si fa
``una piccola spiegazione teorica'' e poi ``partiamo direttamente con la pratica''.

\paragraph{I due obiettivi.} Il docente sposta il \emph{Player Start} e mostra il
livello. A sinistra c'è ``questo bellissimo \textbf{mulino}'', fatto ``facendo ruotare
le pale''. Poi c'è la \textbf{torretta} che gli studenti hanno costruito la volta prima
con Davide (``avete fatto il \emph{merge}, giusto?''), che ``adesso ci segue con lo
sguardo''. O meglio, dovrebbe: nella dimostrazione non lo fa. ``Abbiamo scoperto un
piccolo bug. Poi adesso lo aggiustiamo. E giustamente i bug ogni volta appaiono
praticamente subito.'' Il programma della giornata è questo: le pale del mulino, e la
torretta che segue il giocatore.

% ---------------------------------------------------------------------
\section{I vettori}
\label{ue4:sec:vettori}

``Avete sentito mai parlare di vettori? Matematica, i grafici cartesiani?'' Sì.

\paragraph{Che cos'è un vettore.} ``In generale, un vettore può essere semplificato, in
uno spazio, come letteralmente una \textbf{freccia}'' che ha tre caratteristiche: una
\textbf{direzione}, una \textbf{magnitudine} (``la sua ampiezza'', cioè la lunghezza) e
un \textbf{verso}. In Unreal i vettori si sono già visti senza nominarli: nella
\textbf{location} e nella \textbf{rotation} di ogni oggetto, ``in particolare molto di
più sulla location''. La location è fatta di tre coordinate, \textbf{X}, \textbf{Y} e
\textbf{Z}, e ``queste particolari coordinate possono formare tra di loro, nello spazio
3D, un vettore''. Il docente parte dalla parola ``velocità di movimento'' e si corregge:
``più che una velocità, è più una direzione'', con la sua ampiezza e il suo verso.
Con i vettori si possono ``fare tante cose in Unreal'', per esempio muovere degli
oggetti; le pale del mulino girano in un verso, e anche se in quel caso si parla di
rotazione, e quindi di \emph{rotator}, ``sono sempre formate da delle coordinate''.

\paragraph{Un punto non è ancora un vettore.} Più avanti, alla lavagna, il docente
aggiunge una precisazione importante: ``ricordatevi che quello che vediamo è un
\textbf{punto} nello spazio, non è una freccia, quindi non è già un vettore''. Un
punto, per esempio $(5, 4, 1)$, è solo una posizione; il vettore è la freccia che
collega due punti, per esempio dall'origine $(0,0,0)$ a $(5,4,1)$\footnote{A lezione il
docente dice che ``un vettore è formato dalla somma di due punti''. La relazione
corretta è la \emph{differenza}: il vettore che va da A a B è B $-$ A (si veda la
sottrazione, poco più sotto). Unreal, peraltro, usa la stessa struttura
\emph{Vector} sia per i punti sia per gli spostamenti.}.

\paragraph{Le operazioni.} Con i vettori ``possiamo fare varie operazioni matematiche,
che sono praticamente quelle classiche'': \textbf{addizione}, \textbf{sottrazione},
\textbf{moltiplicazione} e \textbf{prodotto scalare} (\emph{dot product}), ``che è
un'altra cosa, che poi spiegheremo''. La slide le mostra con due vettori, A e B
(figura~\ref{ue4:fig:vettori}).
\begin{itemize}
  \item \textbf{Addizione.} ``Se andiamo a prendere due vettori e li addizioniamo,
        avremo un vettore unico.'' Il vettore A più il vettore B ``dà come risultato
        il vettore C, che è $(5,4)$'': si mette la coda di B sulla punta di A, e C va
        dall'inizio di A alla fine di B. ``Come fare l'ipotenusa, se vi ricordate i
        triangoli.''
  \item \textbf{Sottrazione.} Dà ``la \textbf{differenza} nello spazio tra due
        punti''. Il risultato dipende dall'ordine: A $-$ B è il vettore che va da B
        ad A, B $-$ A quello che va da A a B. La sua lunghezza è la distanza tra i due
        punti.
  \item \textbf{Moltiplicazione.} ``La useremo tantissimo'': si moltiplica tutto ciò
        che fa parte del vettore, X, Y e Z, e quindi il vettore ``viene allungato nella
        sua direzione''.
  \item \textbf{Prodotto scalare.} ``Statemi un attimo a sentire: è l'unica operazione
        che ci dà un valore \textbf{scalare}'', cioè ``non ci ridà un vettore indietro,
        ma un valore in \emph{float}''. (``Chi si ricorda che è il float?'' ``I numeri
        decimali.'' ``Perfetto.'') Lo si vede nella prossima sezione.
\end{itemize}

\begin{figure}[htbp]
\centering
\begin{tikzpicture}[font=\small, scale=0.8,
  ue4vec/.style={-{Stealth[length=2.5mm]}, very thick},
  ue4grid/.style={gray!25, very thin},
  ue4ax/.style={-{Stealth[length=2mm]}, gray!70}]
  % --- addizione
  \begin{scope}
    \draw[ue4grid] (0,0) grid (6,5);
    \draw[ue4ax] (0,0) -- (6.3,0) node[right, font=\scriptsize]{X};
    \draw[ue4ax] (0,0) -- (0,5.3) node[above, font=\scriptsize]{Y};
    \draw[ue4vec, blue!70!black] (0,0) -- node[left, pos=0.55]{A} (2,3);
    \draw[ue4vec, red!70!black] (2,3) -- node[above, pos=0.5]{B} (5,4);
    \draw[ue4vec, green!45!black] (0,0) -- node[below right, pos=0.6]{C $=$ A $+$ B}
      (5,4);
    \fill (5,4) circle (1.5pt) node[above right, font=\scriptsize]{$(5,4)$};
    \node[font=\bfseries] at (3,-0.9) {addizione};
  \end{scope}
  % --- sottrazione
  \begin{scope}[shift={(8.5,0)}]
    \draw[ue4grid] (0,0) grid (6,5);
    \draw[ue4ax] (0,0) -- (6.3,0) node[right, font=\scriptsize]{X};
    \draw[ue4ax] (0,0) -- (0,5.3) node[above, font=\scriptsize]{Y};
    \fill[blue!70!black] (1,1) circle (2.5pt) node[below left]{A};
    \fill[red!70!black] (5,4) circle (2.5pt) node[above right]{B};
    \draw[ue4vec, green!45!black] (1.15,1.11) -- node[above left, pos=0.5]
      {B $-$ A: da A a B} (4.85,3.89);
    \draw[ue4vec, orange!80!black, dashed] (4.9,3.6) -- node[below right, pos=0.5]
      {A $-$ B: da B ad A} (1.2,0.85);
    \node[font=\bfseries] at (3,-0.9) {sottrazione};
  \end{scope}
  % --- moltiplicazione
  \begin{scope}[shift={(0,-7.2)}]
    \draw[ue4grid] (0,0) grid (6,5);
    \draw[ue4ax] (0,0) -- (6.3,0) node[right, font=\scriptsize]{X};
    \draw[ue4ax] (0,0) -- (0,5.3) node[above, font=\scriptsize]{Y};
    \draw[ue4vec, green!45!black, line width=3pt, opacity=0.35] (0,0) -- (5,4);
    \draw[ue4vec, blue!70!black] (0,0) -- (2.5,2) node[pos=0.5, above left]{V};
    \node[green!45!black, font=\small] at (4.6,4.55) {$2 \times$ V};
    \node[font=\bfseries] at (3,-0.9) {moltiplicazione per uno scalare};
  \end{scope}
  % --- normalizzazione
  \begin{scope}[shift={(8.5,-7.2)}]
    \draw[ue4grid] (0,0) grid (6,5);
    \draw[ue4ax] (0,0) -- (6.3,0) node[right, font=\scriptsize]{X};
    \draw[ue4ax] (0,0) -- (0,5.3) node[above, font=\scriptsize]{Y};
    \draw[gray!60, dashed] (0,0) circle (1);
    \node[gray, font=\scriptsize] at (1.35,-0.3) {lunghezza 1};
    \draw[ue4vec, blue!70!black, opacity=0.45] (0,0) -- (4,3)
      node[above right, font=\scriptsize]{V, lunghezza 5};
    \draw[ue4vec, red!70!black] (0,0) -- (0.8,0.6);
    \node[red!70!black, font=\scriptsize, align=left] at (2.7,0.45)
      {Normalize(V)\\$= (0{,}8;\ 0{,}6)$};
    \node[font=\bfseries] at (3,-0.9) {normalizzazione};
  \end{scope}
\end{tikzpicture}
\caption{Le operazioni sui vettori viste a lezione, nel solo piano XY per semplicità.
In alto a sinistra la somma della slide: A seguito da B porta al punto $(5,4)$. In
alto a destra la sottrazione, il cui verso dipende dall'ordine dei fattori. In basso,
la moltiplicazione per un numero, che allunga il vettore senza cambiarne la direzione,
e la normalizzazione, che lo riporta a lunghezza 1 (sezione~\ref{ue4:sec:dot}). I
valori di A e B della slide non si leggono nella trascrizione e sono scelti per
l'esempio; il risultato $(5,4)$ è quello detto a lezione.}
\label{ue4:fig:vettori}
\end{figure}

\paragraph{Approfondimento: i vettori nel libro di Romero e Sewell.} Il
quattordicesimo capitolo di \emph{Blueprints Visual Scripting for Unreal Engine 5},
dedicato ai nodi matematici e alle \emph{trace}, presenta i vettori con un esempio
molto vicino a quello della lezione. In Unreal la struttura \textbf{Vector} contiene
tre \emph{float}, X, Y e Z, e può rappresentare sia un \textbf{punto} nello spazio
(una \emph{location}) sia uno \textbf{spostamento} o una velocità. Un personaggio sta
in $(50, 0, 20)$, un divano in $(450, 0, 20)$, e un'unità di Unreal vale un
centimetro. Per arrivare al divano il personaggio deve spostarsi di $(400, 0, 0)$: 400
centimetri lungo X. Le due regole che ne derivano sono quelle della lezione, scritte
in formula:
\[
  \text{destinazione} = \text{partenza} + \text{spostamento},
  \qquad
  \text{spostamento} = \text{destinazione} - \text{partenza}.
\]
Ogni operazione si fa componente per componente: $(3,5,0) + (5,2,0) = (8,7,0)$ e
$(6,8,0) - (1,4,0) = (5,4,0)$. Il libro aggiunge due esempi pratici: se V1 è la
posizione nel mondo di un personaggio e V2 la posizione \emph{relativa} della sua arma,
V1 $+$ V2 è la posizione nel mondo dell'arma; e la \textbf{lunghezza} della differenza
tra due posizioni (il nodo \emph{Vector Length}) è la loro distanza. Moltiplicare per un
numero negativo inverte il verso: moltiplicando per $-1$ il vettore ``in avanti'' di un
attore si ottiene quello ``all'indietro'', ed è lo stesso trucco usato nel quinto
capitolo per far cambiare direzione a un bersaglio mobile. Il nodo \textbf{Multiply},
avvertono gli autori, è un nodo \emph{wildcard}: se gli si collega un vettore diventa
un \emph{Vector $\times$ Vector}, e per moltiplicare un vettore per un \emph{float}
bisogna fare tasto destro sul secondo pin e scegliere \emph{Convert Pin\ldots} e poi
\emph{Float}.

\begin{table}[htbp]
\centering
\small
\begin{tabular}{@{}p{3.1cm}p{3.4cm}p{2.0cm}p{5.6cm}@{}}
\toprule
\textbf{Operazione} & \textbf{In formula} & \textbf{Risultato} &
\textbf{A che cosa serve in un gioco} \\
\midrule
Addizione & $(a_x+b_x,\ a_y+b_y,\ a_z+b_z)$ & vettore & spostare un punto: posizione
$+$ spostamento $=$ nuova posizione \\
Sottrazione & $(b_x-a_x,\ b_y-a_y,\ b_z-a_z)$ & vettore & la direzione e la distanza
da A a B (``dove sta il player rispetto a me?'') \\
Moltiplicazione per uno scalare & $(k\,a_x,\ k\,a_y,\ k\,a_z)$ & vettore & allungare o
accorciare; con $-1$ invertire il verso; velocità $\times$ \emph{delta seconds} \\
Lunghezza (\emph{Vector Length}) & $\sqrt{a_x^2+a_y^2+a_z^2}$ & \emph{float} & la
magnitudine; della differenza tra due punti, la distanza \\
Normalizzazione (\emph{Normalize}) & $\vec a \,/\, |\vec a|$ & vettore lungo 1 & tenere
solo la direzione, da moltiplicare poi per una velocità \\
Prodotto scalare (\emph{Dot Product}) & $a_x b_x + a_y b_y + a_z b_z$ & \emph{float} &
capire se due direzioni sono concordi, perpendicolari od opposte \\
\bottomrule
\end{tabular}
\caption{Le operazioni sui vettori della lezione, con il nome del nodo Blueprint
corrispondente dove serve e un uso tipico in un gioco.}
\label{ue4:tab:vettori}
\end{table}

% ---------------------------------------------------------------------
\section{Il prodotto scalare e i vettori normalizzati}
\label{ue4:sec:dot}

Il \textbf{prodotto scalare} è l'operazione che restituisce un \emph{float} invece di un
vettore. ``Che cosa ci dà?'' Ci dice come sono orientati due vettori l'uno rispetto
all'altro. Con i vettori A e B della slide (figura~\ref{ue4:fig:dot}):
\begin{itemize}
  \item se il prodotto scalare è \textbf{maggiore di 0}, i due vettori sono
        ``direzionati nella stessa direzione'', grosso modo concordi; se è
        \textbf{1}, ``sono proprio uguali'', cioè puntano esattamente nello stesso
        verso;
  \item se è \textbf{uguale a 0}, sono ``perfettamente \textbf{perpendicolari}'',
        formano un angolo retto;
  \item se è \textbf{minore di 0}, formano un angolo ottuso, ``al verso opposto'': per
        due vettori come quelli della slide ``magari $-0{,}4$, più o meno''; ``se
        proprio sono così'', uno contro l'altro, è $-1$.
\end{itemize}
Questi valori ``puliti'', tra $-1$ e 1, valgono per i \textbf{vettori normalizzati}, ed
è il motivo per cui la slide successiva li introduce.

\paragraph{I vettori normalizzati.} ``Qualcuno di voi sa cos'è un vettore
normalizzato?'' Uno studente: ``un vettore tra 0 e 1''. ``Esatto'': normalizzando un
vettore, ``indipendentemente dalla sua lunghezza, lo ridurremo a 1''. Un vettore ``che
è lungo 5'', normalizzato, ``lo portiamo a 1''\footnote{Più precisamente: un vettore
normalizzato ha lunghezza \emph{esattamente} 1; sono le sue componenti, prese una a
una, a stare tra $-1$ e 1. Il vettore $(4,3,0)$, lungo 5, normalizzato diventa
$(0{,}8;\ 0{,}6;\ 0)$, come nella figura~\ref{ue4:fig:vettori}.}. La direzione resta la
stessa, cambia solo la lunghezza.

\begin{figure}[htbp]
\centering
\begin{tikzpicture}[font=\small,
  ue4vec/.style={-{Stealth[length=2.5mm]}, very thick}]
  \foreach \xs/\ang/\val/\txt in {0/0/{$= 1$}/{stessa direzione},
                                  3.4/40/{$\approx 0{,}77$}/{concordi},
                                  6.8/90/{$= 0$}/{perpendicolari},
                                  10.2/140/{$\approx -0{,}77$}/{angolo ottuso},
                                  13.6/180/{$= -1$}/{opposti}} {
    \begin{scope}[shift={(\xs,0)}]
      \draw[ue4vec, blue!70!black] (0,0) -- (1.5,0) node[below, pos=0.8,
        font=\scriptsize]{A};
      \draw[ue4vec, red!70!black] (0,0) -- (\ang:1.35) node[above,
        font=\scriptsize]{B};
      \fill (0,0) circle (1.5pt);
      \node[font=\small] at (0.3,-1.0) {A $\cdot$ B \val};
      \node[font=\scriptsize, gray!50!black] at (0.3,-1.5) {\txt};
    \end{scope}
  }
\end{tikzpicture}
\caption{Il prodotto scalare tra due vettori normalizzati A e B, al variare
dell'angolo tra di loro: 1 se puntano nello stesso verso, positivo se sono concordi,
0 se sono perpendicolari, negativo se formano un angolo ottuso, $-1$ se sono opposti.
I valori intermedi sono il coseno dell'angolo ($\cos 40^\circ \approx 0{,}77$).}
\label{ue4:fig:dot}
\end{figure}

\paragraph{Approfondimento: il prodotto scalare è un coseno.} Per due vettori qualsiasi
vale
\[
  \vec a \cdot \vec b \;=\; a_x b_x + a_y b_y + a_z b_z \;=\; |\vec a|\,|\vec b|\,
  \cos\theta ,
\]
dove $\theta$ è l'angolo tra i due. Se i vettori sono normalizzati, le lunghezze valgono
1 e il prodotto scalare \emph{è} il coseno dell'angolo: per questo varia tra $-1$ e 1,
come dicono Romero e Sewell nel quattordicesimo capitolo (``il prodotto scalare di due
vettori normalizzati è uguale al coseno dell'angolo che formano''), e per questo i casi
della slide sono 1, 0 e $-1$. Con vettori non normalizzati il \emph{segno} resta
affidabile (positivo, zero, negativo) ma il valore non è più confrontabile. Nei giochi
il prodotto scalare risponde a domande come ``il player è davanti o dietro al
nemico?'': si normalizza la differenza tra la posizione del player e quella del nemico
e la si confronta con il \emph{Forward Vector} del nemico, che Unreal restituisce già
normalizzato (nodi \emph{Get Actor Forward Vector}, \emph{Right Vector}, \emph{Up
Vector}). Un risultato maggiore di $\cos 45^\circ \approx 0{,}71$, per esempio,
significa ``il player è dentro un cono di visione di 90 gradi''. Il quinto capitolo del
libro usa invece la normalizzazione per il movimento: si normalizza una direzione, la si
moltiplica per una velocità e per il \emph{delta time} (con il nodo \textbf{Get World
Delta Seconds}, che tornerà nella torretta) e la si somma alla posizione a ogni
\emph{tick} con \emph{Add Actor World Offset}. ``Normalizzare'', scrivono, ``è una
procedura comune che assicura che il vettore abbia lunghezza 1, il che lo rende
compatibile con il resto dei calcoli''.

% ---------------------------------------------------------------------
\section{I rotator, il gimbal lock e i quaternioni}
\label{ue4:sec:rotator}

``Andiamo avanti: poi abbiamo i \textbf{rotator}.'' Per spiegarli il docente usa
l'immagine del \textbf{giroscopio}: tre ``anelli che girano tra di loro e formano
proprio quella che è la rotazione''. Ogni anello corrisponde a un asse, con i colori di
sempre: la X in rosso, la Y in verde, la Z in blu (figura~\ref{ue4:fig:rotator}).

\paragraph{Il gimbal lock.} Con le ``regole classiche di Eulero'', cioè descrivendo la
rotazione come tre angoli applicati uno dopo l'altro, c'è ``un possibile problema'': il
\textbf{gimbal lock}. ``Non vi sto a spiegare troppo, perché fondamentalmente andremo a
parlare di quattro dimensioni invece di tre, ed è tutta una cosa molto, molto
complessa.'' In breve, ``potrebbe succedere che due assi si blocchino tra di loro'' e
comincino a ruotare ``sempre dipendenti l'uno dall'altro''. Si risolve con i
\textbf{quaternioni}, che usano quattro valori invece di tre, ma ``parliamo di
matematica molto complessa, e non andiamo molto avanti su questa cosa''. Quello che
conta, per ora: per descrivere una rotazione nello spazio servono sempre la X, la Y e
la Z, ``perché abbiamo tre giroscopi che girano di per sé''. Nell'editor, su un
qualsiasi \emph{asset}, si vedono proprio questi tre valori.

\paragraph{L'obiettivo.} ``Cosa vogliamo fare oggi? Vogliamo creare qualcosa che giri'':
qualcosa che usi il \emph{rotator} per far girare le pale del mulino.

\begin{figure}[htbp]
\centering
\begin{tikzpicture}[font=\small,
  ue4ax/.style={-{Stealth[length=2.5mm]}, very thick},
  ue4rot/.style={-{Stealth[length=2mm]}, thick}]
  % assi in proiezione
  \coordinate (O) at (0,0);
  \draw[ue4ax, red!75!black] (O) -- (-1.9,-1.1) node[below left]{X};
  \draw[ue4ax, green!50!black] (O) -- (2.6,0) node[right]{Y};
  \draw[ue4ax, blue!75!black] (O) -- (0,2.6) node[above]{Z};
  % anelli di rotazione
  \draw[ue4rot, red!75!black] (-1.25,-0.25) arc[x radius=0.35, y radius=0.6,
    start angle=90, end angle=400];
  \node[red!75!black, font=\scriptsize, align=center] at (-2.9,-0.1)
    {\textbf{Roll}\\attorno a X};
  \draw[ue4rot, green!50!black] (1.7,0.55) arc[x radius=0.3, y radius=0.55,
    start angle=90, end angle=400];
  \node[green!50!black, font=\scriptsize, align=center] at (2.4,1.05)
    {\textbf{Pitch}\\attorno a Y};
  \draw[ue4rot, blue!75!black] (0.55,1.75) arc[x radius=0.55, y radius=0.2,
    start angle=0, end angle=310];
  \node[blue!75!black, font=\scriptsize, align=center] at (1.65,2.15)
    {\textbf{Yaw}\\attorno a Z};
  % giroscopio stilizzato
  \begin{scope}[shift={(7.0,0.6)}]
    \draw[very thick, red!75!black] (0,0) ellipse (1.9 and 1.9);
    \draw[very thick, green!50!black] (0,0) ellipse (1.45 and 0.55);
    \draw[very thick, blue!75!black] (0,0) ellipse (0.45 and 1.05);
    \fill (0,0) circle (2pt);
    \node[font=\scriptsize, align=center] at (0,-2.4) {il giroscopio: tre anelli
      annidati,\\ognuno ruota attorno al proprio asse};
  \end{scope}
\end{tikzpicture}
\caption{I tre assi di Unreal con i loro colori e i nomi delle tre rotazioni di un
\emph{rotator}: \emph{roll} attorno a X, \emph{pitch} attorno a Y, \emph{yaw} attorno a
Z. A destra, l'immagine del giroscopio usata a lezione. Schema ricostruito.}
\label{ue4:fig:rotator}
\end{figure}

\paragraph{Approfondimento: roll, pitch, yaw e il gimbal lock in pratica.} Nella seconda
edizione di Romero e Sewell, l'elenco dei tipi di variabile descrive il
\textbf{Rotator} così: contiene i \emph{float} X (\textbf{Roll}), Y (\textbf{Pitch}) e
Z (\textbf{Yaw}), espressi in gradi. I tre nomi vengono dall'aviazione: il
\emph{roll} è il rollio (l'aereo che inclina le ali), il \emph{pitch} il beccheggio (il
muso che sale e scende), lo \emph{yaw} l'imbardata (il muso che gira a destra e a
sinistra). Una torretta che insegue un bersaglio usa lo \emph{yaw} per girarsi e il
\emph{pitch} per alzare la canna; le pale del mulino, come si vedrà, girano in
\emph{pitch}. Il \textbf{gimbal lock} succede quando l'anello di mezzo ruota di 90
gradi: il primo e il terzo anello si ritrovano allineati, e ruotarli produce lo stesso
movimento, così un grado di libertà va perso. Con gli angoli di Unreal il caso tipico
è il \emph{pitch} a $\pm 90$: con il muso dritto verso l'alto, \emph{yaw} e \emph{roll}
fanno la stessa cosa. I \textbf{quaternioni} (in Unreal la struttura \emph{Quat})
rappresentano invece ogni rotazione come una rotazione unica attorno a un asse
qualsiasi, con quattro numeri, e non hanno questo problema. Unreal li usa
internamente: un \emph{rotator} è il formato comodo per le persone, il quaternione
quello sicuro per i calcoli, e nodi come \emph{Combine Rotators} convertono i
\emph{rotator} in quaternioni, li combinano e li riconvertono. Un effetto collaterale,
che conviene conoscere: dato che il \emph{pitch} di un \emph{rotator} viene sempre
riportato tra $-90$ e $90$ gradi, una pala che ruota in \emph{pitch} oltre i 90 gradi
viene descritta ``capovolta'' (\emph{yaw} e \emph{roll} a 180), e chi stampasse il suo
\emph{rotator} vedrebbe numeri che saltano. La rotazione sullo schermo, però, resta
corretta.

% ---------------------------------------------------------------------
\section{Il Rotating Movement Component}
\label{ue4:sec:rotating}

``Allora, avete aperto l'editor di Unreal? Sì? Perfetto.'' Prima di costruire il
mulino, il docente mostra ``quello più semplice'' tra i ``vari modi per far ruotare un
actor in generale''; un modo che però, avverte, ``è difficile usarlo su vari
componenti''.

\paragraph{Un componente aggiunto dai Details.} Si trascina in scena un \textbf{cubo}
(con il tasto \textbf{Fine}, \emph{End}, l'oggetto ``va diretto sul
terreno''\footnote{La trascrizione dice ``il tasto AND'': è il tasto \emph{End} (Fine),
che in Unreal appoggia l'oggetto selezionato sulla superficie sottostante.}). Ora,
``qui sui \textbf{Details} sulla destra, potete aggiungere dei componenti anche non
avendo una schermata di Blueprint'': con il pulsante \textbf{Add} dei Details si può
mettere un componente direttamente sull'actor in scena. Si scrive \emph{rotating
movement} e si sceglie il \textbf{Rotating Movement Component}. ``Trovato? Sì?
Perfetto.''

\paragraph{I parametri.} Il componente ha una serie di dati, tra cui il più importante
è il \textbf{Rotation Rate}, con i tre assi X, Y e Z: la velocità di rotazione su
ciascun asse. ``Di default il Rotating Movement, come potete vedere, è messo a
\textbf{180}'' sulla Z: mezzo giro al secondo attorno all'asse verticale.

\paragraph{La regola della mobilità.} ``Intanto vi dico, ragazzi: come abbiamo visto,
se un \emph{asset} si deve muovere all'interno di uno spazio, la sua mobilità dovrà
essere \textbf{Movable}. Ricordatevelo sempre.'' Se lo si lascia \textbf{Static} e si
va in Play, ``prima di tutto non si muove'', e poi l'engine restituisce un errore,
che il docente legge: \emph{Static Mesh Component has to be movable if you'd like to
move it with Rotating Movement Component}. Con la mobilità corretta, in Play il cubo
gira su sé stesso: ``abbiamo fatto le munizioni di \emph{Metal Gear Solid}'', scherza
il docente, pensando agli oggetti da raccogliere che ruotano a mezz'aria. ``Questo è un
modo per far ruotare'': Unreal ``ci dà la possibilità di poter ruotare direttamente un
nostro bellissimo actor tramite un componente che loro hanno già creato''. ``Un oggetto
ruotante che ruota. Perché `ruotante' è brutto, sembra.''

\paragraph{Esercitazione: i tre cubi.} ``Piccola esercitazione velocissima, faccio
queste botte e risposte con le esercitazioni.'' Si copia il cubo (``tanto vi ricordate
come si copiano'') e si fa in modo che \textbf{un cubo ruoti sull'asse X} e
\textbf{un altro sull'asse Y}, alla velocità che si preferisce. ``Ricordatevi di
metterli Movable. Vi do cinque minuti.'' La correzione: sul secondo cubo si azzera la Z
e si mette 180 sulla X; sul terzo, 180 sulla Y. Si controlla che siano tutti e tre
Movable, si va in Play, ``ed ecco le nostre tre rotazioni''. ``Possiamo farlo con
qualsiasi \emph{asset}.''

\paragraph{Approfondimento: la sedia che gira.} Il Rotating Movement Component è il
primo componente che Romero e Sewell fanno usare nel libro, già nel primo capitolo: un
Blueprint \emph{BP\_RotatingChair} con una \emph{Static Mesh} (la sedia dello
\emph{Starter Content}) e un Rotating Movement. ``Di default'', scrivono, ``questo
componente ruoterà il Blueprint attorno all'asse Z, quindi non serve cambiarne le
proprietà'': si compila, si trascina la sedia nel livello e la si guarda girare in Play.
Il componente fa parte della famiglia dei \emph{movement component}, componenti che non
hanno una forma propria ma muovono un altro componente, quello indicato come
\emph{Updated Component}, che di default è la \textbf{radice} dell'actor. È qui la
ragione del limite che il docente scoprirà tra poco con il mulino: il Rotating Movement
fa girare la radice, e con essa tutto ciò che vi è attaccato. Il componente ha anche
un'opzione \emph{Rotation in Local Space} e un \emph{Pivot Translation}, che sposta il
centro di rotazione; ma per far girare una parte sola di un actor la strada più chiara
resta quella che segue, il Blueprint.

% ---------------------------------------------------------------------
\section{Il Blueprint del mulino}
\label{ue4:sec:mulino}

``Proviamo a fare intanto il nostro \emph{windmill}.'' Il mulino del livello di esempio
è formato da più cose, si vede nei Details. ``Chi si ricorda come si chiamano queste
cose qua, che sono tutte insieme in un singolo \emph{asset}?'' ``Componenti.''
``Sono i \textbf{componenti}'': come si è visto nella seconda lezione, un actor ``può
essere, scusate la ridondanza, composto da più componenti''. In questo caso un
\textbf{body}, il corpo, ``la nostra \emph{mesh} principale, il cubone'', e
\textbf{due pale}.

\paragraph{La classe.} Si crea un nuovo \textbf{Blueprint Class} di tipo
\textbf{Actor}, ``perché non sarà disponibile al giocatore, ma sarà solamente là come
sfondo'': non va controllato, non è un \emph{Pawn}. Il docente lo chiama
\textbf{BP\_Windmill}\footnote{La trascrizione dice ``Windmill, io in questo caso lo
chiamerò 1'': probabilmente \emph{BP\_Windmill1}, per non confonderlo con il mulino già
presente nel progetto.}.

\paragraph{Il corpo.} Nella \textbf{Viewport} del Blueprint si aggiunge una
\textbf{Static Mesh} (Add, in alto a sinistra, \emph{Static Mesh}) e la si chiama
\textbf{Body}. Nei dettagli, alla voce \emph{Static Mesh}, si sceglie un cubo, per
esempio \emph{SM\_Cube}. Il primo cubo scelto dal docente ``ha un \emph{pivot}
particolarmente strano''; ne sceglie un altro, ``molto meglio, perché c'ha il
\emph{pivot} al centro''. Lo solleva di circa 40 sulla Z e gli dà un materiale: il
docente mette la \emph{prototype grid}, ``voi mettete quello che volete''. Poi lo scala
per farne il corpo del mulino: ``3 metri per 5'' di altezza, ``di larghezza facciamo 2
metri'', cioè una scala di circa 3, 2 e 5 sui tre assi\footnote{La trascrizione dice
``x32 sulla y5 sull'altezza'' e ``3:25'': la lettura più probabile è X 3, Y 2, Z 5. Il
cubo di base di Unreal misura un metro di lato, quindi una scala di 3 corrisponde a 3
metri.}.

\paragraph{La viewport che non si trova.} Uno studente non ricorda come si apre la
viewport del Blueprint. ``O ce l'hai qua, o ce l'hai qua, oppure dovresti avere
\textbf{Window} e dovresti vedere \textbf{Viewport}.'' La voce è già spuntata, ma la
finestra non c'è: forse è ``nascosta da qualche schermata di Windows'' (``prova a fare
Alt+Tab''). Si passa alla condivisione dello schermo, e si scopre che il problema è un
altro: lo studente ha creato il Blueprint ma non lo ritrova. La soluzione è il
\textbf{filtro} del Content Browser: le ``tre stanghettine'' vicino alla ricerca,
categoria \emph{Blueprint}, ed ecco il file; doppio clic e, se si apre la versione
ridotta, \emph{Open Full Blueprint Editor}. Poi Add, ma nell'editor in italiano la voce
non si trova subito (è ``Componente Mesh Statica''). ``Effettivamente, ragazzi, vi
consiglio di mettervi l'editor in \textbf{inglese}, perché è molto più semplice. Anzi,
ve lo consiglio proprio'': i nomi dei nodi e dei componenti nei tutorial, nei libri e
nella documentazione sono quasi sempre in inglese.

\paragraph{Le pale.} ``Adesso dobbiamo creare le nostre palette.'' Si copia e incolla il
corpo (Ctrl+C, Ctrl+V): le copie finiscono dentro il corpo, ``quindi cliccatele e
tiratele fuori, perché sono uguali''. Si scalano a forma di pala, una orizzontale e una
verticale, si dà loro un materiale diverso (il docente usa \emph{M\_Grid}, ``questa
qua violacea'') e le si rinomina \textbf{Pala1} e \textbf{Pala2}.

\paragraph{Attaccare le pale al corpo.} ``Importante, ragazzi: queste pale devono essere
attaccate al corpo. Come facciamo? Abbiamo visto l'ereditarietà.'' È la
\textbf{gerarchia} dei componenti: ``il fatto che io posso mettere degli \emph{asset}
attaccati l'uno con l'altro''. Si prende una pala e la si \textbf{trascina sopra} Body,
nell'elenco dei componenti: ora è sua figlia. ``Scusami, non ho capito come si fa.''
``Se hai questa pala, la prendi e la trascini sopra Body.''

\paragraph{Il Rotating Movement non basta.} ``Adesso dobbiamo dare una rotazione a
queste pale. Potremmo provare con il Rotating Movement.'' Il docente lo aggiunge al
Blueprint, mette il mulino in scena (``molto più grosso'') e va in Play. ``Guardate che
succede'': il Rotating Movement è impostato sulla Z, e \textbf{ruota tutto il corpo}.
``Una particolarità del Rotating Movement è che non può essere attaccato a
componenti: o ruota tutto, o non ruota niente.'' Quindi bisogna programmare il
movimento delle pale. Il Rotating Movement va tolto: ``se l'avete messo e avete adesso
il vostro mulino che ruota su sé stesso, lo dovete eliminare''.

% ---------------------------------------------------------------------
\section{Far girare una pala: Event Tick e Combine Rotators}
\label{ue4:sec:pala}

``Domanda. Questo movimento deve essere fatto per ogni \emph{frame}. Quindi deve essere
un evento che riprende ogni \emph{frame}. Chi si ricorda qual è l'evento che lo fa? L'abbiamo
usato per la piattaforma dell'altra volta.'' ``L'\textbf{Event Tick}.'' ``Ottimo.''
L'Event Tick, ``che sarà il nostro protagonista'', è già presente nell'Event Graph di un
nuovo Blueprint.

\paragraph{I delta seconds, di nuovo.} ``Chi si ricorda che sono i \textbf{Delta
Seconds}?'' Uno studente: ``determina quanti tick al minuto\ldots''. Un altro: ``i
secondi effettivi, e non quelli basati su \emph{frame}''. Il docente riprende: è ``lo
spazio che intercorre tra un \emph{frame} e l'altro'', la differenza di tempo tra due
\emph{frame}, ``in questo caso si parla di millisecondi''. Si moltiplica un valore per i
\emph{delta seconds} perché la durata di un \emph{frame} cambia (a 30 \emph{frame} al
secondo è molto più lunga che a 120), e così gli spostamenti ``non diventano più lenti
se abbiamo meno \emph{frame}, o più veloci se ne abbiamo di più''. (Il docente prova a
disegnarlo sullo schermo, ma lo strumento ``a quanto pare non mi funziona benissimo''.)

\paragraph{La velocità.} Si tira fuori un cavo da Delta Seconds e si crea un
\textbf{Multiply}: o scrivendo \emph{multiply}, o digitando semplicemente
l'\textbf{asterisco}, ``e vi dà subito il Multiply''. Il secondo valore è la
\textbf{velocità}: ``più è alto, più è veloce''. Per ora, \textbf{90}.

\paragraph{Make Rotator.} ``Adesso vi dico una cosa.'' Il nodo \textbf{Make Rotator}
costruisce un \emph{rotator}, ``quello che muove a livello di rotazione un
\emph{asset}'', a partire da tre \emph{float}: X, Y e Z. Quale asse serve? ``Prendiamo
la nostra bellissima pala e guardiamo subito in che asse la vogliamo ruotare. La X non
mi sembra, la Z neanche: è la nostra carissima \textbf{Y}'', cioè il \textbf{pitch}.
Si collega il risultato del Multiply al pin Y (Pitch): il \emph{rotator} creato ha X e
Z a 0, ``perché non si devono muovere'', e la Y uguale a velocità per \emph{delta
seconds}.

\paragraph{Combine Rotators.} ``Possiamo farlo in tanti modi, ma oggi ne facciamo tipo
tre.'' Il primo è \textbf{Combine Rotators}, che ``combina due \emph{rotator} per
darne uno in uscita''. Il primo dei due è l'\textbf{origine}: ``a livello concettuale,
qual è l'origine che dobbiamo prendere?'' ``Le coordinate del componente.'' ``Sì:
quale componente?'' Per ora una sola pala, la Pala1. Siccome serve un componente e non
l'intero actor, si trascina la pala dall'elenco dei componenti nel grafo: si ottiene
una \textbf{referenza} a quel componente, ``un \emph{get}'' che ``si riferisce a quel
componente preciso''.

\paragraph{Relativa o del mondo?} Da quella referenza si chiede la rotazione. Ma quale?
``Abbiamo visto la differenza tra \textbf{rotazione relativa} e \textbf{rotazione del
mondo}.'' La pala è attaccata al corpo, ma deve muoversi da sola, indipendentemente dal
corpo. ``Relativa al corpo, secondo me'', propone uno studente. ``Relativa a sé
stessa, quindi la sua asse singola. Se prendiamo quella del mondo, girerà in modo
strano.'' Si usa dunque \textbf{Get Relative Rotation}, la si collega al primo
ingresso di Combine Rotators, e il \emph{rotator} di Make Rotator al secondo.

\paragraph{Set Relative Rotation.} ``Chi si ricorda qual è il blocco che serve per
settare la rotazione? Tanto è un \emph{set}.'' È \textbf{Set Relative Rotation},
``perché, ovviamente, abbiamo preso la relativa''. Lo si crea tirando fuori un cavo da
Pala1 (così il \emph{target} è già la pala), lo si collega all'Event Tick e gli si dà
in ingresso la nuova rotazione, l'uscita di Combine Rotators
(figura~\ref{ue4:fig:pala}). ``E mi raccomando, deve essere \textbf{Movable}.'' Si
mette il mulino in scena, si va in Play: ``ed ecco che una pala ruota''.

\begin{figure}[htbp]
\centering
\begin{tikzpicture}[font=\small,
  ue4node/.style={draw, thick, rounded corners=4pt, align=center, fill=gray!8,
                  minimum height=0.9cm},
  ue4event/.style={ue4node, fill=red!12},
  ue4exec/.style={-{Stealth[length=2.5mm]}, very thick, black!80},
  ue4flt/.style={-{Stealth[length=2mm]}, thick, green!50!black},
  ue4rot/.style={-{Stealth[length=2mm]}, thick, violet!70!black},
  ue4obj/.style={-{Stealth[length=2mm]}, thick, blue!60!black}]
  \node[ue4event, text width=2.1cm] (tick) at (0,0)
    {\textbf{Event Tick}\\\footnotesize Delta Seconds};
  \node[ue4node, text width=1.5cm] (mul) at (2.9,-1.6)
    {\textbf{Multiply}\\\footnotesize $\times$ 90};
  \node[ue4node, text width=1.9cm] (make) at (5.6,-1.6)
    {\textbf{Make Rotator}\\\footnotesize X 0, \textbf{Y}, Z 0};
  \node[ue4node, fill=blue!8, text width=1.3cm] (pala) at (2.9,-3.5)
    {\textbf{Pala1}};
  \node[ue4node, text width=2.1cm] (get) at (5.6,-3.5)
    {\textbf{Get Relative}\\\textbf{Rotation}};
  \node[ue4node, text width=2.0cm, minimum height=1.5cm] (comb) at (8.6,-2.5)
    {\textbf{Combine}\\\textbf{Rotators}\\\footnotesize A, B};
  \node[ue4node, text width=2.2cm, minimum height=1.6cm] (set) at (11.6,0)
    {\textbf{Set Relative}\\\textbf{Rotation}\\\footnotesize Target: Pala1};
  \draw[ue4exec] (tick.east) -- (set.west |- tick.east);
  \draw[ue4flt] ([yshift=-3mm]tick.east) -- ++(0.4,0) |- (mul.west);
  \draw[ue4flt] (mul.east) -- node[above, font=\scriptsize]{Pitch} (make.west);
  \draw[ue4obj] (pala.east) -- (get.west);
  \draw[ue4rot] (get.east) -- ++(0.4,0) |- node[pos=0.75, below,
    font=\scriptsize]{A} ([yshift=-3mm]comb.west);
  \draw[ue4rot] (make.east) -- ++(0.4,0) |- node[pos=0.75, above,
    font=\scriptsize]{B} ([yshift=3mm]comb.west);
  \draw[ue4rot] (comb.east) -- node[above, font=\scriptsize]{New Rotation}
    ++(1.6,0) |- ([yshift=-4mm]set.west);
  \draw[ue4obj, dashed] (pala.south) -- ++(0,-0.5) -| node[pos=0.25, below,
    font=\scriptsize]{Target} (set.south);
\end{tikzpicture}
\caption{Il grafo della prima pala. A ogni \emph{frame}, l'Event Tick esegue Set
Relative Rotation su Pala1; la nuova rotazione è la rotazione relativa attuale della
pala combinata con un piccolo incremento in \emph{pitch}, pari a 90 gradi al secondo
moltiplicati per i \emph{delta seconds}. Il cavo nero è l'esecuzione, gli altri sono
dati (verde i \emph{float}, viola i \emph{rotator}, blu le referenze). Schema
ricostruito.}
\label{ue4:fig:pala}
\end{figure}

\paragraph{I problemi degli studenti.} ``Avete la vostra paletta che ruota?'' Non
tutti. Una studentessa chiede di rivedere il codice e non trova il blocco Get Relative
Rotation: ``tira fuori da Pala1 e scrivi \emph{get relative rotation}''. Un altro
studente vede la pala ``staccarsi da tutto l'asset'' quando preme Play, e in
condivisione si scoprono due errori. Il primo: ha calcolato la rotazione con Combine
Rotators ma non l'ha collegata al Set (``hai calcolato la rotazione, ma non
gliel'hai assegnata''). Il secondo: anche dopo, la pala gira attorno a un punto
sbagliato, perché la \emph{mesh} che ha scelto ha il \textbf{pivot} spostato. ``Vedi che
il pivot sta quasi lì? Cambia la \emph{mesh}'': con un cubo che ha il pivot al centro, la
pala gira correttamente. (Nel frattempo il docente accenna anche a un \textbf{Set World
Rotation}, per mostrare la differenza con la rotazione relativa.)

\paragraph{Approfondimento: relativo e mondo nel libro.} Il quattordicesimo capitolo di
Romero e Sewell spiega la differenza con un esempio molto simile al mulino. Un
Blueprint ha come radice un \emph{DefaultSceneRoot}; sotto c'è un componente
\emph{Table}, un tavolo, e sotto il tavolo un componente \emph{Statue}, una statua. Il
\emph{transform} di Table è \textbf{relativo} a DefaultSceneRoot, e quello di Statue è
relativo a Table: ``se muovete il componente Table nella Viewport, anche il componente
Statue si muoverà; ma se cambiate il \emph{transform} relativo di Statue, Table resterà
dov'è''. Lo stesso vale per il mulino: la rotazione relativa di Pala1 è misurata
rispetto a Body. Se il mulino viene ruotato nel livello, la rotazione relativa della
pala non cambia, e il grafo continua a farla girare attorno al proprio asse; usando la
rotazione del mondo, invece, l'incremento in \emph{pitch} verrebbe applicato attorno
all'asse Y del \emph{mondo}, che coincide con quello della pala solo se il mulino non è
ruotato: è il ``girerà in modo strano'' del docente. Nei vettori lo si vede ancora più
chiaramente: il punto iniziale di ogni posizione nel mondo è sempre $(0,0,0)$, mentre
quello di una posizione relativa è il punto in cui sta il genitore.

\paragraph{Approfondimento: una scorciatoia, Add Local Rotation.} Il trio Get Relative
Rotation, Combine Rotators e Set Relative Rotation è istruttivo perché mostra i tre
passaggi di ogni animazione fatta a mano: \emph{leggi} lo stato attuale, \emph{calcola}
il nuovo, \emph{scrivi} il risultato. Unreal offre anche nodi che fanno tutto in uno:
\textbf{Add Local Rotation} (o \emph{Add Relative Rotation}) su un componente aggiunge
un \emph{Delta Rotation} alla rotazione attuale, così come \emph{Add Actor World
Rotation}, usato nella seconda lezione per la piattaforma, lo fa per un intero actor.
Collegando l'Event Tick a un Add Local Rotation su Pala1, con un Delta Rotation fatto
con Make Rotator come sopra, si ottiene lo stesso mulino con un nodo in meno. Il
risultato è lo stesso; cambia la chiarezza di ciò che succede, ed è per questo che a
lezione conviene vedere prima la versione lunga.

% ---------------------------------------------------------------------
\section{La gerarchia: la seconda pala e il cubetto}
\label{ue4:sec:gerarchia}

``Possiamo fare una cosa bella: invece di dover fare la copia esatta per quanto
riguarda la nostra prima pala, possiamo vedere quello che la \textbf{gerarchia} ci
permette di fare in modo molto facile.'' La Pala2 deve seguire in ogni caso il movimento
della Pala1. Invece di programmare anche lei, la si mette ``in eredità'' (``non quella
con Carlo Conti, ragazzi'') alla Pala1: la si trascina sopra Pala1, e diventa sua
\textbf{figlia}. Play: ``non abbiamo programmato niente'', ma la Pala2 segue la
rotazione della Pala1. ``È la Pala1 che gira, non la Pala2; la Pala2 segue solamente.''

\paragraph{Esercitazione: il cubetto sulla pala.} ``Se noi volessimo attaccare un
cubettino su un estremo della pala\ldots anzi, sapete che c'è? Lo fate voi.'' La
consegna: aggiungere una \textbf{Static Mesh} qualsiasi (``un cubo, una \emph{mesh},
quello che vi pare''), piccola, su un'estremità della pala, e attaccarla alla
\textbf{Pala2}, in modo che segua anche lei il movimento. ``Vi do cinque minuti
classici. Ricordatevi l'ereditarietà: per spostarlo dovete prenderlo proprio e metterlo
sopra.''

\paragraph{La domanda sul pivot.} Uno studente l'ha fatto, ma è sorpreso: ``mi
aspettavo di dover spostare il pivot al centro della rotazione. Invece,
automaticamente, è come se si agganciasse alla pala.'' ``Esatto. Perché comunque
\textbf{mantiene la sua posizione relativa}'': il cubetto ha il suo pivot al proprio
centro, ma la sua posizione è espressa rispetto al genitore, e quando il genitore
ruota, il figlio ruota con lui attorno al pivot del genitore. ``Se lo dovessimo
spostare, andrebbe insieme: manterrebbe la sua origine, il fatto che sta sulla parte
destra, sopra, giù, quello che è, ma si muoverebbe relativo all'altro.'' Play: ``eccolo
che gira insieme''. ``Così abbiamo capito come si ereditano piano piano gli oggetti tra
di loro'' (figura~\ref{ue4:fig:gerarchia}).

\begin{figure}[htbp]
\centering
\begin{tikzpicture}[font=\small,
  ue4comp/.style={draw, thick, rounded corners=3pt, minimum width=2.6cm,
                  minimum height=0.7cm, align=center},
  ue4tree/.style={thick, gray!60!black},
  ue4arr/.style={-{Stealth[length=2.5mm]}, thick}]
  % albero dei componenti
  \node[ue4comp, fill=gray!10] (root) at (0,0) {BP\_Windmill (actor)};
  \node[ue4comp, fill=orange!15] (body) at (0.6,-1.0) {Body};
  \node[ue4comp, fill=violet!12] (p1) at (1.2,-2.0) {Pala1 \footnotesize(ruota)};
  \node[ue4comp, fill=violet!6] (p2) at (1.8,-3.0) {Pala2 \footnotesize(segue)};
  \node[ue4comp, fill=green!10] (cub) at (2.4,-4.0) {Cubetto \footnotesize(segue)};
  \foreach \a/\b in {root/body, body/p1, p1/p2, p2/cub}
    \draw[ue4tree] ([xshift=-1.0cm]\a.south) |- (\b.west);
  \node[font=\footnotesize, align=center] at (1.2,-5.0) {l'elenco dei componenti};
  % mulino stilizzato
  \begin{scope}[shift={(8.2,-2.2)}]
    \draw[thick, fill=orange!15] (-0.75,-2.5) rectangle (0.75,1.0);
    \node[font=\scriptsize] at (0,-2.2) {Body};
    \begin{scope}[rotate=25]
      \draw[thick, fill=violet!25] (-2.2,-0.18) rectangle (2.2,0.18);
      \draw[thick, fill=violet!12] (-0.18,-2.2) rectangle (0.18,2.2);
      \draw[thick, fill=green!35] (-0.25,1.85) rectangle (0.25,2.35);
    \end{scope}
    \fill (0,0) circle (2.2pt);
    \draw[ue4arr, red!70!black] (2.5,0.4) arc[start angle=10, end angle=70,
      radius=2.5];
    \node[font=\scriptsize, align=left] at (3.6,-0.4) {Pala1 orizzontale\\Pala2
      verticale\\cubetto in cima};
  \end{scope}
\end{tikzpicture}
\caption{La gerarchia del mulino. Solo Pala1 è programmata per ruotare; Pala2, sua
figlia, e il cubetto, figlio di Pala2, ne seguono la rotazione senza alcun nodo,
mantenendo la loro posizione relativa. Schema ricostruito (nel mulino della lezione le
pale sono di fronte al corpo).}
\label{ue4:fig:gerarchia}
\end{figure}

% ---------------------------------------------------------------------
\section{Il mulino che guarda il player: Find Look at Rotation}
\label{ue4:sec:lookat}

``Adesso vi farò vedere un altro metodino'' per creare una rotazione, sempre
sull'Event Tick, ``così vi faccio vedere anche un altro blocchettino interessante,
perché sono tanti''. Nei Blueprint, come nella programmazione in generale, ``non esiste
un qualcosa che sia fisso'': ``essendo che è un mondo di spazi, punti, vettori, ce lo
possiamo giocare come ci pare''. Il docente ricorda una persona che una volta ha fatto
``un materiale con la matematica senza usare texture. Non so come abbia fatto, ma si
può fare''.

\paragraph{Staccare i cavi.} Si stacca l'Event Tick dal grafo precedente. ``Ricordatevi:
per staccare i cavettini, tenete premuto \textbf{Alt} e cliccate sui puntini.''

\paragraph{Find Look at Rotation.} ``Un blocco chiamato \textbf{Find Look at Rotation}.
Interessante, molto interessante. Questo è un sacco utile, ve lo dico.'' Il docente
legge la descrizione del nodo: \emph{find a rotation for an object at Start location to
point at Target location}. Alla lavagna (``scrittura incredibile, mi dispiace'') disegna
due punti nello spazio: lo \textbf{Start}, ``da dove parto'', e il \textbf{Target}, il
punto da raggiungere. Il nodo restituisce ``la rotazione verso quel punto'': la
rotazione che un oggetto posto in Start deve avere per guardare Target.

\paragraph{Un falso inizio.} Il docente parte con un esempio sbagliato (lo Start e un
Target ottenuto sommando 360, poi 359, alla posizione), e si ferma: ``scusate, ragazzi,
perdonatemi: il Find Look at Rotation lo facciamo per un esercizio che era diverso.
Ho sbagliato io, perdonate, perché non ho gli appunti.'' L'esercizio vero: far ruotare
il mulino \textbf{in direzione del character}, ``un mulino molto magico'' che si gira
verso il giocatore.

\paragraph{Il grafo.} I due punti sono:
\begin{itemize}
  \item lo \textbf{Start}: la posizione del mulino, con \textbf{Get Actor Location}
        (senza \emph{target}, cioè il Blueprint stesso);
  \item il \textbf{Target}: la posizione del giocatore, con \textbf{Get Player
        Character} e, da questo, Get Actor Location.
\end{itemize}
Il Find Look at Rotation ``ci darà il vettore verso la posizione del nostro player'',
o meglio la rotazione che lo punta.

\paragraph{RInterp To: l'interpolazione.} ``Possiamo usare anche un altro blocco, che ci
darà un'\textbf{interpolazione}.'' Si chiama \textbf{RInterp To} (R come
\emph{rotator}). ``Qualcuno sa cosa vuol dire interpolazione? No.'' ``È letteralmente
il passaggio tra due punti'': da un valore A a un valore B. I nodi di questo tipo
(\emph{Lerp}, \emph{VLerp}, \emph{Interp}) sono ovunque: ``avete presente le animazioni?
Un'animazione è un passaggio da un punto A a un punto B'', e di base è fatta a livello
di valori, perché ruotare un braccio è cambiare i valori di mano, braccio, gamba. Il
nodo ha quattro ingressi:
\begin{itemize}
  \item \textbf{Current}, la rotazione attuale del mulino (Get Actor Rotation, o la
        rotazione relativa del corpo);
  \item \textbf{Target}, la rotazione appena trovata da Find Look at Rotation;
  \item \textbf{Delta Time}, che ``ci dice in quanto tempo'' si fa l'interpolazione;
  \item \textbf{Interp Speed}, ``la velocità in cui lo fa''.
\end{itemize}

\paragraph{Solo lo yaw.} ``Però una cosa importante, ragazzi: non vogliamo che si muova
tutto sulla stessa asse.'' Il mulino deve girare ``al massimo così'', in orizzontale,
non inclinarsi verso l'alto o verso il basso. ``Quale asse dovrò influenzare?''
``La Z?'' ``Esatto, perfetto.'' Si fa \textbf{Split Struct Pin} (tasto destro sul pin)
sull'uscita di Find Look at Rotation e sull'ingresso Target, e si collega solo la Z,
lo \emph{yaw}. Il risultato dell'interpolazione si applica con un nodo di rotazione
(nella trascrizione, \emph{Set Relative Rotation})\footnote{La trascrizione dice ``set
relative rotation di pala 1'', ma a girare verso il player è tutto il mulino. Con ogni
probabilità il nodo è un Set Relative Rotation sul corpo (Body), a cui sono attaccate le
pale, o un Set Actor Rotation sull'intero Blueprint: il risultato visibile è lo
stesso.}, collegato all'Event Tick (figura~\ref{ue4:fig:lookat}).

\paragraph{L'offset di 90 gradi.} In Play il mulino si gira verso il player, ma di
lato. ``Possiamo mettere più 90'': si somma \textbf{90} alla Z prima di collegarla, ``in
modo tale che sia\ldots ecco''. Ora il mulino guarda il giocatore. ``Quello che sta
facendo è continuare a calcolare, visto che lo fa sull'Event Tick, il vettore che
indirizza in un punto, che in questo caso è il player, in modo continuo.'' ``Questo è un
blocco che serve moltissimo, ragazzi'': serve ogni volta che bisogna fare ``delle robe a
\emph{target}'', e servirà per la torretta. ``Vedete che segue continuamente la nostra
torretta? Questo è fatto con il Find Look at Rotation.''

\begin{figure}[htbp]
\centering
\begin{tikzpicture}[font=\small,
  ue4node/.style={draw, thick, rounded corners=4pt, align=center, fill=gray!8,
                  minimum height=0.9cm},
  ue4event/.style={ue4node, fill=red!12},
  ue4exec/.style={-{Stealth[length=2.5mm]}, very thick, black!80},
  ue4flt/.style={-{Stealth[length=2mm]}, thick, green!50!black},
  ue4rot/.style={-{Stealth[length=2mm]}, thick, violet!70!black},
  ue4vecw/.style={-{Stealth[length=2mm]}, thick, yellow!50!black},
  ue4obj/.style={-{Stealth[length=2mm]}, thick, blue!60!black}]
  \node[ue4event, text width=2.0cm] (tick) at (0,0)
    {\textbf{Event Tick}\\\footnotesize Delta Seconds};
  \node[ue4node, text width=2.0cm] (self) at (0,-1.9)
    {\textbf{Get Actor}\\\textbf{Location}\\\footnotesize (self)};
  \node[ue4node, fill=blue!8, text width=2.0cm] (pc) at (0,-3.9)
    {\textbf{Get Player}\\\textbf{Character}};
  \node[ue4node, text width=2.0cm] (pcl) at (2.8,-3.9)
    {\textbf{Get Actor}\\\textbf{Location}};
  \node[ue4node, text width=2.1cm, minimum height=1.4cm] (look) at (6.0,-2.6)
    {\textbf{Find Look at}\\\textbf{Rotation}};
  \node[ue4node, text width=1.2cm] (add) at (8.7,-2.6) {\textbf{$+$ 90}\\
    \footnotesize su Z};
  \node[ue4node, text width=2.0cm] (cur) at (8.7,-4.5)
    {\textbf{Get Actor}\\\textbf{Rotation}};
  \node[ue4node, text width=2.6cm, minimum height=2.2cm] (interp) at (11.5,-2.9)
    {\textbf{RInterp To}\\[2pt]\footnotesize Current\\\footnotesize Target (Z)\\
    \footnotesize Delta Time (Tick)\\\footnotesize Interp Speed};
  \node[ue4node, text width=1.5cm, minimum height=1.3cm] (set) at (14.3,0)
    {\textbf{Set}\\\textbf{Rotation}\\\footnotesize (Z)};
  \draw[ue4exec] (tick.east) -- (set.west |- tick.east);
  \draw[ue4vecw] (self.east) -- node[above, font=\scriptsize]{Start} (4.3,-1.9)
    |- ([yshift=3mm]look.west);
  \draw[ue4obj] (pc.east) -- (pcl.west);
  \draw[ue4vecw] (pcl.east) -- node[below, font=\scriptsize, pos=0.6]{Target}
    (4.5,-3.9) |- ([yshift=-3mm]look.west);
  \draw[ue4flt] (look.east) -- node[above, font=\scriptsize]{Z} (add.west);
  \draw[ue4flt] (add.east) -- ++(0.5,0) |- ([yshift=2mm]interp.west);
  \draw[ue4rot] (cur.east) -- ++(0.5,0) |- ([yshift=-5mm]interp.west);
  \draw[ue4rot] (interp.east) -- ++(0.4,0) |- ([yshift=-4mm]set.west);
\end{tikzpicture}
\caption{Il mulino che guarda il player. A ogni \emph{frame}, Find Look at Rotation
calcola la rotazione che dal mulino (Start) punta al giocatore (Target); se ne tiene
solo lo \emph{yaw} (Z), gli si sommano 90 gradi per compensare l'orientamento della
\emph{mesh}, e RInterp To fa avvicinare gradualmente la rotazione attuale a quella
cercata. Il Delta Time può venire dall'Event Tick, come qui (un cavo non disegnato per
chiarezza), o essere un numero fisso (a lezione, 1); si veda l'approfondimento. Schema ricostruito.}
\label{ue4:fig:lookat}
\end{figure}

\paragraph{``A me non va.''} Seguono tre studenti in condivisione.
\begin{itemize}
  \item Al primo il mulino resta fermo: ``ti manca il Delta Time e
        l'Interp Speed''. Li riempie (``metti uno e due''), compila (``non ho
        compilato, vabbè'') e aggiunge l'offset: ``tira fuori dalla Z, fai più, metti
        90''. ``Perfetto, adesso ti segue.''
  \item Al secondo manca anche a lui il Delta Time: ``metti uno''.
  \item Il terzo, collegando il nodo a Find Look at Rotation, si ritrova un pin
        diviso che non voleva: ha preso \textbf{Get Actor Rotation} invece di
        \textbf{Get Actor Location}. Poi non riesce a collegare l'uscita alla Z:
        ``lo devi splittare: tasto destro sul pinnetto, quello violaceo, quello
        lilla, Split Struct Pin''. E infine l'offset: da Return Value di Find Look at
        Rotation, la Z, un ``più'', 90, e al Target.
\end{itemize}

\paragraph{Approfondimento: che cosa calcola Find Look at Rotation.} Il nodo fa una cosa
semplice: sottrae la posizione di partenza da quella di arrivo, ottenendo il vettore
direzione $\vec d = T - S$ (la sottrazione della sezione~\ref{ue4:sec:vettori}), e
trasforma quella direzione in un \emph{rotator} con \emph{roll} nullo. Per la sola
rotazione orizzontale, quella che interessa al mulino e alla torretta, il conto è
\[
  \text{yaw} \;=\; \operatorname{atan2}(d_y,\ d_x) \quad \text{(in gradi)},
\]
cioè l'angolo, misurato nel piano XY a partire dall'asse X, sotto cui il bersaglio è
visto. Il risultato presuppone che la ``faccia'' dell'oggetto guardi lungo il suo asse
\textbf{X} locale, che in Unreal è la direzione ``in avanti'' (il \emph{Forward
Vector}). Se la \emph{mesh} è stata modellata guardando lungo la Y, come le pale del
mulino della lezione, la rotazione calcolata è sfasata di 90 gradi, e il ``più 90'' del
docente non fa che correggere questo sfasamento. Lo stesso problema, come si vedrà, si
ripresenta con la torretta.

\paragraph{Approfondimento: come funziona RInterp To.} A ogni chiamata, RInterp To
calcola quanto manca per arrivare al Target e ne percorre una frazione, pari
(semplificando) a
\[
  \alpha \;=\; \min\bigl(1,\ \text{Delta Time} \times \text{Interp Speed}\bigr),
  \qquad
  \text{nuova rotazione} = \text{Current} + \alpha \cdot (\text{Target} - \text{Current}),
\]
con la differenza presa sempre ``dal lato più corto'' (da 350 a 10 gradi si passa per
0, non per 180). Ne seguono tre cose utili. Primo: il movimento è veloce quando il
bersaglio è lontano e rallenta avvicinandosi, perché a ogni \emph{frame} si percorre
una frazione di una distanza sempre più piccola; è l'effetto ``morbido'' che si vede in
Play. Secondo: il nodo è pensato per ricevere come \textbf{Delta Time} proprio i
\emph{delta seconds} (dall'Event Tick o da \emph{Get World Delta Seconds}) e come
\textbf{Interp Speed} un numero intorno a 1--10: così la velocità non dipende dal
\emph{frame rate}. Con un Delta Time fisso, come l'1 messo a lezione, la frazione
percorsa a ogni \emph{frame} è sempre la stessa, e il mulino gira più veloce su un PC
che fa più \emph{frame}. Terzo: se il prodotto supera 1, la frazione è 1 e la rotazione
\textbf{scatta} subito al Target; e se l'Interp Speed è 0 o negativo, Unreal restituisce
direttamente il Target. Con Delta Time 1 e Interp Speed 2 il prodotto è 2: il mulino,
in pratica, ``snappa''. Il fenomeno tornerà con la torretta
(sezione~\ref{ue4:sec:problemi}).

% ---------------------------------------------------------------------
\section{Il Clamp, il debug con Print String e la pausa}
\label{ue4:sec:clamp}

``Abbiamo fatto il nostro bellissimo Find Look at Rotation. Potremmo anche
\textbf{clamparlo}.'' ``Si ricorda il Clamp qualcuno di voi?'' ``È quello per dare un
valore minimo'', e massimo. ``Esatto, bravissimo.'' Lo scopo: limitare la rotazione del
mulino, in modo che non possa girarsi del tutto.

\paragraph{Il grafo.} Si splitta il pin con Split Struct Pin, si inserisce un
\textbf{Clamp (float)}, ``perché è un \emph{float}'', sulla Z, e la si ricompone con un
\textbf{Make Rotator} prima di collegarla. Il primo tentativo è un intervallo stretto,
da $-25$ a $25$ gradi, ma qualcosa non torna: i valori in gioco sono molto più grandi
(``mi dice: Pala1 è un botto di valori'').

\paragraph{Il debug con Print String.} ``Questo come ho fatto, ragazzi? Questo è un modo
per \textbf{debuggare} le cose.'' Non sapendo di quanto ruotasse il mulino e quali
limiti dargli, il docente collega un \textbf{Print String} al valore e lo stampa a
schermo mentre si muove attorno al mulino: ``abbiamo stampato letteralmente il
valore''. I numeri che compaiono vanno più o meno da $-160$ a $160$. Il docente prova
quei limiti, poi li allarga: ``snappa un pochino. Facciamo duecento massimo. Molto
meglio. Non dovrebbe andare oltre''\footnote{La trascrizione di questo passaggio è
molto confusa (``facciamo zero, no, mettere meno duecento\ldots{} facciamo duecento
massimo''), e non permette di ricostruire con certezza i valori finali del Clamp.
Quello che conta è il metodo: stampare i valori, osservarli, scegliere i limiti.}.

\paragraph{Approfondimento: Clamp e gli angoli.} Il Clamp restituisce il valore stesso
se è compreso tra Min e Max, altrimenti il limite più vicino:
$\operatorname{clamp}(v, m, M) = \min(\max(v, m), M)$. Con gli angoli c'è una
sottigliezza che spiega i valori ``strani'' della lezione: lo \emph{yaw} restituito da
Find Look at Rotation sta tra $-180$ e $180$, e passando ``dietro'' al mulino salta da
$180$ a $-180$; sommandoci 90 gradi di offset, l'intervallo si sposta (da $-90$ a
$270$). Un limite come $-25 \ldots 25$, applicato a questi numeri, non corrisponde a
``25 gradi a destra e a sinistra della direzione di riposo'', ma a una fetta del cerchio
che dipende da dove sta il player e da come è orientato il mulino nel livello. Il
modo pulito è ragionare sulla differenza tra l'angolo voluto e quello di riposo,
riportarla tra $-180$ e $180$ (Unreal ha il nodo \emph{Normalize Axis}, o \emph{Delta
(Rotator)} per due \emph{rotator}), e solo allora fare il Clamp. Il metodo del docente,
stampare i numeri prima di decidere, resta comunque la prima cosa da fare.

\paragraph{La pausa.} ``Ci siete? C'è qualche cosa che non avete capito? C'è qualche
blocco che vi risulta ancora un pochino nebuloso?'' Tutto chiaro. ``Sono le 17 e 28.
Facciamo pausa, così poi riprendiamo e cominciamo a fare la nostra bellissima
torrettina. Ci vediamo tra 20 minuti.''

% ---------------------------------------------------------------------
\section{Dopo la pausa: la torretta in tre pezzi}
\label{ue4:sec:torretta}

``Eccoci qua. Allora, ragazzi, adesso andremo a fare una bellissima torrettina.'' Il
docente mostra la versione finita. Attorno alla torretta c'è una \textbf{sfera} disegnata
con linee rosse: è il suo raggio d'azione. Quando il giocatore ci entra, la torretta
``fa lo stesso processo che abbiamo visto'', cioè ruota verso il \emph{character}, e
``in questo caso spara dei proiettili''. Non ci sono effetti speciali (``non abbiamo
Niagara''\footnote{La trascrizione dice ``BevoFX'': con ogni probabilità
\emph{Niagara VFX}, il sistema di effetti visivi di Unreal Engine, di cui Davide aveva
parlato nella lezione precedente.}), ma si costruisce ``la logica che ci servirà con la
torretta che avete \emph{mergiato}''. Si vedono anche delle linee che partono dal
cannone: sono un \textbf{line trace}, ``che permette a noi di avere una corrispondenza
dipendente da cosa incontra''. Oggi si fa la rotazione; proiettili e \emph{line trace}
saranno la prossima volta.

\paragraph{Il Blueprint.} ``La prima cosa che ovviamente dobbiamo fare è creare un
nuovo Blueprint, sempre \textbf{Actor}'', chiamato \textbf{BP\_Turret}. Ma come hanno
composto la torretta gli studenti, con Davide? ``Abbiamo preso i tre pezzi forniti e li
abbiamo uniti'' con \textbf{Merge Actors}: una \emph{mesh} unica. ``Allora facciamo una
cosa: la facciamo a \textbf{tre pezzi}.'' Per far girare il supporto e il cannone
indipendentemente dalla base servono tre componenti distinti, quindi si torna ai tre
pezzi importati con Davide: tre \textbf{Static Mesh} nel Blueprint, una per la
\textbf{base}, una per il supporto (``l'\emph{holder}, chiamiamolo'', poi rinominato
\textbf{Pivot}) e una per il \textbf{cannone} (``la \emph{gun}'', poi
\textbf{Cannon}).

\paragraph{Chi segue chi.} I tre pezzi vanno ``uno dentro l'altro''. Il docente parte
con l'idea che il supporto segua il cannone, poi si corregge: ``fate il contrario, che
è meglio: fate che la \emph{gun} segua l'\emph{holder}''. La gerarchia giusta è quindi:
la \textbf{base} come radice, il \textbf{Pivot} figlio della base, il \textbf{Cannon}
figlio del Pivot (figura~\ref{ue4:fig:turretbp}). Così, quando il Pivot ruota
orizzontalmente, il cannone gira con lui; e un giorno il cannone potrà alzarsi e
abbassarsi da solo, senza trascinarsi dietro il supporto.

\paragraph{Un trucco: Browse to Asset.} Per assegnare le \emph{mesh} ai componenti non
serve cercarle a mano nel Content Browser ``con l'elenchino e tutto quanto''. ``Andate
qua sull'\emph{asset} che avete in scena, cliccate tasto destro e fate \textbf{Browse to
Asset}'': il Content Browser salta direttamente alla posizione del file, e da lì lo si
può \textbf{trascinare} (\emph{drag and drop}) nello slot \emph{Static Mesh} del
componente. In alternativa, come sempre, si clicca sullo slot e si scrive il nome:
\emph{turret base}, \emph{turret pivot}, \emph{cannon}. Poi si aggiustano le posizioni
(``un buchettino sta un po' più giù'') finché i pezzi combaciano.

\paragraph{Un problema di Add.} Una studentessa, premendo Add, vede solo la \emph{mesh}
intera della torretta (\emph{SM\_\allowbreak Torretta}) e nessun'altra opzione. In condivisione si
prova a compilare e a riaprire il Blueprint, e alla fine il componente si riesce ad
aggiungere: ``va benissimo così, abbiamo risolto. A volte Unreal può fare dei giochi
simpatici.''

% ---------------------------------------------------------------------
\section{Il raggio d'azione: la Sphere Collision}
\label{ue4:sec:sfera}

``Adesso dobbiamo creare quello che abbiamo visto qua: vedi questo cerchio rosso intorno
alla torretta? Questo è una \textbf{Sphere Collision}.'' È il \textbf{range}: entrandoci
dentro, la torretta si attiva.

\paragraph{Aggiungerla.} In BP\_Turret si aggiunge, tra i componenti, una
\textbf{Sphere Collision}, ``che tranquillamente può stare anche esterna alla base'',
cioè non deve essere figlia di un pezzo che ruota. Appena aggiunta non si vede quasi,
``semplicemente perché è molto piccola''. Cliccando su \emph{Sphere}, nei dettagli c'è
lo \textbf{Sphere Radius}: di default è 32; il docente lo porta a \textbf{800}. C'è anche
la \textbf{Line Thickness}, lo spessore delle linee con cui la sfera è disegnata
nell'editor, ``che serve magari solo per vederla un po' più spessa''.

\paragraph{Hidden in Game.} In Play, però, la sfera non si vede. ``Di base, questa è una
\textbf{collisione}'', una cosa interna agli \emph{asset}, come le collisioni della
\emph{mesh} viste con Davide: esiste per la logica, non per essere guardata. Per
vederla mentre si prova il gioco (``per avere facilità nel vedere quando entriamo e
usciamo'') si cerca nei dettagli \textbf{Hidden in Game} e lo si toglie. ``Parlando di
\emph{rendering}, Hidden in Game dice: so che logicamente esiste all'interno del mondo,
ma se non è \emph{hidden in game}, la renderizzo pure.'' Tolta la spunta, il valore è
\emph{false}, ``quindi non è nascosto in gioco'', e in Play ecco il range. È la stessa
opzione usata nella seconda lezione, al contrario, per rendere invisibile il soffitto
della vaschetta.

\paragraph{Ripetere i passaggi.} Uno studente ha perso questa parte mentre cercava i
pezzi della torretta, e chiede di rivedere come si mettono le \emph{mesh} nei tre
componenti. Il docente lo rifà: si clicca sul componente, si scrive il nome della
\emph{mesh} nello slot Static Mesh (\emph{turret base}, \emph{turret pivot}), si
rinominano i componenti in modo che siano coerenti (\textbf{Pivot}, \textbf{Cannon},
``come a me è chiamata Cannon nella nostra \emph{mesh}'').

% ---------------------------------------------------------------------
\section{Gli eventi dei componenti e l'overlap}
\label{ue4:sec:overlap}

``Adesso andiamo a vedere un evento in particolare.'' Oltre agli eventi già visti, che
appartengono all'intero actor (l'Event Begin Play, l'Event Actor Begin Overlap,
l'Event Tick), ``possiamo richiamare degli eventi anche all'interno dei
\textbf{componenti}''. Selezionando \emph{Sphere} e scorrendo in fondo ai dettagli, c'è
una sezione \textbf{Events} piena di pulsanti: \emph{On Component Begin Overlap},
\emph{On Component End Overlap}, \emph{On Component Wake}, \emph{On Component Sleep},
gli eventi del cursore e altri. ``Tutta una serie di eventi'' che si possono usare
``per far sì che avvenga qualcosa in un preciso momento''. Qui la torretta deve partire
quando il player entra nella sfera: si clicca sul \textbf{+} di \textbf{On Component
Begin Overlap}, e nel grafo compare l'evento, già legato alla Sphere.

\paragraph{Che cos'è un overlap.} ``Sapete cosa si intende per \emph{overlap}?''
Uno studente: ``tipo appena si passa un confine, qualcosa del genere''. ``In un certo
senso.'' Alla lavagna (``perdonate le mie doti molto brutte di disegno'') il docente
disegna la sfera con la torretta al centro: un \emph{overlap} avviene quando il player,
``o un qualsiasi actor'', \textbf{tocca o entra} in quel volume. Un actor fuori non sta
sovrapponendosi; uno dentro sì. E uno a cavallo del bordo? ``Sta overlappando o non sta
overlappando?'' ``Sì.'' ``Qualcun altro ha un'opinione differente?'' Sta
sovrapponendosi, ``perché anche un qualsiasi poligono da cui è formato è parte
integrante di esso'': basta un pezzo (figura~\ref{ue4:fig:overlap}). Ci sono quindi tre
situazioni: non si sovrappone, \textbf{comincia} a sovrapporsi (Begin Overlap),
\textbf{smette} di sovrapporsi (End Overlap, ``quando ci esce''). Per ora serve solo la
prima: ``vogliamo che quando lui entra in questo volume la torretta faccia tipo: ti ho
visto, mo' comincio a ruotare verso di te''.

\begin{figure}[htbp]
\centering
\begin{tikzpicture}[font=\small,
  ue4arr/.style={-{Stealth[length=2.5mm]}, thick}]
  \draw[red!70!black, thick, fill=red!4] (0,0) circle (2.4);
  \node[red!70!black, font=\scriptsize] at (0,2.65) {Sphere Collision (raggio 800)};
  \draw[thick, fill=orange!35] (-0.35,-0.3) rectangle (0.35,0.3);
  \draw[thick, fill=orange!20] (0.35,-0.08) rectangle (0.9,0.08);
  \node[font=\scriptsize] at (0,-0.55) {torretta};
  % fuori
  \draw[thick, fill=blue!15] (-4.3,0.6) circle (0.32);
  \node[font=\scriptsize, align=center] at (-4.3,-0.2) {fuori:\\nessun overlap};
  % a cavallo
  \draw[thick, fill=blue!15] (1.75,1.65) circle (0.32);
  \node[font=\scriptsize, align=left, anchor=west] at (2.3,1.9) {a cavallo del bordo:\\
    \textbf{overlap} (basta un pezzo)};
  % dentro
  \draw[thick, fill=blue!15] (-1.0,-1.25) circle (0.32);
  \node[font=\scriptsize, align=center] at (-1.0,-1.95) {dentro: overlap};
  % frecce degli eventi
  \draw[ue4arr, green!45!black] (-3.9,0.45) -- (-2.55,0.25);
  \node[green!45!black, font=\scriptsize, align=center] at (-3.4,1.45)
    {entra:\\\textbf{Begin Overlap}};
  \draw[ue4arr, orange!80!black] (1.95,-1.55) -- (3.2,-2.3);
  \node[orange!80!black, font=\scriptsize, align=left, anchor=west] at (3.3,-2.3)
    {esce: \textbf{End Overlap}};
\end{tikzpicture}
\caption{Le situazioni dell'\emph{overlap} disegnate a lezione. Un actor si sovrappone
alla sfera appena una parte qualsiasi della sua forma di collisione ne entra; Begin
Overlap scatta una volta, nel momento in cui entra, End Overlap una volta, quando ne
esce del tutto. Schema ricostruito.}
\label{ue4:fig:overlap}
\end{figure}

\paragraph{Other Actor.} ``Come facciamo a capire quale actor ci dà indietro?'' L'evento
ha una serie di uscite: \emph{Overlapped Component}, \textbf{Other Actor}, \emph{Other
Comp}, \emph{Other Body Index}, \emph{From Sweep}, \emph{Sweep Result}. ``Sono tutti
input che sono risultanti dall'overlap'', ``in parole povere'' le informazioni su chi si
è sovrapposto a che cosa. Per mostrarlo, il docente collega Other Actor a un
\textbf{Print String} (Unreal inserisce da solo un \emph{Get Display Name}, per
trasformare l'actor in testo), va in Play ed entra nella sfera: in alto a sinistra
compare \emph{BP\_ThirdPersonCharacter\_0}\footnote{La trascrizione dice ``BP ha perso
un character 0''. Il suffisso \_0 è il numero che Unreal assegna a ogni istanza di una
classe presente nel livello.}. ``Mi ridà cosa sta overlappando, che in questo caso è il
\emph{character}.'' Esistono anche gli \emph{overlapped component} e gli \emph{other
component}, ma quello che interessa è l'Other Actor, ``perché il nostro
\emph{character} è un actor''.

\paragraph{Approfondimento: perché l'overlap funziona ``da solo''.} Perché Unreal generi
un evento di \emph{overlap} servono tre condizioni: entrambi i componenti devono avere
attivo \emph{Generate Overlap Events}; le loro risposte di collisione reciproche devono
essere \emph{Overlap} (e non \emph{Ignore}); e almeno uno dei due deve muoversi. La
Sphere Collision nasce con il profilo di collisione \emph{OverlapAllDynamic}, che si
sovrappone a tutto ciò che si muove, e la capsula del \emph{character} è un oggetto
dinamico di tipo \emph{Pawn}: per questo, nella lezione, non c'è stato bisogno di
toccare nulla. Se un giorno un evento di \emph{overlap} non scatta, è la prima cosa da
controllare, nei dettagli di entrambi i componenti, sezione \emph{Collision}. La seconda
edizione di Romero e Sewell elenca proprio gli eventi di collisione (``eseguiti quando
due actor collidono o si sovrappongono'') tra le famiglie principali di eventi, accanto
a quelli di input e all'Event Begin Play.

% ---------------------------------------------------------------------
\section{Is Valid: esiste davvero?}
\label{ue4:sec:isvalid}

``Vogliamo essere sicuri che il nostro actor in questione sia proprio il player. Perché
un actor può essere anche un cubo, volendo.'' Prima, però, un controllo più semplice. Si
tira fuori da Other Actor un \textbf{Is Valid}: ``ce ne avete tanti, ma quello che ci
interessa è quello col \textbf{punto interrogativo}'', la \emph{macro} con due uscite di
esecuzione, \textbf{Is Valid} e \textbf{Is Not Valid}.

\paragraph{Che cosa controlla.} ``Is Valid, statemi a sentire, ci dà una risultante: è
valido o non è valido?'' Qui chiede: ``hai overlappato qualcosa? È vero, esiste questa
cosa?'' In generale, Is Valid serve a dire se esiste o no qualcosa, ``che può essere
anche una \textbf{variabile}, quindi una variabile vuota''.

\paragraph{L'esempio della variabile.} Per farlo vedere il docente costruisce un piccolo
esperimento:
\begin{enumerate}
  \item da Other Actor crea una variabile con \textbf{Promote to Variable};
  \item mette un Is Valid su quella variabile, all'inizio dell'evento di \emph{overlap};
  \item sul ramo \textbf{Is Not Valid} assegna alla variabile l'actor appena
        sovrapposto (il \emph{Third Person Character});
  \item sul ramo \textbf{Is Valid} stampa con un Print String ``valore registrato'' e il
        nome contenuto nella variabile.
\end{enumerate}
Prima entrata nella sfera: non succede nulla di visibile, ``è normale''. Il Blueprint ha
chiesto: ``questa variabile ha qualcosa dentro? No, non c'è niente. Quindi non è
valida'', e allora le ha assegnato l'actor. Si esce e si rientra: ``ah, me l'hai fatto
prima, quindi questa cosa c'ha qualcosa dentro'', e il ramo è Is Valid: compare
\emph{valore registrato: BP\_ThirdPersonCharacter\_0}. ``Qualcuno ha qualche dubbio
riguardo all'Is Valid? Abbastanza chiaro?'' L'esperimento si cancella: ``ci serve
solamente a capire cosa fa l'Is Valid''.

\paragraph{Approfondimento: la regola di Romero e Sewell.} Il quarto capitolo di
\emph{Blueprints Visual Scripting for Unreal Engine 5}, dedicato alla comunicazione tra
Blueprint, mette in evidenza una nota: ``usate sempre la macro Is Valid prima di
eseguire una funzione usando una variabile che contiene una referenza a un oggetto. La
variabile può non essere valida per diversi motivi, ed eseguire una funzione con una
variabile non valida causerà un errore durante l'esecuzione''. I motivi tipici sono
due: la variabile non è mai stata riempita (come nella prima entrata dell'esempio), o
l'oggetto a cui puntava è stato distrutto nel frattempo. In entrambi i casi un nodo che
prova a leggere la posizione di ``nessuno'' produce nel registro il classico errore
\emph{Accessed None}. Il libro usa lo stesso schema in molti punti: un Is Valid prima di
distruggere un actor, prima di applicare un danno a un bersaglio, prima di leggere un
salvataggio (se non è valido, lo si crea). La torretta lo userà due volte, nella
prossima sezione e nell'evento Move Turret.

% ---------------------------------------------------------------------
\section{I tag e Actor Has Tag}
\label{ue4:sec:tag}

Is Valid dice che \emph{qualcosa} è entrato. Per sapere \emph{chi}, si tira fuori da
Other Actor un \textbf{Actor Has Tag}, ``un blocco che va a esaminare una cosa chiamata
\textbf{tag}''. I tag ``sono dei \emph{label}, delle \textbf{etichette} che noi possiamo
mettere all'interno di vari \emph{asset} del gioco''.

\paragraph{Mettere il tag al player.} Si apre il \textbf{Third Person Character} e,
``non dobbiamo fare niente di particolare'', si scrive \emph{tag} nella ricerca dei
dettagli: compare \textbf{Actor} $\rightarrow$ \textbf{Tags}, un elenco vuoto con un
\textbf{+}. Si aggiunge un elemento e lo si chiama \textbf{Player}. ``Cosa abbiamo
fatto? Abbiamo letteralmente etichettato il player con un tag Player'': è come avere un
cartellino di riconoscimento sulla schiena, con scritto sopra ``tu sei il player''.

\paragraph{Il controllo.} Tornati nella torretta, Actor Has Tag chiede: il
\emph{target}, cioè l'actor sovrapposto (Other Actor), ha il tag \textbf{Player}? Il
nodo restituisce un valore \textbf{booleano}: ``sì, questo actor che ho overlappato ha
il tag Player'', e si va nel \textbf{True} di un \textbf{Branch}; ``se no, invece andrà
in False''. Il docente lo mostra con due Print String, ``Player trovato'' sul ramo True e
``Player non trovato'' sul ramo False: entrando nella sfera, ``Player trovato''.

\paragraph{Attenzione all'ortografia.} ``Non lo fate, ragazzi, faccio solo vedere per
esempio.'' Il docente modifica il tag, togliendo la R finale, e rientra: ``Player non
trovato''. ``Questo vi fa vedere anche che dovete essere precisi nel dire qual è il
tag.'' Il consiglio: ``fate Ctrl+C, Ctrl+V del tag che avete scritto'', copiandolo dal
\emph{character} nel nodo invece di riscriverlo a mano.

\paragraph{Approfondimento: tag, Cast o confronto?} Ci sono almeno tre modi per
rispondere alla domanda ``l'actor che è entrato è il player?''. Il primo è il
\textbf{tag}, quello della lezione: semplice, e indipendente dalla classe (lo stesso tag
\emph{Player} può stare su personaggi diversi, e un tag \emph{Enemy} su nemici di classi
diverse); esiste anche \emph{Component Has Tag}, per etichettare singoli componenti. Il
secondo è il \textbf{Cast To} (per esempio \emph{Cast To BP\_ThirdPersonCharacter}),
che verifica la classe e, se riesce, restituisce l'actor già ``convertito'', con
accesso alle sue variabili e funzioni. Il terzo è il confronto diretto tra oggetti: nel
capitolo sull'intelligenza artificiale della seconda edizione di Romero e Sewell, un
nemico che pattuglia usa un \emph{Event Actor Begin Overlap} seguito da un Branch su un
nodo \textbf{Equal (Object)}, che controlla se l'Other Actor è proprio il punto di
pattuglia salvato in una variabile. Una curiosità sui tag: sono di tipo \emph{Name}, e
in Unreal il confronto tra \emph{Name} \textbf{non distingue maiuscole e minuscole},
per cui \emph{player} e \emph{Player} risultano uguali; quello che conta, come ha
mostrato il docente, è che le lettere siano le stesse. Copiare e incollare il tag
resta comunque la via più sicura.

% ---------------------------------------------------------------------
\section{La variabile del bersaglio e i Custom Event}
\label{ue4:sec:custom}

Ora che il player è stato trovato, bisogna ricordarselo: ``andiamo a crearci una
variabile in cui mettiamo il nostro player che è stato trovato'', di tipo
\textbf{Actor}. Il modo più veloce: tirare fuori un cavo da Other Actor e scegliere
\textbf{Promote to Variable}. Si crea un nodo \emph{Set} di una variabile nuova, che si
rinomina \textbf{PlayerTarget}. Per rinominare, ``o vai qua e fai \textbf{F2}, che è più
facile, oppure fai tasto destro e Rinomina''.

\paragraph{Un Custom Event.} ``Manca mezz'ora'', e il docente si chiede se si riuscirà a
finire (``questo lo continueremo la prossima volta''), ma va avanti. ``Adesso qua dentro
andremo a fare un \textbf{Custom Event}.'' Si scrive \emph{custom event} nel grafo, si
sceglie \emph{Add Custom Event} e gli si dà un nome: \textbf{TrackTargetTimer}. ``Come
potete vedere, questo bellissimo evento è un Custom Event: è uguale agli altri'', come
il Begin Overlap, il Tick o il Begin Play, ``ma è un evento che abbiamo creato noi'',
``che possiamo richiamare quando vogliamo, con delle regole specifiche''. Gli eventi di
Unreal sono ``prefissati'': l'\emph{overlap} scatta perché l'editor di Unreal lo
controlla già per conto suo. I Custom Event, invece, partono quando li si chiama.

\paragraph{Chiamarlo.} Per mostrarlo, il docente collega al Custom Event un Print String
che stampa ``Hello'' e lo chiama dal ramo del Begin Overlap, subito dopo il Set di
PlayerTarget: si scrive il nome dell'evento (\emph{TrackTargetTimer}) e compare il nodo
che lo invoca. In Play, entrando nella sfera: ``Hello''. ``Ho creato un evento
customizzato che richiamo quando entro nella Sphere Collision. Compreso?'' ``Sappiate
che di Custom Event ne potete creare a milioni'': basta scrivere \emph{custom event} e
dare un nome.

\paragraph{Approfondimento: eventi, funzioni e macro.} La seconda edizione di Romero e
Sewell ricorda che nell'Event Graph ogni evento di Unreal può comparire una sola volta,
mentre i \textbf{Custom Event} si possono creare a volontà; e offre una tabella per
scegliere tra \emph{macro}, \emph{funzioni} ed \emph{eventi}. In sintesi: tutti e tre
accettano parametri in ingresso; solo macro e funzioni restituiscono valori; le funzioni
e gli eventi possono essere chiamati da altri Blueprint, le macro no; e solo macro ed
eventi possono contenere azioni \emph{latenti}, cioè che aspettano (come un
\emph{Delay}), mentre le \emph{Timeline} stanno solo negli eventi. Il capitolo sui
vincoli di gioco della stessa edizione usa proprio un Custom Event, \emph{Sprint Drain},
richiamato da un timer, per consumare la stamina del giocatore mentre corre: ``un Custom
Event permette di attivare i nodi collegati a lui ogni volta che un altro nodo lo
chiama; così gruppi di nodi dello stesso Blueprint possono comunicare tra loro anche
quando non sono collegati direttamente''. È lo schema esatto della torretta.

% ---------------------------------------------------------------------
\section{Il timer: Set Timer by Event}
\label{ue4:sec:timer}

``Cosa vogliamo fare con questo Custom Event?'' Durante tutto il tempo, ``come un Event
Tick'', la torretta deve cercare la posizione verso cui ruotare per seguire il player.
``Lo potremmo fare tranquillamente nell'Event Tick, volendo, ma possiamo fare una cosa
interessante: possiamo fare un \textbf{timer}. L'abbiamo chiamato timer apposta.''

\paragraph{Il nodo.} Dal Custom Event si tira fuori \textbf{Set Timer by Event}, un
blocco ``che ci permette di far partire un timer che ci chiede ogni quanto''. I suoi
parametri, così come li imposta il docente:
\begin{itemize}
  \item \textbf{Time}: ogni quanti secondi scatta. ``Vogliamo che lo faccia allo stesso
        tempo del Tick.'' Si potrebbe portare fin qui il cavo dei \emph{delta seconds}
        dall'Event Tick (``ma è un po' brutto''); c'è invece un nodo apposta,
        \textbf{Get World Delta Seconds}, ``un blocchetto semplicemente per comodità'',
        ``per non far sì di filare il Tick così, a cavolo'';
  \item \textbf{Looping}: spuntato, ``vogliamo che sia un timer che \emph{loopi}'';
  \item \textbf{Max Once Per Frame}: spuntato anche questo, ``massimo una volta per
        ogni \emph{frame}''. ``Ci permette almeno di dire: ok, per ogni \emph{frame} si
        calcola, tanto effettivamente si muoverà per ogni \emph{frame}. Questa è una
        cosa un po' più nel micro, però noi mettiamo \emph{max}, in modo che sia più
        ottimizzato.''
\end{itemize}

\paragraph{Quale evento?} ``Set Timer by Event: quale evento?'' Sul nodo c'è un
\textbf{quadratino rosso}, il pin \textbf{Event}, ``che guarda caso è uguale a quello del
Custom Event'': un pin di tipo \emph{delegato}. Vuol dire che lo si può tirare fuori e
creare da lì un \textbf{altro Custom Event}, quello che il timer chiamerà a ogni giro:
lo si chiama \textbf{MoveTurret}. È lì che si farà muovere la torretta.

\paragraph{Come ragiona.} Per farlo vedere, a MoveTurret si collega il solito Print
String, ``Hello''. La catena completa è: Begin Overlap, Is Valid, Branch su Actor Has
Tag, Set PlayerTarget, TrackTargetTimer; e TrackTargetTimer fa partire un timer, alla
velocità dei \emph{delta seconds}, in \emph{loop}, al massimo una volta per
\emph{frame}, che chiama MoveTurret, che stampa ``Hello''. ``Domanda: quando mi
stamperà Hello?'' ``Da quando entri in range.'' ``Ottimo, esatto.'' Appena il player
entra, Hello a ogni \emph{frame}. ``Esco fuori: si cancellano\ldots{} ah, no, scusate,
non si stanno cancellando, è vero. Perché non stoppiamo mai il timer.'' Dalla prima
entrata in poi, il timer continua all'infinito: ``ciò vuol dire che il timer a un certo
punto va disattivato; lo vedremo dopo''.

\paragraph{Approfondimento: timer o Event Tick?} Romero e Sewell presentano i timer, già
nel secondo capitolo, come un'alternativa all'Event Tick: tutto ciò che è collegato al
Tick viene eseguito a ogni \emph{frame}, e ciò che non ha bisogno di aggiornarsi così
spesso può girare a intervalli più lunghi. Nel quinto capitolo un Custom Event
\emph{ChangeDirection}, chiamato da un Set Timer by Event con Looping attivo, fa
invertire periodicamente la direzione di un bersaglio mobile; nella seconda edizione un
timer chiama ripetutamente \emph{Sprint Drain}, e un \emph{Clear Timer by Function Name}
lo ferma quando il giocatore smette di correre. Nella torretta della lezione l'intervallo
è pari ai \emph{delta seconds}, quindi il timer scatta (circa) una volta per
\emph{frame}, come un Tick: il vantaggio non è la frequenza, ma il fatto che il timer
\textbf{parte solo quando serve} (all'ingresso del player) e \textbf{si può fermare}.
Tre dettagli utili. Primo: il valore di Get World Delta Seconds viene letto una sola
volta, quando il timer viene impostato, e resta fisso anche se poi il \emph{frame rate}
cambia. Secondo: se l'intervallo di un timer è più breve di un \emph{frame}, Unreal può
eseguirlo più volte nello stesso \emph{frame} per ``recuperare''; \textbf{Max Once Per
Frame}, un'opzione delle versioni recenti di Unreal Engine 5, impedisce proprio questo.
Terzo: per fermarlo, Set Timer by Event restituisce un \textbf{Timer Handle}, una
``maniglia'' che conviene salvare in una variabile (Promote to Variable) e passare,
quando serve, a \textbf{Clear and Invalidate Timer by Handle}. Il posto naturale per
farlo è l'evento \textbf{On Component End Overlap} della sfera, insieme a
\emph{svuotare} la variabile PlayerTarget (un Set senza nulla collegato): altrimenti
l'Is Valid di MoveTurret continuerà a trovare un bersaglio valido anche dopo che il
player è uscito dal range.

% ---------------------------------------------------------------------
\section{Move Turret: la torretta che segue il player}
\label{ue4:sec:move}

``Cosa dobbiamo fare? L'abbiamo fatto prima, praticamente'': un \textbf{Find Look at
Rotation}, ``perché dobbiamo ruotare il nostro Pivot in direzione del player''.

\paragraph{Prima, un Is Valid.} ``Intanto facciamo un Is Valid. Perché? Perché diciamo:
ma ce l'ho un target?'' Si trascina nel grafo la variabile \textbf{PlayerTarget} e la si
collega a un Is Valid, in testa a MoveTurret. ``Ce l'ho un target effettivo, o non ce
l'ho più, l'ho perso, sono uscito fuori? Se ce l'ho, lo faccio e mi muovo. Se non ce
l'ho, non lo faccio'': senza questo controllo la torretta ``continuerebbe a ruotare a
cavolo''.

\paragraph{Start e Target.} ``Del Find Look at Rotation ci servono due punti nello
spazio. Quali sono?'' ``Lo Start è l'\emph{holder}, che tanto la \emph{gun} è collegata
all'\emph{holder}, e il Target è il player.'' ``Esatto.'' Per lo Start si usa
\textbf{Get Actor Location} della torretta stessa, ``la locazione generale
dell'actor''. Per il Target, il docente comincia con Get Player Character e Get Actor
Location, come per il mulino, ma uno studente fa notare che si può usare
\textbf{PlayerTarget}: ``bravo, mi piace che l'hai notato, effettivamente sarebbe stato
il secondo metodo''. Da PlayerTarget si tira fuori Get Actor Location, ``una sua
particolarità'', e la si collega al Target.

\paragraph{Lo yaw e RInterp To.} Il Pivot deve ruotare attorno alla \textbf{Z}:
``guardate un po' dove si modifica: sulla Z''. Si splitta l'uscita di Find Look at
Rotation e si fa ``quello che abbiamo fatto prima'', un \textbf{RInterp To}:
\begin{itemize}
  \item \textbf{Current}: la rotazione attuale del Pivot, con \textbf{Get Relative
        Rotation} sul componente Pivot;
  \item \textbf{Target}: splittato, con la sola Z collegata, ``perché non ci interessa
        la X e la Y'';
  \item \textbf{Delta Time}: 5;
  \item \textbf{Interp Speed}: \textbf{0,04}, ``perché, essendo che sta sui
        \emph{frame}, è molto veloce. Se voi mettete uno, questa \emph{snappa}
        diretto''.
\end{itemize}
Infine \textbf{Set Relative Rotation} sul Pivot (``prendete Pivot, Set Relative
Rotation''), con i pin splittati ``per sicurezza'', collegato al ramo Is Valid; e un
Print String per il debug (figura~\ref{ue4:fig:turretbp}).

\begin{figure}[htbp]
\centering
\begin{tikzpicture}[font=\footnotesize,
  ue4node/.style={draw, thick, rounded corners=4pt, align=center, fill=gray!8,
                  minimum height=0.85cm, text width=2.3cm},
  ue4event/.style={ue4node, fill=red!12},
  ue4exec/.style={-{Stealth[length=2.5mm]}, very thick, black!80},
  ue4data/.style={-{Stealth[length=2mm]}, thick, blue!60!black},
  ue4flt/.style={-{Stealth[length=2mm]}, thick, green!50!black},
  ue4dlg/.style={-{Stealth[length=2mm]}, thick, red!70!black, dashed},
  ue4comp/.style={draw, thick, rounded corners=3pt, minimum width=2.3cm,
                  minimum height=0.55cm, align=center, font=\scriptsize}]
  % --- riga 1: overlap
  \node[ue4event] (ov) at (0,0)
    {\textbf{On Component Begin Overlap}\\(Sphere)};
  \node[ue4node] (iv1) at (2.95,0) {\textbf{Is Valid}\\Other Actor};
  \node[ue4node] (br) at (5.9,0) {\textbf{Branch}\\Actor Has Tag ``Player''};
  \node[ue4node] (set) at (8.85,0) {\textbf{Set}\\\textbf{PlayerTarget}};
  \node[ue4node, text width=2.8cm] (call) at (12.05,0) {\textbf{TrackTargetTimer}\\(chiamata)};
  \draw[ue4exec] (ov) -- (iv1);
  \draw[ue4exec] (iv1) -- (br);
  \draw[ue4exec] (br) -- node[above, font=\scriptsize]{True} (set);
  \draw[ue4exec] (set) -- (call);
  \draw[ue4data] (ov.south) -- (0,-1.15) -| node[pos=0.2, below,
    font=\scriptsize]{Other Actor} (set.south);
  % --- riga 2: timer
  \node[ue4event, text width=2.8cm] (cev) at (0,-2.9)
    {\textbf{Custom Event}\\\textbf{TrackTargetTimer}};
  \node[ue4node, text width=3.0cm] (timer) at (3.65,-2.9)
    {\textbf{Set Timer by Event}\\Time: Get World Delta Seconds\\Looping \checkmark\\
    Max Once Per Frame \checkmark};
  \node[ue4event] (mv) at (7.1,-2.9) {\textbf{Custom Event}\\\textbf{MoveTurret}};
  \draw[ue4exec] (cev) -- (timer);
  \draw[ue4dlg] (timer.east) -- node[above, font=\scriptsize]{Event} (mv.west);
  % --- riga 3: move turret
  \node[ue4node] (iv2) at (0,-5.7) {\textbf{Is Valid}\\PlayerTarget};
  \node[ue4node] (look) at (3.1,-5.7) {\textbf{Find Look at Rotation}\\
    Start: self\\Target: PlayerTarget};
  \node[ue4node] (interp) at (6.3,-5.7) {\textbf{RInterp To}\\
    Current: Pivot\\Target: solo Z\\Speed 0,04--0,05};
  \node[ue4node] (srr) at (9.5,-5.7) {\textbf{Set Relative Rotation}\\Target: Pivot};
  \draw[ue4exec] (mv.south) -- (7.1,-4.35) -| (iv2.north);
  \draw[ue4exec] (iv2.south) -- (0,-7.1) -| node[pos=0.2, below,
    font=\scriptsize]{Is Valid} (srr.south);
  \draw[ue4flt] (look) -- node[above, font=\scriptsize]{Z} (interp);
  \draw[ue4flt] (interp) -- (srr);
  % --- gerarchia componenti
  \node[font=\scriptsize] at (11.6,-1.95) {i componenti};
  \node[ue4comp, fill=orange!25] (cb) at (11.4,-2.5) {Base (radice)};
  \node[ue4comp, fill=gray!20] (cp) at (11.9,-3.2) {Pivot};
  \node[ue4comp, fill=orange!12] (cc) at (12.4,-3.9) {Cannon};
  \node[ue4comp, fill=red!8] (cs) at (11.9,-4.6) {Sphere (r. 800)};
  \draw[thick, gray!60!black] ([xshift=-0.85cm]cb.south) |- (cp.west);
  \draw[thick, gray!60!black] ([xshift=-0.85cm]cp.south) |- (cc.west);
  \draw[thick, gray!60!black] ([xshift=-0.85cm]cb.south) |- (cs.west);
\end{tikzpicture}
\caption{La torretta della lezione. In alto: quando qualcosa entra nella sfera, se
esiste ed è etichettato \emph{Player}, lo si salva in PlayerTarget e si chiama
TrackTargetTimer. In mezzo: TrackTargetTimer fa partire un timer che, a ogni
\emph{frame}, chiama MoveTurret (il pin rosso tratteggiato è il delegato Event). In
basso: MoveTurret, se c'è ancora un bersaglio, calcola la rotazione verso di lui, ne
tiene lo \emph{yaw}, la interpola e la applica al Pivot (la Current dell'interpolazione
viene da Get Relative Rotation del Pivot; l'eventuale offset è discusso nel testo). A
destra la gerarchia dei componenti: il Cannon segue il Pivot, la sfera non ruota. Schema
ricostruito.}
\label{ue4:fig:turretbp}
\end{figure}

\paragraph{L'offset, e il pivot reimportato.} In Play la torretta ruota, ma ``ha
bisogno di un offset, come abbiamo visto con il nostro mulino''. Il docente prova con
$-90$, poi con $+90$: ``eccolo qua. Non è esattamente precisissimo''. Uno studente
osserva: ``noi questa cosa non la dobbiamo fare, perché l'altra volta abbiamo fatto un
\textbf{reimport} modificando i valori''. Ed è la chiave: senza alcun offset, la
torretta del docente ``guarda direttamente in faccia'', mentre quelle degli studenti no,
perché con Davide avevano reimportato il \textbf{Pivot} e il \textbf{Cannon} con una
\textbf{rotazione di 113} gradi sulla Z (la base no). Il docente ricostruisce i numeri
con la classe (si parla anche di $-67$ e di $-271$, ``proprio perfetti''), reimporta
prima il pivot, poi il cannone, prova con il solo pivot, e alla fine: ``ecco a voi.
Adesso la nostra torretta ci segue perfettamente''. Per gli studenti la regola pratica,
ripetuta più volte nel finale, è: se la torretta ha ``un \emph{offsettino}'',
\textbf{reimportare il Pivot mettendo $-113$} sulla rotazione Z\footnote{La trascrizione
di questo passaggio è tra le più confuse della lezione: i valori citati sono 113, $-113$,
$-67$, 118, $-270$ e $-271$, e non è chiaro in che ordine siano stati provati. Il senso
è comunque chiaro: la rotazione di 113 gradi aggiunta in import nella terza lezione
sposta la direzione ``in avanti'' della \emph{mesh}, e Find Look at Rotation, che
presuppone che l'oggetto guardi lungo la sua X, ne risulta sfasato. Le due correzioni
possibili sono equivalenti: annullare la rotazione in import, o sommare un offset fisso
allo \emph{yaw} nel grafo, come per il mulino.}.

\paragraph{Approfondimento: Find Look at Rotation e lo spazio relativo.} Il grafo della
lezione mescola due spazi: Find Look at Rotation lavora nello spazio del
\textbf{mondo} (le posizioni di torretta e player sono posizioni nel mondo), mentre
Set Relative Rotation imposta la rotazione del Pivot \textbf{relativa} alla base. Le
due cose coincidono finché la torretta nel livello non è ruotata: se la si piazza
girata di 30 gradi, il Pivot punterà 30 gradi a lato del player. Le soluzioni robuste
sono due: usare \textbf{Set World Rotation} sul Pivot (la rotazione calcolata è già
quella del mondo), oppure sottrarre allo \emph{yaw} calcolato quello dell'actor (Get
Actor Rotation). Anche la Current dell'interpolazione andrebbe presa nello stesso
spazio del Target: Get World Rotation con Set World Rotation, Get Relative Rotation con
Set Relative Rotation. Lo Start, infine, è la posizione dell'actor, cioè della base; per
una torretta alta conviene usare la posizione del Pivot o del cannone (\emph{Get World
Location} sul componente), specie quando si vorrà calcolare anche il \emph{pitch} per
alzare la canna.

% ---------------------------------------------------------------------
\section{Il tuning e i problemi degli studenti}
\label{ue4:sec:problemi}

\paragraph{La velocità di inseguimento.} ``Giratele intorno.'' Per aumentare o diminuire
la velocità con cui la torretta ``ci mette a cercare il target'' si cambia l'\textbf{Interp
Speed}. A 0,07 è ``molto più reattiva''; a 2 ``sarà proprio quasi uno \emph{snap}''
(``pam, \emph{snappa} proprio''). Il consiglio: ``mettetela tipo 0,05''. E, volendo, la
si può trasformare in variabile (Promote to Variable) ``e la potete settare
esternamente'', nei dettagli dell'istanza in scena, come le variabili \emph{Instance
Editable} della seconda lezione.

\paragraph{``A me snappa lo stesso.''} Uno studente ha Delta Time 5 e Interp Speed 0,05,
e la torretta scatta comunque. Si prova 0,02 (``vedo leggerissima differenza''), poi
Delta Time 7: ``stampa sempre la stessa velocità''. In condivisione il docente controlla
Get World Delta Seconds, Looping e Max Once Per Frame: tutto giusto. ``Prova ad andare di
\emph{tweak}, di \emph{tuning}: vedi un po' tu i valori.'' Abbassando ancora l'Interp
Speed, verso lo zero, la torretta rallenta: ``ora sì, già un po' più così''. (Nel
frattempo lo studente che lavora deve andare: ``ciao, ciao''.)

\paragraph{Approfondimento: perché ``snappa''.} La formula di RInterp To vista per il
mulino spiega tutto. Con Delta Time fisso a 5 e Interp Speed 0,04, la frazione di
strada percorsa a ogni chiamata è $5 \times 0{,}04 = 0{,}2$: il 20\% di quel che manca.
Dopo dieci chiamate manca $0{,}8^{10} \approx 11\%$ dell'angolo iniziale; a 100
\emph{frame} al secondo, sono un decimo di secondo. Con 0,05 la frazione è 0,25; con 7
e 0,02 è 0,14; con Interp Speed 2 il prodotto è 10, cioè più di 1, e la torretta arriva
in un solo \emph{frame}: lo \emph{snap}. In tutti questi casi, inoltre, la velocità
dipende dai \emph{frame}: con il doppio dei \emph{frame} la torretta è due volte più
rapida. L'uso previsto del nodo è il contrario di quello della lezione: \textbf{Delta
Time} collegato a \textbf{Get World Delta Seconds} (circa 0,01 a 100 \emph{fps}) e
\textbf{Interp Speed} nell'ordine delle unità (da 2 a 5 per un inseguimento morbido),
così che il prodotto resti piccolo e la velocità sia la stessa su ogni macchina. È lo
stesso principio dei \emph{delta seconds} della seconda lezione, applicato a
un'interpolazione. E un'avvertenza: un Interp Speed esattamente 0 non ferma la torretta,
ma la fa scattare subito sul bersaglio, perché Unreal lo interpreta come ``nessuna
interpolazione''.

\paragraph{La lista degli errori.} Il resto della lezione è una carrellata di
condivisioni dello schermo, molto istruttiva perché mostra gli errori più frequenti. La
tabella~\ref{ue4:tab:errori} li raccoglie. Il caso più ricco è quello di uno studente
a cui ``non va perché probabilmente ho mancato delle parti'':
\begin{itemize}
  \item per la Current di RInterp To aveva usato un \emph{Set} invece di \textbf{Get
        Relative Rotation} (``no, no, no, non \emph{set}: \emph{get}'');
  \item l'Is Valid di MoveTurret non era collegato alla variabile \textbf{PlayerTarget}
        (``a sinistra, proprio sulle variabili, non sul grafo'');
  \item MoveTurret non era un evento separato. Per riorganizzare il grafo il docente
        gli fa selezionare tutto tranne TrackTargetTimer, Set Timer by Event e Get
        World Delta Seconds (``tieni premuto Ctrl e clicca su quelli che vuoi
        escludere''), spostare in basso la selezione, tirare fuori dal pin Event del
        timer un nuovo Custom Event chiamato \textbf{MoveTurret} e collegarlo all'Is
        Valid;
  \item il timer non era in \textbf{Looping};
  \item il \textbf{Target} di Find Look at Rotation era scollegato: da PlayerTarget,
        Get Actor Location, al Target;
  \item il \emph{character} non aveva il \textbf{tag}: lo si aggiunge nel Third Person
        Character, ``\emph{Player} esattamente come l'hai scritto, quindi con la P
        maiuscola'';
  \item la sfera aveva ancora Hidden in Game attivo e raggio \textbf{32}: ``devi
        metterla a 800'';
  \item in scena c'era solo la \textbf{mesh} della torretta, non il
        \textbf{Blueprint}: ``questa è solo la \emph{mesh}, è solo il modello 3D che hai
        messo in scena''; si cancella e si trascina nel livello BP\_Turret;
  \item il player partiva già dentro il range, e quindi il Begin Overlap non scattava:
        ``il player parte da fuori, mettilo fuori'';
  \item in Play \textbf{girava la base}: nella gerarchia, la base era figlia del Pivot.
        ``Devi staccare la base'': trascinando il Pivot dentro la base, ``perfetto, ce
        l'hai perfettamente'';
  \item infine, l'\emph{offset} del pivot reimportato: ``devi reimportare il pivot a
        $-113$''.
\end{itemize}
Un altro studente ha la torretta ferma: ``Other Actor devi attaccarlo a PlayerTarget,
perché così non gli stai impostando il \emph{target} che è overlappato: non avrà un
\emph{target} da seguire''. Un terzo non vede la sfera in gioco, solo in modalità
\emph{spectator}: aveva saltato il passaggio di Hidden in Game.

\begin{table}[htbp]
\centering
\small
\begin{tabular}{@{}p{5.4cm}p{8.6cm}@{}}
\toprule
\textbf{Sintomo} & \textbf{Causa e rimedio} \\
\midrule
La pala gira attorno a un punto sbagliato & la \emph{mesh} ha il pivot fuori
centro: scegliere un cubo con il pivot al centro (o correggere il pivot) \\
La pala non gira & il risultato di Combine Rotators non è collegato a Set
Relative Rotation; oppure il componente non è Movable \\
Il mulino o la torretta restano fermi & mancano Delta Time e Interp Speed in
RInterp To; oppure non si è compilato \\
Non si riesce a collegare un valore a una sola componente & fare Split Struct Pin
(tasto destro sul pin) \\
Gira tutto l'actor, non solo una parte & il Rotating Movement fa girare la radice:
programmare la rotazione del componente \\
La torretta non si accorge del player & manca il tag \emph{Player} (o è scritto
diversamente); Other Actor non collegato al Set di PlayerTarget; il player parte già
dentro la sfera \\
La sfera non si vede in Play & Hidden in Game ancora attivo; raggio di default 32 \\
Non succede niente in Play & in scena c'è la \emph{mesh}, non il Blueprint \\
Gira la base & gerarchia sbagliata: la base deve essere la radice, il Pivot suo
figlio, il Cannon figlio del Pivot \\
La torretta guarda di lato & la \emph{mesh} non guarda lungo X (Pivot reimportato
a 113): reimportare a $-113$ o sommare un offset allo \emph{yaw} \\
La torretta scatta (\emph{snap}) & prodotto Delta Time $\times$ Interp Speed troppo
grande: usare Get World Delta Seconds e un Interp Speed di qualche unità \\
La torretta continua a seguire dopo l'uscita & il timer non viene mai fermato e
PlayerTarget resta pieno: End Overlap, Clear and Invalidate Timer, svuotare la
variabile \\
\bottomrule
\end{tabular}
\caption{I problemi incontrati durante la lezione, con la causa e il rimedio. Le ultime
due righe riguardano punti che il docente ha rimandato alla lezione successiva.}
\label{ue4:tab:errori}
\end{table}

% ---------------------------------------------------------------------
\section{La chiusura e la prossima lezione}
\label{ue4:sec:chiusura}

``Ok, perfetto, ragazzi.'' Per questa volta, ``nessun particolare compito da casa o
dimostrazione''. Chi vuole può fare una \textbf{torretta di colore diverso}, con un
materiale suo, verde o blu, ripassando quello che si è fatto con Davide.

\paragraph{La prossima volta.} ``La prossima volta faremo la nostra bellissima
\textbf{torretta che spara}'': \textbf{spawnerà} dei \textbf{proiettili}. Si vedrà
``come funziona il sistema di \emph{spawning} di Unreal'': come si basa una cosa
\emph{spawnata} ``con le sue regole'', come si fa apparire un oggetto in generale, e
come si fa il \textbf{line trace}. ``Altre domande? Se no, andiamo a cena tutti,
ragazzi.'' ``Vi auguro buona serata. E a martedì.''

\begin{table}[htbp]
\centering
\small
\begin{tabular}{@{}p{4.0cm}p{10.0cm}@{}}
\toprule
\textbf{Nodo o componente} & \textbf{Che cosa fa} \\
\midrule
Rotating Movement Component & fa ruotare la radice dell'actor a velocità costante
(Rotation Rate, default Z 180) \\
Event Tick, Delta Seconds & esegue a ogni \emph{frame}; i \emph{delta seconds} sono il
tempo tra due \emph{frame} \\
Get World Delta Seconds & gli stessi \emph{delta seconds}, ovunque nel grafo \\
Make Rotator & costruisce un \emph{rotator} da tre \emph{float} (Roll X, Pitch Y,
Yaw Z) \\
Combine Rotators & combina due rotazioni in una \\
Get / Set Relative Rotation & legge e scrive la rotazione di un componente rispetto al
genitore \\
Get Actor Location & la posizione dell'actor nel mondo \\
Get Player Character & una referenza al personaggio del giocatore \\
Find Look at Rotation & la rotazione per guardare da Start verso Target \\
RInterp To & avvicina gradualmente un \emph{rotator} a un altro \\
Clamp (float) & limita un valore tra un minimo e un massimo \\
Split Struct Pin & divide un pin di struttura (vettore, \emph{rotator}) nelle sue
componenti \\
Sphere Collision & un volume sferico per rilevare \emph{overlap} (Sphere Radius,
Hidden in Game) \\
On Component Begin Overlap & scatta quando qualcosa entra nel volume del componente;
Other Actor è chi è entrato \\
Is Valid (macro con il ``?'') & controlla che una referenza punti a qualcosa che
esiste \\
Actor Has Tag & controlla se un actor porta un'etichetta (Actor $\rightarrow$ Tags) \\
Promote to Variable & crea una variabile del tipo giusto da un pin \\
Custom Event & un evento creato da noi, che parte quando lo si chiama \\
Set Timer by Event & chiama un evento dopo un tempo, una volta o in \emph{loop} (Max
Once Per Frame) \\
Print String & stampa a schermo: lo strumento base del debug \\
\bottomrule
\end{tabular}
\caption{I nodi e i componenti della quarta lezione.}
\label{ue4:tab:nodi}
\end{table}

% ---------------------------------------------------------------------
\section*{Punti chiave del capitolo}
\addcontentsline{toc}{section}{Punti chiave del capitolo}

\begin{itemize}
  \item \textbf{Quixel Megascans} è un ``calderone'' gratuito di materiali, texture e
        \emph{mesh} scansionati, da usare con l'account Epic. Per la pallina che
        attraversa il piano: il \textbf{CCD} (\emph{Continuous Collision Detection}).
  \item Un \textbf{vettore} ha direzione, magnitudine e verso; in Unreal la struttura
        Vector (X, Y, Z) rappresenta sia punti sia spostamenti. Somma: si spostano i
        punti; \textbf{sottrazione}: B $-$ A va da A a B; \textbf{moltiplicazione}:
        allunga; \textbf{prodotto scalare}: un \emph{float}, 1 se concordi, 0 se
        perpendicolari, $-1$ se opposti (per vettori normalizzati).
  \item Un vettore \textbf{normalizzato} ha lunghezza 1: tiene solo la direzione.
  \item Un \textbf{rotator} ha tre angoli: \textbf{Roll} (X), \textbf{Pitch} (Y),
        \textbf{Yaw} (Z). Gli angoli di Eulero soffrono di \textbf{gimbal lock}; i
        \textbf{quaternioni}, usati internamente da Unreal, no.
  \item Il \textbf{Rotating Movement Component} fa girare un actor senza programmare,
        ma solo per intero. Ogni oggetto che si muove deve essere \textbf{Movable}.
  \item Per far girare un componente: \textbf{Event Tick} $\rightarrow$ velocità
        $\times$ \emph{delta seconds} $\rightarrow$ \textbf{Make Rotator}
        $\rightarrow$ \textbf{Combine Rotators} con \textbf{Get Relative Rotation}
        $\rightarrow$ \textbf{Set Relative Rotation}. Le rotazioni dei componenti si
        lavorano in spazio \textbf{relativo}.
  \item La \textbf{gerarchia} fa il resto: un componente figlio segue il genitore,
        mantenendo la propria posizione relativa (la seconda pala, il cubetto).
  \item \textbf{Find Look at Rotation} calcola la rotazione per guardare un punto;
        dello \emph{yaw} si tiene la Z (\textbf{Split Struct Pin}), con un
        \textbf{offset} se la \emph{mesh} non guarda lungo X. \textbf{RInterp To}
        rende il movimento graduale: Delta Time dai \emph{delta seconds}, Interp Speed
        di qualche unità, altrimenti ``snappa''.
  \item Il \textbf{Clamp} limita un valore; prima di sceglierne i limiti, si
        \textbf{stampano} i valori con Print String.
  \item La torretta è un Blueprint in \textbf{tre pezzi}: Base (radice), Pivot
        (figlio, ruota in \emph{yaw}), Cannon (figlio del Pivot). Browse to Asset aiuta
        a trovare le \emph{mesh}.
  \item Il range è una \textbf{Sphere Collision} (raggio 800, Hidden in Game tolto per
        vederla). \textbf{On Component Begin Overlap} scatta quando qualcosa ci entra;
        \textbf{Other Actor} dice chi.
  \item \textbf{Is Valid} prima di usare una referenza; \textbf{Actor Has Tag} con il tag
        \emph{Player} (scritto identico) per riconoscere il giocatore;
        \textbf{Promote to Variable} per ricordarlo in \emph{PlayerTarget}.
  \item I \textbf{Custom Event} sono eventi nostri, chiamati quando vogliamo;
        \textbf{Set Timer by Event} (Looping, Max Once Per Frame) chiama a ripetizione
        l'evento \textbf{MoveTurret}, che ruota il Pivot verso il bersaglio. Il timer
        va fermato (End Overlap, Clear and Invalidate Timer): sarà il prossimo passo.
  \item Prossima lezione: la torretta che \textbf{spara}, lo \textbf{spawn} dei
        proiettili e il \textbf{line trace}.
\end{itemize}

%
%  Tutte le etichette di questo capitolo hanno il prefisso "ue4:"
%  (ue4:cap:..., ue4:sec:..., ue4:fig:..., ue4:tab:...), e tutti gli
%  stili TikZ hanno il prefisso "ue4" (ue4node, ue4exec, ue4data, ...)
%  per non collidere con chiavi native di TikZ (come "step") né con
%  quelli delle lezioni precedenti. Il capitolo è autonomo: non contiene
%  \ref verso altri capitoli (i rimandi alle lezioni precedenti sono a
%  parole).
%  Le figure sono ricostruzioni schematiche di quanto mostrato e detto a
%  lezione, non copie delle slide o delle schermate del docente. Dove la
%  trascrizione è troppo rovinata per ricostruire un dettaglio, il
%  capitolo lo dice invece di inventarlo.
% =====================================================================

\chapter[Vettori, rotazioni e la torretta]{Vettori, rotazioni e la torretta: il
mulino che gira e la torretta che ti segue}
\label{ue4:cap:torretta}

Questa è la \textbf{quarta lezione} del corso breve di \textbf{introduzione a Unreal
Engine}, e segna il ritorno del docente titolare dopo la lezione di Davide sulla
grafica. Se la terza lezione guardava a ciò che ``fa il grafico all'interno di Unreal''
(texture, materiali, l'import della torretta), questa torna alla \textbf{logica}, cioè
ai \textbf{Blueprint}, e lo fa con un obiettivo preciso, mostrato in apertura: un
\textbf{mulino} le cui pale girano, e la \textbf{torretta} importata con Davide che
``adesso ci segue con lo sguardo''. Per arrivarci serve un po' di matematica dello
spazio 3D: i \textbf{vettori}, le loro operazioni, i \textbf{rotator}.

La lezione ha due metà separate da una pausa di venti minuti, presa alle 17:28. La
prima metà parte dalla \textbf{teoria}: che cos'è un \textbf{vettore} (``una freccia
con una direzione, una magnitudine e un verso''), le quattro operazioni
(\textbf{addizione}, \textbf{sottrazione}, \textbf{moltiplicazione} e \textbf{prodotto
scalare}, il \emph{dot product}), i \textbf{vettori normalizzati}, i
\textbf{rotator} e il problema del \textbf{gimbal lock}. Poi si passa all'engine con un
mulino: prima il modo più semplice di far ruotare un oggetto, il \textbf{Rotating
Movement Component}, poi un Blueprint vero, \textbf{BP\_Windmill}, in cui una pala
gira grazie all'\textbf{Event Tick}, ai \emph{delta seconds} e ai nodi \textbf{Make
Rotator}, \textbf{Combine Rotators} e \textbf{Set Relative Rotation}; la seconda pala
gira ``gratis'' grazie alla \textbf{gerarchia} dei componenti. Infine il mulino
diventa ``magico'' e si gira verso il giocatore con \textbf{Find Look at Rotation} e
\textbf{RInterp To}, con un \textbf{Clamp} per limitarne la rotazione. La seconda metà
costruisce la \textbf{torretta}: un Blueprint in tre pezzi (\textbf{Base},
\textbf{Pivot}, \textbf{Cannon}), una \textbf{Sphere Collision} come raggio d'azione,
l'evento \textbf{On Component Begin Overlap}, i controlli \textbf{Is Valid} e
\textbf{Actor Has Tag}, una variabile per ricordare il bersaglio, un \textbf{Custom
Event} e un \textbf{timer} che, a ogni \emph{frame}, fa ruotare il pivot verso il
player.

La lezione è \textbf{online}, su Zoom, come le precedenti. Il docente alterna
spiegazione e piccole esercitazioni ``botta e risposta'', da cinque minuti l'una, e
passa gran parte della seconda metà a risolvere i problemi degli studenti in
condivisione dello schermo. Finisce verso le 19, ``andiamo a cena tutti'', con
l'appuntamento a \textbf{martedì}.

Tre avvertenze di lettura, le stesse delle lezioni precedenti. La prima: la
trascrizione è automatica e storpia moltissimi nomi. Unreal è ``a Real'', ``Arreal'',
``Arry''; l'\textbf{Event Tick} è ``Leventic'', ``l'eventic'', ``eventHIC'' e
``l'event kick''; i \textbf{vettori} sono una volta ``vittori''; i \emph{rotator} sono
``rotatempi'' e ``rotatori''; \textbf{Find Look at Rotation} è ``findLookAtRotation'',
``find all the rotation'' e ``Find, look, rotation''; \textbf{RInterp To} è ``Er inter
tu'', ``R Inter2'' e ``r interp tu''; \textbf{Get Player Character} è
``getLayerCharacter''; \textbf{Get Actor Location} è ``Get. Hector. Location'' e ``get
active location''; il \textbf{Rotating Movement} è a un certo punto ``vetro in
movement''; \textbf{Is Valid} è ``Svalid'' e ``isvalid''; \textbf{Other Actor} è
``OderActor'' e ``Oder Actor''; la \textbf{gun} della torretta è ``Ganna'' e ``GAN'';
il \textbf{Pivot} è spesso ``pivo''; il \emph{Third Person Character} è ``ter person
character'' e ``Peter per personacarte''; \textbf{Quixel Megascans} è ``QuickSell'' e
``Quixel megaschanico''. Dove la ricostruzione è sicura il capitolo usa la forma
corretta e, la prima volta, lo segnala in nota; dove non lo è riporta la forma della
trascrizione e lo dice. La seconda: all'inizio della registrazione, in un punto della
prima esercitazione e durante tutta la pausa, il riconoscimento vocale produce lunghi
blocchi di testo senza senso e ripetuto (``qual è la risposta che ci può essere?'',
``grazie a tutti per l'attenzione'', discorsi sul ``mercato del lavoro'', sul ``governo
di Dalai Lama'' e su ``la migliore dal 2005''): sono artefatti, non parole di nessuno,
e il capitolo li ignora. La terza: gran parte della lezione è una dimostrazione dal
vivo nell'editor dei Blueprint, spesso costruita a voce (``tirate fuori'',
``collegate'', ``splittate''). Il capitolo ricostruisce i grafi da ciò che il docente
dice mentre li costruisce, e li completa, dove serve, con i nomi esatti dei nodi e dei
pin di Unreal Engine 5.

\paragraph{Come leggere questo capitolo.} Il capitolo segue l'ordine della lezione.
Le figure ricostruiscono gli schemi spiegati a voce o alla lavagna (``perdonate le mie
doti molto brutte di disegno''): le operazioni tra vettori, i casi del prodotto
scalare, gli assi di rotazione, la gerarchia del mulino, i grafi della pala, del
mulino che guarda il player e della torretta, i casi dell'\emph{overlap}. Le tabelle
riassumono le operazioni sui vettori, i nodi della lezione e gli errori più comuni
incontrati dagli studenti. Qua e là, dei paragrafi di \textbf{approfondimento}
collegano la lezione ai libri di Marcos Romero e Brenden Sewell (\emph{Blueprints
Visual Scripting for Unreal Engine 5} e la seconda edizione dello stesso testo, scritta
per Unreal Engine 4), i testi di riferimento per questa parte del libro, e a qualche
nozione tecnica generale su Unreal Engine 5; sono segnalati come tali e non fanno parte
della lezione.

Il percorso del capitolo:
\begin{enumerate}
  \item l'attesa: \textbf{Quixel Megascans} e il materiale della lava; il trucco del
        \textbf{CCD}; il programma della giornata (sezioni da~\ref{ue4:sec:attesa}
        a~\ref{ue4:sec:programma});
  \item la teoria: i \textbf{vettori} e le loro operazioni, il \textbf{prodotto
        scalare}, la \textbf{normalizzazione}, i \textbf{rotator} e il \textbf{gimbal
        lock} (sezioni da~\ref{ue4:sec:vettori} a~\ref{ue4:sec:rotator});
  \item il \textbf{Rotating Movement Component} e l'esercizio dei tre cubi (sezione
        \ref{ue4:sec:rotating});
  \item il Blueprint del \textbf{mulino}: i componenti, la pala che gira con
        l'Event Tick, la \textbf{gerarchia} (sezioni da~\ref{ue4:sec:mulino}
        a~\ref{ue4:sec:gerarchia});
  \item il mulino che guarda il player: \textbf{Find Look at Rotation},
        \textbf{RInterp To}, il \textbf{Clamp} e il debug (sezioni
        \ref{ue4:sec:lookat} e~\ref{ue4:sec:clamp});
  \item dopo la pausa, la \textbf{torretta}: i tre pezzi, la \textbf{Sphere
        Collision}, l'\textbf{overlap} (sezioni da~\ref{ue4:sec:torretta}
        a~\ref{ue4:sec:overlap});
  \item \textbf{Is Valid}, i \textbf{tag}, la variabile del bersaglio e i
        \textbf{Custom Event} (sezioni da~\ref{ue4:sec:isvalid}
        a~\ref{ue4:sec:custom});
  \item il \textbf{timer} e l'evento \textbf{Move Turret}: la torretta che segue il
        player (sezioni \ref{ue4:sec:timer} e~\ref{ue4:sec:move});
  \item il \emph{tuning}, i problemi degli studenti e la chiusura (sezioni
        \ref{ue4:sec:problemi} e~\ref{ue4:sec:chiusura}).
\end{enumerate}

% ---------------------------------------------------------------------
\section{L'attesa: Quixel Megascans e la lava}
\label{ue4:sec:attesa}

La registrazione parte mentre il docente fa entrare gli studenti nella stanza di Zoom
(``ho aperto io la prima a far entrare un po' tutti'') e chiacchiera con il primo
arrivato. ``Come va con Unreal? Hai sperimentato un po', l'hai toccato un pochetto?''
Lo studente ci ha giocato con un collega, che gli ha insegnato le texture, ma non ha
trovato online ``alcuna fonte per prendere cose utili, file utili'': per l'esercizio
della lava assegnato da Davide nella lezione precedente, un materiale di lava animata,
ha usato quello che c'era nella cartella del Drive.

\paragraph{Quixel Megascans.} La risposta del docente è \textbf{Quixel
Megascans}\footnote{La trascrizione dice ``QuickSell'', ``Quixel megaschanico'' e
``Materials pixel mega scan''. Il nome corretto è Quixel Megascans: la libreria di
asset scansionati di Quixel, società acquisita da Epic Games, gratuita per l'uso con
Unreal Engine.}, di cui passa il link in chat. Quixel ``ti permette di avere proprio
dei materiali'': ``una marea di materiali di qualsiasi tipo'', ma anche rocce,
\emph{mesh} e texture di alta qualità. Il docente fatica un po' a ritrovare la pagina
giusta, ``perché hanno cambiato il sito'', e ricorda che serve accedere con l'account
di Unreal (cioè di Epic Games). ``Ormai fanno la cooperazione tra di loro'': c'è un
collegamento diretto tra Quixel e Unreal, ``proprio tipo \emph{face to face}''. Cercando
\emph{lava} escono diverse risorse; una, \emph{lava story} o qualcosa del genere,
``sembra carina''.

\paragraph{Come sono fatte.} Le texture di Megascans ``sono proprio \emph{scan}'': c'è
``gente che si mette là con queste macchine che costano una vita umana'' a
scansionare le rocce. ``Quixel è un grande, grande calderone di risorse.'' Da lì si
può partire anche per la lava senza trovare la texture perfetta: si prende una
\textbf{roccia} e, ``come avete visto con Davide, con la matematica'', cioè con un
\textbf{Multiply} e con il \textbf{Lerp} pilotato da un'\textbf{alfa}, si inserisce
``il rosso della lava all'interno di queste rocce''. Un'altra strada sono i video su
YouTube: alcuni tutorial lasciano in descrizione le texture usate (altri le fanno
pagare, ``il maledetto''), e ``ti interessano le texture, il materiale poi lo fai tu''.

\paragraph{La lava degli studenti.} Arrivano gli altri (``buongiorno, Filippo, Marco,
Patrick''), e il docente chiede chi ha provato il materiale della lava ``che vi è stato
dato da Davide''\footnote{La trascrizione dice ``da Dario'', ma subito dopo e più
avanti il nome è Davide, il docente della terza lezione.}. Uno studente l'ha mandato sul
Discord, e il docente lo commenta: ``bello il movimento che hai fatto''. L'hai messo
su un piano, in engine, per vedere come reagiva? No. ``Quello vedilo un po' tu'',
perché a volte il risultato in \emph{viewport} cambia a seconda di come è costruito il
materiale (``che so, hai messo il Multiply sbagliato''): il consiglio è di provarlo
sempre in \emph{viewport}, su un oggetto vero. E ``sicuramente una texture differente
potrebbe essere più carina'': torna utile Megascans.

\paragraph{L'appello.} Il docente aspetta ``2--3 minuti'' che entrino gli altri. Uno
studente avvisa che anche oggi dovrà staccarsi venti minuti prima per andare a
lavorare: ``assolutamente nessun problema''.

% ---------------------------------------------------------------------
\section{Le domande e il trucco del CCD}
\label{ue4:sec:ccd}

Prima di cominciare, il docente apre lo spazio delle domande: ``se qualcuno di voi ha
avuto qualche particolare cosa di cui vuole parlare, domandare già prima che magari
iniziamo, così facciamo il Q\&A già da subito''. Nessuno: ``mi sembrate tutti molto
tranquilli''.

\paragraph{La pallina che entra nel piano.} Allora è il docente a recuperare ``un
piccolo trucchetto che mi sono dimenticato di dirvi'' per il \textbf{gioco della
pallina} della seconda lezione, quello della piattaforma inclinabile: a volte la
pallina, con la fisica attiva, \textbf{penetra} nel piano, ``come abbiamo visto''. Nei
dettagli del componente, nella sezione della fisica (vicino a \textbf{Simulate
Physics}), c'è un'opzione da spuntare: il \textbf{CCD}, \emph{Continuous Collision
Detection}\footnote{La trascrizione dice ``Trucco Collision Collision Detection''.}.
``Fondamentalmente, è un \emph{tick} che, se lo mettiamo su questo particolare
\emph{asset}, lui in modo continuo cerca di vedere le collisioni'': le rende ``più
precise'', ed evita che la pallina attraversi il terreno. ``Se fate questo, vi risolve
tranquillamente il problema.'' Chi vuole, può sistemare il gioco per la prossima volta.

\paragraph{Approfondimento: perché la pallina attraversa il piano.} La fisica di un
gioco non è continua: il motore fisico controlla le collisioni a passi discreti, una
volta per \emph{frame} (o per \emph{substep}). Se in un passo la pallina si sposta più
dello spessore del piano, può trovarsi ``prima'' sopra e ``dopo'' già dentro o sotto,
senza che nessun controllo l'abbia mai vista toccare la superficie: è il fenomeno del
\emph{tunneling}, tanto più probabile quanto più l'oggetto è piccolo e veloce e il piano
sottile. Il CCD (in Unreal la casella \emph{Use CCD} nella sezione \emph{Collision} del
componente, nelle impostazioni avanzate) fa sì che il motore fisico consideri anche il
percorso tra le due posizioni, la ``scia'' dell'oggetto, e non solo le posizioni
finali. Costa un po' di più, ed è per questo che non è attivo di default: lo si accende
solo sugli oggetti che ne hanno bisogno, come proiettili e palline veloci. Un'alternativa
complementare è dare al piano uno spessore maggiore.

% ---------------------------------------------------------------------
\section{Il programma: il mulino e la torretta}
\label{ue4:sec:programma}

``Oggi vi dico già da subito cosa faremo.'' Gli studenti hanno scaricato il
\textbf{materiale didattico}: dentro il file estratto c'è il progetto
\emph{UE5 Intro Start}\footnote{La trascrizione dice ``Hive UE5 intro start'', la stessa
parola iniziale incerta già incontrata nella terza lezione.}, con il suo file di
progetto. Lo si apre, ``così almeno gli \emph{shader} si cominciano a buildare'', si fa
``una piccola spiegazione teorica'' e poi ``partiamo direttamente con la pratica''.

\paragraph{I due obiettivi.} Il docente sposta il \emph{Player Start} e mostra il
livello. A sinistra c'è ``questo bellissimo \textbf{mulino}'', fatto ``facendo ruotare
le pale''. Poi c'è la \textbf{torretta} che gli studenti hanno costruito la volta prima
con Davide (``avete fatto il \emph{merge}, giusto?''), che ``adesso ci segue con lo
sguardo''. O meglio, dovrebbe: nella dimostrazione non lo fa. ``Abbiamo scoperto un
piccolo bug. Poi adesso lo aggiustiamo. E giustamente i bug ogni volta appaiono
praticamente subito.'' Il programma della giornata è questo: le pale del mulino, e la
torretta che segue il giocatore.

% ---------------------------------------------------------------------
\section{I vettori}
\label{ue4:sec:vettori}

``Avete sentito mai parlare di vettori? Matematica, i grafici cartesiani?'' Sì.

\paragraph{Che cos'è un vettore.} ``In generale, un vettore può essere semplificato, in
uno spazio, come letteralmente una \textbf{freccia}'' che ha tre caratteristiche: una
\textbf{direzione}, una \textbf{magnitudine} (``la sua ampiezza'', cioè la lunghezza) e
un \textbf{verso}. In Unreal i vettori si sono già visti senza nominarli: nella
\textbf{location} e nella \textbf{rotation} di ogni oggetto, ``in particolare molto di
più sulla location''. La location è fatta di tre coordinate, \textbf{X}, \textbf{Y} e
\textbf{Z}, e ``queste particolari coordinate possono formare tra di loro, nello spazio
3D, un vettore''. Il docente parte dalla parola ``velocità di movimento'' e si corregge:
``più che una velocità, è più una direzione'', con la sua ampiezza e il suo verso.
Con i vettori si possono ``fare tante cose in Unreal'', per esempio muovere degli
oggetti; le pale del mulino girano in un verso, e anche se in quel caso si parla di
rotazione, e quindi di \emph{rotator}, ``sono sempre formate da delle coordinate''.

\paragraph{Un punto non è ancora un vettore.} Più avanti, alla lavagna, il docente
aggiunge una precisazione importante: ``ricordatevi che quello che vediamo è un
\textbf{punto} nello spazio, non è una freccia, quindi non è già un vettore''. Un
punto, per esempio $(5, 4, 1)$, è solo una posizione; il vettore è la freccia che
collega due punti, per esempio dall'origine $(0,0,0)$ a $(5,4,1)$\footnote{A lezione il
docente dice che ``un vettore è formato dalla somma di due punti''. La relazione
corretta è la \emph{differenza}: il vettore che va da A a B è B $-$ A (si veda la
sottrazione, poco più sotto). Unreal, peraltro, usa la stessa struttura
\emph{Vector} sia per i punti sia per gli spostamenti.}.

\paragraph{Le operazioni.} Con i vettori ``possiamo fare varie operazioni matematiche,
che sono praticamente quelle classiche'': \textbf{addizione}, \textbf{sottrazione},
\textbf{moltiplicazione} e \textbf{prodotto scalare} (\emph{dot product}), ``che è
un'altra cosa, che poi spiegheremo''. La slide le mostra con due vettori, A e B
(figura~\ref{ue4:fig:vettori}).
\begin{itemize}
  \item \textbf{Addizione.} ``Se andiamo a prendere due vettori e li addizioniamo,
        avremo un vettore unico.'' Il vettore A più il vettore B ``dà come risultato
        il vettore C, che è $(5,4)$'': si mette la coda di B sulla punta di A, e C va
        dall'inizio di A alla fine di B. ``Come fare l'ipotenusa, se vi ricordate i
        triangoli.''
  \item \textbf{Sottrazione.} Dà ``la \textbf{differenza} nello spazio tra due
        punti''. Il risultato dipende dall'ordine: A $-$ B è il vettore che va da B
        ad A, B $-$ A quello che va da A a B. La sua lunghezza è la distanza tra i due
        punti.
  \item \textbf{Moltiplicazione.} ``La useremo tantissimo'': si moltiplica tutto ciò
        che fa parte del vettore, X, Y e Z, e quindi il vettore ``viene allungato nella
        sua direzione''.
  \item \textbf{Prodotto scalare.} ``Statemi un attimo a sentire: è l'unica operazione
        che ci dà un valore \textbf{scalare}'', cioè ``non ci ridà un vettore indietro,
        ma un valore in \emph{float}''. (``Chi si ricorda che è il float?'' ``I numeri
        decimali.'' ``Perfetto.'') Lo si vede nella prossima sezione.
\end{itemize}

\begin{figure}[htbp]
\centering
\begin{tikzpicture}[font=\small, scale=0.8,
  ue4vec/.style={-{Stealth[length=2.5mm]}, very thick},
  ue4grid/.style={gray!25, very thin},
  ue4ax/.style={-{Stealth[length=2mm]}, gray!70}]
  % --- addizione
  \begin{scope}
    \draw[ue4grid] (0,0) grid (6,5);
    \draw[ue4ax] (0,0) -- (6.3,0) node[right, font=\scriptsize]{X};
    \draw[ue4ax] (0,0) -- (0,5.3) node[above, font=\scriptsize]{Y};
    \draw[ue4vec, blue!70!black] (0,0) -- node[left, pos=0.55]{A} (2,3);
    \draw[ue4vec, red!70!black] (2,3) -- node[above, pos=0.5]{B} (5,4);
    \draw[ue4vec, green!45!black] (0,0) -- node[below right, pos=0.6]{C $=$ A $+$ B}
      (5,4);
    \fill (5,4) circle (1.5pt) node[above right, font=\scriptsize]{$(5,4)$};
    \node[font=\bfseries] at (3,-0.9) {addizione};
  \end{scope}
  % --- sottrazione
  \begin{scope}[shift={(8.5,0)}]
    \draw[ue4grid] (0,0) grid (6,5);
    \draw[ue4ax] (0,0) -- (6.3,0) node[right, font=\scriptsize]{X};
    \draw[ue4ax] (0,0) -- (0,5.3) node[above, font=\scriptsize]{Y};
    \fill[blue!70!black] (1,1) circle (2.5pt) node[below left]{A};
    \fill[red!70!black] (5,4) circle (2.5pt) node[above right]{B};
    \draw[ue4vec, green!45!black] (1.15,1.11) -- node[above left, pos=0.5]
      {B $-$ A: da A a B} (4.85,3.89);
    \draw[ue4vec, orange!80!black, dashed] (4.9,3.6) -- node[below right, pos=0.5]
      {A $-$ B: da B ad A} (1.2,0.85);
    \node[font=\bfseries] at (3,-0.9) {sottrazione};
  \end{scope}
  % --- moltiplicazione
  \begin{scope}[shift={(0,-7.2)}]
    \draw[ue4grid] (0,0) grid (6,5);
    \draw[ue4ax] (0,0) -- (6.3,0) node[right, font=\scriptsize]{X};
    \draw[ue4ax] (0,0) -- (0,5.3) node[above, font=\scriptsize]{Y};
    \draw[ue4vec, green!45!black, line width=3pt, opacity=0.35] (0,0) -- (5,4);
    \draw[ue4vec, blue!70!black] (0,0) -- (2.5,2) node[pos=0.5, above left]{V};
    \node[green!45!black, font=\small] at (4.6,4.55) {$2 \times$ V};
    \node[font=\bfseries] at (3,-0.9) {moltiplicazione per uno scalare};
  \end{scope}
  % --- normalizzazione
  \begin{scope}[shift={(8.5,-7.2)}]
    \draw[ue4grid] (0,0) grid (6,5);
    \draw[ue4ax] (0,0) -- (6.3,0) node[right, font=\scriptsize]{X};
    \draw[ue4ax] (0,0) -- (0,5.3) node[above, font=\scriptsize]{Y};
    \draw[gray!60, dashed] (0,0) circle (1);
    \node[gray, font=\scriptsize] at (1.35,-0.3) {lunghezza 1};
    \draw[ue4vec, blue!70!black, opacity=0.45] (0,0) -- (4,3)
      node[above right, font=\scriptsize]{V, lunghezza 5};
    \draw[ue4vec, red!70!black] (0,0) -- (0.8,0.6);
    \node[red!70!black, font=\scriptsize, align=left] at (2.7,0.45)
      {Normalize(V)\\$= (0{,}8;\ 0{,}6)$};
    \node[font=\bfseries] at (3,-0.9) {normalizzazione};
  \end{scope}
\end{tikzpicture}
\caption{Le operazioni sui vettori viste a lezione, nel solo piano XY per semplicità.
In alto a sinistra la somma della slide: A seguito da B porta al punto $(5,4)$. In
alto a destra la sottrazione, il cui verso dipende dall'ordine dei fattori. In basso,
la moltiplicazione per un numero, che allunga il vettore senza cambiarne la direzione,
e la normalizzazione, che lo riporta a lunghezza 1 (sezione~\ref{ue4:sec:dot}). I
valori di A e B della slide non si leggono nella trascrizione e sono scelti per
l'esempio; il risultato $(5,4)$ è quello detto a lezione.}
\label{ue4:fig:vettori}
\end{figure}

\paragraph{Approfondimento: i vettori nel libro di Romero e Sewell.} Il
quattordicesimo capitolo di \emph{Blueprints Visual Scripting for Unreal Engine 5},
dedicato ai nodi matematici e alle \emph{trace}, presenta i vettori con un esempio
molto vicino a quello della lezione. In Unreal la struttura \textbf{Vector} contiene
tre \emph{float}, X, Y e Z, e può rappresentare sia un \textbf{punto} nello spazio
(una \emph{location}) sia uno \textbf{spostamento} o una velocità. Un personaggio sta
in $(50, 0, 20)$, un divano in $(450, 0, 20)$, e un'unità di Unreal vale un
centimetro. Per arrivare al divano il personaggio deve spostarsi di $(400, 0, 0)$: 400
centimetri lungo X. Le due regole che ne derivano sono quelle della lezione, scritte
in formula:
\[
  \text{destinazione} = \text{partenza} + \text{spostamento},
  \qquad
  \text{spostamento} = \text{destinazione} - \text{partenza}.
\]
Ogni operazione si fa componente per componente: $(3,5,0) + (5,2,0) = (8,7,0)$ e
$(6,8,0) - (1,4,0) = (5,4,0)$. Il libro aggiunge due esempi pratici: se V1 è la
posizione nel mondo di un personaggio e V2 la posizione \emph{relativa} della sua arma,
V1 $+$ V2 è la posizione nel mondo dell'arma; e la \textbf{lunghezza} della differenza
tra due posizioni (il nodo \emph{Vector Length}) è la loro distanza. Moltiplicare per un
numero negativo inverte il verso: moltiplicando per $-1$ il vettore ``in avanti'' di un
attore si ottiene quello ``all'indietro'', ed è lo stesso trucco usato nel quinto
capitolo per far cambiare direzione a un bersaglio mobile. Il nodo \textbf{Multiply},
avvertono gli autori, è un nodo \emph{wildcard}: se gli si collega un vettore diventa
un \emph{Vector $\times$ Vector}, e per moltiplicare un vettore per un \emph{float}
bisogna fare tasto destro sul secondo pin e scegliere \emph{Convert Pin\ldots} e poi
\emph{Float}.

\begin{table}[htbp]
\centering
\small
\begin{tabular}{@{}p{3.1cm}p{3.4cm}p{2.0cm}p{5.6cm}@{}}
\toprule
\textbf{Operazione} & \textbf{In formula} & \textbf{Risultato} &
\textbf{A che cosa serve in un gioco} \\
\midrule
Addizione & $(a_x+b_x,\ a_y+b_y,\ a_z+b_z)$ & vettore & spostare un punto: posizione
$+$ spostamento $=$ nuova posizione \\
Sottrazione & $(b_x-a_x,\ b_y-a_y,\ b_z-a_z)$ & vettore & la direzione e la distanza
da A a B (``dove sta il player rispetto a me?'') \\
Moltiplicazione per uno scalare & $(k\,a_x,\ k\,a_y,\ k\,a_z)$ & vettore & allungare o
accorciare; con $-1$ invertire il verso; velocità $\times$ \emph{delta seconds} \\
Lunghezza (\emph{Vector Length}) & $\sqrt{a_x^2+a_y^2+a_z^2}$ & \emph{float} & la
magnitudine; della differenza tra due punti, la distanza \\
Normalizzazione (\emph{Normalize}) & $\vec a \,/\, |\vec a|$ & vettore lungo 1 & tenere
solo la direzione, da moltiplicare poi per una velocità \\
Prodotto scalare (\emph{Dot Product}) & $a_x b_x + a_y b_y + a_z b_z$ & \emph{float} &
capire se due direzioni sono concordi, perpendicolari od opposte \\
\bottomrule
\end{tabular}
\caption{Le operazioni sui vettori della lezione, con il nome del nodo Blueprint
corrispondente dove serve e un uso tipico in un gioco.}
\label{ue4:tab:vettori}
\end{table}

% ---------------------------------------------------------------------
\section{Il prodotto scalare e i vettori normalizzati}
\label{ue4:sec:dot}

Il \textbf{prodotto scalare} è l'operazione che restituisce un \emph{float} invece di un
vettore. ``Che cosa ci dà?'' Ci dice come sono orientati due vettori l'uno rispetto
all'altro. Con i vettori A e B della slide (figura~\ref{ue4:fig:dot}):
\begin{itemize}
  \item se il prodotto scalare è \textbf{maggiore di 0}, i due vettori sono
        ``direzionati nella stessa direzione'', grosso modo concordi; se è
        \textbf{1}, ``sono proprio uguali'', cioè puntano esattamente nello stesso
        verso;
  \item se è \textbf{uguale a 0}, sono ``perfettamente \textbf{perpendicolari}'',
        formano un angolo retto;
  \item se è \textbf{minore di 0}, formano un angolo ottuso, ``al verso opposto'': per
        due vettori come quelli della slide ``magari $-0{,}4$, più o meno''; ``se
        proprio sono così'', uno contro l'altro, è $-1$.
\end{itemize}
Questi valori ``puliti'', tra $-1$ e 1, valgono per i \textbf{vettori normalizzati}, ed
è il motivo per cui la slide successiva li introduce.

\paragraph{I vettori normalizzati.} ``Qualcuno di voi sa cos'è un vettore
normalizzato?'' Uno studente: ``un vettore tra 0 e 1''. ``Esatto'': normalizzando un
vettore, ``indipendentemente dalla sua lunghezza, lo ridurremo a 1''. Un vettore ``che
è lungo 5'', normalizzato, ``lo portiamo a 1''\footnote{Più precisamente: un vettore
normalizzato ha lunghezza \emph{esattamente} 1; sono le sue componenti, prese una a
una, a stare tra $-1$ e 1. Il vettore $(4,3,0)$, lungo 5, normalizzato diventa
$(0{,}8;\ 0{,}6;\ 0)$, come nella figura~\ref{ue4:fig:vettori}.}. La direzione resta la
stessa, cambia solo la lunghezza.

\begin{figure}[htbp]
\centering
\begin{tikzpicture}[font=\small,
  ue4vec/.style={-{Stealth[length=2.5mm]}, very thick}]
  \foreach \xs/\ang/\val/\txt in {0/0/{$= 1$}/{stessa direzione},
                                  3.4/40/{$\approx 0{,}77$}/{concordi},
                                  6.8/90/{$= 0$}/{perpendicolari},
                                  10.2/140/{$\approx -0{,}77$}/{angolo ottuso},
                                  13.6/180/{$= -1$}/{opposti}} {
    \begin{scope}[shift={(\xs,0)}]
      \draw[ue4vec, blue!70!black] (0,0) -- (1.5,0) node[below, pos=0.8,
        font=\scriptsize]{A};
      \draw[ue4vec, red!70!black] (0,0) -- (\ang:1.35) node[above,
        font=\scriptsize]{B};
      \fill (0,0) circle (1.5pt);
      \node[font=\small] at (0.3,-1.0) {A $\cdot$ B \val};
      \node[font=\scriptsize, gray!50!black] at (0.3,-1.5) {\txt};
    \end{scope}
  }
\end{tikzpicture}
\caption{Il prodotto scalare tra due vettori normalizzati A e B, al variare
dell'angolo tra di loro: 1 se puntano nello stesso verso, positivo se sono concordi,
0 se sono perpendicolari, negativo se formano un angolo ottuso, $-1$ se sono opposti.
I valori intermedi sono il coseno dell'angolo ($\cos 40^\circ \approx 0{,}77$).}
\label{ue4:fig:dot}
\end{figure}

\paragraph{Approfondimento: il prodotto scalare è un coseno.} Per due vettori qualsiasi
vale
\[
  \vec a \cdot \vec b \;=\; a_x b_x + a_y b_y + a_z b_z \;=\; |\vec a|\,|\vec b|\,
  \cos\theta ,
\]
dove $\theta$ è l'angolo tra i due. Se i vettori sono normalizzati, le lunghezze valgono
1 e il prodotto scalare \emph{è} il coseno dell'angolo: per questo varia tra $-1$ e 1,
come dicono Romero e Sewell nel quattordicesimo capitolo (``il prodotto scalare di due
vettori normalizzati è uguale al coseno dell'angolo che formano''), e per questo i casi
della slide sono 1, 0 e $-1$. Con vettori non normalizzati il \emph{segno} resta
affidabile (positivo, zero, negativo) ma il valore non è più confrontabile. Nei giochi
il prodotto scalare risponde a domande come ``il player è davanti o dietro al
nemico?'': si normalizza la differenza tra la posizione del player e quella del nemico
e la si confronta con il \emph{Forward Vector} del nemico, che Unreal restituisce già
normalizzato (nodi \emph{Get Actor Forward Vector}, \emph{Right Vector}, \emph{Up
Vector}). Un risultato maggiore di $\cos 45^\circ \approx 0{,}71$, per esempio,
significa ``il player è dentro un cono di visione di 90 gradi''. Il quinto capitolo del
libro usa invece la normalizzazione per il movimento: si normalizza una direzione, la si
moltiplica per una velocità e per il \emph{delta time} (con il nodo \textbf{Get World
Delta Seconds}, che tornerà nella torretta) e la si somma alla posizione a ogni
\emph{tick} con \emph{Add Actor World Offset}. ``Normalizzare'', scrivono, ``è una
procedura comune che assicura che il vettore abbia lunghezza 1, il che lo rende
compatibile con il resto dei calcoli''.

% ---------------------------------------------------------------------
\section{I rotator, il gimbal lock e i quaternioni}
\label{ue4:sec:rotator}

``Andiamo avanti: poi abbiamo i \textbf{rotator}.'' Per spiegarli il docente usa
l'immagine del \textbf{giroscopio}: tre ``anelli che girano tra di loro e formano
proprio quella che è la rotazione''. Ogni anello corrisponde a un asse, con i colori di
sempre: la X in rosso, la Y in verde, la Z in blu (figura~\ref{ue4:fig:rotator}).

\paragraph{Il gimbal lock.} Con le ``regole classiche di Eulero'', cioè descrivendo la
rotazione come tre angoli applicati uno dopo l'altro, c'è ``un possibile problema'': il
\textbf{gimbal lock}. ``Non vi sto a spiegare troppo, perché fondamentalmente andremo a
parlare di quattro dimensioni invece di tre, ed è tutta una cosa molto, molto
complessa.'' In breve, ``potrebbe succedere che due assi si blocchino tra di loro'' e
comincino a ruotare ``sempre dipendenti l'uno dall'altro''. Si risolve con i
\textbf{quaternioni}, che usano quattro valori invece di tre, ma ``parliamo di
matematica molto complessa, e non andiamo molto avanti su questa cosa''. Quello che
conta, per ora: per descrivere una rotazione nello spazio servono sempre la X, la Y e
la Z, ``perché abbiamo tre giroscopi che girano di per sé''. Nell'editor, su un
qualsiasi \emph{asset}, si vedono proprio questi tre valori.

\paragraph{L'obiettivo.} ``Cosa vogliamo fare oggi? Vogliamo creare qualcosa che giri'':
qualcosa che usi il \emph{rotator} per far girare le pale del mulino.

\begin{figure}[htbp]
\centering
\begin{tikzpicture}[font=\small,
  ue4ax/.style={-{Stealth[length=2.5mm]}, very thick},
  ue4rot/.style={-{Stealth[length=2mm]}, thick}]
  % assi in proiezione
  \coordinate (O) at (0,0);
  \draw[ue4ax, red!75!black] (O) -- (-1.9,-1.1) node[below left]{X};
  \draw[ue4ax, green!50!black] (O) -- (2.6,0) node[right]{Y};
  \draw[ue4ax, blue!75!black] (O) -- (0,2.6) node[above]{Z};
  % anelli di rotazione
  \draw[ue4rot, red!75!black] (-1.25,-0.25) arc[x radius=0.35, y radius=0.6,
    start angle=90, end angle=400];
  \node[red!75!black, font=\scriptsize, align=center] at (-2.9,-0.1)
    {\textbf{Roll}\\attorno a X};
  \draw[ue4rot, green!50!black] (1.7,0.55) arc[x radius=0.3, y radius=0.55,
    start angle=90, end angle=400];
  \node[green!50!black, font=\scriptsize, align=center] at (2.4,1.05)
    {\textbf{Pitch}\\attorno a Y};
  \draw[ue4rot, blue!75!black] (0.55,1.75) arc[x radius=0.55, y radius=0.2,
    start angle=0, end angle=310];
  \node[blue!75!black, font=\scriptsize, align=center] at (1.65,2.15)
    {\textbf{Yaw}\\attorno a Z};
  % giroscopio stilizzato
  \begin{scope}[shift={(7.0,0.6)}]
    \draw[very thick, red!75!black] (0,0) ellipse (1.9 and 1.9);
    \draw[very thick, green!50!black] (0,0) ellipse (1.45 and 0.55);
    \draw[very thick, blue!75!black] (0,0) ellipse (0.45 and 1.05);
    \fill (0,0) circle (2pt);
    \node[font=\scriptsize, align=center] at (0,-2.4) {il giroscopio: tre anelli
      annidati,\\ognuno ruota attorno al proprio asse};
  \end{scope}
\end{tikzpicture}
\caption{I tre assi di Unreal con i loro colori e i nomi delle tre rotazioni di un
\emph{rotator}: \emph{roll} attorno a X, \emph{pitch} attorno a Y, \emph{yaw} attorno a
Z. A destra, l'immagine del giroscopio usata a lezione. Schema ricostruito.}
\label{ue4:fig:rotator}
\end{figure}

\paragraph{Approfondimento: roll, pitch, yaw e il gimbal lock in pratica.} Nella seconda
edizione di Romero e Sewell, l'elenco dei tipi di variabile descrive il
\textbf{Rotator} così: contiene i \emph{float} X (\textbf{Roll}), Y (\textbf{Pitch}) e
Z (\textbf{Yaw}), espressi in gradi. I tre nomi vengono dall'aviazione: il
\emph{roll} è il rollio (l'aereo che inclina le ali), il \emph{pitch} il beccheggio (il
muso che sale e scende), lo \emph{yaw} l'imbardata (il muso che gira a destra e a
sinistra). Una torretta che insegue un bersaglio usa lo \emph{yaw} per girarsi e il
\emph{pitch} per alzare la canna; le pale del mulino, come si vedrà, girano in
\emph{pitch}. Il \textbf{gimbal lock} succede quando l'anello di mezzo ruota di 90
gradi: il primo e il terzo anello si ritrovano allineati, e ruotarli produce lo stesso
movimento, così un grado di libertà va perso. Con gli angoli di Unreal il caso tipico
è il \emph{pitch} a $\pm 90$: con il muso dritto verso l'alto, \emph{yaw} e \emph{roll}
fanno la stessa cosa. I \textbf{quaternioni} (in Unreal la struttura \emph{Quat})
rappresentano invece ogni rotazione come una rotazione unica attorno a un asse
qualsiasi, con quattro numeri, e non hanno questo problema. Unreal li usa
internamente: un \emph{rotator} è il formato comodo per le persone, il quaternione
quello sicuro per i calcoli, e nodi come \emph{Combine Rotators} convertono i
\emph{rotator} in quaternioni, li combinano e li riconvertono. Un effetto collaterale,
che conviene conoscere: dato che il \emph{pitch} di un \emph{rotator} viene sempre
riportato tra $-90$ e $90$ gradi, una pala che ruota in \emph{pitch} oltre i 90 gradi
viene descritta ``capovolta'' (\emph{yaw} e \emph{roll} a 180), e chi stampasse il suo
\emph{rotator} vedrebbe numeri che saltano. La rotazione sullo schermo, però, resta
corretta.

% ---------------------------------------------------------------------
\section{Il Rotating Movement Component}
\label{ue4:sec:rotating}

``Allora, avete aperto l'editor di Unreal? Sì? Perfetto.'' Prima di costruire il
mulino, il docente mostra ``quello più semplice'' tra i ``vari modi per far ruotare un
actor in generale''; un modo che però, avverte, ``è difficile usarlo su vari
componenti''.

\paragraph{Un componente aggiunto dai Details.} Si trascina in scena un \textbf{cubo}
(con il tasto \textbf{Fine}, \emph{End}, l'oggetto ``va diretto sul
terreno''\footnote{La trascrizione dice ``il tasto AND'': è il tasto \emph{End} (Fine),
che in Unreal appoggia l'oggetto selezionato sulla superficie sottostante.}). Ora,
``qui sui \textbf{Details} sulla destra, potete aggiungere dei componenti anche non
avendo una schermata di Blueprint'': con il pulsante \textbf{Add} dei Details si può
mettere un componente direttamente sull'actor in scena. Si scrive \emph{rotating
movement} e si sceglie il \textbf{Rotating Movement Component}. ``Trovato? Sì?
Perfetto.''

\paragraph{I parametri.} Il componente ha una serie di dati, tra cui il più importante
è il \textbf{Rotation Rate}, con i tre assi X, Y e Z: la velocità di rotazione su
ciascun asse. ``Di default il Rotating Movement, come potete vedere, è messo a
\textbf{180}'' sulla Z: mezzo giro al secondo attorno all'asse verticale.

\paragraph{La regola della mobilità.} ``Intanto vi dico, ragazzi: come abbiamo visto,
se un \emph{asset} si deve muovere all'interno di uno spazio, la sua mobilità dovrà
essere \textbf{Movable}. Ricordatevelo sempre.'' Se lo si lascia \textbf{Static} e si
va in Play, ``prima di tutto non si muove'', e poi l'engine restituisce un errore,
che il docente legge: \emph{Static Mesh Component has to be movable if you'd like to
move it with Rotating Movement Component}. Con la mobilità corretta, in Play il cubo
gira su sé stesso: ``abbiamo fatto le munizioni di \emph{Metal Gear Solid}'', scherza
il docente, pensando agli oggetti da raccogliere che ruotano a mezz'aria. ``Questo è un
modo per far ruotare'': Unreal ``ci dà la possibilità di poter ruotare direttamente un
nostro bellissimo actor tramite un componente che loro hanno già creato''. ``Un oggetto
ruotante che ruota. Perché `ruotante' è brutto, sembra.''

\paragraph{Esercitazione: i tre cubi.} ``Piccola esercitazione velocissima, faccio
queste botte e risposte con le esercitazioni.'' Si copia il cubo (``tanto vi ricordate
come si copiano'') e si fa in modo che \textbf{un cubo ruoti sull'asse X} e
\textbf{un altro sull'asse Y}, alla velocità che si preferisce. ``Ricordatevi di
metterli Movable. Vi do cinque minuti.'' La correzione: sul secondo cubo si azzera la Z
e si mette 180 sulla X; sul terzo, 180 sulla Y. Si controlla che siano tutti e tre
Movable, si va in Play, ``ed ecco le nostre tre rotazioni''. ``Possiamo farlo con
qualsiasi \emph{asset}.''

\paragraph{Approfondimento: la sedia che gira.} Il Rotating Movement Component è il
primo componente che Romero e Sewell fanno usare nel libro, già nel primo capitolo: un
Blueprint \emph{BP\_RotatingChair} con una \emph{Static Mesh} (la sedia dello
\emph{Starter Content}) e un Rotating Movement. ``Di default'', scrivono, ``questo
componente ruoterà il Blueprint attorno all'asse Z, quindi non serve cambiarne le
proprietà'': si compila, si trascina la sedia nel livello e la si guarda girare in Play.
Il componente fa parte della famiglia dei \emph{movement component}, componenti che non
hanno una forma propria ma muovono un altro componente, quello indicato come
\emph{Updated Component}, che di default è la \textbf{radice} dell'actor. È qui la
ragione del limite che il docente scoprirà tra poco con il mulino: il Rotating Movement
fa girare la radice, e con essa tutto ciò che vi è attaccato. Il componente ha anche
un'opzione \emph{Rotation in Local Space} e un \emph{Pivot Translation}, che sposta il
centro di rotazione; ma per far girare una parte sola di un actor la strada più chiara
resta quella che segue, il Blueprint.

% ---------------------------------------------------------------------
\section{Il Blueprint del mulino}
\label{ue4:sec:mulino}

``Proviamo a fare intanto il nostro \emph{windmill}.'' Il mulino del livello di esempio
è formato da più cose, si vede nei Details. ``Chi si ricorda come si chiamano queste
cose qua, che sono tutte insieme in un singolo \emph{asset}?'' ``Componenti.''
``Sono i \textbf{componenti}'': come si è visto nella seconda lezione, un actor ``può
essere, scusate la ridondanza, composto da più componenti''. In questo caso un
\textbf{body}, il corpo, ``la nostra \emph{mesh} principale, il cubone'', e
\textbf{due pale}.

\paragraph{La classe.} Si crea un nuovo \textbf{Blueprint Class} di tipo
\textbf{Actor}, ``perché non sarà disponibile al giocatore, ma sarà solamente là come
sfondo'': non va controllato, non è un \emph{Pawn}. Il docente lo chiama
\textbf{BP\_Windmill}\footnote{La trascrizione dice ``Windmill, io in questo caso lo
chiamerò 1'': probabilmente \emph{BP\_Windmill1}, per non confonderlo con il mulino già
presente nel progetto.}.

\paragraph{Il corpo.} Nella \textbf{Viewport} del Blueprint si aggiunge una
\textbf{Static Mesh} (Add, in alto a sinistra, \emph{Static Mesh}) e la si chiama
\textbf{Body}. Nei dettagli, alla voce \emph{Static Mesh}, si sceglie un cubo, per
esempio \emph{SM\_Cube}. Il primo cubo scelto dal docente ``ha un \emph{pivot}
particolarmente strano''; ne sceglie un altro, ``molto meglio, perché c'ha il
\emph{pivot} al centro''. Lo solleva di circa 40 sulla Z e gli dà un materiale: il
docente mette la \emph{prototype grid}, ``voi mettete quello che volete''. Poi lo scala
per farne il corpo del mulino: ``3 metri per 5'' di altezza, ``di larghezza facciamo 2
metri'', cioè una scala di circa 3, 2 e 5 sui tre assi\footnote{La trascrizione dice
``x32 sulla y5 sull'altezza'' e ``3:25'': la lettura più probabile è X 3, Y 2, Z 5. Il
cubo di base di Unreal misura un metro di lato, quindi una scala di 3 corrisponde a 3
metri.}.

\paragraph{La viewport che non si trova.} Uno studente non ricorda come si apre la
viewport del Blueprint. ``O ce l'hai qua, o ce l'hai qua, oppure dovresti avere
\textbf{Window} e dovresti vedere \textbf{Viewport}.'' La voce è già spuntata, ma la
finestra non c'è: forse è ``nascosta da qualche schermata di Windows'' (``prova a fare
Alt+Tab''). Si passa alla condivisione dello schermo, e si scopre che il problema è un
altro: lo studente ha creato il Blueprint ma non lo ritrova. La soluzione è il
\textbf{filtro} del Content Browser: le ``tre stanghettine'' vicino alla ricerca,
categoria \emph{Blueprint}, ed ecco il file; doppio clic e, se si apre la versione
ridotta, \emph{Open Full Blueprint Editor}. Poi Add, ma nell'editor in italiano la voce
non si trova subito (è ``Componente Mesh Statica''). ``Effettivamente, ragazzi, vi
consiglio di mettervi l'editor in \textbf{inglese}, perché è molto più semplice. Anzi,
ve lo consiglio proprio'': i nomi dei nodi e dei componenti nei tutorial, nei libri e
nella documentazione sono quasi sempre in inglese.

\paragraph{Le pale.} ``Adesso dobbiamo creare le nostre palette.'' Si copia e incolla il
corpo (Ctrl+C, Ctrl+V): le copie finiscono dentro il corpo, ``quindi cliccatele e
tiratele fuori, perché sono uguali''. Si scalano a forma di pala, una orizzontale e una
verticale, si dà loro un materiale diverso (il docente usa \emph{M\_Grid}, ``questa
qua violacea'') e le si rinomina \textbf{Pala1} e \textbf{Pala2}.

\paragraph{Attaccare le pale al corpo.} ``Importante, ragazzi: queste pale devono essere
attaccate al corpo. Come facciamo? Abbiamo visto l'ereditarietà.'' È la
\textbf{gerarchia} dei componenti: ``il fatto che io posso mettere degli \emph{asset}
attaccati l'uno con l'altro''. Si prende una pala e la si \textbf{trascina sopra} Body,
nell'elenco dei componenti: ora è sua figlia. ``Scusami, non ho capito come si fa.''
``Se hai questa pala, la prendi e la trascini sopra Body.''

\paragraph{Il Rotating Movement non basta.} ``Adesso dobbiamo dare una rotazione a
queste pale. Potremmo provare con il Rotating Movement.'' Il docente lo aggiunge al
Blueprint, mette il mulino in scena (``molto più grosso'') e va in Play. ``Guardate che
succede'': il Rotating Movement è impostato sulla Z, e \textbf{ruota tutto il corpo}.
``Una particolarità del Rotating Movement è che non può essere attaccato a
componenti: o ruota tutto, o non ruota niente.'' Quindi bisogna programmare il
movimento delle pale. Il Rotating Movement va tolto: ``se l'avete messo e avete adesso
il vostro mulino che ruota su sé stesso, lo dovete eliminare''.

% ---------------------------------------------------------------------
\section{Far girare una pala: Event Tick e Combine Rotators}
\label{ue4:sec:pala}

``Domanda. Questo movimento deve essere fatto per ogni \emph{frame}. Quindi deve essere
un evento che riprende ogni \emph{frame}. Chi si ricorda qual è l'evento che lo fa? L'abbiamo
usato per la piattaforma dell'altra volta.'' ``L'\textbf{Event Tick}.'' ``Ottimo.''
L'Event Tick, ``che sarà il nostro protagonista'', è già presente nell'Event Graph di un
nuovo Blueprint.

\paragraph{I delta seconds, di nuovo.} ``Chi si ricorda che sono i \textbf{Delta
Seconds}?'' Uno studente: ``determina quanti tick al minuto\ldots''. Un altro: ``i
secondi effettivi, e non quelli basati su \emph{frame}''. Il docente riprende: è ``lo
spazio che intercorre tra un \emph{frame} e l'altro'', la differenza di tempo tra due
\emph{frame}, ``in questo caso si parla di millisecondi''. Si moltiplica un valore per i
\emph{delta seconds} perché la durata di un \emph{frame} cambia (a 30 \emph{frame} al
secondo è molto più lunga che a 120), e così gli spostamenti ``non diventano più lenti
se abbiamo meno \emph{frame}, o più veloci se ne abbiamo di più''. (Il docente prova a
disegnarlo sullo schermo, ma lo strumento ``a quanto pare non mi funziona benissimo''.)

\paragraph{La velocità.} Si tira fuori un cavo da Delta Seconds e si crea un
\textbf{Multiply}: o scrivendo \emph{multiply}, o digitando semplicemente
l'\textbf{asterisco}, ``e vi dà subito il Multiply''. Il secondo valore è la
\textbf{velocità}: ``più è alto, più è veloce''. Per ora, \textbf{90}.

\paragraph{Make Rotator.} ``Adesso vi dico una cosa.'' Il nodo \textbf{Make Rotator}
costruisce un \emph{rotator}, ``quello che muove a livello di rotazione un
\emph{asset}'', a partire da tre \emph{float}: X, Y e Z. Quale asse serve? ``Prendiamo
la nostra bellissima pala e guardiamo subito in che asse la vogliamo ruotare. La X non
mi sembra, la Z neanche: è la nostra carissima \textbf{Y}'', cioè il \textbf{pitch}.
Si collega il risultato del Multiply al pin Y (Pitch): il \emph{rotator} creato ha X e
Z a 0, ``perché non si devono muovere'', e la Y uguale a velocità per \emph{delta
seconds}.

\paragraph{Combine Rotators.} ``Possiamo farlo in tanti modi, ma oggi ne facciamo tipo
tre.'' Il primo è \textbf{Combine Rotators}, che ``combina due \emph{rotator} per
darne uno in uscita''. Il primo dei due è l'\textbf{origine}: ``a livello concettuale,
qual è l'origine che dobbiamo prendere?'' ``Le coordinate del componente.'' ``Sì:
quale componente?'' Per ora una sola pala, la Pala1. Siccome serve un componente e non
l'intero actor, si trascina la pala dall'elenco dei componenti nel grafo: si ottiene
una \textbf{referenza} a quel componente, ``un \emph{get}'' che ``si riferisce a quel
componente preciso''.

\paragraph{Relativa o del mondo?} Da quella referenza si chiede la rotazione. Ma quale?
``Abbiamo visto la differenza tra \textbf{rotazione relativa} e \textbf{rotazione del
mondo}.'' La pala è attaccata al corpo, ma deve muoversi da sola, indipendentemente dal
corpo. ``Relativa al corpo, secondo me'', propone uno studente. ``Relativa a sé
stessa, quindi la sua asse singola. Se prendiamo quella del mondo, girerà in modo
strano.'' Si usa dunque \textbf{Get Relative Rotation}, la si collega al primo
ingresso di Combine Rotators, e il \emph{rotator} di Make Rotator al secondo.

\paragraph{Set Relative Rotation.} ``Chi si ricorda qual è il blocco che serve per
settare la rotazione? Tanto è un \emph{set}.'' È \textbf{Set Relative Rotation},
``perché, ovviamente, abbiamo preso la relativa''. Lo si crea tirando fuori un cavo da
Pala1 (così il \emph{target} è già la pala), lo si collega all'Event Tick e gli si dà
in ingresso la nuova rotazione, l'uscita di Combine Rotators
(figura~\ref{ue4:fig:pala}). ``E mi raccomando, deve essere \textbf{Movable}.'' Si
mette il mulino in scena, si va in Play: ``ed ecco che una pala ruota''.

\begin{figure}[htbp]
\centering
\begin{tikzpicture}[font=\small,
  ue4node/.style={draw, thick, rounded corners=4pt, align=center, fill=gray!8,
                  minimum height=0.9cm},
  ue4event/.style={ue4node, fill=red!12},
  ue4exec/.style={-{Stealth[length=2.5mm]}, very thick, black!80},
  ue4flt/.style={-{Stealth[length=2mm]}, thick, green!50!black},
  ue4rot/.style={-{Stealth[length=2mm]}, thick, violet!70!black},
  ue4obj/.style={-{Stealth[length=2mm]}, thick, blue!60!black}]
  \node[ue4event, text width=2.1cm] (tick) at (0,0)
    {\textbf{Event Tick}\\\footnotesize Delta Seconds};
  \node[ue4node, text width=1.5cm] (mul) at (2.9,-1.6)
    {\textbf{Multiply}\\\footnotesize $\times$ 90};
  \node[ue4node, text width=1.9cm] (make) at (5.6,-1.6)
    {\textbf{Make Rotator}\\\footnotesize X 0, \textbf{Y}, Z 0};
  \node[ue4node, fill=blue!8, text width=1.3cm] (pala) at (2.9,-3.5)
    {\textbf{Pala1}};
  \node[ue4node, text width=2.1cm] (get) at (5.6,-3.5)
    {\textbf{Get Relative}\\\textbf{Rotation}};
  \node[ue4node, text width=2.0cm, minimum height=1.5cm] (comb) at (8.6,-2.5)
    {\textbf{Combine}\\\textbf{Rotators}\\\footnotesize A, B};
  \node[ue4node, text width=2.2cm, minimum height=1.6cm] (set) at (11.6,0)
    {\textbf{Set Relative}\\\textbf{Rotation}\\\footnotesize Target: Pala1};
  \draw[ue4exec] (tick.east) -- (set.west |- tick.east);
  \draw[ue4flt] ([yshift=-3mm]tick.east) -- ++(0.4,0) |- (mul.west);
  \draw[ue4flt] (mul.east) -- node[above, font=\scriptsize]{Pitch} (make.west);
  \draw[ue4obj] (pala.east) -- (get.west);
  \draw[ue4rot] (get.east) -- ++(0.4,0) |- node[pos=0.75, below,
    font=\scriptsize]{A} ([yshift=-3mm]comb.west);
  \draw[ue4rot] (make.east) -- ++(0.4,0) |- node[pos=0.75, above,
    font=\scriptsize]{B} ([yshift=3mm]comb.west);
  \draw[ue4rot] (comb.east) -- node[above, font=\scriptsize]{New Rotation}
    ++(1.6,0) |- ([yshift=-4mm]set.west);
  \draw[ue4obj, dashed] (pala.south) -- ++(0,-0.5) -| node[pos=0.25, below,
    font=\scriptsize]{Target} (set.south);
\end{tikzpicture}
\caption{Il grafo della prima pala. A ogni \emph{frame}, l'Event Tick esegue Set
Relative Rotation su Pala1; la nuova rotazione è la rotazione relativa attuale della
pala combinata con un piccolo incremento in \emph{pitch}, pari a 90 gradi al secondo
moltiplicati per i \emph{delta seconds}. Il cavo nero è l'esecuzione, gli altri sono
dati (verde i \emph{float}, viola i \emph{rotator}, blu le referenze). Schema
ricostruito.}
\label{ue4:fig:pala}
\end{figure}

\paragraph{I problemi degli studenti.} ``Avete la vostra paletta che ruota?'' Non
tutti. Una studentessa chiede di rivedere il codice e non trova il blocco Get Relative
Rotation: ``tira fuori da Pala1 e scrivi \emph{get relative rotation}''. Un altro
studente vede la pala ``staccarsi da tutto l'asset'' quando preme Play, e in
condivisione si scoprono due errori. Il primo: ha calcolato la rotazione con Combine
Rotators ma non l'ha collegata al Set (``hai calcolato la rotazione, ma non
gliel'hai assegnata''). Il secondo: anche dopo, la pala gira attorno a un punto
sbagliato, perché la \emph{mesh} che ha scelto ha il \textbf{pivot} spostato. ``Vedi che
il pivot sta quasi lì? Cambia la \emph{mesh}'': con un cubo che ha il pivot al centro, la
pala gira correttamente. (Nel frattempo il docente accenna anche a un \textbf{Set World
Rotation}, per mostrare la differenza con la rotazione relativa.)

\paragraph{Approfondimento: relativo e mondo nel libro.} Il quattordicesimo capitolo di
Romero e Sewell spiega la differenza con un esempio molto simile al mulino. Un
Blueprint ha come radice un \emph{DefaultSceneRoot}; sotto c'è un componente
\emph{Table}, un tavolo, e sotto il tavolo un componente \emph{Statue}, una statua. Il
\emph{transform} di Table è \textbf{relativo} a DefaultSceneRoot, e quello di Statue è
relativo a Table: ``se muovete il componente Table nella Viewport, anche il componente
Statue si muoverà; ma se cambiate il \emph{transform} relativo di Statue, Table resterà
dov'è''. Lo stesso vale per il mulino: la rotazione relativa di Pala1 è misurata
rispetto a Body. Se il mulino viene ruotato nel livello, la rotazione relativa della
pala non cambia, e il grafo continua a farla girare attorno al proprio asse; usando la
rotazione del mondo, invece, l'incremento in \emph{pitch} verrebbe applicato attorno
all'asse Y del \emph{mondo}, che coincide con quello della pala solo se il mulino non è
ruotato: è il ``girerà in modo strano'' del docente. Nei vettori lo si vede ancora più
chiaramente: il punto iniziale di ogni posizione nel mondo è sempre $(0,0,0)$, mentre
quello di una posizione relativa è il punto in cui sta il genitore.

\paragraph{Approfondimento: una scorciatoia, Add Local Rotation.} Il trio Get Relative
Rotation, Combine Rotators e Set Relative Rotation è istruttivo perché mostra i tre
passaggi di ogni animazione fatta a mano: \emph{leggi} lo stato attuale, \emph{calcola}
il nuovo, \emph{scrivi} il risultato. Unreal offre anche nodi che fanno tutto in uno:
\textbf{Add Local Rotation} (o \emph{Add Relative Rotation}) su un componente aggiunge
un \emph{Delta Rotation} alla rotazione attuale, così come \emph{Add Actor World
Rotation}, usato nella seconda lezione per la piattaforma, lo fa per un intero actor.
Collegando l'Event Tick a un Add Local Rotation su Pala1, con un Delta Rotation fatto
con Make Rotator come sopra, si ottiene lo stesso mulino con un nodo in meno. Il
risultato è lo stesso; cambia la chiarezza di ciò che succede, ed è per questo che a
lezione conviene vedere prima la versione lunga.

% ---------------------------------------------------------------------
\section{La gerarchia: la seconda pala e il cubetto}
\label{ue4:sec:gerarchia}

``Possiamo fare una cosa bella: invece di dover fare la copia esatta per quanto
riguarda la nostra prima pala, possiamo vedere quello che la \textbf{gerarchia} ci
permette di fare in modo molto facile.'' La Pala2 deve seguire in ogni caso il movimento
della Pala1. Invece di programmare anche lei, la si mette ``in eredità'' (``non quella
con Carlo Conti, ragazzi'') alla Pala1: la si trascina sopra Pala1, e diventa sua
\textbf{figlia}. Play: ``non abbiamo programmato niente'', ma la Pala2 segue la
rotazione della Pala1. ``È la Pala1 che gira, non la Pala2; la Pala2 segue solamente.''

\paragraph{Esercitazione: il cubetto sulla pala.} ``Se noi volessimo attaccare un
cubettino su un estremo della pala\ldots anzi, sapete che c'è? Lo fate voi.'' La
consegna: aggiungere una \textbf{Static Mesh} qualsiasi (``un cubo, una \emph{mesh},
quello che vi pare''), piccola, su un'estremità della pala, e attaccarla alla
\textbf{Pala2}, in modo che segua anche lei il movimento. ``Vi do cinque minuti
classici. Ricordatevi l'ereditarietà: per spostarlo dovete prenderlo proprio e metterlo
sopra.''

\paragraph{La domanda sul pivot.} Uno studente l'ha fatto, ma è sorpreso: ``mi
aspettavo di dover spostare il pivot al centro della rotazione. Invece,
automaticamente, è come se si agganciasse alla pala.'' ``Esatto. Perché comunque
\textbf{mantiene la sua posizione relativa}'': il cubetto ha il suo pivot al proprio
centro, ma la sua posizione è espressa rispetto al genitore, e quando il genitore
ruota, il figlio ruota con lui attorno al pivot del genitore. ``Se lo dovessimo
spostare, andrebbe insieme: manterrebbe la sua origine, il fatto che sta sulla parte
destra, sopra, giù, quello che è, ma si muoverebbe relativo all'altro.'' Play: ``eccolo
che gira insieme''. ``Così abbiamo capito come si ereditano piano piano gli oggetti tra
di loro'' (figura~\ref{ue4:fig:gerarchia}).

\begin{figure}[htbp]
\centering
\begin{tikzpicture}[font=\small,
  ue4comp/.style={draw, thick, rounded corners=3pt, minimum width=2.6cm,
                  minimum height=0.7cm, align=center},
  ue4tree/.style={thick, gray!60!black},
  ue4arr/.style={-{Stealth[length=2.5mm]}, thick}]
  % albero dei componenti
  \node[ue4comp, fill=gray!10] (root) at (0,0) {BP\_Windmill (actor)};
  \node[ue4comp, fill=orange!15] (body) at (0.6,-1.0) {Body};
  \node[ue4comp, fill=violet!12] (p1) at (1.2,-2.0) {Pala1 \footnotesize(ruota)};
  \node[ue4comp, fill=violet!6] (p2) at (1.8,-3.0) {Pala2 \footnotesize(segue)};
  \node[ue4comp, fill=green!10] (cub) at (2.4,-4.0) {Cubetto \footnotesize(segue)};
  \foreach \a/\b in {root/body, body/p1, p1/p2, p2/cub}
    \draw[ue4tree] ([xshift=-1.0cm]\a.south) |- (\b.west);
  \node[font=\footnotesize, align=center] at (1.2,-5.0) {l'elenco dei componenti};
  % mulino stilizzato
  \begin{scope}[shift={(8.2,-2.2)}]
    \draw[thick, fill=orange!15] (-0.75,-2.5) rectangle (0.75,1.0);
    \node[font=\scriptsize] at (0,-2.2) {Body};
    \begin{scope}[rotate=25]
      \draw[thick, fill=violet!25] (-2.2,-0.18) rectangle (2.2,0.18);
      \draw[thick, fill=violet!12] (-0.18,-2.2) rectangle (0.18,2.2);
      \draw[thick, fill=green!35] (-0.25,1.85) rectangle (0.25,2.35);
    \end{scope}
    \fill (0,0) circle (2.2pt);
    \draw[ue4arr, red!70!black] (2.5,0.4) arc[start angle=10, end angle=70,
      radius=2.5];
    \node[font=\scriptsize, align=left] at (3.6,-0.4) {Pala1 orizzontale\\Pala2
      verticale\\cubetto in cima};
  \end{scope}
\end{tikzpicture}
\caption{La gerarchia del mulino. Solo Pala1 è programmata per ruotare; Pala2, sua
figlia, e il cubetto, figlio di Pala2, ne seguono la rotazione senza alcun nodo,
mantenendo la loro posizione relativa. Schema ricostruito (nel mulino della lezione le
pale sono di fronte al corpo).}
\label{ue4:fig:gerarchia}
\end{figure}

% ---------------------------------------------------------------------
\section{Il mulino che guarda il player: Find Look at Rotation}
\label{ue4:sec:lookat}

``Adesso vi farò vedere un altro metodino'' per creare una rotazione, sempre
sull'Event Tick, ``così vi faccio vedere anche un altro blocchettino interessante,
perché sono tanti''. Nei Blueprint, come nella programmazione in generale, ``non esiste
un qualcosa che sia fisso'': ``essendo che è un mondo di spazi, punti, vettori, ce lo
possiamo giocare come ci pare''. Il docente ricorda una persona che una volta ha fatto
``un materiale con la matematica senza usare texture. Non so come abbia fatto, ma si
può fare''.

\paragraph{Staccare i cavi.} Si stacca l'Event Tick dal grafo precedente. ``Ricordatevi:
per staccare i cavettini, tenete premuto \textbf{Alt} e cliccate sui puntini.''

\paragraph{Find Look at Rotation.} ``Un blocco chiamato \textbf{Find Look at Rotation}.
Interessante, molto interessante. Questo è un sacco utile, ve lo dico.'' Il docente
legge la descrizione del nodo: \emph{find a rotation for an object at Start location to
point at Target location}. Alla lavagna (``scrittura incredibile, mi dispiace'') disegna
due punti nello spazio: lo \textbf{Start}, ``da dove parto'', e il \textbf{Target}, il
punto da raggiungere. Il nodo restituisce ``la rotazione verso quel punto'': la
rotazione che un oggetto posto in Start deve avere per guardare Target.

\paragraph{Un falso inizio.} Il docente parte con un esempio sbagliato (lo Start e un
Target ottenuto sommando 360, poi 359, alla posizione), e si ferma: ``scusate, ragazzi,
perdonatemi: il Find Look at Rotation lo facciamo per un esercizio che era diverso.
Ho sbagliato io, perdonate, perché non ho gli appunti.'' L'esercizio vero: far ruotare
il mulino \textbf{in direzione del character}, ``un mulino molto magico'' che si gira
verso il giocatore.

\paragraph{Il grafo.} I due punti sono:
\begin{itemize}
  \item lo \textbf{Start}: la posizione del mulino, con \textbf{Get Actor Location}
        (senza \emph{target}, cioè il Blueprint stesso);
  \item il \textbf{Target}: la posizione del giocatore, con \textbf{Get Player
        Character} e, da questo, Get Actor Location.
\end{itemize}
Il Find Look at Rotation ``ci darà il vettore verso la posizione del nostro player'',
o meglio la rotazione che lo punta.

\paragraph{RInterp To: l'interpolazione.} ``Possiamo usare anche un altro blocco, che ci
darà un'\textbf{interpolazione}.'' Si chiama \textbf{RInterp To} (R come
\emph{rotator}). ``Qualcuno sa cosa vuol dire interpolazione? No.'' ``È letteralmente
il passaggio tra due punti'': da un valore A a un valore B. I nodi di questo tipo
(\emph{Lerp}, \emph{VLerp}, \emph{Interp}) sono ovunque: ``avete presente le animazioni?
Un'animazione è un passaggio da un punto A a un punto B'', e di base è fatta a livello
di valori, perché ruotare un braccio è cambiare i valori di mano, braccio, gamba. Il
nodo ha quattro ingressi:
\begin{itemize}
  \item \textbf{Current}, la rotazione attuale del mulino (Get Actor Rotation, o la
        rotazione relativa del corpo);
  \item \textbf{Target}, la rotazione appena trovata da Find Look at Rotation;
  \item \textbf{Delta Time}, che ``ci dice in quanto tempo'' si fa l'interpolazione;
  \item \textbf{Interp Speed}, ``la velocità in cui lo fa''.
\end{itemize}

\paragraph{Solo lo yaw.} ``Però una cosa importante, ragazzi: non vogliamo che si muova
tutto sulla stessa asse.'' Il mulino deve girare ``al massimo così'', in orizzontale,
non inclinarsi verso l'alto o verso il basso. ``Quale asse dovrò influenzare?''
``La Z?'' ``Esatto, perfetto.'' Si fa \textbf{Split Struct Pin} (tasto destro sul pin)
sull'uscita di Find Look at Rotation e sull'ingresso Target, e si collega solo la Z,
lo \emph{yaw}. Il risultato dell'interpolazione si applica con un nodo di rotazione
(nella trascrizione, \emph{Set Relative Rotation})\footnote{La trascrizione dice ``set
relative rotation di pala 1'', ma a girare verso il player è tutto il mulino. Con ogni
probabilità il nodo è un Set Relative Rotation sul corpo (Body), a cui sono attaccate le
pale, o un Set Actor Rotation sull'intero Blueprint: il risultato visibile è lo
stesso.}, collegato all'Event Tick (figura~\ref{ue4:fig:lookat}).

\paragraph{L'offset di 90 gradi.} In Play il mulino si gira verso il player, ma di
lato. ``Possiamo mettere più 90'': si somma \textbf{90} alla Z prima di collegarla, ``in
modo tale che sia\ldots ecco''. Ora il mulino guarda il giocatore. ``Quello che sta
facendo è continuare a calcolare, visto che lo fa sull'Event Tick, il vettore che
indirizza in un punto, che in questo caso è il player, in modo continuo.'' ``Questo è un
blocco che serve moltissimo, ragazzi'': serve ogni volta che bisogna fare ``delle robe a
\emph{target}'', e servirà per la torretta. ``Vedete che segue continuamente la nostra
torretta? Questo è fatto con il Find Look at Rotation.''

\begin{figure}[htbp]
\centering
\begin{tikzpicture}[font=\small,
  ue4node/.style={draw, thick, rounded corners=4pt, align=center, fill=gray!8,
                  minimum height=0.9cm},
  ue4event/.style={ue4node, fill=red!12},
  ue4exec/.style={-{Stealth[length=2.5mm]}, very thick, black!80},
  ue4flt/.style={-{Stealth[length=2mm]}, thick, green!50!black},
  ue4rot/.style={-{Stealth[length=2mm]}, thick, violet!70!black},
  ue4vecw/.style={-{Stealth[length=2mm]}, thick, yellow!50!black},
  ue4obj/.style={-{Stealth[length=2mm]}, thick, blue!60!black}]
  \node[ue4event, text width=2.0cm] (tick) at (0,0)
    {\textbf{Event Tick}\\\footnotesize Delta Seconds};
  \node[ue4node, text width=2.0cm] (self) at (0,-1.9)
    {\textbf{Get Actor}\\\textbf{Location}\\\footnotesize (self)};
  \node[ue4node, fill=blue!8, text width=2.0cm] (pc) at (0,-3.9)
    {\textbf{Get Player}\\\textbf{Character}};
  \node[ue4node, text width=2.0cm] (pcl) at (2.8,-3.9)
    {\textbf{Get Actor}\\\textbf{Location}};
  \node[ue4node, text width=2.1cm, minimum height=1.4cm] (look) at (6.0,-2.6)
    {\textbf{Find Look at}\\\textbf{Rotation}};
  \node[ue4node, text width=1.2cm] (add) at (8.7,-2.6) {\textbf{$+$ 90}\\
    \footnotesize su Z};
  \node[ue4node, text width=2.0cm] (cur) at (8.7,-4.5)
    {\textbf{Get Actor}\\\textbf{Rotation}};
  \node[ue4node, text width=2.6cm, minimum height=2.2cm] (interp) at (11.5,-2.9)
    {\textbf{RInterp To}\\[2pt]\footnotesize Current\\\footnotesize Target (Z)\\
    \footnotesize Delta Time (Tick)\\\footnotesize Interp Speed};
  \node[ue4node, text width=1.5cm, minimum height=1.3cm] (set) at (14.3,0)
    {\textbf{Set}\\\textbf{Rotation}\\\footnotesize (Z)};
  \draw[ue4exec] (tick.east) -- (set.west |- tick.east);
  \draw[ue4vecw] (self.east) -- node[above, font=\scriptsize]{Start} (4.3,-1.9)
    |- ([yshift=3mm]look.west);
  \draw[ue4obj] (pc.east) -- (pcl.west);
  \draw[ue4vecw] (pcl.east) -- node[below, font=\scriptsize, pos=0.6]{Target}
    (4.5,-3.9) |- ([yshift=-3mm]look.west);
  \draw[ue4flt] (look.east) -- node[above, font=\scriptsize]{Z} (add.west);
  \draw[ue4flt] (add.east) -- ++(0.5,0) |- ([yshift=2mm]interp.west);
  \draw[ue4rot] (cur.east) -- ++(0.5,0) |- ([yshift=-5mm]interp.west);
  \draw[ue4rot] (interp.east) -- ++(0.4,0) |- ([yshift=-4mm]set.west);
\end{tikzpicture}
\caption{Il mulino che guarda il player. A ogni \emph{frame}, Find Look at Rotation
calcola la rotazione che dal mulino (Start) punta al giocatore (Target); se ne tiene
solo lo \emph{yaw} (Z), gli si sommano 90 gradi per compensare l'orientamento della
\emph{mesh}, e RInterp To fa avvicinare gradualmente la rotazione attuale a quella
cercata. Il Delta Time può venire dall'Event Tick, come qui (un cavo non disegnato per
chiarezza), o essere un numero fisso (a lezione, 1); si veda l'approfondimento. Schema ricostruito.}
\label{ue4:fig:lookat}
\end{figure}

\paragraph{``A me non va.''} Seguono tre studenti in condivisione.
\begin{itemize}
  \item Al primo il mulino resta fermo: ``ti manca il Delta Time e
        l'Interp Speed''. Li riempie (``metti uno e due''), compila (``non ho
        compilato, vabbè'') e aggiunge l'offset: ``tira fuori dalla Z, fai più, metti
        90''. ``Perfetto, adesso ti segue.''
  \item Al secondo manca anche a lui il Delta Time: ``metti uno''.
  \item Il terzo, collegando il nodo a Find Look at Rotation, si ritrova un pin
        diviso che non voleva: ha preso \textbf{Get Actor Rotation} invece di
        \textbf{Get Actor Location}. Poi non riesce a collegare l'uscita alla Z:
        ``lo devi splittare: tasto destro sul pinnetto, quello violaceo, quello
        lilla, Split Struct Pin''. E infine l'offset: da Return Value di Find Look at
        Rotation, la Z, un ``più'', 90, e al Target.
\end{itemize}

\paragraph{Approfondimento: che cosa calcola Find Look at Rotation.} Il nodo fa una cosa
semplice: sottrae la posizione di partenza da quella di arrivo, ottenendo il vettore
direzione $\vec d = T - S$ (la sottrazione della sezione~\ref{ue4:sec:vettori}), e
trasforma quella direzione in un \emph{rotator} con \emph{roll} nullo. Per la sola
rotazione orizzontale, quella che interessa al mulino e alla torretta, il conto è
\[
  \text{yaw} \;=\; \operatorname{atan2}(d_y,\ d_x) \quad \text{(in gradi)},
\]
cioè l'angolo, misurato nel piano XY a partire dall'asse X, sotto cui il bersaglio è
visto. Il risultato presuppone che la ``faccia'' dell'oggetto guardi lungo il suo asse
\textbf{X} locale, che in Unreal è la direzione ``in avanti'' (il \emph{Forward
Vector}). Se la \emph{mesh} è stata modellata guardando lungo la Y, come le pale del
mulino della lezione, la rotazione calcolata è sfasata di 90 gradi, e il ``più 90'' del
docente non fa che correggere questo sfasamento. Lo stesso problema, come si vedrà, si
ripresenta con la torretta.

\paragraph{Approfondimento: come funziona RInterp To.} A ogni chiamata, RInterp To
calcola quanto manca per arrivare al Target e ne percorre una frazione, pari
(semplificando) a
\[
  \alpha \;=\; \min\bigl(1,\ \text{Delta Time} \times \text{Interp Speed}\bigr),
  \qquad
  \text{nuova rotazione} = \text{Current} + \alpha \cdot (\text{Target} - \text{Current}),
\]
con la differenza presa sempre ``dal lato più corto'' (da 350 a 10 gradi si passa per
0, non per 180). Ne seguono tre cose utili. Primo: il movimento è veloce quando il
bersaglio è lontano e rallenta avvicinandosi, perché a ogni \emph{frame} si percorre
una frazione di una distanza sempre più piccola; è l'effetto ``morbido'' che si vede in
Play. Secondo: il nodo è pensato per ricevere come \textbf{Delta Time} proprio i
\emph{delta seconds} (dall'Event Tick o da \emph{Get World Delta Seconds}) e come
\textbf{Interp Speed} un numero intorno a 1--10: così la velocità non dipende dal
\emph{frame rate}. Con un Delta Time fisso, come l'1 messo a lezione, la frazione
percorsa a ogni \emph{frame} è sempre la stessa, e il mulino gira più veloce su un PC
che fa più \emph{frame}. Terzo: se il prodotto supera 1, la frazione è 1 e la rotazione
\textbf{scatta} subito al Target; e se l'Interp Speed è 0 o negativo, Unreal restituisce
direttamente il Target. Con Delta Time 1 e Interp Speed 2 il prodotto è 2: il mulino,
in pratica, ``snappa''. Il fenomeno tornerà con la torretta
(sezione~\ref{ue4:sec:problemi}).

% ---------------------------------------------------------------------
\section{Il Clamp, il debug con Print String e la pausa}
\label{ue4:sec:clamp}

``Abbiamo fatto il nostro bellissimo Find Look at Rotation. Potremmo anche
\textbf{clamparlo}.'' ``Si ricorda il Clamp qualcuno di voi?'' ``È quello per dare un
valore minimo'', e massimo. ``Esatto, bravissimo.'' Lo scopo: limitare la rotazione del
mulino, in modo che non possa girarsi del tutto.

\paragraph{Il grafo.} Si splitta il pin con Split Struct Pin, si inserisce un
\textbf{Clamp (float)}, ``perché è un \emph{float}'', sulla Z, e la si ricompone con un
\textbf{Make Rotator} prima di collegarla. Il primo tentativo è un intervallo stretto,
da $-25$ a $25$ gradi, ma qualcosa non torna: i valori in gioco sono molto più grandi
(``mi dice: Pala1 è un botto di valori'').

\paragraph{Il debug con Print String.} ``Questo come ho fatto, ragazzi? Questo è un modo
per \textbf{debuggare} le cose.'' Non sapendo di quanto ruotasse il mulino e quali
limiti dargli, il docente collega un \textbf{Print String} al valore e lo stampa a
schermo mentre si muove attorno al mulino: ``abbiamo stampato letteralmente il
valore''. I numeri che compaiono vanno più o meno da $-160$ a $160$. Il docente prova
quei limiti, poi li allarga: ``snappa un pochino. Facciamo duecento massimo. Molto
meglio. Non dovrebbe andare oltre''\footnote{La trascrizione di questo passaggio è
molto confusa (``facciamo zero, no, mettere meno duecento\ldots{} facciamo duecento
massimo''), e non permette di ricostruire con certezza i valori finali del Clamp.
Quello che conta è il metodo: stampare i valori, osservarli, scegliere i limiti.}.

\paragraph{Approfondimento: Clamp e gli angoli.} Il Clamp restituisce il valore stesso
se è compreso tra Min e Max, altrimenti il limite più vicino:
$\operatorname{clamp}(v, m, M) = \min(\max(v, m), M)$. Con gli angoli c'è una
sottigliezza che spiega i valori ``strani'' della lezione: lo \emph{yaw} restituito da
Find Look at Rotation sta tra $-180$ e $180$, e passando ``dietro'' al mulino salta da
$180$ a $-180$; sommandoci 90 gradi di offset, l'intervallo si sposta (da $-90$ a
$270$). Un limite come $-25 \ldots 25$, applicato a questi numeri, non corrisponde a
``25 gradi a destra e a sinistra della direzione di riposo'', ma a una fetta del cerchio
che dipende da dove sta il player e da come è orientato il mulino nel livello. Il
modo pulito è ragionare sulla differenza tra l'angolo voluto e quello di riposo,
riportarla tra $-180$ e $180$ (Unreal ha il nodo \emph{Normalize Axis}, o \emph{Delta
(Rotator)} per due \emph{rotator}), e solo allora fare il Clamp. Il metodo del docente,
stampare i numeri prima di decidere, resta comunque la prima cosa da fare.

\paragraph{La pausa.} ``Ci siete? C'è qualche cosa che non avete capito? C'è qualche
blocco che vi risulta ancora un pochino nebuloso?'' Tutto chiaro. ``Sono le 17 e 28.
Facciamo pausa, così poi riprendiamo e cominciamo a fare la nostra bellissima
torrettina. Ci vediamo tra 20 minuti.''

% ---------------------------------------------------------------------
\section{Dopo la pausa: la torretta in tre pezzi}
\label{ue4:sec:torretta}

``Eccoci qua. Allora, ragazzi, adesso andremo a fare una bellissima torrettina.'' Il
docente mostra la versione finita. Attorno alla torretta c'è una \textbf{sfera} disegnata
con linee rosse: è il suo raggio d'azione. Quando il giocatore ci entra, la torretta
``fa lo stesso processo che abbiamo visto'', cioè ruota verso il \emph{character}, e
``in questo caso spara dei proiettili''. Non ci sono effetti speciali (``non abbiamo
Niagara''\footnote{La trascrizione dice ``BevoFX'': con ogni probabilità
\emph{Niagara VFX}, il sistema di effetti visivi di Unreal Engine, di cui Davide aveva
parlato nella lezione precedente.}), ma si costruisce ``la logica che ci servirà con la
torretta che avete \emph{mergiato}''. Si vedono anche delle linee che partono dal
cannone: sono un \textbf{line trace}, ``che permette a noi di avere una corrispondenza
dipendente da cosa incontra''. Oggi si fa la rotazione; proiettili e \emph{line trace}
saranno la prossima volta.

\paragraph{Il Blueprint.} ``La prima cosa che ovviamente dobbiamo fare è creare un
nuovo Blueprint, sempre \textbf{Actor}'', chiamato \textbf{BP\_Turret}. Ma come hanno
composto la torretta gli studenti, con Davide? ``Abbiamo preso i tre pezzi forniti e li
abbiamo uniti'' con \textbf{Merge Actors}: una \emph{mesh} unica. ``Allora facciamo una
cosa: la facciamo a \textbf{tre pezzi}.'' Per far girare il supporto e il cannone
indipendentemente dalla base servono tre componenti distinti, quindi si torna ai tre
pezzi importati con Davide: tre \textbf{Static Mesh} nel Blueprint, una per la
\textbf{base}, una per il supporto (``l'\emph{holder}, chiamiamolo'', poi rinominato
\textbf{Pivot}) e una per il \textbf{cannone} (``la \emph{gun}'', poi
\textbf{Cannon}).

\paragraph{Chi segue chi.} I tre pezzi vanno ``uno dentro l'altro''. Il docente parte
con l'idea che il supporto segua il cannone, poi si corregge: ``fate il contrario, che
è meglio: fate che la \emph{gun} segua l'\emph{holder}''. La gerarchia giusta è quindi:
la \textbf{base} come radice, il \textbf{Pivot} figlio della base, il \textbf{Cannon}
figlio del Pivot (figura~\ref{ue4:fig:turretbp}). Così, quando il Pivot ruota
orizzontalmente, il cannone gira con lui; e un giorno il cannone potrà alzarsi e
abbassarsi da solo, senza trascinarsi dietro il supporto.

\paragraph{Un trucco: Browse to Asset.} Per assegnare le \emph{mesh} ai componenti non
serve cercarle a mano nel Content Browser ``con l'elenchino e tutto quanto''. ``Andate
qua sull'\emph{asset} che avete in scena, cliccate tasto destro e fate \textbf{Browse to
Asset}'': il Content Browser salta direttamente alla posizione del file, e da lì lo si
può \textbf{trascinare} (\emph{drag and drop}) nello slot \emph{Static Mesh} del
componente. In alternativa, come sempre, si clicca sullo slot e si scrive il nome:
\emph{turret base}, \emph{turret pivot}, \emph{cannon}. Poi si aggiustano le posizioni
(``un buchettino sta un po' più giù'') finché i pezzi combaciano.

\paragraph{Un problema di Add.} Una studentessa, premendo Add, vede solo la \emph{mesh}
intera della torretta (\emph{SM\_\allowbreak Torretta}) e nessun'altra opzione. In condivisione si
prova a compilare e a riaprire il Blueprint, e alla fine il componente si riesce ad
aggiungere: ``va benissimo così, abbiamo risolto. A volte Unreal può fare dei giochi
simpatici.''

% ---------------------------------------------------------------------
\section{Il raggio d'azione: la Sphere Collision}
\label{ue4:sec:sfera}

``Adesso dobbiamo creare quello che abbiamo visto qua: vedi questo cerchio rosso intorno
alla torretta? Questo è una \textbf{Sphere Collision}.'' È il \textbf{range}: entrandoci
dentro, la torretta si attiva.

\paragraph{Aggiungerla.} In BP\_Turret si aggiunge, tra i componenti, una
\textbf{Sphere Collision}, ``che tranquillamente può stare anche esterna alla base'',
cioè non deve essere figlia di un pezzo che ruota. Appena aggiunta non si vede quasi,
``semplicemente perché è molto piccola''. Cliccando su \emph{Sphere}, nei dettagli c'è
lo \textbf{Sphere Radius}: di default è 32; il docente lo porta a \textbf{800}. C'è anche
la \textbf{Line Thickness}, lo spessore delle linee con cui la sfera è disegnata
nell'editor, ``che serve magari solo per vederla un po' più spessa''.

\paragraph{Hidden in Game.} In Play, però, la sfera non si vede. ``Di base, questa è una
\textbf{collisione}'', una cosa interna agli \emph{asset}, come le collisioni della
\emph{mesh} viste con Davide: esiste per la logica, non per essere guardata. Per
vederla mentre si prova il gioco (``per avere facilità nel vedere quando entriamo e
usciamo'') si cerca nei dettagli \textbf{Hidden in Game} e lo si toglie. ``Parlando di
\emph{rendering}, Hidden in Game dice: so che logicamente esiste all'interno del mondo,
ma se non è \emph{hidden in game}, la renderizzo pure.'' Tolta la spunta, il valore è
\emph{false}, ``quindi non è nascosto in gioco'', e in Play ecco il range. È la stessa
opzione usata nella seconda lezione, al contrario, per rendere invisibile il soffitto
della vaschetta.

\paragraph{Ripetere i passaggi.} Uno studente ha perso questa parte mentre cercava i
pezzi della torretta, e chiede di rivedere come si mettono le \emph{mesh} nei tre
componenti. Il docente lo rifà: si clicca sul componente, si scrive il nome della
\emph{mesh} nello slot Static Mesh (\emph{turret base}, \emph{turret pivot}), si
rinominano i componenti in modo che siano coerenti (\textbf{Pivot}, \textbf{Cannon},
``come a me è chiamata Cannon nella nostra \emph{mesh}'').

% ---------------------------------------------------------------------
\section{Gli eventi dei componenti e l'overlap}
\label{ue4:sec:overlap}

``Adesso andiamo a vedere un evento in particolare.'' Oltre agli eventi già visti, che
appartengono all'intero actor (l'Event Begin Play, l'Event Actor Begin Overlap,
l'Event Tick), ``possiamo richiamare degli eventi anche all'interno dei
\textbf{componenti}''. Selezionando \emph{Sphere} e scorrendo in fondo ai dettagli, c'è
una sezione \textbf{Events} piena di pulsanti: \emph{On Component Begin Overlap},
\emph{On Component End Overlap}, \emph{On Component Wake}, \emph{On Component Sleep},
gli eventi del cursore e altri. ``Tutta una serie di eventi'' che si possono usare
``per far sì che avvenga qualcosa in un preciso momento''. Qui la torretta deve partire
quando il player entra nella sfera: si clicca sul \textbf{+} di \textbf{On Component
Begin Overlap}, e nel grafo compare l'evento, già legato alla Sphere.

\paragraph{Che cos'è un overlap.} ``Sapete cosa si intende per \emph{overlap}?''
Uno studente: ``tipo appena si passa un confine, qualcosa del genere''. ``In un certo
senso.'' Alla lavagna (``perdonate le mie doti molto brutte di disegno'') il docente
disegna la sfera con la torretta al centro: un \emph{overlap} avviene quando il player,
``o un qualsiasi actor'', \textbf{tocca o entra} in quel volume. Un actor fuori non sta
sovrapponendosi; uno dentro sì. E uno a cavallo del bordo? ``Sta overlappando o non sta
overlappando?'' ``Sì.'' ``Qualcun altro ha un'opinione differente?'' Sta
sovrapponendosi, ``perché anche un qualsiasi poligono da cui è formato è parte
integrante di esso'': basta un pezzo (figura~\ref{ue4:fig:overlap}). Ci sono quindi tre
situazioni: non si sovrappone, \textbf{comincia} a sovrapporsi (Begin Overlap),
\textbf{smette} di sovrapporsi (End Overlap, ``quando ci esce''). Per ora serve solo la
prima: ``vogliamo che quando lui entra in questo volume la torretta faccia tipo: ti ho
visto, mo' comincio a ruotare verso di te''.

\begin{figure}[htbp]
\centering
\begin{tikzpicture}[font=\small,
  ue4arr/.style={-{Stealth[length=2.5mm]}, thick}]
  \draw[red!70!black, thick, fill=red!4] (0,0) circle (2.4);
  \node[red!70!black, font=\scriptsize] at (0,2.65) {Sphere Collision (raggio 800)};
  \draw[thick, fill=orange!35] (-0.35,-0.3) rectangle (0.35,0.3);
  \draw[thick, fill=orange!20] (0.35,-0.08) rectangle (0.9,0.08);
  \node[font=\scriptsize] at (0,-0.55) {torretta};
  % fuori
  \draw[thick, fill=blue!15] (-4.3,0.6) circle (0.32);
  \node[font=\scriptsize, align=center] at (-4.3,-0.2) {fuori:\\nessun overlap};
  % a cavallo
  \draw[thick, fill=blue!15] (1.75,1.65) circle (0.32);
  \node[font=\scriptsize, align=left, anchor=west] at (2.3,1.9) {a cavallo del bordo:\\
    \textbf{overlap} (basta un pezzo)};
  % dentro
  \draw[thick, fill=blue!15] (-1.0,-1.25) circle (0.32);
  \node[font=\scriptsize, align=center] at (-1.0,-1.95) {dentro: overlap};
  % frecce degli eventi
  \draw[ue4arr, green!45!black] (-3.9,0.45) -- (-2.55,0.25);
  \node[green!45!black, font=\scriptsize, align=center] at (-3.4,1.45)
    {entra:\\\textbf{Begin Overlap}};
  \draw[ue4arr, orange!80!black] (1.95,-1.55) -- (3.2,-2.3);
  \node[orange!80!black, font=\scriptsize, align=left, anchor=west] at (3.3,-2.3)
    {esce: \textbf{End Overlap}};
\end{tikzpicture}
\caption{Le situazioni dell'\emph{overlap} disegnate a lezione. Un actor si sovrappone
alla sfera appena una parte qualsiasi della sua forma di collisione ne entra; Begin
Overlap scatta una volta, nel momento in cui entra, End Overlap una volta, quando ne
esce del tutto. Schema ricostruito.}
\label{ue4:fig:overlap}
\end{figure}

\paragraph{Other Actor.} ``Come facciamo a capire quale actor ci dà indietro?'' L'evento
ha una serie di uscite: \emph{Overlapped Component}, \textbf{Other Actor}, \emph{Other
Comp}, \emph{Other Body Index}, \emph{From Sweep}, \emph{Sweep Result}. ``Sono tutti
input che sono risultanti dall'overlap'', ``in parole povere'' le informazioni su chi si
è sovrapposto a che cosa. Per mostrarlo, il docente collega Other Actor a un
\textbf{Print String} (Unreal inserisce da solo un \emph{Get Display Name}, per
trasformare l'actor in testo), va in Play ed entra nella sfera: in alto a sinistra
compare \emph{BP\_ThirdPersonCharacter\_0}\footnote{La trascrizione dice ``BP ha perso
un character 0''. Il suffisso \_0 è il numero che Unreal assegna a ogni istanza di una
classe presente nel livello.}. ``Mi ridà cosa sta overlappando, che in questo caso è il
\emph{character}.'' Esistono anche gli \emph{overlapped component} e gli \emph{other
component}, ma quello che interessa è l'Other Actor, ``perché il nostro
\emph{character} è un actor''.

\paragraph{Approfondimento: perché l'overlap funziona ``da solo''.} Perché Unreal generi
un evento di \emph{overlap} servono tre condizioni: entrambi i componenti devono avere
attivo \emph{Generate Overlap Events}; le loro risposte di collisione reciproche devono
essere \emph{Overlap} (e non \emph{Ignore}); e almeno uno dei due deve muoversi. La
Sphere Collision nasce con il profilo di collisione \emph{OverlapAllDynamic}, che si
sovrappone a tutto ciò che si muove, e la capsula del \emph{character} è un oggetto
dinamico di tipo \emph{Pawn}: per questo, nella lezione, non c'è stato bisogno di
toccare nulla. Se un giorno un evento di \emph{overlap} non scatta, è la prima cosa da
controllare, nei dettagli di entrambi i componenti, sezione \emph{Collision}. La seconda
edizione di Romero e Sewell elenca proprio gli eventi di collisione (``eseguiti quando
due actor collidono o si sovrappongono'') tra le famiglie principali di eventi, accanto
a quelli di input e all'Event Begin Play.

% ---------------------------------------------------------------------
\section{Is Valid: esiste davvero?}
\label{ue4:sec:isvalid}

``Vogliamo essere sicuri che il nostro actor in questione sia proprio il player. Perché
un actor può essere anche un cubo, volendo.'' Prima, però, un controllo più semplice. Si
tira fuori da Other Actor un \textbf{Is Valid}: ``ce ne avete tanti, ma quello che ci
interessa è quello col \textbf{punto interrogativo}'', la \emph{macro} con due uscite di
esecuzione, \textbf{Is Valid} e \textbf{Is Not Valid}.

\paragraph{Che cosa controlla.} ``Is Valid, statemi a sentire, ci dà una risultante: è
valido o non è valido?'' Qui chiede: ``hai overlappato qualcosa? È vero, esiste questa
cosa?'' In generale, Is Valid serve a dire se esiste o no qualcosa, ``che può essere
anche una \textbf{variabile}, quindi una variabile vuota''.

\paragraph{L'esempio della variabile.} Per farlo vedere il docente costruisce un piccolo
esperimento:
\begin{enumerate}
  \item da Other Actor crea una variabile con \textbf{Promote to Variable};
  \item mette un Is Valid su quella variabile, all'inizio dell'evento di \emph{overlap};
  \item sul ramo \textbf{Is Not Valid} assegna alla variabile l'actor appena
        sovrapposto (il \emph{Third Person Character});
  \item sul ramo \textbf{Is Valid} stampa con un Print String ``valore registrato'' e il
        nome contenuto nella variabile.
\end{enumerate}
Prima entrata nella sfera: non succede nulla di visibile, ``è normale''. Il Blueprint ha
chiesto: ``questa variabile ha qualcosa dentro? No, non c'è niente. Quindi non è
valida'', e allora le ha assegnato l'actor. Si esce e si rientra: ``ah, me l'hai fatto
prima, quindi questa cosa c'ha qualcosa dentro'', e il ramo è Is Valid: compare
\emph{valore registrato: BP\_ThirdPersonCharacter\_0}. ``Qualcuno ha qualche dubbio
riguardo all'Is Valid? Abbastanza chiaro?'' L'esperimento si cancella: ``ci serve
solamente a capire cosa fa l'Is Valid''.

\paragraph{Approfondimento: la regola di Romero e Sewell.} Il quarto capitolo di
\emph{Blueprints Visual Scripting for Unreal Engine 5}, dedicato alla comunicazione tra
Blueprint, mette in evidenza una nota: ``usate sempre la macro Is Valid prima di
eseguire una funzione usando una variabile che contiene una referenza a un oggetto. La
variabile può non essere valida per diversi motivi, ed eseguire una funzione con una
variabile non valida causerà un errore durante l'esecuzione''. I motivi tipici sono
due: la variabile non è mai stata riempita (come nella prima entrata dell'esempio), o
l'oggetto a cui puntava è stato distrutto nel frattempo. In entrambi i casi un nodo che
prova a leggere la posizione di ``nessuno'' produce nel registro il classico errore
\emph{Accessed None}. Il libro usa lo stesso schema in molti punti: un Is Valid prima di
distruggere un actor, prima di applicare un danno a un bersaglio, prima di leggere un
salvataggio (se non è valido, lo si crea). La torretta lo userà due volte, nella
prossima sezione e nell'evento Move Turret.

% ---------------------------------------------------------------------
\section{I tag e Actor Has Tag}
\label{ue4:sec:tag}

Is Valid dice che \emph{qualcosa} è entrato. Per sapere \emph{chi}, si tira fuori da
Other Actor un \textbf{Actor Has Tag}, ``un blocco che va a esaminare una cosa chiamata
\textbf{tag}''. I tag ``sono dei \emph{label}, delle \textbf{etichette} che noi possiamo
mettere all'interno di vari \emph{asset} del gioco''.

\paragraph{Mettere il tag al player.} Si apre il \textbf{Third Person Character} e,
``non dobbiamo fare niente di particolare'', si scrive \emph{tag} nella ricerca dei
dettagli: compare \textbf{Actor} $\rightarrow$ \textbf{Tags}, un elenco vuoto con un
\textbf{+}. Si aggiunge un elemento e lo si chiama \textbf{Player}. ``Cosa abbiamo
fatto? Abbiamo letteralmente etichettato il player con un tag Player'': è come avere un
cartellino di riconoscimento sulla schiena, con scritto sopra ``tu sei il player''.

\paragraph{Il controllo.} Tornati nella torretta, Actor Has Tag chiede: il
\emph{target}, cioè l'actor sovrapposto (Other Actor), ha il tag \textbf{Player}? Il
nodo restituisce un valore \textbf{booleano}: ``sì, questo actor che ho overlappato ha
il tag Player'', e si va nel \textbf{True} di un \textbf{Branch}; ``se no, invece andrà
in False''. Il docente lo mostra con due Print String, ``Player trovato'' sul ramo True e
``Player non trovato'' sul ramo False: entrando nella sfera, ``Player trovato''.

\paragraph{Attenzione all'ortografia.} ``Non lo fate, ragazzi, faccio solo vedere per
esempio.'' Il docente modifica il tag, togliendo la R finale, e rientra: ``Player non
trovato''. ``Questo vi fa vedere anche che dovete essere precisi nel dire qual è il
tag.'' Il consiglio: ``fate Ctrl+C, Ctrl+V del tag che avete scritto'', copiandolo dal
\emph{character} nel nodo invece di riscriverlo a mano.

\paragraph{Approfondimento: tag, Cast o confronto?} Ci sono almeno tre modi per
rispondere alla domanda ``l'actor che è entrato è il player?''. Il primo è il
\textbf{tag}, quello della lezione: semplice, e indipendente dalla classe (lo stesso tag
\emph{Player} può stare su personaggi diversi, e un tag \emph{Enemy} su nemici di classi
diverse); esiste anche \emph{Component Has Tag}, per etichettare singoli componenti. Il
secondo è il \textbf{Cast To} (per esempio \emph{Cast To BP\_ThirdPersonCharacter}),
che verifica la classe e, se riesce, restituisce l'actor già ``convertito'', con
accesso alle sue variabili e funzioni. Il terzo è il confronto diretto tra oggetti: nel
capitolo sull'intelligenza artificiale della seconda edizione di Romero e Sewell, un
nemico che pattuglia usa un \emph{Event Actor Begin Overlap} seguito da un Branch su un
nodo \textbf{Equal (Object)}, che controlla se l'Other Actor è proprio il punto di
pattuglia salvato in una variabile. Una curiosità sui tag: sono di tipo \emph{Name}, e
in Unreal il confronto tra \emph{Name} \textbf{non distingue maiuscole e minuscole},
per cui \emph{player} e \emph{Player} risultano uguali; quello che conta, come ha
mostrato il docente, è che le lettere siano le stesse. Copiare e incollare il tag
resta comunque la via più sicura.

% ---------------------------------------------------------------------
\section{La variabile del bersaglio e i Custom Event}
\label{ue4:sec:custom}

Ora che il player è stato trovato, bisogna ricordarselo: ``andiamo a crearci una
variabile in cui mettiamo il nostro player che è stato trovato'', di tipo
\textbf{Actor}. Il modo più veloce: tirare fuori un cavo da Other Actor e scegliere
\textbf{Promote to Variable}. Si crea un nodo \emph{Set} di una variabile nuova, che si
rinomina \textbf{PlayerTarget}. Per rinominare, ``o vai qua e fai \textbf{F2}, che è più
facile, oppure fai tasto destro e Rinomina''.

\paragraph{Un Custom Event.} ``Manca mezz'ora'', e il docente si chiede se si riuscirà a
finire (``questo lo continueremo la prossima volta''), ma va avanti. ``Adesso qua dentro
andremo a fare un \textbf{Custom Event}.'' Si scrive \emph{custom event} nel grafo, si
sceglie \emph{Add Custom Event} e gli si dà un nome: \textbf{TrackTargetTimer}. ``Come
potete vedere, questo bellissimo evento è un Custom Event: è uguale agli altri'', come
il Begin Overlap, il Tick o il Begin Play, ``ma è un evento che abbiamo creato noi'',
``che possiamo richiamare quando vogliamo, con delle regole specifiche''. Gli eventi di
Unreal sono ``prefissati'': l'\emph{overlap} scatta perché l'editor di Unreal lo
controlla già per conto suo. I Custom Event, invece, partono quando li si chiama.

\paragraph{Chiamarlo.} Per mostrarlo, il docente collega al Custom Event un Print String
che stampa ``Hello'' e lo chiama dal ramo del Begin Overlap, subito dopo il Set di
PlayerTarget: si scrive il nome dell'evento (\emph{TrackTargetTimer}) e compare il nodo
che lo invoca. In Play, entrando nella sfera: ``Hello''. ``Ho creato un evento
customizzato che richiamo quando entro nella Sphere Collision. Compreso?'' ``Sappiate
che di Custom Event ne potete creare a milioni'': basta scrivere \emph{custom event} e
dare un nome.

\paragraph{Approfondimento: eventi, funzioni e macro.} La seconda edizione di Romero e
Sewell ricorda che nell'Event Graph ogni evento di Unreal può comparire una sola volta,
mentre i \textbf{Custom Event} si possono creare a volontà; e offre una tabella per
scegliere tra \emph{macro}, \emph{funzioni} ed \emph{eventi}. In sintesi: tutti e tre
accettano parametri in ingresso; solo macro e funzioni restituiscono valori; le funzioni
e gli eventi possono essere chiamati da altri Blueprint, le macro no; e solo macro ed
eventi possono contenere azioni \emph{latenti}, cioè che aspettano (come un
\emph{Delay}), mentre le \emph{Timeline} stanno solo negli eventi. Il capitolo sui
vincoli di gioco della stessa edizione usa proprio un Custom Event, \emph{Sprint Drain},
richiamato da un timer, per consumare la stamina del giocatore mentre corre: ``un Custom
Event permette di attivare i nodi collegati a lui ogni volta che un altro nodo lo
chiama; così gruppi di nodi dello stesso Blueprint possono comunicare tra loro anche
quando non sono collegati direttamente''. È lo schema esatto della torretta.

% ---------------------------------------------------------------------
\section{Il timer: Set Timer by Event}
\label{ue4:sec:timer}

``Cosa vogliamo fare con questo Custom Event?'' Durante tutto il tempo, ``come un Event
Tick'', la torretta deve cercare la posizione verso cui ruotare per seguire il player.
``Lo potremmo fare tranquillamente nell'Event Tick, volendo, ma possiamo fare una cosa
interessante: possiamo fare un \textbf{timer}. L'abbiamo chiamato timer apposta.''

\paragraph{Il nodo.} Dal Custom Event si tira fuori \textbf{Set Timer by Event}, un
blocco ``che ci permette di far partire un timer che ci chiede ogni quanto''. I suoi
parametri, così come li imposta il docente:
\begin{itemize}
  \item \textbf{Time}: ogni quanti secondi scatta. ``Vogliamo che lo faccia allo stesso
        tempo del Tick.'' Si potrebbe portare fin qui il cavo dei \emph{delta seconds}
        dall'Event Tick (``ma è un po' brutto''); c'è invece un nodo apposta,
        \textbf{Get World Delta Seconds}, ``un blocchetto semplicemente per comodità'',
        ``per non far sì di filare il Tick così, a cavolo'';
  \item \textbf{Looping}: spuntato, ``vogliamo che sia un timer che \emph{loopi}'';
  \item \textbf{Max Once Per Frame}: spuntato anche questo, ``massimo una volta per
        ogni \emph{frame}''. ``Ci permette almeno di dire: ok, per ogni \emph{frame} si
        calcola, tanto effettivamente si muoverà per ogni \emph{frame}. Questa è una
        cosa un po' più nel micro, però noi mettiamo \emph{max}, in modo che sia più
        ottimizzato.''
\end{itemize}

\paragraph{Quale evento?} ``Set Timer by Event: quale evento?'' Sul nodo c'è un
\textbf{quadratino rosso}, il pin \textbf{Event}, ``che guarda caso è uguale a quello del
Custom Event'': un pin di tipo \emph{delegato}. Vuol dire che lo si può tirare fuori e
creare da lì un \textbf{altro Custom Event}, quello che il timer chiamerà a ogni giro:
lo si chiama \textbf{MoveTurret}. È lì che si farà muovere la torretta.

\paragraph{Come ragiona.} Per farlo vedere, a MoveTurret si collega il solito Print
String, ``Hello''. La catena completa è: Begin Overlap, Is Valid, Branch su Actor Has
Tag, Set PlayerTarget, TrackTargetTimer; e TrackTargetTimer fa partire un timer, alla
velocità dei \emph{delta seconds}, in \emph{loop}, al massimo una volta per
\emph{frame}, che chiama MoveTurret, che stampa ``Hello''. ``Domanda: quando mi
stamperà Hello?'' ``Da quando entri in range.'' ``Ottimo, esatto.'' Appena il player
entra, Hello a ogni \emph{frame}. ``Esco fuori: si cancellano\ldots{} ah, no, scusate,
non si stanno cancellando, è vero. Perché non stoppiamo mai il timer.'' Dalla prima
entrata in poi, il timer continua all'infinito: ``ciò vuol dire che il timer a un certo
punto va disattivato; lo vedremo dopo''.

\paragraph{Approfondimento: timer o Event Tick?} Romero e Sewell presentano i timer, già
nel secondo capitolo, come un'alternativa all'Event Tick: tutto ciò che è collegato al
Tick viene eseguito a ogni \emph{frame}, e ciò che non ha bisogno di aggiornarsi così
spesso può girare a intervalli più lunghi. Nel quinto capitolo un Custom Event
\emph{ChangeDirection}, chiamato da un Set Timer by Event con Looping attivo, fa
invertire periodicamente la direzione di un bersaglio mobile; nella seconda edizione un
timer chiama ripetutamente \emph{Sprint Drain}, e un \emph{Clear Timer by Function Name}
lo ferma quando il giocatore smette di correre. Nella torretta della lezione l'intervallo
è pari ai \emph{delta seconds}, quindi il timer scatta (circa) una volta per
\emph{frame}, come un Tick: il vantaggio non è la frequenza, ma il fatto che il timer
\textbf{parte solo quando serve} (all'ingresso del player) e \textbf{si può fermare}.
Tre dettagli utili. Primo: il valore di Get World Delta Seconds viene letto una sola
volta, quando il timer viene impostato, e resta fisso anche se poi il \emph{frame rate}
cambia. Secondo: se l'intervallo di un timer è più breve di un \emph{frame}, Unreal può
eseguirlo più volte nello stesso \emph{frame} per ``recuperare''; \textbf{Max Once Per
Frame}, un'opzione delle versioni recenti di Unreal Engine 5, impedisce proprio questo.
Terzo: per fermarlo, Set Timer by Event restituisce un \textbf{Timer Handle}, una
``maniglia'' che conviene salvare in una variabile (Promote to Variable) e passare,
quando serve, a \textbf{Clear and Invalidate Timer by Handle}. Il posto naturale per
farlo è l'evento \textbf{On Component End Overlap} della sfera, insieme a
\emph{svuotare} la variabile PlayerTarget (un Set senza nulla collegato): altrimenti
l'Is Valid di MoveTurret continuerà a trovare un bersaglio valido anche dopo che il
player è uscito dal range.

% ---------------------------------------------------------------------
\section{Move Turret: la torretta che segue il player}
\label{ue4:sec:move}

``Cosa dobbiamo fare? L'abbiamo fatto prima, praticamente'': un \textbf{Find Look at
Rotation}, ``perché dobbiamo ruotare il nostro Pivot in direzione del player''.

\paragraph{Prima, un Is Valid.} ``Intanto facciamo un Is Valid. Perché? Perché diciamo:
ma ce l'ho un target?'' Si trascina nel grafo la variabile \textbf{PlayerTarget} e la si
collega a un Is Valid, in testa a MoveTurret. ``Ce l'ho un target effettivo, o non ce
l'ho più, l'ho perso, sono uscito fuori? Se ce l'ho, lo faccio e mi muovo. Se non ce
l'ho, non lo faccio'': senza questo controllo la torretta ``continuerebbe a ruotare a
cavolo''.

\paragraph{Start e Target.} ``Del Find Look at Rotation ci servono due punti nello
spazio. Quali sono?'' ``Lo Start è l'\emph{holder}, che tanto la \emph{gun} è collegata
all'\emph{holder}, e il Target è il player.'' ``Esatto.'' Per lo Start si usa
\textbf{Get Actor Location} della torretta stessa, ``la locazione generale
dell'actor''. Per il Target, il docente comincia con Get Player Character e Get Actor
Location, come per il mulino, ma uno studente fa notare che si può usare
\textbf{PlayerTarget}: ``bravo, mi piace che l'hai notato, effettivamente sarebbe stato
il secondo metodo''. Da PlayerTarget si tira fuori Get Actor Location, ``una sua
particolarità'', e la si collega al Target.

\paragraph{Lo yaw e RInterp To.} Il Pivot deve ruotare attorno alla \textbf{Z}:
``guardate un po' dove si modifica: sulla Z''. Si splitta l'uscita di Find Look at
Rotation e si fa ``quello che abbiamo fatto prima'', un \textbf{RInterp To}:
\begin{itemize}
  \item \textbf{Current}: la rotazione attuale del Pivot, con \textbf{Get Relative
        Rotation} sul componente Pivot;
  \item \textbf{Target}: splittato, con la sola Z collegata, ``perché non ci interessa
        la X e la Y'';
  \item \textbf{Delta Time}: 5;
  \item \textbf{Interp Speed}: \textbf{0,04}, ``perché, essendo che sta sui
        \emph{frame}, è molto veloce. Se voi mettete uno, questa \emph{snappa}
        diretto''.
\end{itemize}
Infine \textbf{Set Relative Rotation} sul Pivot (``prendete Pivot, Set Relative
Rotation''), con i pin splittati ``per sicurezza'', collegato al ramo Is Valid; e un
Print String per il debug (figura~\ref{ue4:fig:turretbp}).

\begin{figure}[htbp]
\centering
\begin{tikzpicture}[font=\footnotesize,
  ue4node/.style={draw, thick, rounded corners=4pt, align=center, fill=gray!8,
                  minimum height=0.85cm, text width=2.3cm},
  ue4event/.style={ue4node, fill=red!12},
  ue4exec/.style={-{Stealth[length=2.5mm]}, very thick, black!80},
  ue4data/.style={-{Stealth[length=2mm]}, thick, blue!60!black},
  ue4flt/.style={-{Stealth[length=2mm]}, thick, green!50!black},
  ue4dlg/.style={-{Stealth[length=2mm]}, thick, red!70!black, dashed},
  ue4comp/.style={draw, thick, rounded corners=3pt, minimum width=2.3cm,
                  minimum height=0.55cm, align=center, font=\scriptsize}]
  % --- riga 1: overlap
  \node[ue4event] (ov) at (0,0)
    {\textbf{On Component Begin Overlap}\\(Sphere)};
  \node[ue4node] (iv1) at (2.95,0) {\textbf{Is Valid}\\Other Actor};
  \node[ue4node] (br) at (5.9,0) {\textbf{Branch}\\Actor Has Tag ``Player''};
  \node[ue4node] (set) at (8.85,0) {\textbf{Set}\\\textbf{PlayerTarget}};
  \node[ue4node, text width=2.8cm] (call) at (12.05,0) {\textbf{TrackTargetTimer}\\(chiamata)};
  \draw[ue4exec] (ov) -- (iv1);
  \draw[ue4exec] (iv1) -- (br);
  \draw[ue4exec] (br) -- node[above, font=\scriptsize]{True} (set);
  \draw[ue4exec] (set) -- (call);
  \draw[ue4data] (ov.south) -- (0,-1.15) -| node[pos=0.2, below,
    font=\scriptsize]{Other Actor} (set.south);
  % --- riga 2: timer
  \node[ue4event, text width=2.8cm] (cev) at (0,-2.9)
    {\textbf{Custom Event}\\\textbf{TrackTargetTimer}};
  \node[ue4node, text width=3.0cm] (timer) at (3.65,-2.9)
    {\textbf{Set Timer by Event}\\Time: Get World Delta Seconds\\Looping \checkmark\\
    Max Once Per Frame \checkmark};
  \node[ue4event] (mv) at (7.1,-2.9) {\textbf{Custom Event}\\\textbf{MoveTurret}};
  \draw[ue4exec] (cev) -- (timer);
  \draw[ue4dlg] (timer.east) -- node[above, font=\scriptsize]{Event} (mv.west);
  % --- riga 3: move turret
  \node[ue4node] (iv2) at (0,-5.7) {\textbf{Is Valid}\\PlayerTarget};
  \node[ue4node] (look) at (3.1,-5.7) {\textbf{Find Look at Rotation}\\
    Start: self\\Target: PlayerTarget};
  \node[ue4node] (interp) at (6.3,-5.7) {\textbf{RInterp To}\\
    Current: Pivot\\Target: solo Z\\Speed 0,04--0,05};
  \node[ue4node] (srr) at (9.5,-5.7) {\textbf{Set Relative Rotation}\\Target: Pivot};
  \draw[ue4exec] (mv.south) -- (7.1,-4.35) -| (iv2.north);
  \draw[ue4exec] (iv2.south) -- (0,-7.1) -| node[pos=0.2, below,
    font=\scriptsize]{Is Valid} (srr.south);
  \draw[ue4flt] (look) -- node[above, font=\scriptsize]{Z} (interp);
  \draw[ue4flt] (interp) -- (srr);
  % --- gerarchia componenti
  \node[font=\scriptsize] at (11.6,-1.95) {i componenti};
  \node[ue4comp, fill=orange!25] (cb) at (11.4,-2.5) {Base (radice)};
  \node[ue4comp, fill=gray!20] (cp) at (11.9,-3.2) {Pivot};
  \node[ue4comp, fill=orange!12] (cc) at (12.4,-3.9) {Cannon};
  \node[ue4comp, fill=red!8] (cs) at (11.9,-4.6) {Sphere (r. 800)};
  \draw[thick, gray!60!black] ([xshift=-0.85cm]cb.south) |- (cp.west);
  \draw[thick, gray!60!black] ([xshift=-0.85cm]cp.south) |- (cc.west);
  \draw[thick, gray!60!black] ([xshift=-0.85cm]cb.south) |- (cs.west);
\end{tikzpicture}
\caption{La torretta della lezione. In alto: quando qualcosa entra nella sfera, se
esiste ed è etichettato \emph{Player}, lo si salva in PlayerTarget e si chiama
TrackTargetTimer. In mezzo: TrackTargetTimer fa partire un timer che, a ogni
\emph{frame}, chiama MoveTurret (il pin rosso tratteggiato è il delegato Event). In
basso: MoveTurret, se c'è ancora un bersaglio, calcola la rotazione verso di lui, ne
tiene lo \emph{yaw}, la interpola e la applica al Pivot (la Current dell'interpolazione
viene da Get Relative Rotation del Pivot; l'eventuale offset è discusso nel testo). A
destra la gerarchia dei componenti: il Cannon segue il Pivot, la sfera non ruota. Schema
ricostruito.}
\label{ue4:fig:turretbp}
\end{figure}

\paragraph{L'offset, e il pivot reimportato.} In Play la torretta ruota, ma ``ha
bisogno di un offset, come abbiamo visto con il nostro mulino''. Il docente prova con
$-90$, poi con $+90$: ``eccolo qua. Non è esattamente precisissimo''. Uno studente
osserva: ``noi questa cosa non la dobbiamo fare, perché l'altra volta abbiamo fatto un
\textbf{reimport} modificando i valori''. Ed è la chiave: senza alcun offset, la
torretta del docente ``guarda direttamente in faccia'', mentre quelle degli studenti no,
perché con Davide avevano reimportato il \textbf{Pivot} e il \textbf{Cannon} con una
\textbf{rotazione di 113} gradi sulla Z (la base no). Il docente ricostruisce i numeri
con la classe (si parla anche di $-67$ e di $-271$, ``proprio perfetti''), reimporta
prima il pivot, poi il cannone, prova con il solo pivot, e alla fine: ``ecco a voi.
Adesso la nostra torretta ci segue perfettamente''. Per gli studenti la regola pratica,
ripetuta più volte nel finale, è: se la torretta ha ``un \emph{offsettino}'',
\textbf{reimportare il Pivot mettendo $-113$} sulla rotazione Z\footnote{La trascrizione
di questo passaggio è tra le più confuse della lezione: i valori citati sono 113, $-113$,
$-67$, 118, $-270$ e $-271$, e non è chiaro in che ordine siano stati provati. Il senso
è comunque chiaro: la rotazione di 113 gradi aggiunta in import nella terza lezione
sposta la direzione ``in avanti'' della \emph{mesh}, e Find Look at Rotation, che
presuppone che l'oggetto guardi lungo la sua X, ne risulta sfasato. Le due correzioni
possibili sono equivalenti: annullare la rotazione in import, o sommare un offset fisso
allo \emph{yaw} nel grafo, come per il mulino.}.

\paragraph{Approfondimento: Find Look at Rotation e lo spazio relativo.} Il grafo della
lezione mescola due spazi: Find Look at Rotation lavora nello spazio del
\textbf{mondo} (le posizioni di torretta e player sono posizioni nel mondo), mentre
Set Relative Rotation imposta la rotazione del Pivot \textbf{relativa} alla base. Le
due cose coincidono finché la torretta nel livello non è ruotata: se la si piazza
girata di 30 gradi, il Pivot punterà 30 gradi a lato del player. Le soluzioni robuste
sono due: usare \textbf{Set World Rotation} sul Pivot (la rotazione calcolata è già
quella del mondo), oppure sottrarre allo \emph{yaw} calcolato quello dell'actor (Get
Actor Rotation). Anche la Current dell'interpolazione andrebbe presa nello stesso
spazio del Target: Get World Rotation con Set World Rotation, Get Relative Rotation con
Set Relative Rotation. Lo Start, infine, è la posizione dell'actor, cioè della base; per
una torretta alta conviene usare la posizione del Pivot o del cannone (\emph{Get World
Location} sul componente), specie quando si vorrà calcolare anche il \emph{pitch} per
alzare la canna.

% ---------------------------------------------------------------------
\section{Il tuning e i problemi degli studenti}
\label{ue4:sec:problemi}

\paragraph{La velocità di inseguimento.} ``Giratele intorno.'' Per aumentare o diminuire
la velocità con cui la torretta ``ci mette a cercare il target'' si cambia l'\textbf{Interp
Speed}. A 0,07 è ``molto più reattiva''; a 2 ``sarà proprio quasi uno \emph{snap}''
(``pam, \emph{snappa} proprio''). Il consiglio: ``mettetela tipo 0,05''. E, volendo, la
si può trasformare in variabile (Promote to Variable) ``e la potete settare
esternamente'', nei dettagli dell'istanza in scena, come le variabili \emph{Instance
Editable} della seconda lezione.

\paragraph{``A me snappa lo stesso.''} Uno studente ha Delta Time 5 e Interp Speed 0,05,
e la torretta scatta comunque. Si prova 0,02 (``vedo leggerissima differenza''), poi
Delta Time 7: ``stampa sempre la stessa velocità''. In condivisione il docente controlla
Get World Delta Seconds, Looping e Max Once Per Frame: tutto giusto. ``Prova ad andare di
\emph{tweak}, di \emph{tuning}: vedi un po' tu i valori.'' Abbassando ancora l'Interp
Speed, verso lo zero, la torretta rallenta: ``ora sì, già un po' più così''. (Nel
frattempo lo studente che lavora deve andare: ``ciao, ciao''.)

\paragraph{Approfondimento: perché ``snappa''.} La formula di RInterp To vista per il
mulino spiega tutto. Con Delta Time fisso a 5 e Interp Speed 0,04, la frazione di
strada percorsa a ogni chiamata è $5 \times 0{,}04 = 0{,}2$: il 20\% di quel che manca.
Dopo dieci chiamate manca $0{,}8^{10} \approx 11\%$ dell'angolo iniziale; a 100
\emph{frame} al secondo, sono un decimo di secondo. Con 0,05 la frazione è 0,25; con 7
e 0,02 è 0,14; con Interp Speed 2 il prodotto è 10, cioè più di 1, e la torretta arriva
in un solo \emph{frame}: lo \emph{snap}. In tutti questi casi, inoltre, la velocità
dipende dai \emph{frame}: con il doppio dei \emph{frame} la torretta è due volte più
rapida. L'uso previsto del nodo è il contrario di quello della lezione: \textbf{Delta
Time} collegato a \textbf{Get World Delta Seconds} (circa 0,01 a 100 \emph{fps}) e
\textbf{Interp Speed} nell'ordine delle unità (da 2 a 5 per un inseguimento morbido),
così che il prodotto resti piccolo e la velocità sia la stessa su ogni macchina. È lo
stesso principio dei \emph{delta seconds} della seconda lezione, applicato a
un'interpolazione. E un'avvertenza: un Interp Speed esattamente 0 non ferma la torretta,
ma la fa scattare subito sul bersaglio, perché Unreal lo interpreta come ``nessuna
interpolazione''.

\paragraph{La lista degli errori.} Il resto della lezione è una carrellata di
condivisioni dello schermo, molto istruttiva perché mostra gli errori più frequenti. La
tabella~\ref{ue4:tab:errori} li raccoglie. Il caso più ricco è quello di uno studente
a cui ``non va perché probabilmente ho mancato delle parti'':
\begin{itemize}
  \item per la Current di RInterp To aveva usato un \emph{Set} invece di \textbf{Get
        Relative Rotation} (``no, no, no, non \emph{set}: \emph{get}'');
  \item l'Is Valid di MoveTurret non era collegato alla variabile \textbf{PlayerTarget}
        (``a sinistra, proprio sulle variabili, non sul grafo'');
  \item MoveTurret non era un evento separato. Per riorganizzare il grafo il docente
        gli fa selezionare tutto tranne TrackTargetTimer, Set Timer by Event e Get
        World Delta Seconds (``tieni premuto Ctrl e clicca su quelli che vuoi
        escludere''), spostare in basso la selezione, tirare fuori dal pin Event del
        timer un nuovo Custom Event chiamato \textbf{MoveTurret} e collegarlo all'Is
        Valid;
  \item il timer non era in \textbf{Looping};
  \item il \textbf{Target} di Find Look at Rotation era scollegato: da PlayerTarget,
        Get Actor Location, al Target;
  \item il \emph{character} non aveva il \textbf{tag}: lo si aggiunge nel Third Person
        Character, ``\emph{Player} esattamente come l'hai scritto, quindi con la P
        maiuscola'';
  \item la sfera aveva ancora Hidden in Game attivo e raggio \textbf{32}: ``devi
        metterla a 800'';
  \item in scena c'era solo la \textbf{mesh} della torretta, non il
        \textbf{Blueprint}: ``questa è solo la \emph{mesh}, è solo il modello 3D che hai
        messo in scena''; si cancella e si trascina nel livello BP\_Turret;
  \item il player partiva già dentro il range, e quindi il Begin Overlap non scattava:
        ``il player parte da fuori, mettilo fuori'';
  \item in Play \textbf{girava la base}: nella gerarchia, la base era figlia del Pivot.
        ``Devi staccare la base'': trascinando il Pivot dentro la base, ``perfetto, ce
        l'hai perfettamente'';
  \item infine, l'\emph{offset} del pivot reimportato: ``devi reimportare il pivot a
        $-113$''.
\end{itemize}
Un altro studente ha la torretta ferma: ``Other Actor devi attaccarlo a PlayerTarget,
perché così non gli stai impostando il \emph{target} che è overlappato: non avrà un
\emph{target} da seguire''. Un terzo non vede la sfera in gioco, solo in modalità
\emph{spectator}: aveva saltato il passaggio di Hidden in Game.

\begin{table}[htbp]
\centering
\small
\begin{tabular}{@{}p{5.4cm}p{8.6cm}@{}}
\toprule
\textbf{Sintomo} & \textbf{Causa e rimedio} \\
\midrule
La pala gira attorno a un punto sbagliato & la \emph{mesh} ha il pivot fuori
centro: scegliere un cubo con il pivot al centro (o correggere il pivot) \\
La pala non gira & il risultato di Combine Rotators non è collegato a Set
Relative Rotation; oppure il componente non è Movable \\
Il mulino o la torretta restano fermi & mancano Delta Time e Interp Speed in
RInterp To; oppure non si è compilato \\
Non si riesce a collegare un valore a una sola componente & fare Split Struct Pin
(tasto destro sul pin) \\
Gira tutto l'actor, non solo una parte & il Rotating Movement fa girare la radice:
programmare la rotazione del componente \\
La torretta non si accorge del player & manca il tag \emph{Player} (o è scritto
diversamente); Other Actor non collegato al Set di PlayerTarget; il player parte già
dentro la sfera \\
La sfera non si vede in Play & Hidden in Game ancora attivo; raggio di default 32 \\
Non succede niente in Play & in scena c'è la \emph{mesh}, non il Blueprint \\
Gira la base & gerarchia sbagliata: la base deve essere la radice, il Pivot suo
figlio, il Cannon figlio del Pivot \\
La torretta guarda di lato & la \emph{mesh} non guarda lungo X (Pivot reimportato
a 113): reimportare a $-113$ o sommare un offset allo \emph{yaw} \\
La torretta scatta (\emph{snap}) & prodotto Delta Time $\times$ Interp Speed troppo
grande: usare Get World Delta Seconds e un Interp Speed di qualche unità \\
La torretta continua a seguire dopo l'uscita & il timer non viene mai fermato e
PlayerTarget resta pieno: End Overlap, Clear and Invalidate Timer, svuotare la
variabile \\
\bottomrule
\end{tabular}
\caption{I problemi incontrati durante la lezione, con la causa e il rimedio. Le ultime
due righe riguardano punti che il docente ha rimandato alla lezione successiva.}
\label{ue4:tab:errori}
\end{table}

% ---------------------------------------------------------------------
\section{La chiusura e la prossima lezione}
\label{ue4:sec:chiusura}

``Ok, perfetto, ragazzi.'' Per questa volta, ``nessun particolare compito da casa o
dimostrazione''. Chi vuole può fare una \textbf{torretta di colore diverso}, con un
materiale suo, verde o blu, ripassando quello che si è fatto con Davide.

\paragraph{La prossima volta.} ``La prossima volta faremo la nostra bellissima
\textbf{torretta che spara}'': \textbf{spawnerà} dei \textbf{proiettili}. Si vedrà
``come funziona il sistema di \emph{spawning} di Unreal'': come si basa una cosa
\emph{spawnata} ``con le sue regole'', come si fa apparire un oggetto in generale, e
come si fa il \textbf{line trace}. ``Altre domande? Se no, andiamo a cena tutti,
ragazzi.'' ``Vi auguro buona serata. E a martedì.''

\begin{table}[htbp]
\centering
\small
\begin{tabular}{@{}p{4.0cm}p{10.0cm}@{}}
\toprule
\textbf{Nodo o componente} & \textbf{Che cosa fa} \\
\midrule
Rotating Movement Component & fa ruotare la radice dell'actor a velocità costante
(Rotation Rate, default Z 180) \\
Event Tick, Delta Seconds & esegue a ogni \emph{frame}; i \emph{delta seconds} sono il
tempo tra due \emph{frame} \\
Get World Delta Seconds & gli stessi \emph{delta seconds}, ovunque nel grafo \\
Make Rotator & costruisce un \emph{rotator} da tre \emph{float} (Roll X, Pitch Y,
Yaw Z) \\
Combine Rotators & combina due rotazioni in una \\
Get / Set Relative Rotation & legge e scrive la rotazione di un componente rispetto al
genitore \\
Get Actor Location & la posizione dell'actor nel mondo \\
Get Player Character & una referenza al personaggio del giocatore \\
Find Look at Rotation & la rotazione per guardare da Start verso Target \\
RInterp To & avvicina gradualmente un \emph{rotator} a un altro \\
Clamp (float) & limita un valore tra un minimo e un massimo \\
Split Struct Pin & divide un pin di struttura (vettore, \emph{rotator}) nelle sue
componenti \\
Sphere Collision & un volume sferico per rilevare \emph{overlap} (Sphere Radius,
Hidden in Game) \\
On Component Begin Overlap & scatta quando qualcosa entra nel volume del componente;
Other Actor è chi è entrato \\
Is Valid (macro con il ``?'') & controlla che una referenza punti a qualcosa che
esiste \\
Actor Has Tag & controlla se un actor porta un'etichetta (Actor $\rightarrow$ Tags) \\
Promote to Variable & crea una variabile del tipo giusto da un pin \\
Custom Event & un evento creato da noi, che parte quando lo si chiama \\
Set Timer by Event & chiama un evento dopo un tempo, una volta o in \emph{loop} (Max
Once Per Frame) \\
Print String & stampa a schermo: lo strumento base del debug \\
\bottomrule
\end{tabular}
\caption{I nodi e i componenti della quarta lezione.}
\label{ue4:tab:nodi}
\end{table}

% ---------------------------------------------------------------------
\section*{Punti chiave del capitolo}
\addcontentsline{toc}{section}{Punti chiave del capitolo}

\begin{itemize}
  \item \textbf{Quixel Megascans} è un ``calderone'' gratuito di materiali, texture e
        \emph{mesh} scansionati, da usare con l'account Epic. Per la pallina che
        attraversa il piano: il \textbf{CCD} (\emph{Continuous Collision Detection}).
  \item Un \textbf{vettore} ha direzione, magnitudine e verso; in Unreal la struttura
        Vector (X, Y, Z) rappresenta sia punti sia spostamenti. Somma: si spostano i
        punti; \textbf{sottrazione}: B $-$ A va da A a B; \textbf{moltiplicazione}:
        allunga; \textbf{prodotto scalare}: un \emph{float}, 1 se concordi, 0 se
        perpendicolari, $-1$ se opposti (per vettori normalizzati).
  \item Un vettore \textbf{normalizzato} ha lunghezza 1: tiene solo la direzione.
  \item Un \textbf{rotator} ha tre angoli: \textbf{Roll} (X), \textbf{Pitch} (Y),
        \textbf{Yaw} (Z). Gli angoli di Eulero soffrono di \textbf{gimbal lock}; i
        \textbf{quaternioni}, usati internamente da Unreal, no.
  \item Il \textbf{Rotating Movement Component} fa girare un actor senza programmare,
        ma solo per intero. Ogni oggetto che si muove deve essere \textbf{Movable}.
  \item Per far girare un componente: \textbf{Event Tick} $\rightarrow$ velocità
        $\times$ \emph{delta seconds} $\rightarrow$ \textbf{Make Rotator}
        $\rightarrow$ \textbf{Combine Rotators} con \textbf{Get Relative Rotation}
        $\rightarrow$ \textbf{Set Relative Rotation}. Le rotazioni dei componenti si
        lavorano in spazio \textbf{relativo}.
  \item La \textbf{gerarchia} fa il resto: un componente figlio segue il genitore,
        mantenendo la propria posizione relativa (la seconda pala, il cubetto).
  \item \textbf{Find Look at Rotation} calcola la rotazione per guardare un punto;
        dello \emph{yaw} si tiene la Z (\textbf{Split Struct Pin}), con un
        \textbf{offset} se la \emph{mesh} non guarda lungo X. \textbf{RInterp To}
        rende il movimento graduale: Delta Time dai \emph{delta seconds}, Interp Speed
        di qualche unità, altrimenti ``snappa''.
  \item Il \textbf{Clamp} limita un valore; prima di sceglierne i limiti, si
        \textbf{stampano} i valori con Print String.
  \item La torretta è un Blueprint in \textbf{tre pezzi}: Base (radice), Pivot
        (figlio, ruota in \emph{yaw}), Cannon (figlio del Pivot). Browse to Asset aiuta
        a trovare le \emph{mesh}.
  \item Il range è una \textbf{Sphere Collision} (raggio 800, Hidden in Game tolto per
        vederla). \textbf{On Component Begin Overlap} scatta quando qualcosa ci entra;
        \textbf{Other Actor} dice chi.
  \item \textbf{Is Valid} prima di usare una referenza; \textbf{Actor Has Tag} con il tag
        \emph{Player} (scritto identico) per riconoscere il giocatore;
        \textbf{Promote to Variable} per ricordarlo in \emph{PlayerTarget}.
  \item I \textbf{Custom Event} sono eventi nostri, chiamati quando vogliamo;
        \textbf{Set Timer by Event} (Looping, Max Once Per Frame) chiama a ripetizione
        l'evento \textbf{MoveTurret}, che ruota il Pivot verso il bersaglio. Il timer
        va fermato (End Overlap, Clear and Invalidate Timer): sarà il prossimo passo.
  \item Prossima lezione: la torretta che \textbf{spara}, lo \textbf{spawn} dei
        proiettili e il \textbf{line trace}.
\end{itemize}

%
%  Tutte le etichette di questo capitolo hanno il prefisso "ue4:"
%  (ue4:cap:..., ue4:sec:..., ue4:fig:..., ue4:tab:...), e tutti gli
%  stili TikZ hanno il prefisso "ue4" (ue4node, ue4exec, ue4data, ...)
%  per non collidere con chiavi native di TikZ (come "step") né con
%  quelli delle lezioni precedenti. Il capitolo è autonomo: non contiene
%  \ref verso altri capitoli (i rimandi alle lezioni precedenti sono a
%  parole).
%  Le figure sono ricostruzioni schematiche di quanto mostrato e detto a
%  lezione, non copie delle slide o delle schermate del docente. Dove la
%  trascrizione è troppo rovinata per ricostruire un dettaglio, il
%  capitolo lo dice invece di inventarlo.
% =====================================================================

\chapter[Vettori, rotazioni e la torretta]{Vettori, rotazioni e la torretta: il
mulino che gira e la torretta che ti segue}
\label{ue4:cap:torretta}

Questa è la \textbf{quarta lezione} del corso breve di \textbf{introduzione a Unreal
Engine}, e segna il ritorno del docente titolare dopo la lezione di Davide sulla
grafica. Se la terza lezione guardava a ciò che ``fa il grafico all'interno di Unreal''
(texture, materiali, l'import della torretta), questa torna alla \textbf{logica}, cioè
ai \textbf{Blueprint}, e lo fa con un obiettivo preciso, mostrato in apertura: un
\textbf{mulino} le cui pale girano, e la \textbf{torretta} importata con Davide che
``adesso ci segue con lo sguardo''. Per arrivarci serve un po' di matematica dello
spazio 3D: i \textbf{vettori}, le loro operazioni, i \textbf{rotator}.

La lezione ha due metà separate da una pausa di venti minuti, presa alle 17:28. La
prima metà parte dalla \textbf{teoria}: che cos'è un \textbf{vettore} (``una freccia
con una direzione, una magnitudine e un verso''), le quattro operazioni
(\textbf{addizione}, \textbf{sottrazione}, \textbf{moltiplicazione} e \textbf{prodotto
scalare}, il \emph{dot product}), i \textbf{vettori normalizzati}, i
\textbf{rotator} e il problema del \textbf{gimbal lock}. Poi si passa all'engine con un
mulino: prima il modo più semplice di far ruotare un oggetto, il \textbf{Rotating
Movement Component}, poi un Blueprint vero, \textbf{BP\_Windmill}, in cui una pala
gira grazie all'\textbf{Event Tick}, ai \emph{delta seconds} e ai nodi \textbf{Make
Rotator}, \textbf{Combine Rotators} e \textbf{Set Relative Rotation}; la seconda pala
gira ``gratis'' grazie alla \textbf{gerarchia} dei componenti. Infine il mulino
diventa ``magico'' e si gira verso il giocatore con \textbf{Find Look at Rotation} e
\textbf{RInterp To}, con un \textbf{Clamp} per limitarne la rotazione. La seconda metà
costruisce la \textbf{torretta}: un Blueprint in tre pezzi (\textbf{Base},
\textbf{Pivot}, \textbf{Cannon}), una \textbf{Sphere Collision} come raggio d'azione,
l'evento \textbf{On Component Begin Overlap}, i controlli \textbf{Is Valid} e
\textbf{Actor Has Tag}, una variabile per ricordare il bersaglio, un \textbf{Custom
Event} e un \textbf{timer} che, a ogni \emph{frame}, fa ruotare il pivot verso il
player.

La lezione è \textbf{online}, su Zoom, come le precedenti. Il docente alterna
spiegazione e piccole esercitazioni ``botta e risposta'', da cinque minuti l'una, e
passa gran parte della seconda metà a risolvere i problemi degli studenti in
condivisione dello schermo. Finisce verso le 19, ``andiamo a cena tutti'', con
l'appuntamento a \textbf{martedì}.

Tre avvertenze di lettura, le stesse delle lezioni precedenti. La prima: la
trascrizione è automatica e storpia moltissimi nomi. Unreal è ``a Real'', ``Arreal'',
``Arry''; l'\textbf{Event Tick} è ``Leventic'', ``l'eventic'', ``eventHIC'' e
``l'event kick''; i \textbf{vettori} sono una volta ``vittori''; i \emph{rotator} sono
``rotatempi'' e ``rotatori''; \textbf{Find Look at Rotation} è ``findLookAtRotation'',
``find all the rotation'' e ``Find, look, rotation''; \textbf{RInterp To} è ``Er inter
tu'', ``R Inter2'' e ``r interp tu''; \textbf{Get Player Character} è
``getLayerCharacter''; \textbf{Get Actor Location} è ``Get. Hector. Location'' e ``get
active location''; il \textbf{Rotating Movement} è a un certo punto ``vetro in
movement''; \textbf{Is Valid} è ``Svalid'' e ``isvalid''; \textbf{Other Actor} è
``OderActor'' e ``Oder Actor''; la \textbf{gun} della torretta è ``Ganna'' e ``GAN'';
il \textbf{Pivot} è spesso ``pivo''; il \emph{Third Person Character} è ``ter person
character'' e ``Peter per personacarte''; \textbf{Quixel Megascans} è ``QuickSell'' e
``Quixel megaschanico''. Dove la ricostruzione è sicura il capitolo usa la forma
corretta e, la prima volta, lo segnala in nota; dove non lo è riporta la forma della
trascrizione e lo dice. La seconda: all'inizio della registrazione, in un punto della
prima esercitazione e durante tutta la pausa, il riconoscimento vocale produce lunghi
blocchi di testo senza senso e ripetuto (``qual è la risposta che ci può essere?'',
``grazie a tutti per l'attenzione'', discorsi sul ``mercato del lavoro'', sul ``governo
di Dalai Lama'' e su ``la migliore dal 2005''): sono artefatti, non parole di nessuno,
e il capitolo li ignora. La terza: gran parte della lezione è una dimostrazione dal
vivo nell'editor dei Blueprint, spesso costruita a voce (``tirate fuori'',
``collegate'', ``splittate''). Il capitolo ricostruisce i grafi da ciò che il docente
dice mentre li costruisce, e li completa, dove serve, con i nomi esatti dei nodi e dei
pin di Unreal Engine 5.

\paragraph{Come leggere questo capitolo.} Il capitolo segue l'ordine della lezione.
Le figure ricostruiscono gli schemi spiegati a voce o alla lavagna (``perdonate le mie
doti molto brutte di disegno''): le operazioni tra vettori, i casi del prodotto
scalare, gli assi di rotazione, la gerarchia del mulino, i grafi della pala, del
mulino che guarda il player e della torretta, i casi dell'\emph{overlap}. Le tabelle
riassumono le operazioni sui vettori, i nodi della lezione e gli errori più comuni
incontrati dagli studenti. Qua e là, dei paragrafi di \textbf{approfondimento}
collegano la lezione ai libri di Marcos Romero e Brenden Sewell (\emph{Blueprints
Visual Scripting for Unreal Engine 5} e la seconda edizione dello stesso testo, scritta
per Unreal Engine 4), i testi di riferimento per questa parte del libro, e a qualche
nozione tecnica generale su Unreal Engine 5; sono segnalati come tali e non fanno parte
della lezione.

Il percorso del capitolo:
\begin{enumerate}
  \item l'attesa: \textbf{Quixel Megascans} e il materiale della lava; il trucco del
        \textbf{CCD}; il programma della giornata (sezioni da~\ref{ue4:sec:attesa}
        a~\ref{ue4:sec:programma});
  \item la teoria: i \textbf{vettori} e le loro operazioni, il \textbf{prodotto
        scalare}, la \textbf{normalizzazione}, i \textbf{rotator} e il \textbf{gimbal
        lock} (sezioni da~\ref{ue4:sec:vettori} a~\ref{ue4:sec:rotator});
  \item il \textbf{Rotating Movement Component} e l'esercizio dei tre cubi (sezione
        \ref{ue4:sec:rotating});
  \item il Blueprint del \textbf{mulino}: i componenti, la pala che gira con
        l'Event Tick, la \textbf{gerarchia} (sezioni da~\ref{ue4:sec:mulino}
        a~\ref{ue4:sec:gerarchia});
  \item il mulino che guarda il player: \textbf{Find Look at Rotation},
        \textbf{RInterp To}, il \textbf{Clamp} e il debug (sezioni
        \ref{ue4:sec:lookat} e~\ref{ue4:sec:clamp});
  \item dopo la pausa, la \textbf{torretta}: i tre pezzi, la \textbf{Sphere
        Collision}, l'\textbf{overlap} (sezioni da~\ref{ue4:sec:torretta}
        a~\ref{ue4:sec:overlap});
  \item \textbf{Is Valid}, i \textbf{tag}, la variabile del bersaglio e i
        \textbf{Custom Event} (sezioni da~\ref{ue4:sec:isvalid}
        a~\ref{ue4:sec:custom});
  \item il \textbf{timer} e l'evento \textbf{Move Turret}: la torretta che segue il
        player (sezioni \ref{ue4:sec:timer} e~\ref{ue4:sec:move});
  \item il \emph{tuning}, i problemi degli studenti e la chiusura (sezioni
        \ref{ue4:sec:problemi} e~\ref{ue4:sec:chiusura}).
\end{enumerate}

% ---------------------------------------------------------------------
\section{L'attesa: Quixel Megascans e la lava}
\label{ue4:sec:attesa}

La registrazione parte mentre il docente fa entrare gli studenti nella stanza di Zoom
(``ho aperto io la prima a far entrare un po' tutti'') e chiacchiera con il primo
arrivato. ``Come va con Unreal? Hai sperimentato un po', l'hai toccato un pochetto?''
Lo studente ci ha giocato con un collega, che gli ha insegnato le texture, ma non ha
trovato online ``alcuna fonte per prendere cose utili, file utili'': per l'esercizio
della lava assegnato da Davide nella lezione precedente, un materiale di lava animata,
ha usato quello che c'era nella cartella del Drive.

\paragraph{Quixel Megascans.} La risposta del docente è \textbf{Quixel
Megascans}\footnote{La trascrizione dice ``QuickSell'', ``Quixel megaschanico'' e
``Materials pixel mega scan''. Il nome corretto è Quixel Megascans: la libreria di
asset scansionati di Quixel, società acquisita da Epic Games, gratuita per l'uso con
Unreal Engine.}, di cui passa il link in chat. Quixel ``ti permette di avere proprio
dei materiali'': ``una marea di materiali di qualsiasi tipo'', ma anche rocce,
\emph{mesh} e texture di alta qualità. Il docente fatica un po' a ritrovare la pagina
giusta, ``perché hanno cambiato il sito'', e ricorda che serve accedere con l'account
di Unreal (cioè di Epic Games). ``Ormai fanno la cooperazione tra di loro'': c'è un
collegamento diretto tra Quixel e Unreal, ``proprio tipo \emph{face to face}''. Cercando
\emph{lava} escono diverse risorse; una, \emph{lava story} o qualcosa del genere,
``sembra carina''.

\paragraph{Come sono fatte.} Le texture di Megascans ``sono proprio \emph{scan}'': c'è
``gente che si mette là con queste macchine che costano una vita umana'' a
scansionare le rocce. ``Quixel è un grande, grande calderone di risorse.'' Da lì si
può partire anche per la lava senza trovare la texture perfetta: si prende una
\textbf{roccia} e, ``come avete visto con Davide, con la matematica'', cioè con un
\textbf{Multiply} e con il \textbf{Lerp} pilotato da un'\textbf{alfa}, si inserisce
``il rosso della lava all'interno di queste rocce''. Un'altra strada sono i video su
YouTube: alcuni tutorial lasciano in descrizione le texture usate (altri le fanno
pagare, ``il maledetto''), e ``ti interessano le texture, il materiale poi lo fai tu''.

\paragraph{La lava degli studenti.} Arrivano gli altri (``buongiorno, Filippo, Marco,
Patrick''), e il docente chiede chi ha provato il materiale della lava ``che vi è stato
dato da Davide''\footnote{La trascrizione dice ``da Dario'', ma subito dopo e più
avanti il nome è Davide, il docente della terza lezione.}. Uno studente l'ha mandato sul
Discord, e il docente lo commenta: ``bello il movimento che hai fatto''. L'hai messo
su un piano, in engine, per vedere come reagiva? No. ``Quello vedilo un po' tu'',
perché a volte il risultato in \emph{viewport} cambia a seconda di come è costruito il
materiale (``che so, hai messo il Multiply sbagliato''): il consiglio è di provarlo
sempre in \emph{viewport}, su un oggetto vero. E ``sicuramente una texture differente
potrebbe essere più carina'': torna utile Megascans.

\paragraph{L'appello.} Il docente aspetta ``2--3 minuti'' che entrino gli altri. Uno
studente avvisa che anche oggi dovrà staccarsi venti minuti prima per andare a
lavorare: ``assolutamente nessun problema''.

% ---------------------------------------------------------------------
\section{Le domande e il trucco del CCD}
\label{ue4:sec:ccd}

Prima di cominciare, il docente apre lo spazio delle domande: ``se qualcuno di voi ha
avuto qualche particolare cosa di cui vuole parlare, domandare già prima che magari
iniziamo, così facciamo il Q\&A già da subito''. Nessuno: ``mi sembrate tutti molto
tranquilli''.

\paragraph{La pallina che entra nel piano.} Allora è il docente a recuperare ``un
piccolo trucchetto che mi sono dimenticato di dirvi'' per il \textbf{gioco della
pallina} della seconda lezione, quello della piattaforma inclinabile: a volte la
pallina, con la fisica attiva, \textbf{penetra} nel piano, ``come abbiamo visto''. Nei
dettagli del componente, nella sezione della fisica (vicino a \textbf{Simulate
Physics}), c'è un'opzione da spuntare: il \textbf{CCD}, \emph{Continuous Collision
Detection}\footnote{La trascrizione dice ``Trucco Collision Collision Detection''.}.
``Fondamentalmente, è un \emph{tick} che, se lo mettiamo su questo particolare
\emph{asset}, lui in modo continuo cerca di vedere le collisioni'': le rende ``più
precise'', ed evita che la pallina attraversi il terreno. ``Se fate questo, vi risolve
tranquillamente il problema.'' Chi vuole, può sistemare il gioco per la prossima volta.

\paragraph{Approfondimento: perché la pallina attraversa il piano.} La fisica di un
gioco non è continua: il motore fisico controlla le collisioni a passi discreti, una
volta per \emph{frame} (o per \emph{substep}). Se in un passo la pallina si sposta più
dello spessore del piano, può trovarsi ``prima'' sopra e ``dopo'' già dentro o sotto,
senza che nessun controllo l'abbia mai vista toccare la superficie: è il fenomeno del
\emph{tunneling}, tanto più probabile quanto più l'oggetto è piccolo e veloce e il piano
sottile. Il CCD (in Unreal la casella \emph{Use CCD} nella sezione \emph{Collision} del
componente, nelle impostazioni avanzate) fa sì che il motore fisico consideri anche il
percorso tra le due posizioni, la ``scia'' dell'oggetto, e non solo le posizioni
finali. Costa un po' di più, ed è per questo che non è attivo di default: lo si accende
solo sugli oggetti che ne hanno bisogno, come proiettili e palline veloci. Un'alternativa
complementare è dare al piano uno spessore maggiore.

% ---------------------------------------------------------------------
\section{Il programma: il mulino e la torretta}
\label{ue4:sec:programma}

``Oggi vi dico già da subito cosa faremo.'' Gli studenti hanno scaricato il
\textbf{materiale didattico}: dentro il file estratto c'è il progetto
\emph{UE5 Intro Start}\footnote{La trascrizione dice ``Hive UE5 intro start'', la stessa
parola iniziale incerta già incontrata nella terza lezione.}, con il suo file di
progetto. Lo si apre, ``così almeno gli \emph{shader} si cominciano a buildare'', si fa
``una piccola spiegazione teorica'' e poi ``partiamo direttamente con la pratica''.

\paragraph{I due obiettivi.} Il docente sposta il \emph{Player Start} e mostra il
livello. A sinistra c'è ``questo bellissimo \textbf{mulino}'', fatto ``facendo ruotare
le pale''. Poi c'è la \textbf{torretta} che gli studenti hanno costruito la volta prima
con Davide (``avete fatto il \emph{merge}, giusto?''), che ``adesso ci segue con lo
sguardo''. O meglio, dovrebbe: nella dimostrazione non lo fa. ``Abbiamo scoperto un
piccolo bug. Poi adesso lo aggiustiamo. E giustamente i bug ogni volta appaiono
praticamente subito.'' Il programma della giornata è questo: le pale del mulino, e la
torretta che segue il giocatore.

% ---------------------------------------------------------------------
\section{I vettori}
\label{ue4:sec:vettori}

``Avete sentito mai parlare di vettori? Matematica, i grafici cartesiani?'' Sì.

\paragraph{Che cos'è un vettore.} ``In generale, un vettore può essere semplificato, in
uno spazio, come letteralmente una \textbf{freccia}'' che ha tre caratteristiche: una
\textbf{direzione}, una \textbf{magnitudine} (``la sua ampiezza'', cioè la lunghezza) e
un \textbf{verso}. In Unreal i vettori si sono già visti senza nominarli: nella
\textbf{location} e nella \textbf{rotation} di ogni oggetto, ``in particolare molto di
più sulla location''. La location è fatta di tre coordinate, \textbf{X}, \textbf{Y} e
\textbf{Z}, e ``queste particolari coordinate possono formare tra di loro, nello spazio
3D, un vettore''. Il docente parte dalla parola ``velocità di movimento'' e si corregge:
``più che una velocità, è più una direzione'', con la sua ampiezza e il suo verso.
Con i vettori si possono ``fare tante cose in Unreal'', per esempio muovere degli
oggetti; le pale del mulino girano in un verso, e anche se in quel caso si parla di
rotazione, e quindi di \emph{rotator}, ``sono sempre formate da delle coordinate''.

\paragraph{Un punto non è ancora un vettore.} Più avanti, alla lavagna, il docente
aggiunge una precisazione importante: ``ricordatevi che quello che vediamo è un
\textbf{punto} nello spazio, non è una freccia, quindi non è già un vettore''. Un
punto, per esempio $(5, 4, 1)$, è solo una posizione; il vettore è la freccia che
collega due punti, per esempio dall'origine $(0,0,0)$ a $(5,4,1)$\footnote{A lezione il
docente dice che ``un vettore è formato dalla somma di due punti''. La relazione
corretta è la \emph{differenza}: il vettore che va da A a B è B $-$ A (si veda la
sottrazione, poco più sotto). Unreal, peraltro, usa la stessa struttura
\emph{Vector} sia per i punti sia per gli spostamenti.}.

\paragraph{Le operazioni.} Con i vettori ``possiamo fare varie operazioni matematiche,
che sono praticamente quelle classiche'': \textbf{addizione}, \textbf{sottrazione},
\textbf{moltiplicazione} e \textbf{prodotto scalare} (\emph{dot product}), ``che è
un'altra cosa, che poi spiegheremo''. La slide le mostra con due vettori, A e B
(figura~\ref{ue4:fig:vettori}).
\begin{itemize}
  \item \textbf{Addizione.} ``Se andiamo a prendere due vettori e li addizioniamo,
        avremo un vettore unico.'' Il vettore A più il vettore B ``dà come risultato
        il vettore C, che è $(5,4)$'': si mette la coda di B sulla punta di A, e C va
        dall'inizio di A alla fine di B. ``Come fare l'ipotenusa, se vi ricordate i
        triangoli.''
  \item \textbf{Sottrazione.} Dà ``la \textbf{differenza} nello spazio tra due
        punti''. Il risultato dipende dall'ordine: A $-$ B è il vettore che va da B
        ad A, B $-$ A quello che va da A a B. La sua lunghezza è la distanza tra i due
        punti.
  \item \textbf{Moltiplicazione.} ``La useremo tantissimo'': si moltiplica tutto ciò
        che fa parte del vettore, X, Y e Z, e quindi il vettore ``viene allungato nella
        sua direzione''.
  \item \textbf{Prodotto scalare.} ``Statemi un attimo a sentire: è l'unica operazione
        che ci dà un valore \textbf{scalare}'', cioè ``non ci ridà un vettore indietro,
        ma un valore in \emph{float}''. (``Chi si ricorda che è il float?'' ``I numeri
        decimali.'' ``Perfetto.'') Lo si vede nella prossima sezione.
\end{itemize}

\begin{figure}[htbp]
\centering
\begin{tikzpicture}[font=\small, scale=0.8,
  ue4vec/.style={-{Stealth[length=2.5mm]}, very thick},
  ue4grid/.style={gray!25, very thin},
  ue4ax/.style={-{Stealth[length=2mm]}, gray!70}]
  % --- addizione
  \begin{scope}
    \draw[ue4grid] (0,0) grid (6,5);
    \draw[ue4ax] (0,0) -- (6.3,0) node[right, font=\scriptsize]{X};
    \draw[ue4ax] (0,0) -- (0,5.3) node[above, font=\scriptsize]{Y};
    \draw[ue4vec, blue!70!black] (0,0) -- node[left, pos=0.55]{A} (2,3);
    \draw[ue4vec, red!70!black] (2,3) -- node[above, pos=0.5]{B} (5,4);
    \draw[ue4vec, green!45!black] (0,0) -- node[below right, pos=0.6]{C $=$ A $+$ B}
      (5,4);
    \fill (5,4) circle (1.5pt) node[above right, font=\scriptsize]{$(5,4)$};
    \node[font=\bfseries] at (3,-0.9) {addizione};
  \end{scope}
  % --- sottrazione
  \begin{scope}[shift={(8.5,0)}]
    \draw[ue4grid] (0,0) grid (6,5);
    \draw[ue4ax] (0,0) -- (6.3,0) node[right, font=\scriptsize]{X};
    \draw[ue4ax] (0,0) -- (0,5.3) node[above, font=\scriptsize]{Y};
    \fill[blue!70!black] (1,1) circle (2.5pt) node[below left]{A};
    \fill[red!70!black] (5,4) circle (2.5pt) node[above right]{B};
    \draw[ue4vec, green!45!black] (1.15,1.11) -- node[above left, pos=0.5]
      {B $-$ A: da A a B} (4.85,3.89);
    \draw[ue4vec, orange!80!black, dashed] (4.9,3.6) -- node[below right, pos=0.5]
      {A $-$ B: da B ad A} (1.2,0.85);
    \node[font=\bfseries] at (3,-0.9) {sottrazione};
  \end{scope}
  % --- moltiplicazione
  \begin{scope}[shift={(0,-7.2)}]
    \draw[ue4grid] (0,0) grid (6,5);
    \draw[ue4ax] (0,0) -- (6.3,0) node[right, font=\scriptsize]{X};
    \draw[ue4ax] (0,0) -- (0,5.3) node[above, font=\scriptsize]{Y};
    \draw[ue4vec, green!45!black, line width=3pt, opacity=0.35] (0,0) -- (5,4);
    \draw[ue4vec, blue!70!black] (0,0) -- (2.5,2) node[pos=0.5, above left]{V};
    \node[green!45!black, font=\small] at (4.6,4.55) {$2 \times$ V};
    \node[font=\bfseries] at (3,-0.9) {moltiplicazione per uno scalare};
  \end{scope}
  % --- normalizzazione
  \begin{scope}[shift={(8.5,-7.2)}]
    \draw[ue4grid] (0,0) grid (6,5);
    \draw[ue4ax] (0,0) -- (6.3,0) node[right, font=\scriptsize]{X};
    \draw[ue4ax] (0,0) -- (0,5.3) node[above, font=\scriptsize]{Y};
    \draw[gray!60, dashed] (0,0) circle (1);
    \node[gray, font=\scriptsize] at (1.35,-0.3) {lunghezza 1};
    \draw[ue4vec, blue!70!black, opacity=0.45] (0,0) -- (4,3)
      node[above right, font=\scriptsize]{V, lunghezza 5};
    \draw[ue4vec, red!70!black] (0,0) -- (0.8,0.6);
    \node[red!70!black, font=\scriptsize, align=left] at (2.7,0.45)
      {Normalize(V)\\$= (0{,}8;\ 0{,}6)$};
    \node[font=\bfseries] at (3,-0.9) {normalizzazione};
  \end{scope}
\end{tikzpicture}
\caption{Le operazioni sui vettori viste a lezione, nel solo piano XY per semplicità.
In alto a sinistra la somma della slide: A seguito da B porta al punto $(5,4)$. In
alto a destra la sottrazione, il cui verso dipende dall'ordine dei fattori. In basso,
la moltiplicazione per un numero, che allunga il vettore senza cambiarne la direzione,
e la normalizzazione, che lo riporta a lunghezza 1 (sezione~\ref{ue4:sec:dot}). I
valori di A e B della slide non si leggono nella trascrizione e sono scelti per
l'esempio; il risultato $(5,4)$ è quello detto a lezione.}
\label{ue4:fig:vettori}
\end{figure}

\paragraph{Approfondimento: i vettori nel libro di Romero e Sewell.} Il
quattordicesimo capitolo di \emph{Blueprints Visual Scripting for Unreal Engine 5},
dedicato ai nodi matematici e alle \emph{trace}, presenta i vettori con un esempio
molto vicino a quello della lezione. In Unreal la struttura \textbf{Vector} contiene
tre \emph{float}, X, Y e Z, e può rappresentare sia un \textbf{punto} nello spazio
(una \emph{location}) sia uno \textbf{spostamento} o una velocità. Un personaggio sta
in $(50, 0, 20)$, un divano in $(450, 0, 20)$, e un'unità di Unreal vale un
centimetro. Per arrivare al divano il personaggio deve spostarsi di $(400, 0, 0)$: 400
centimetri lungo X. Le due regole che ne derivano sono quelle della lezione, scritte
in formula:
\[
  \text{destinazione} = \text{partenza} + \text{spostamento},
  \qquad
  \text{spostamento} = \text{destinazione} - \text{partenza}.
\]
Ogni operazione si fa componente per componente: $(3,5,0) + (5,2,0) = (8,7,0)$ e
$(6,8,0) - (1,4,0) = (5,4,0)$. Il libro aggiunge due esempi pratici: se V1 è la
posizione nel mondo di un personaggio e V2 la posizione \emph{relativa} della sua arma,
V1 $+$ V2 è la posizione nel mondo dell'arma; e la \textbf{lunghezza} della differenza
tra due posizioni (il nodo \emph{Vector Length}) è la loro distanza. Moltiplicare per un
numero negativo inverte il verso: moltiplicando per $-1$ il vettore ``in avanti'' di un
attore si ottiene quello ``all'indietro'', ed è lo stesso trucco usato nel quinto
capitolo per far cambiare direzione a un bersaglio mobile. Il nodo \textbf{Multiply},
avvertono gli autori, è un nodo \emph{wildcard}: se gli si collega un vettore diventa
un \emph{Vector $\times$ Vector}, e per moltiplicare un vettore per un \emph{float}
bisogna fare tasto destro sul secondo pin e scegliere \emph{Convert Pin\ldots} e poi
\emph{Float}.

\begin{table}[htbp]
\centering
\small
\begin{tabular}{@{}p{3.1cm}p{3.4cm}p{2.0cm}p{5.6cm}@{}}
\toprule
\textbf{Operazione} & \textbf{In formula} & \textbf{Risultato} &
\textbf{A che cosa serve in un gioco} \\
\midrule
Addizione & $(a_x+b_x,\ a_y+b_y,\ a_z+b_z)$ & vettore & spostare un punto: posizione
$+$ spostamento $=$ nuova posizione \\
Sottrazione & $(b_x-a_x,\ b_y-a_y,\ b_z-a_z)$ & vettore & la direzione e la distanza
da A a B (``dove sta il player rispetto a me?'') \\
Moltiplicazione per uno scalare & $(k\,a_x,\ k\,a_y,\ k\,a_z)$ & vettore & allungare o
accorciare; con $-1$ invertire il verso; velocità $\times$ \emph{delta seconds} \\
Lunghezza (\emph{Vector Length}) & $\sqrt{a_x^2+a_y^2+a_z^2}$ & \emph{float} & la
magnitudine; della differenza tra due punti, la distanza \\
Normalizzazione (\emph{Normalize}) & $\vec a \,/\, |\vec a|$ & vettore lungo 1 & tenere
solo la direzione, da moltiplicare poi per una velocità \\
Prodotto scalare (\emph{Dot Product}) & $a_x b_x + a_y b_y + a_z b_z$ & \emph{float} &
capire se due direzioni sono concordi, perpendicolari od opposte \\
\bottomrule
\end{tabular}
\caption{Le operazioni sui vettori della lezione, con il nome del nodo Blueprint
corrispondente dove serve e un uso tipico in un gioco.}
\label{ue4:tab:vettori}
\end{table}

% ---------------------------------------------------------------------
\section{Il prodotto scalare e i vettori normalizzati}
\label{ue4:sec:dot}

Il \textbf{prodotto scalare} è l'operazione che restituisce un \emph{float} invece di un
vettore. ``Che cosa ci dà?'' Ci dice come sono orientati due vettori l'uno rispetto
all'altro. Con i vettori A e B della slide (figura~\ref{ue4:fig:dot}):
\begin{itemize}
  \item se il prodotto scalare è \textbf{maggiore di 0}, i due vettori sono
        ``direzionati nella stessa direzione'', grosso modo concordi; se è
        \textbf{1}, ``sono proprio uguali'', cioè puntano esattamente nello stesso
        verso;
  \item se è \textbf{uguale a 0}, sono ``perfettamente \textbf{perpendicolari}'',
        formano un angolo retto;
  \item se è \textbf{minore di 0}, formano un angolo ottuso, ``al verso opposto'': per
        due vettori come quelli della slide ``magari $-0{,}4$, più o meno''; ``se
        proprio sono così'', uno contro l'altro, è $-1$.
\end{itemize}
Questi valori ``puliti'', tra $-1$ e 1, valgono per i \textbf{vettori normalizzati}, ed
è il motivo per cui la slide successiva li introduce.

\paragraph{I vettori normalizzati.} ``Qualcuno di voi sa cos'è un vettore
normalizzato?'' Uno studente: ``un vettore tra 0 e 1''. ``Esatto'': normalizzando un
vettore, ``indipendentemente dalla sua lunghezza, lo ridurremo a 1''. Un vettore ``che
è lungo 5'', normalizzato, ``lo portiamo a 1''\footnote{Più precisamente: un vettore
normalizzato ha lunghezza \emph{esattamente} 1; sono le sue componenti, prese una a
una, a stare tra $-1$ e 1. Il vettore $(4,3,0)$, lungo 5, normalizzato diventa
$(0{,}8;\ 0{,}6;\ 0)$, come nella figura~\ref{ue4:fig:vettori}.}. La direzione resta la
stessa, cambia solo la lunghezza.

\begin{figure}[htbp]
\centering
\begin{tikzpicture}[font=\small,
  ue4vec/.style={-{Stealth[length=2.5mm]}, very thick}]
  \foreach \xs/\ang/\val/\txt in {0/0/{$= 1$}/{stessa direzione},
                                  3.4/40/{$\approx 0{,}77$}/{concordi},
                                  6.8/90/{$= 0$}/{perpendicolari},
                                  10.2/140/{$\approx -0{,}77$}/{angolo ottuso},
                                  13.6/180/{$= -1$}/{opposti}} {
    \begin{scope}[shift={(\xs,0)}]
      \draw[ue4vec, blue!70!black] (0,0) -- (1.5,0) node[below, pos=0.8,
        font=\scriptsize]{A};
      \draw[ue4vec, red!70!black] (0,0) -- (\ang:1.35) node[above,
        font=\scriptsize]{B};
      \fill (0,0) circle (1.5pt);
      \node[font=\small] at (0.3,-1.0) {A $\cdot$ B \val};
      \node[font=\scriptsize, gray!50!black] at (0.3,-1.5) {\txt};
    \end{scope}
  }
\end{tikzpicture}
\caption{Il prodotto scalare tra due vettori normalizzati A e B, al variare
dell'angolo tra di loro: 1 se puntano nello stesso verso, positivo se sono concordi,
0 se sono perpendicolari, negativo se formano un angolo ottuso, $-1$ se sono opposti.
I valori intermedi sono il coseno dell'angolo ($\cos 40^\circ \approx 0{,}77$).}
\label{ue4:fig:dot}
\end{figure}

\paragraph{Approfondimento: il prodotto scalare è un coseno.} Per due vettori qualsiasi
vale
\[
  \vec a \cdot \vec b \;=\; a_x b_x + a_y b_y + a_z b_z \;=\; |\vec a|\,|\vec b|\,
  \cos\theta ,
\]
dove $\theta$ è l'angolo tra i due. Se i vettori sono normalizzati, le lunghezze valgono
1 e il prodotto scalare \emph{è} il coseno dell'angolo: per questo varia tra $-1$ e 1,
come dicono Romero e Sewell nel quattordicesimo capitolo (``il prodotto scalare di due
vettori normalizzati è uguale al coseno dell'angolo che formano''), e per questo i casi
della slide sono 1, 0 e $-1$. Con vettori non normalizzati il \emph{segno} resta
affidabile (positivo, zero, negativo) ma il valore non è più confrontabile. Nei giochi
il prodotto scalare risponde a domande come ``il player è davanti o dietro al
nemico?'': si normalizza la differenza tra la posizione del player e quella del nemico
e la si confronta con il \emph{Forward Vector} del nemico, che Unreal restituisce già
normalizzato (nodi \emph{Get Actor Forward Vector}, \emph{Right Vector}, \emph{Up
Vector}). Un risultato maggiore di $\cos 45^\circ \approx 0{,}71$, per esempio,
significa ``il player è dentro un cono di visione di 90 gradi''. Il quinto capitolo del
libro usa invece la normalizzazione per il movimento: si normalizza una direzione, la si
moltiplica per una velocità e per il \emph{delta time} (con il nodo \textbf{Get World
Delta Seconds}, che tornerà nella torretta) e la si somma alla posizione a ogni
\emph{tick} con \emph{Add Actor World Offset}. ``Normalizzare'', scrivono, ``è una
procedura comune che assicura che il vettore abbia lunghezza 1, il che lo rende
compatibile con il resto dei calcoli''.

% ---------------------------------------------------------------------
\section{I rotator, il gimbal lock e i quaternioni}
\label{ue4:sec:rotator}

``Andiamo avanti: poi abbiamo i \textbf{rotator}.'' Per spiegarli il docente usa
l'immagine del \textbf{giroscopio}: tre ``anelli che girano tra di loro e formano
proprio quella che è la rotazione''. Ogni anello corrisponde a un asse, con i colori di
sempre: la X in rosso, la Y in verde, la Z in blu (figura~\ref{ue4:fig:rotator}).

\paragraph{Il gimbal lock.} Con le ``regole classiche di Eulero'', cioè descrivendo la
rotazione come tre angoli applicati uno dopo l'altro, c'è ``un possibile problema'': il
\textbf{gimbal lock}. ``Non vi sto a spiegare troppo, perché fondamentalmente andremo a
parlare di quattro dimensioni invece di tre, ed è tutta una cosa molto, molto
complessa.'' In breve, ``potrebbe succedere che due assi si blocchino tra di loro'' e
comincino a ruotare ``sempre dipendenti l'uno dall'altro''. Si risolve con i
\textbf{quaternioni}, che usano quattro valori invece di tre, ma ``parliamo di
matematica molto complessa, e non andiamo molto avanti su questa cosa''. Quello che
conta, per ora: per descrivere una rotazione nello spazio servono sempre la X, la Y e
la Z, ``perché abbiamo tre giroscopi che girano di per sé''. Nell'editor, su un
qualsiasi \emph{asset}, si vedono proprio questi tre valori.

\paragraph{L'obiettivo.} ``Cosa vogliamo fare oggi? Vogliamo creare qualcosa che giri'':
qualcosa che usi il \emph{rotator} per far girare le pale del mulino.

\begin{figure}[htbp]
\centering
\begin{tikzpicture}[font=\small,
  ue4ax/.style={-{Stealth[length=2.5mm]}, very thick},
  ue4rot/.style={-{Stealth[length=2mm]}, thick}]
  % assi in proiezione
  \coordinate (O) at (0,0);
  \draw[ue4ax, red!75!black] (O) -- (-1.9,-1.1) node[below left]{X};
  \draw[ue4ax, green!50!black] (O) -- (2.6,0) node[right]{Y};
  \draw[ue4ax, blue!75!black] (O) -- (0,2.6) node[above]{Z};
  % anelli di rotazione
  \draw[ue4rot, red!75!black] (-1.25,-0.25) arc[x radius=0.35, y radius=0.6,
    start angle=90, end angle=400];
  \node[red!75!black, font=\scriptsize, align=center] at (-2.9,-0.1)
    {\textbf{Roll}\\attorno a X};
  \draw[ue4rot, green!50!black] (1.7,0.55) arc[x radius=0.3, y radius=0.55,
    start angle=90, end angle=400];
  \node[green!50!black, font=\scriptsize, align=center] at (2.4,1.05)
    {\textbf{Pitch}\\attorno a Y};
  \draw[ue4rot, blue!75!black] (0.55,1.75) arc[x radius=0.55, y radius=0.2,
    start angle=0, end angle=310];
  \node[blue!75!black, font=\scriptsize, align=center] at (1.65,2.15)
    {\textbf{Yaw}\\attorno a Z};
  % giroscopio stilizzato
  \begin{scope}[shift={(7.0,0.6)}]
    \draw[very thick, red!75!black] (0,0) ellipse (1.9 and 1.9);
    \draw[very thick, green!50!black] (0,0) ellipse (1.45 and 0.55);
    \draw[very thick, blue!75!black] (0,0) ellipse (0.45 and 1.05);
    \fill (0,0) circle (2pt);
    \node[font=\scriptsize, align=center] at (0,-2.4) {il giroscopio: tre anelli
      annidati,\\ognuno ruota attorno al proprio asse};
  \end{scope}
\end{tikzpicture}
\caption{I tre assi di Unreal con i loro colori e i nomi delle tre rotazioni di un
\emph{rotator}: \emph{roll} attorno a X, \emph{pitch} attorno a Y, \emph{yaw} attorno a
Z. A destra, l'immagine del giroscopio usata a lezione. Schema ricostruito.}
\label{ue4:fig:rotator}
\end{figure}

\paragraph{Approfondimento: roll, pitch, yaw e il gimbal lock in pratica.} Nella seconda
edizione di Romero e Sewell, l'elenco dei tipi di variabile descrive il
\textbf{Rotator} così: contiene i \emph{float} X (\textbf{Roll}), Y (\textbf{Pitch}) e
Z (\textbf{Yaw}), espressi in gradi. I tre nomi vengono dall'aviazione: il
\emph{roll} è il rollio (l'aereo che inclina le ali), il \emph{pitch} il beccheggio (il
muso che sale e scende), lo \emph{yaw} l'imbardata (il muso che gira a destra e a
sinistra). Una torretta che insegue un bersaglio usa lo \emph{yaw} per girarsi e il
\emph{pitch} per alzare la canna; le pale del mulino, come si vedrà, girano in
\emph{pitch}. Il \textbf{gimbal lock} succede quando l'anello di mezzo ruota di 90
gradi: il primo e il terzo anello si ritrovano allineati, e ruotarli produce lo stesso
movimento, così un grado di libertà va perso. Con gli angoli di Unreal il caso tipico
è il \emph{pitch} a $\pm 90$: con il muso dritto verso l'alto, \emph{yaw} e \emph{roll}
fanno la stessa cosa. I \textbf{quaternioni} (in Unreal la struttura \emph{Quat})
rappresentano invece ogni rotazione come una rotazione unica attorno a un asse
qualsiasi, con quattro numeri, e non hanno questo problema. Unreal li usa
internamente: un \emph{rotator} è il formato comodo per le persone, il quaternione
quello sicuro per i calcoli, e nodi come \emph{Combine Rotators} convertono i
\emph{rotator} in quaternioni, li combinano e li riconvertono. Un effetto collaterale,
che conviene conoscere: dato che il \emph{pitch} di un \emph{rotator} viene sempre
riportato tra $-90$ e $90$ gradi, una pala che ruota in \emph{pitch} oltre i 90 gradi
viene descritta ``capovolta'' (\emph{yaw} e \emph{roll} a 180), e chi stampasse il suo
\emph{rotator} vedrebbe numeri che saltano. La rotazione sullo schermo, però, resta
corretta.

% ---------------------------------------------------------------------
\section{Il Rotating Movement Component}
\label{ue4:sec:rotating}

``Allora, avete aperto l'editor di Unreal? Sì? Perfetto.'' Prima di costruire il
mulino, il docente mostra ``quello più semplice'' tra i ``vari modi per far ruotare un
actor in generale''; un modo che però, avverte, ``è difficile usarlo su vari
componenti''.

\paragraph{Un componente aggiunto dai Details.} Si trascina in scena un \textbf{cubo}
(con il tasto \textbf{Fine}, \emph{End}, l'oggetto ``va diretto sul
terreno''\footnote{La trascrizione dice ``il tasto AND'': è il tasto \emph{End} (Fine),
che in Unreal appoggia l'oggetto selezionato sulla superficie sottostante.}). Ora,
``qui sui \textbf{Details} sulla destra, potete aggiungere dei componenti anche non
avendo una schermata di Blueprint'': con il pulsante \textbf{Add} dei Details si può
mettere un componente direttamente sull'actor in scena. Si scrive \emph{rotating
movement} e si sceglie il \textbf{Rotating Movement Component}. ``Trovato? Sì?
Perfetto.''

\paragraph{I parametri.} Il componente ha una serie di dati, tra cui il più importante
è il \textbf{Rotation Rate}, con i tre assi X, Y e Z: la velocità di rotazione su
ciascun asse. ``Di default il Rotating Movement, come potete vedere, è messo a
\textbf{180}'' sulla Z: mezzo giro al secondo attorno all'asse verticale.

\paragraph{La regola della mobilità.} ``Intanto vi dico, ragazzi: come abbiamo visto,
se un \emph{asset} si deve muovere all'interno di uno spazio, la sua mobilità dovrà
essere \textbf{Movable}. Ricordatevelo sempre.'' Se lo si lascia \textbf{Static} e si
va in Play, ``prima di tutto non si muove'', e poi l'engine restituisce un errore,
che il docente legge: \emph{Static Mesh Component has to be movable if you'd like to
move it with Rotating Movement Component}. Con la mobilità corretta, in Play il cubo
gira su sé stesso: ``abbiamo fatto le munizioni di \emph{Metal Gear Solid}'', scherza
il docente, pensando agli oggetti da raccogliere che ruotano a mezz'aria. ``Questo è un
modo per far ruotare'': Unreal ``ci dà la possibilità di poter ruotare direttamente un
nostro bellissimo actor tramite un componente che loro hanno già creato''. ``Un oggetto
ruotante che ruota. Perché `ruotante' è brutto, sembra.''

\paragraph{Esercitazione: i tre cubi.} ``Piccola esercitazione velocissima, faccio
queste botte e risposte con le esercitazioni.'' Si copia il cubo (``tanto vi ricordate
come si copiano'') e si fa in modo che \textbf{un cubo ruoti sull'asse X} e
\textbf{un altro sull'asse Y}, alla velocità che si preferisce. ``Ricordatevi di
metterli Movable. Vi do cinque minuti.'' La correzione: sul secondo cubo si azzera la Z
e si mette 180 sulla X; sul terzo, 180 sulla Y. Si controlla che siano tutti e tre
Movable, si va in Play, ``ed ecco le nostre tre rotazioni''. ``Possiamo farlo con
qualsiasi \emph{asset}.''

\paragraph{Approfondimento: la sedia che gira.} Il Rotating Movement Component è il
primo componente che Romero e Sewell fanno usare nel libro, già nel primo capitolo: un
Blueprint \emph{BP\_RotatingChair} con una \emph{Static Mesh} (la sedia dello
\emph{Starter Content}) e un Rotating Movement. ``Di default'', scrivono, ``questo
componente ruoterà il Blueprint attorno all'asse Z, quindi non serve cambiarne le
proprietà'': si compila, si trascina la sedia nel livello e la si guarda girare in Play.
Il componente fa parte della famiglia dei \emph{movement component}, componenti che non
hanno una forma propria ma muovono un altro componente, quello indicato come
\emph{Updated Component}, che di default è la \textbf{radice} dell'actor. È qui la
ragione del limite che il docente scoprirà tra poco con il mulino: il Rotating Movement
fa girare la radice, e con essa tutto ciò che vi è attaccato. Il componente ha anche
un'opzione \emph{Rotation in Local Space} e un \emph{Pivot Translation}, che sposta il
centro di rotazione; ma per far girare una parte sola di un actor la strada più chiara
resta quella che segue, il Blueprint.

% ---------------------------------------------------------------------
\section{Il Blueprint del mulino}
\label{ue4:sec:mulino}

``Proviamo a fare intanto il nostro \emph{windmill}.'' Il mulino del livello di esempio
è formato da più cose, si vede nei Details. ``Chi si ricorda come si chiamano queste
cose qua, che sono tutte insieme in un singolo \emph{asset}?'' ``Componenti.''
``Sono i \textbf{componenti}'': come si è visto nella seconda lezione, un actor ``può
essere, scusate la ridondanza, composto da più componenti''. In questo caso un
\textbf{body}, il corpo, ``la nostra \emph{mesh} principale, il cubone'', e
\textbf{due pale}.

\paragraph{La classe.} Si crea un nuovo \textbf{Blueprint Class} di tipo
\textbf{Actor}, ``perché non sarà disponibile al giocatore, ma sarà solamente là come
sfondo'': non va controllato, non è un \emph{Pawn}. Il docente lo chiama
\textbf{BP\_Windmill}\footnote{La trascrizione dice ``Windmill, io in questo caso lo
chiamerò 1'': probabilmente \emph{BP\_Windmill1}, per non confonderlo con il mulino già
presente nel progetto.}.

\paragraph{Il corpo.} Nella \textbf{Viewport} del Blueprint si aggiunge una
\textbf{Static Mesh} (Add, in alto a sinistra, \emph{Static Mesh}) e la si chiama
\textbf{Body}. Nei dettagli, alla voce \emph{Static Mesh}, si sceglie un cubo, per
esempio \emph{SM\_Cube}. Il primo cubo scelto dal docente ``ha un \emph{pivot}
particolarmente strano''; ne sceglie un altro, ``molto meglio, perché c'ha il
\emph{pivot} al centro''. Lo solleva di circa 40 sulla Z e gli dà un materiale: il
docente mette la \emph{prototype grid}, ``voi mettete quello che volete''. Poi lo scala
per farne il corpo del mulino: ``3 metri per 5'' di altezza, ``di larghezza facciamo 2
metri'', cioè una scala di circa 3, 2 e 5 sui tre assi\footnote{La trascrizione dice
``x32 sulla y5 sull'altezza'' e ``3:25'': la lettura più probabile è X 3, Y 2, Z 5. Il
cubo di base di Unreal misura un metro di lato, quindi una scala di 3 corrisponde a 3
metri.}.

\paragraph{La viewport che non si trova.} Uno studente non ricorda come si apre la
viewport del Blueprint. ``O ce l'hai qua, o ce l'hai qua, oppure dovresti avere
\textbf{Window} e dovresti vedere \textbf{Viewport}.'' La voce è già spuntata, ma la
finestra non c'è: forse è ``nascosta da qualche schermata di Windows'' (``prova a fare
Alt+Tab''). Si passa alla condivisione dello schermo, e si scopre che il problema è un
altro: lo studente ha creato il Blueprint ma non lo ritrova. La soluzione è il
\textbf{filtro} del Content Browser: le ``tre stanghettine'' vicino alla ricerca,
categoria \emph{Blueprint}, ed ecco il file; doppio clic e, se si apre la versione
ridotta, \emph{Open Full Blueprint Editor}. Poi Add, ma nell'editor in italiano la voce
non si trova subito (è ``Componente Mesh Statica''). ``Effettivamente, ragazzi, vi
consiglio di mettervi l'editor in \textbf{inglese}, perché è molto più semplice. Anzi,
ve lo consiglio proprio'': i nomi dei nodi e dei componenti nei tutorial, nei libri e
nella documentazione sono quasi sempre in inglese.

\paragraph{Le pale.} ``Adesso dobbiamo creare le nostre palette.'' Si copia e incolla il
corpo (Ctrl+C, Ctrl+V): le copie finiscono dentro il corpo, ``quindi cliccatele e
tiratele fuori, perché sono uguali''. Si scalano a forma di pala, una orizzontale e una
verticale, si dà loro un materiale diverso (il docente usa \emph{M\_Grid}, ``questa
qua violacea'') e le si rinomina \textbf{Pala1} e \textbf{Pala2}.

\paragraph{Attaccare le pale al corpo.} ``Importante, ragazzi: queste pale devono essere
attaccate al corpo. Come facciamo? Abbiamo visto l'ereditarietà.'' È la
\textbf{gerarchia} dei componenti: ``il fatto che io posso mettere degli \emph{asset}
attaccati l'uno con l'altro''. Si prende una pala e la si \textbf{trascina sopra} Body,
nell'elenco dei componenti: ora è sua figlia. ``Scusami, non ho capito come si fa.''
``Se hai questa pala, la prendi e la trascini sopra Body.''

\paragraph{Il Rotating Movement non basta.} ``Adesso dobbiamo dare una rotazione a
queste pale. Potremmo provare con il Rotating Movement.'' Il docente lo aggiunge al
Blueprint, mette il mulino in scena (``molto più grosso'') e va in Play. ``Guardate che
succede'': il Rotating Movement è impostato sulla Z, e \textbf{ruota tutto il corpo}.
``Una particolarità del Rotating Movement è che non può essere attaccato a
componenti: o ruota tutto, o non ruota niente.'' Quindi bisogna programmare il
movimento delle pale. Il Rotating Movement va tolto: ``se l'avete messo e avete adesso
il vostro mulino che ruota su sé stesso, lo dovete eliminare''.

% ---------------------------------------------------------------------
\section{Far girare una pala: Event Tick e Combine Rotators}
\label{ue4:sec:pala}

``Domanda. Questo movimento deve essere fatto per ogni \emph{frame}. Quindi deve essere
un evento che riprende ogni \emph{frame}. Chi si ricorda qual è l'evento che lo fa? L'abbiamo
usato per la piattaforma dell'altra volta.'' ``L'\textbf{Event Tick}.'' ``Ottimo.''
L'Event Tick, ``che sarà il nostro protagonista'', è già presente nell'Event Graph di un
nuovo Blueprint.

\paragraph{I delta seconds, di nuovo.} ``Chi si ricorda che sono i \textbf{Delta
Seconds}?'' Uno studente: ``determina quanti tick al minuto\ldots''. Un altro: ``i
secondi effettivi, e non quelli basati su \emph{frame}''. Il docente riprende: è ``lo
spazio che intercorre tra un \emph{frame} e l'altro'', la differenza di tempo tra due
\emph{frame}, ``in questo caso si parla di millisecondi''. Si moltiplica un valore per i
\emph{delta seconds} perché la durata di un \emph{frame} cambia (a 30 \emph{frame} al
secondo è molto più lunga che a 120), e così gli spostamenti ``non diventano più lenti
se abbiamo meno \emph{frame}, o più veloci se ne abbiamo di più''. (Il docente prova a
disegnarlo sullo schermo, ma lo strumento ``a quanto pare non mi funziona benissimo''.)

\paragraph{La velocità.} Si tira fuori un cavo da Delta Seconds e si crea un
\textbf{Multiply}: o scrivendo \emph{multiply}, o digitando semplicemente
l'\textbf{asterisco}, ``e vi dà subito il Multiply''. Il secondo valore è la
\textbf{velocità}: ``più è alto, più è veloce''. Per ora, \textbf{90}.

\paragraph{Make Rotator.} ``Adesso vi dico una cosa.'' Il nodo \textbf{Make Rotator}
costruisce un \emph{rotator}, ``quello che muove a livello di rotazione un
\emph{asset}'', a partire da tre \emph{float}: X, Y e Z. Quale asse serve? ``Prendiamo
la nostra bellissima pala e guardiamo subito in che asse la vogliamo ruotare. La X non
mi sembra, la Z neanche: è la nostra carissima \textbf{Y}'', cioè il \textbf{pitch}.
Si collega il risultato del Multiply al pin Y (Pitch): il \emph{rotator} creato ha X e
Z a 0, ``perché non si devono muovere'', e la Y uguale a velocità per \emph{delta
seconds}.

\paragraph{Combine Rotators.} ``Possiamo farlo in tanti modi, ma oggi ne facciamo tipo
tre.'' Il primo è \textbf{Combine Rotators}, che ``combina due \emph{rotator} per
darne uno in uscita''. Il primo dei due è l'\textbf{origine}: ``a livello concettuale,
qual è l'origine che dobbiamo prendere?'' ``Le coordinate del componente.'' ``Sì:
quale componente?'' Per ora una sola pala, la Pala1. Siccome serve un componente e non
l'intero actor, si trascina la pala dall'elenco dei componenti nel grafo: si ottiene
una \textbf{referenza} a quel componente, ``un \emph{get}'' che ``si riferisce a quel
componente preciso''.

\paragraph{Relativa o del mondo?} Da quella referenza si chiede la rotazione. Ma quale?
``Abbiamo visto la differenza tra \textbf{rotazione relativa} e \textbf{rotazione del
mondo}.'' La pala è attaccata al corpo, ma deve muoversi da sola, indipendentemente dal
corpo. ``Relativa al corpo, secondo me'', propone uno studente. ``Relativa a sé
stessa, quindi la sua asse singola. Se prendiamo quella del mondo, girerà in modo
strano.'' Si usa dunque \textbf{Get Relative Rotation}, la si collega al primo
ingresso di Combine Rotators, e il \emph{rotator} di Make Rotator al secondo.

\paragraph{Set Relative Rotation.} ``Chi si ricorda qual è il blocco che serve per
settare la rotazione? Tanto è un \emph{set}.'' È \textbf{Set Relative Rotation},
``perché, ovviamente, abbiamo preso la relativa''. Lo si crea tirando fuori un cavo da
Pala1 (così il \emph{target} è già la pala), lo si collega all'Event Tick e gli si dà
in ingresso la nuova rotazione, l'uscita di Combine Rotators
(figura~\ref{ue4:fig:pala}). ``E mi raccomando, deve essere \textbf{Movable}.'' Si
mette il mulino in scena, si va in Play: ``ed ecco che una pala ruota''.

\begin{figure}[htbp]
\centering
\begin{tikzpicture}[font=\small,
  ue4node/.style={draw, thick, rounded corners=4pt, align=center, fill=gray!8,
                  minimum height=0.9cm},
  ue4event/.style={ue4node, fill=red!12},
  ue4exec/.style={-{Stealth[length=2.5mm]}, very thick, black!80},
  ue4flt/.style={-{Stealth[length=2mm]}, thick, green!50!black},
  ue4rot/.style={-{Stealth[length=2mm]}, thick, violet!70!black},
  ue4obj/.style={-{Stealth[length=2mm]}, thick, blue!60!black}]
  \node[ue4event, text width=2.1cm] (tick) at (0,0)
    {\textbf{Event Tick}\\\footnotesize Delta Seconds};
  \node[ue4node, text width=1.5cm] (mul) at (2.9,-1.6)
    {\textbf{Multiply}\\\footnotesize $\times$ 90};
  \node[ue4node, text width=1.9cm] (make) at (5.6,-1.6)
    {\textbf{Make Rotator}\\\footnotesize X 0, \textbf{Y}, Z 0};
  \node[ue4node, fill=blue!8, text width=1.3cm] (pala) at (2.9,-3.5)
    {\textbf{Pala1}};
  \node[ue4node, text width=2.1cm] (get) at (5.6,-3.5)
    {\textbf{Get Relative}\\\textbf{Rotation}};
  \node[ue4node, text width=2.0cm, minimum height=1.5cm] (comb) at (8.6,-2.5)
    {\textbf{Combine}\\\textbf{Rotators}\\\footnotesize A, B};
  \node[ue4node, text width=2.2cm, minimum height=1.6cm] (set) at (11.6,0)
    {\textbf{Set Relative}\\\textbf{Rotation}\\\footnotesize Target: Pala1};
  \draw[ue4exec] (tick.east) -- (set.west |- tick.east);
  \draw[ue4flt] ([yshift=-3mm]tick.east) -- ++(0.4,0) |- (mul.west);
  \draw[ue4flt] (mul.east) -- node[above, font=\scriptsize]{Pitch} (make.west);
  \draw[ue4obj] (pala.east) -- (get.west);
  \draw[ue4rot] (get.east) -- ++(0.4,0) |- node[pos=0.75, below,
    font=\scriptsize]{A} ([yshift=-3mm]comb.west);
  \draw[ue4rot] (make.east) -- ++(0.4,0) |- node[pos=0.75, above,
    font=\scriptsize]{B} ([yshift=3mm]comb.west);
  \draw[ue4rot] (comb.east) -- node[above, font=\scriptsize]{New Rotation}
    ++(1.6,0) |- ([yshift=-4mm]set.west);
  \draw[ue4obj, dashed] (pala.south) -- ++(0,-0.5) -| node[pos=0.25, below,
    font=\scriptsize]{Target} (set.south);
\end{tikzpicture}
\caption{Il grafo della prima pala. A ogni \emph{frame}, l'Event Tick esegue Set
Relative Rotation su Pala1; la nuova rotazione è la rotazione relativa attuale della
pala combinata con un piccolo incremento in \emph{pitch}, pari a 90 gradi al secondo
moltiplicati per i \emph{delta seconds}. Il cavo nero è l'esecuzione, gli altri sono
dati (verde i \emph{float}, viola i \emph{rotator}, blu le referenze). Schema
ricostruito.}
\label{ue4:fig:pala}
\end{figure}

\paragraph{I problemi degli studenti.} ``Avete la vostra paletta che ruota?'' Non
tutti. Una studentessa chiede di rivedere il codice e non trova il blocco Get Relative
Rotation: ``tira fuori da Pala1 e scrivi \emph{get relative rotation}''. Un altro
studente vede la pala ``staccarsi da tutto l'asset'' quando preme Play, e in
condivisione si scoprono due errori. Il primo: ha calcolato la rotazione con Combine
Rotators ma non l'ha collegata al Set (``hai calcolato la rotazione, ma non
gliel'hai assegnata''). Il secondo: anche dopo, la pala gira attorno a un punto
sbagliato, perché la \emph{mesh} che ha scelto ha il \textbf{pivot} spostato. ``Vedi che
il pivot sta quasi lì? Cambia la \emph{mesh}'': con un cubo che ha il pivot al centro, la
pala gira correttamente. (Nel frattempo il docente accenna anche a un \textbf{Set World
Rotation}, per mostrare la differenza con la rotazione relativa.)

\paragraph{Approfondimento: relativo e mondo nel libro.} Il quattordicesimo capitolo di
Romero e Sewell spiega la differenza con un esempio molto simile al mulino. Un
Blueprint ha come radice un \emph{DefaultSceneRoot}; sotto c'è un componente
\emph{Table}, un tavolo, e sotto il tavolo un componente \emph{Statue}, una statua. Il
\emph{transform} di Table è \textbf{relativo} a DefaultSceneRoot, e quello di Statue è
relativo a Table: ``se muovete il componente Table nella Viewport, anche il componente
Statue si muoverà; ma se cambiate il \emph{transform} relativo di Statue, Table resterà
dov'è''. Lo stesso vale per il mulino: la rotazione relativa di Pala1 è misurata
rispetto a Body. Se il mulino viene ruotato nel livello, la rotazione relativa della
pala non cambia, e il grafo continua a farla girare attorno al proprio asse; usando la
rotazione del mondo, invece, l'incremento in \emph{pitch} verrebbe applicato attorno
all'asse Y del \emph{mondo}, che coincide con quello della pala solo se il mulino non è
ruotato: è il ``girerà in modo strano'' del docente. Nei vettori lo si vede ancora più
chiaramente: il punto iniziale di ogni posizione nel mondo è sempre $(0,0,0)$, mentre
quello di una posizione relativa è il punto in cui sta il genitore.

\paragraph{Approfondimento: una scorciatoia, Add Local Rotation.} Il trio Get Relative
Rotation, Combine Rotators e Set Relative Rotation è istruttivo perché mostra i tre
passaggi di ogni animazione fatta a mano: \emph{leggi} lo stato attuale, \emph{calcola}
il nuovo, \emph{scrivi} il risultato. Unreal offre anche nodi che fanno tutto in uno:
\textbf{Add Local Rotation} (o \emph{Add Relative Rotation}) su un componente aggiunge
un \emph{Delta Rotation} alla rotazione attuale, così come \emph{Add Actor World
Rotation}, usato nella seconda lezione per la piattaforma, lo fa per un intero actor.
Collegando l'Event Tick a un Add Local Rotation su Pala1, con un Delta Rotation fatto
con Make Rotator come sopra, si ottiene lo stesso mulino con un nodo in meno. Il
risultato è lo stesso; cambia la chiarezza di ciò che succede, ed è per questo che a
lezione conviene vedere prima la versione lunga.

% ---------------------------------------------------------------------
\section{La gerarchia: la seconda pala e il cubetto}
\label{ue4:sec:gerarchia}

``Possiamo fare una cosa bella: invece di dover fare la copia esatta per quanto
riguarda la nostra prima pala, possiamo vedere quello che la \textbf{gerarchia} ci
permette di fare in modo molto facile.'' La Pala2 deve seguire in ogni caso il movimento
della Pala1. Invece di programmare anche lei, la si mette ``in eredità'' (``non quella
con Carlo Conti, ragazzi'') alla Pala1: la si trascina sopra Pala1, e diventa sua
\textbf{figlia}. Play: ``non abbiamo programmato niente'', ma la Pala2 segue la
rotazione della Pala1. ``È la Pala1 che gira, non la Pala2; la Pala2 segue solamente.''

\paragraph{Esercitazione: il cubetto sulla pala.} ``Se noi volessimo attaccare un
cubettino su un estremo della pala\ldots anzi, sapete che c'è? Lo fate voi.'' La
consegna: aggiungere una \textbf{Static Mesh} qualsiasi (``un cubo, una \emph{mesh},
quello che vi pare''), piccola, su un'estremità della pala, e attaccarla alla
\textbf{Pala2}, in modo che segua anche lei il movimento. ``Vi do cinque minuti
classici. Ricordatevi l'ereditarietà: per spostarlo dovete prenderlo proprio e metterlo
sopra.''

\paragraph{La domanda sul pivot.} Uno studente l'ha fatto, ma è sorpreso: ``mi
aspettavo di dover spostare il pivot al centro della rotazione. Invece,
automaticamente, è come se si agganciasse alla pala.'' ``Esatto. Perché comunque
\textbf{mantiene la sua posizione relativa}'': il cubetto ha il suo pivot al proprio
centro, ma la sua posizione è espressa rispetto al genitore, e quando il genitore
ruota, il figlio ruota con lui attorno al pivot del genitore. ``Se lo dovessimo
spostare, andrebbe insieme: manterrebbe la sua origine, il fatto che sta sulla parte
destra, sopra, giù, quello che è, ma si muoverebbe relativo all'altro.'' Play: ``eccolo
che gira insieme''. ``Così abbiamo capito come si ereditano piano piano gli oggetti tra
di loro'' (figura~\ref{ue4:fig:gerarchia}).

\begin{figure}[htbp]
\centering
\begin{tikzpicture}[font=\small,
  ue4comp/.style={draw, thick, rounded corners=3pt, minimum width=2.6cm,
                  minimum height=0.7cm, align=center},
  ue4tree/.style={thick, gray!60!black},
  ue4arr/.style={-{Stealth[length=2.5mm]}, thick}]
  % albero dei componenti
  \node[ue4comp, fill=gray!10] (root) at (0,0) {BP\_Windmill (actor)};
  \node[ue4comp, fill=orange!15] (body) at (0.6,-1.0) {Body};
  \node[ue4comp, fill=violet!12] (p1) at (1.2,-2.0) {Pala1 \footnotesize(ruota)};
  \node[ue4comp, fill=violet!6] (p2) at (1.8,-3.0) {Pala2 \footnotesize(segue)};
  \node[ue4comp, fill=green!10] (cub) at (2.4,-4.0) {Cubetto \footnotesize(segue)};
  \foreach \a/\b in {root/body, body/p1, p1/p2, p2/cub}
    \draw[ue4tree] ([xshift=-1.0cm]\a.south) |- (\b.west);
  \node[font=\footnotesize, align=center] at (1.2,-5.0) {l'elenco dei componenti};
  % mulino stilizzato
  \begin{scope}[shift={(8.2,-2.2)}]
    \draw[thick, fill=orange!15] (-0.75,-2.5) rectangle (0.75,1.0);
    \node[font=\scriptsize] at (0,-2.2) {Body};
    \begin{scope}[rotate=25]
      \draw[thick, fill=violet!25] (-2.2,-0.18) rectangle (2.2,0.18);
      \draw[thick, fill=violet!12] (-0.18,-2.2) rectangle (0.18,2.2);
      \draw[thick, fill=green!35] (-0.25,1.85) rectangle (0.25,2.35);
    \end{scope}
    \fill (0,0) circle (2.2pt);
    \draw[ue4arr, red!70!black] (2.5,0.4) arc[start angle=10, end angle=70,
      radius=2.5];
    \node[font=\scriptsize, align=left] at (3.6,-0.4) {Pala1 orizzontale\\Pala2
      verticale\\cubetto in cima};
  \end{scope}
\end{tikzpicture}
\caption{La gerarchia del mulino. Solo Pala1 è programmata per ruotare; Pala2, sua
figlia, e il cubetto, figlio di Pala2, ne seguono la rotazione senza alcun nodo,
mantenendo la loro posizione relativa. Schema ricostruito (nel mulino della lezione le
pale sono di fronte al corpo).}
\label{ue4:fig:gerarchia}
\end{figure}

% ---------------------------------------------------------------------
\section{Il mulino che guarda il player: Find Look at Rotation}
\label{ue4:sec:lookat}

``Adesso vi farò vedere un altro metodino'' per creare una rotazione, sempre
sull'Event Tick, ``così vi faccio vedere anche un altro blocchettino interessante,
perché sono tanti''. Nei Blueprint, come nella programmazione in generale, ``non esiste
un qualcosa che sia fisso'': ``essendo che è un mondo di spazi, punti, vettori, ce lo
possiamo giocare come ci pare''. Il docente ricorda una persona che una volta ha fatto
``un materiale con la matematica senza usare texture. Non so come abbia fatto, ma si
può fare''.

\paragraph{Staccare i cavi.} Si stacca l'Event Tick dal grafo precedente. ``Ricordatevi:
per staccare i cavettini, tenete premuto \textbf{Alt} e cliccate sui puntini.''

\paragraph{Find Look at Rotation.} ``Un blocco chiamato \textbf{Find Look at Rotation}.
Interessante, molto interessante. Questo è un sacco utile, ve lo dico.'' Il docente
legge la descrizione del nodo: \emph{find a rotation for an object at Start location to
point at Target location}. Alla lavagna (``scrittura incredibile, mi dispiace'') disegna
due punti nello spazio: lo \textbf{Start}, ``da dove parto'', e il \textbf{Target}, il
punto da raggiungere. Il nodo restituisce ``la rotazione verso quel punto'': la
rotazione che un oggetto posto in Start deve avere per guardare Target.

\paragraph{Un falso inizio.} Il docente parte con un esempio sbagliato (lo Start e un
Target ottenuto sommando 360, poi 359, alla posizione), e si ferma: ``scusate, ragazzi,
perdonatemi: il Find Look at Rotation lo facciamo per un esercizio che era diverso.
Ho sbagliato io, perdonate, perché non ho gli appunti.'' L'esercizio vero: far ruotare
il mulino \textbf{in direzione del character}, ``un mulino molto magico'' che si gira
verso il giocatore.

\paragraph{Il grafo.} I due punti sono:
\begin{itemize}
  \item lo \textbf{Start}: la posizione del mulino, con \textbf{Get Actor Location}
        (senza \emph{target}, cioè il Blueprint stesso);
  \item il \textbf{Target}: la posizione del giocatore, con \textbf{Get Player
        Character} e, da questo, Get Actor Location.
\end{itemize}
Il Find Look at Rotation ``ci darà il vettore verso la posizione del nostro player'',
o meglio la rotazione che lo punta.

\paragraph{RInterp To: l'interpolazione.} ``Possiamo usare anche un altro blocco, che ci
darà un'\textbf{interpolazione}.'' Si chiama \textbf{RInterp To} (R come
\emph{rotator}). ``Qualcuno sa cosa vuol dire interpolazione? No.'' ``È letteralmente
il passaggio tra due punti'': da un valore A a un valore B. I nodi di questo tipo
(\emph{Lerp}, \emph{VLerp}, \emph{Interp}) sono ovunque: ``avete presente le animazioni?
Un'animazione è un passaggio da un punto A a un punto B'', e di base è fatta a livello
di valori, perché ruotare un braccio è cambiare i valori di mano, braccio, gamba. Il
nodo ha quattro ingressi:
\begin{itemize}
  \item \textbf{Current}, la rotazione attuale del mulino (Get Actor Rotation, o la
        rotazione relativa del corpo);
  \item \textbf{Target}, la rotazione appena trovata da Find Look at Rotation;
  \item \textbf{Delta Time}, che ``ci dice in quanto tempo'' si fa l'interpolazione;
  \item \textbf{Interp Speed}, ``la velocità in cui lo fa''.
\end{itemize}

\paragraph{Solo lo yaw.} ``Però una cosa importante, ragazzi: non vogliamo che si muova
tutto sulla stessa asse.'' Il mulino deve girare ``al massimo così'', in orizzontale,
non inclinarsi verso l'alto o verso il basso. ``Quale asse dovrò influenzare?''
``La Z?'' ``Esatto, perfetto.'' Si fa \textbf{Split Struct Pin} (tasto destro sul pin)
sull'uscita di Find Look at Rotation e sull'ingresso Target, e si collega solo la Z,
lo \emph{yaw}. Il risultato dell'interpolazione si applica con un nodo di rotazione
(nella trascrizione, \emph{Set Relative Rotation})\footnote{La trascrizione dice ``set
relative rotation di pala 1'', ma a girare verso il player è tutto il mulino. Con ogni
probabilità il nodo è un Set Relative Rotation sul corpo (Body), a cui sono attaccate le
pale, o un Set Actor Rotation sull'intero Blueprint: il risultato visibile è lo
stesso.}, collegato all'Event Tick (figura~\ref{ue4:fig:lookat}).

\paragraph{L'offset di 90 gradi.} In Play il mulino si gira verso il player, ma di
lato. ``Possiamo mettere più 90'': si somma \textbf{90} alla Z prima di collegarla, ``in
modo tale che sia\ldots ecco''. Ora il mulino guarda il giocatore. ``Quello che sta
facendo è continuare a calcolare, visto che lo fa sull'Event Tick, il vettore che
indirizza in un punto, che in questo caso è il player, in modo continuo.'' ``Questo è un
blocco che serve moltissimo, ragazzi'': serve ogni volta che bisogna fare ``delle robe a
\emph{target}'', e servirà per la torretta. ``Vedete che segue continuamente la nostra
torretta? Questo è fatto con il Find Look at Rotation.''

\begin{figure}[htbp]
\centering
\begin{tikzpicture}[font=\small,
  ue4node/.style={draw, thick, rounded corners=4pt, align=center, fill=gray!8,
                  minimum height=0.9cm},
  ue4event/.style={ue4node, fill=red!12},
  ue4exec/.style={-{Stealth[length=2.5mm]}, very thick, black!80},
  ue4flt/.style={-{Stealth[length=2mm]}, thick, green!50!black},
  ue4rot/.style={-{Stealth[length=2mm]}, thick, violet!70!black},
  ue4vecw/.style={-{Stealth[length=2mm]}, thick, yellow!50!black},
  ue4obj/.style={-{Stealth[length=2mm]}, thick, blue!60!black}]
  \node[ue4event, text width=2.0cm] (tick) at (0,0)
    {\textbf{Event Tick}\\\footnotesize Delta Seconds};
  \node[ue4node, text width=2.0cm] (self) at (0,-1.9)
    {\textbf{Get Actor}\\\textbf{Location}\\\footnotesize (self)};
  \node[ue4node, fill=blue!8, text width=2.0cm] (pc) at (0,-3.9)
    {\textbf{Get Player}\\\textbf{Character}};
  \node[ue4node, text width=2.0cm] (pcl) at (2.8,-3.9)
    {\textbf{Get Actor}\\\textbf{Location}};
  \node[ue4node, text width=2.1cm, minimum height=1.4cm] (look) at (6.0,-2.6)
    {\textbf{Find Look at}\\\textbf{Rotation}};
  \node[ue4node, text width=1.2cm] (add) at (8.7,-2.6) {\textbf{$+$ 90}\\
    \footnotesize su Z};
  \node[ue4node, text width=2.0cm] (cur) at (8.7,-4.5)
    {\textbf{Get Actor}\\\textbf{Rotation}};
  \node[ue4node, text width=2.6cm, minimum height=2.2cm] (interp) at (11.5,-2.9)
    {\textbf{RInterp To}\\[2pt]\footnotesize Current\\\footnotesize Target (Z)\\
    \footnotesize Delta Time (Tick)\\\footnotesize Interp Speed};
  \node[ue4node, text width=1.5cm, minimum height=1.3cm] (set) at (14.3,0)
    {\textbf{Set}\\\textbf{Rotation}\\\footnotesize (Z)};
  \draw[ue4exec] (tick.east) -- (set.west |- tick.east);
  \draw[ue4vecw] (self.east) -- node[above, font=\scriptsize]{Start} (4.3,-1.9)
    |- ([yshift=3mm]look.west);
  \draw[ue4obj] (pc.east) -- (pcl.west);
  \draw[ue4vecw] (pcl.east) -- node[below, font=\scriptsize, pos=0.6]{Target}
    (4.5,-3.9) |- ([yshift=-3mm]look.west);
  \draw[ue4flt] (look.east) -- node[above, font=\scriptsize]{Z} (add.west);
  \draw[ue4flt] (add.east) -- ++(0.5,0) |- ([yshift=2mm]interp.west);
  \draw[ue4rot] (cur.east) -- ++(0.5,0) |- ([yshift=-5mm]interp.west);
  \draw[ue4rot] (interp.east) -- ++(0.4,0) |- ([yshift=-4mm]set.west);
\end{tikzpicture}
\caption{Il mulino che guarda il player. A ogni \emph{frame}, Find Look at Rotation
calcola la rotazione che dal mulino (Start) punta al giocatore (Target); se ne tiene
solo lo \emph{yaw} (Z), gli si sommano 90 gradi per compensare l'orientamento della
\emph{mesh}, e RInterp To fa avvicinare gradualmente la rotazione attuale a quella
cercata. Il Delta Time può venire dall'Event Tick, come qui (un cavo non disegnato per
chiarezza), o essere un numero fisso (a lezione, 1); si veda l'approfondimento. Schema ricostruito.}
\label{ue4:fig:lookat}
\end{figure}

\paragraph{``A me non va.''} Seguono tre studenti in condivisione.
\begin{itemize}
  \item Al primo il mulino resta fermo: ``ti manca il Delta Time e
        l'Interp Speed''. Li riempie (``metti uno e due''), compila (``non ho
        compilato, vabbè'') e aggiunge l'offset: ``tira fuori dalla Z, fai più, metti
        90''. ``Perfetto, adesso ti segue.''
  \item Al secondo manca anche a lui il Delta Time: ``metti uno''.
  \item Il terzo, collegando il nodo a Find Look at Rotation, si ritrova un pin
        diviso che non voleva: ha preso \textbf{Get Actor Rotation} invece di
        \textbf{Get Actor Location}. Poi non riesce a collegare l'uscita alla Z:
        ``lo devi splittare: tasto destro sul pinnetto, quello violaceo, quello
        lilla, Split Struct Pin''. E infine l'offset: da Return Value di Find Look at
        Rotation, la Z, un ``più'', 90, e al Target.
\end{itemize}

\paragraph{Approfondimento: che cosa calcola Find Look at Rotation.} Il nodo fa una cosa
semplice: sottrae la posizione di partenza da quella di arrivo, ottenendo il vettore
direzione $\vec d = T - S$ (la sottrazione della sezione~\ref{ue4:sec:vettori}), e
trasforma quella direzione in un \emph{rotator} con \emph{roll} nullo. Per la sola
rotazione orizzontale, quella che interessa al mulino e alla torretta, il conto è
\[
  \text{yaw} \;=\; \operatorname{atan2}(d_y,\ d_x) \quad \text{(in gradi)},
\]
cioè l'angolo, misurato nel piano XY a partire dall'asse X, sotto cui il bersaglio è
visto. Il risultato presuppone che la ``faccia'' dell'oggetto guardi lungo il suo asse
\textbf{X} locale, che in Unreal è la direzione ``in avanti'' (il \emph{Forward
Vector}). Se la \emph{mesh} è stata modellata guardando lungo la Y, come le pale del
mulino della lezione, la rotazione calcolata è sfasata di 90 gradi, e il ``più 90'' del
docente non fa che correggere questo sfasamento. Lo stesso problema, come si vedrà, si
ripresenta con la torretta.

\paragraph{Approfondimento: come funziona RInterp To.} A ogni chiamata, RInterp To
calcola quanto manca per arrivare al Target e ne percorre una frazione, pari
(semplificando) a
\[
  \alpha \;=\; \min\bigl(1,\ \text{Delta Time} \times \text{Interp Speed}\bigr),
  \qquad
  \text{nuova rotazione} = \text{Current} + \alpha \cdot (\text{Target} - \text{Current}),
\]
con la differenza presa sempre ``dal lato più corto'' (da 350 a 10 gradi si passa per
0, non per 180). Ne seguono tre cose utili. Primo: il movimento è veloce quando il
bersaglio è lontano e rallenta avvicinandosi, perché a ogni \emph{frame} si percorre
una frazione di una distanza sempre più piccola; è l'effetto ``morbido'' che si vede in
Play. Secondo: il nodo è pensato per ricevere come \textbf{Delta Time} proprio i
\emph{delta seconds} (dall'Event Tick o da \emph{Get World Delta Seconds}) e come
\textbf{Interp Speed} un numero intorno a 1--10: così la velocità non dipende dal
\emph{frame rate}. Con un Delta Time fisso, come l'1 messo a lezione, la frazione
percorsa a ogni \emph{frame} è sempre la stessa, e il mulino gira più veloce su un PC
che fa più \emph{frame}. Terzo: se il prodotto supera 1, la frazione è 1 e la rotazione
\textbf{scatta} subito al Target; e se l'Interp Speed è 0 o negativo, Unreal restituisce
direttamente il Target. Con Delta Time 1 e Interp Speed 2 il prodotto è 2: il mulino,
in pratica, ``snappa''. Il fenomeno tornerà con la torretta
(sezione~\ref{ue4:sec:problemi}).

% ---------------------------------------------------------------------
\section{Il Clamp, il debug con Print String e la pausa}
\label{ue4:sec:clamp}

``Abbiamo fatto il nostro bellissimo Find Look at Rotation. Potremmo anche
\textbf{clamparlo}.'' ``Si ricorda il Clamp qualcuno di voi?'' ``È quello per dare un
valore minimo'', e massimo. ``Esatto, bravissimo.'' Lo scopo: limitare la rotazione del
mulino, in modo che non possa girarsi del tutto.

\paragraph{Il grafo.} Si splitta il pin con Split Struct Pin, si inserisce un
\textbf{Clamp (float)}, ``perché è un \emph{float}'', sulla Z, e la si ricompone con un
\textbf{Make Rotator} prima di collegarla. Il primo tentativo è un intervallo stretto,
da $-25$ a $25$ gradi, ma qualcosa non torna: i valori in gioco sono molto più grandi
(``mi dice: Pala1 è un botto di valori'').

\paragraph{Il debug con Print String.} ``Questo come ho fatto, ragazzi? Questo è un modo
per \textbf{debuggare} le cose.'' Non sapendo di quanto ruotasse il mulino e quali
limiti dargli, il docente collega un \textbf{Print String} al valore e lo stampa a
schermo mentre si muove attorno al mulino: ``abbiamo stampato letteralmente il
valore''. I numeri che compaiono vanno più o meno da $-160$ a $160$. Il docente prova
quei limiti, poi li allarga: ``snappa un pochino. Facciamo duecento massimo. Molto
meglio. Non dovrebbe andare oltre''\footnote{La trascrizione di questo passaggio è
molto confusa (``facciamo zero, no, mettere meno duecento\ldots{} facciamo duecento
massimo''), e non permette di ricostruire con certezza i valori finali del Clamp.
Quello che conta è il metodo: stampare i valori, osservarli, scegliere i limiti.}.

\paragraph{Approfondimento: Clamp e gli angoli.} Il Clamp restituisce il valore stesso
se è compreso tra Min e Max, altrimenti il limite più vicino:
$\operatorname{clamp}(v, m, M) = \min(\max(v, m), M)$. Con gli angoli c'è una
sottigliezza che spiega i valori ``strani'' della lezione: lo \emph{yaw} restituito da
Find Look at Rotation sta tra $-180$ e $180$, e passando ``dietro'' al mulino salta da
$180$ a $-180$; sommandoci 90 gradi di offset, l'intervallo si sposta (da $-90$ a
$270$). Un limite come $-25 \ldots 25$, applicato a questi numeri, non corrisponde a
``25 gradi a destra e a sinistra della direzione di riposo'', ma a una fetta del cerchio
che dipende da dove sta il player e da come è orientato il mulino nel livello. Il
modo pulito è ragionare sulla differenza tra l'angolo voluto e quello di riposo,
riportarla tra $-180$ e $180$ (Unreal ha il nodo \emph{Normalize Axis}, o \emph{Delta
(Rotator)} per due \emph{rotator}), e solo allora fare il Clamp. Il metodo del docente,
stampare i numeri prima di decidere, resta comunque la prima cosa da fare.

\paragraph{La pausa.} ``Ci siete? C'è qualche cosa che non avete capito? C'è qualche
blocco che vi risulta ancora un pochino nebuloso?'' Tutto chiaro. ``Sono le 17 e 28.
Facciamo pausa, così poi riprendiamo e cominciamo a fare la nostra bellissima
torrettina. Ci vediamo tra 20 minuti.''

% ---------------------------------------------------------------------
\section{Dopo la pausa: la torretta in tre pezzi}
\label{ue4:sec:torretta}

``Eccoci qua. Allora, ragazzi, adesso andremo a fare una bellissima torrettina.'' Il
docente mostra la versione finita. Attorno alla torretta c'è una \textbf{sfera} disegnata
con linee rosse: è il suo raggio d'azione. Quando il giocatore ci entra, la torretta
``fa lo stesso processo che abbiamo visto'', cioè ruota verso il \emph{character}, e
``in questo caso spara dei proiettili''. Non ci sono effetti speciali (``non abbiamo
Niagara''\footnote{La trascrizione dice ``BevoFX'': con ogni probabilità
\emph{Niagara VFX}, il sistema di effetti visivi di Unreal Engine, di cui Davide aveva
parlato nella lezione precedente.}), ma si costruisce ``la logica che ci servirà con la
torretta che avete \emph{mergiato}''. Si vedono anche delle linee che partono dal
cannone: sono un \textbf{line trace}, ``che permette a noi di avere una corrispondenza
dipendente da cosa incontra''. Oggi si fa la rotazione; proiettili e \emph{line trace}
saranno la prossima volta.

\paragraph{Il Blueprint.} ``La prima cosa che ovviamente dobbiamo fare è creare un
nuovo Blueprint, sempre \textbf{Actor}'', chiamato \textbf{BP\_Turret}. Ma come hanno
composto la torretta gli studenti, con Davide? ``Abbiamo preso i tre pezzi forniti e li
abbiamo uniti'' con \textbf{Merge Actors}: una \emph{mesh} unica. ``Allora facciamo una
cosa: la facciamo a \textbf{tre pezzi}.'' Per far girare il supporto e il cannone
indipendentemente dalla base servono tre componenti distinti, quindi si torna ai tre
pezzi importati con Davide: tre \textbf{Static Mesh} nel Blueprint, una per la
\textbf{base}, una per il supporto (``l'\emph{holder}, chiamiamolo'', poi rinominato
\textbf{Pivot}) e una per il \textbf{cannone} (``la \emph{gun}'', poi
\textbf{Cannon}).

\paragraph{Chi segue chi.} I tre pezzi vanno ``uno dentro l'altro''. Il docente parte
con l'idea che il supporto segua il cannone, poi si corregge: ``fate il contrario, che
è meglio: fate che la \emph{gun} segua l'\emph{holder}''. La gerarchia giusta è quindi:
la \textbf{base} come radice, il \textbf{Pivot} figlio della base, il \textbf{Cannon}
figlio del Pivot (figura~\ref{ue4:fig:turretbp}). Così, quando il Pivot ruota
orizzontalmente, il cannone gira con lui; e un giorno il cannone potrà alzarsi e
abbassarsi da solo, senza trascinarsi dietro il supporto.

\paragraph{Un trucco: Browse to Asset.} Per assegnare le \emph{mesh} ai componenti non
serve cercarle a mano nel Content Browser ``con l'elenchino e tutto quanto''. ``Andate
qua sull'\emph{asset} che avete in scena, cliccate tasto destro e fate \textbf{Browse to
Asset}'': il Content Browser salta direttamente alla posizione del file, e da lì lo si
può \textbf{trascinare} (\emph{drag and drop}) nello slot \emph{Static Mesh} del
componente. In alternativa, come sempre, si clicca sullo slot e si scrive il nome:
\emph{turret base}, \emph{turret pivot}, \emph{cannon}. Poi si aggiustano le posizioni
(``un buchettino sta un po' più giù'') finché i pezzi combaciano.

\paragraph{Un problema di Add.} Una studentessa, premendo Add, vede solo la \emph{mesh}
intera della torretta (\emph{SM\_\allowbreak Torretta}) e nessun'altra opzione. In condivisione si
prova a compilare e a riaprire il Blueprint, e alla fine il componente si riesce ad
aggiungere: ``va benissimo così, abbiamo risolto. A volte Unreal può fare dei giochi
simpatici.''

% ---------------------------------------------------------------------
\section{Il raggio d'azione: la Sphere Collision}
\label{ue4:sec:sfera}

``Adesso dobbiamo creare quello che abbiamo visto qua: vedi questo cerchio rosso intorno
alla torretta? Questo è una \textbf{Sphere Collision}.'' È il \textbf{range}: entrandoci
dentro, la torretta si attiva.

\paragraph{Aggiungerla.} In BP\_Turret si aggiunge, tra i componenti, una
\textbf{Sphere Collision}, ``che tranquillamente può stare anche esterna alla base'',
cioè non deve essere figlia di un pezzo che ruota. Appena aggiunta non si vede quasi,
``semplicemente perché è molto piccola''. Cliccando su \emph{Sphere}, nei dettagli c'è
lo \textbf{Sphere Radius}: di default è 32; il docente lo porta a \textbf{800}. C'è anche
la \textbf{Line Thickness}, lo spessore delle linee con cui la sfera è disegnata
nell'editor, ``che serve magari solo per vederla un po' più spessa''.

\paragraph{Hidden in Game.} In Play, però, la sfera non si vede. ``Di base, questa è una
\textbf{collisione}'', una cosa interna agli \emph{asset}, come le collisioni della
\emph{mesh} viste con Davide: esiste per la logica, non per essere guardata. Per
vederla mentre si prova il gioco (``per avere facilità nel vedere quando entriamo e
usciamo'') si cerca nei dettagli \textbf{Hidden in Game} e lo si toglie. ``Parlando di
\emph{rendering}, Hidden in Game dice: so che logicamente esiste all'interno del mondo,
ma se non è \emph{hidden in game}, la renderizzo pure.'' Tolta la spunta, il valore è
\emph{false}, ``quindi non è nascosto in gioco'', e in Play ecco il range. È la stessa
opzione usata nella seconda lezione, al contrario, per rendere invisibile il soffitto
della vaschetta.

\paragraph{Ripetere i passaggi.} Uno studente ha perso questa parte mentre cercava i
pezzi della torretta, e chiede di rivedere come si mettono le \emph{mesh} nei tre
componenti. Il docente lo rifà: si clicca sul componente, si scrive il nome della
\emph{mesh} nello slot Static Mesh (\emph{turret base}, \emph{turret pivot}), si
rinominano i componenti in modo che siano coerenti (\textbf{Pivot}, \textbf{Cannon},
``come a me è chiamata Cannon nella nostra \emph{mesh}'').

% ---------------------------------------------------------------------
\section{Gli eventi dei componenti e l'overlap}
\label{ue4:sec:overlap}

``Adesso andiamo a vedere un evento in particolare.'' Oltre agli eventi già visti, che
appartengono all'intero actor (l'Event Begin Play, l'Event Actor Begin Overlap,
l'Event Tick), ``possiamo richiamare degli eventi anche all'interno dei
\textbf{componenti}''. Selezionando \emph{Sphere} e scorrendo in fondo ai dettagli, c'è
una sezione \textbf{Events} piena di pulsanti: \emph{On Component Begin Overlap},
\emph{On Component End Overlap}, \emph{On Component Wake}, \emph{On Component Sleep},
gli eventi del cursore e altri. ``Tutta una serie di eventi'' che si possono usare
``per far sì che avvenga qualcosa in un preciso momento''. Qui la torretta deve partire
quando il player entra nella sfera: si clicca sul \textbf{+} di \textbf{On Component
Begin Overlap}, e nel grafo compare l'evento, già legato alla Sphere.

\paragraph{Che cos'è un overlap.} ``Sapete cosa si intende per \emph{overlap}?''
Uno studente: ``tipo appena si passa un confine, qualcosa del genere''. ``In un certo
senso.'' Alla lavagna (``perdonate le mie doti molto brutte di disegno'') il docente
disegna la sfera con la torretta al centro: un \emph{overlap} avviene quando il player,
``o un qualsiasi actor'', \textbf{tocca o entra} in quel volume. Un actor fuori non sta
sovrapponendosi; uno dentro sì. E uno a cavallo del bordo? ``Sta overlappando o non sta
overlappando?'' ``Sì.'' ``Qualcun altro ha un'opinione differente?'' Sta
sovrapponendosi, ``perché anche un qualsiasi poligono da cui è formato è parte
integrante di esso'': basta un pezzo (figura~\ref{ue4:fig:overlap}). Ci sono quindi tre
situazioni: non si sovrappone, \textbf{comincia} a sovrapporsi (Begin Overlap),
\textbf{smette} di sovrapporsi (End Overlap, ``quando ci esce''). Per ora serve solo la
prima: ``vogliamo che quando lui entra in questo volume la torretta faccia tipo: ti ho
visto, mo' comincio a ruotare verso di te''.

\begin{figure}[htbp]
\centering
\begin{tikzpicture}[font=\small,
  ue4arr/.style={-{Stealth[length=2.5mm]}, thick}]
  \draw[red!70!black, thick, fill=red!4] (0,0) circle (2.4);
  \node[red!70!black, font=\scriptsize] at (0,2.65) {Sphere Collision (raggio 800)};
  \draw[thick, fill=orange!35] (-0.35,-0.3) rectangle (0.35,0.3);
  \draw[thick, fill=orange!20] (0.35,-0.08) rectangle (0.9,0.08);
  \node[font=\scriptsize] at (0,-0.55) {torretta};
  % fuori
  \draw[thick, fill=blue!15] (-4.3,0.6) circle (0.32);
  \node[font=\scriptsize, align=center] at (-4.3,-0.2) {fuori:\\nessun overlap};
  % a cavallo
  \draw[thick, fill=blue!15] (1.75,1.65) circle (0.32);
  \node[font=\scriptsize, align=left, anchor=west] at (2.3,1.9) {a cavallo del bordo:\\
    \textbf{overlap} (basta un pezzo)};
  % dentro
  \draw[thick, fill=blue!15] (-1.0,-1.25) circle (0.32);
  \node[font=\scriptsize, align=center] at (-1.0,-1.95) {dentro: overlap};
  % frecce degli eventi
  \draw[ue4arr, green!45!black] (-3.9,0.45) -- (-2.55,0.25);
  \node[green!45!black, font=\scriptsize, align=center] at (-3.4,1.45)
    {entra:\\\textbf{Begin Overlap}};
  \draw[ue4arr, orange!80!black] (1.95,-1.55) -- (3.2,-2.3);
  \node[orange!80!black, font=\scriptsize, align=left, anchor=west] at (3.3,-2.3)
    {esce: \textbf{End Overlap}};
\end{tikzpicture}
\caption{Le situazioni dell'\emph{overlap} disegnate a lezione. Un actor si sovrappone
alla sfera appena una parte qualsiasi della sua forma di collisione ne entra; Begin
Overlap scatta una volta, nel momento in cui entra, End Overlap una volta, quando ne
esce del tutto. Schema ricostruito.}
\label{ue4:fig:overlap}
\end{figure}

\paragraph{Other Actor.} ``Come facciamo a capire quale actor ci dà indietro?'' L'evento
ha una serie di uscite: \emph{Overlapped Component}, \textbf{Other Actor}, \emph{Other
Comp}, \emph{Other Body Index}, \emph{From Sweep}, \emph{Sweep Result}. ``Sono tutti
input che sono risultanti dall'overlap'', ``in parole povere'' le informazioni su chi si
è sovrapposto a che cosa. Per mostrarlo, il docente collega Other Actor a un
\textbf{Print String} (Unreal inserisce da solo un \emph{Get Display Name}, per
trasformare l'actor in testo), va in Play ed entra nella sfera: in alto a sinistra
compare \emph{BP\_ThirdPersonCharacter\_0}\footnote{La trascrizione dice ``BP ha perso
un character 0''. Il suffisso \_0 è il numero che Unreal assegna a ogni istanza di una
classe presente nel livello.}. ``Mi ridà cosa sta overlappando, che in questo caso è il
\emph{character}.'' Esistono anche gli \emph{overlapped component} e gli \emph{other
component}, ma quello che interessa è l'Other Actor, ``perché il nostro
\emph{character} è un actor''.

\paragraph{Approfondimento: perché l'overlap funziona ``da solo''.} Perché Unreal generi
un evento di \emph{overlap} servono tre condizioni: entrambi i componenti devono avere
attivo \emph{Generate Overlap Events}; le loro risposte di collisione reciproche devono
essere \emph{Overlap} (e non \emph{Ignore}); e almeno uno dei due deve muoversi. La
Sphere Collision nasce con il profilo di collisione \emph{OverlapAllDynamic}, che si
sovrappone a tutto ciò che si muove, e la capsula del \emph{character} è un oggetto
dinamico di tipo \emph{Pawn}: per questo, nella lezione, non c'è stato bisogno di
toccare nulla. Se un giorno un evento di \emph{overlap} non scatta, è la prima cosa da
controllare, nei dettagli di entrambi i componenti, sezione \emph{Collision}. La seconda
edizione di Romero e Sewell elenca proprio gli eventi di collisione (``eseguiti quando
due actor collidono o si sovrappongono'') tra le famiglie principali di eventi, accanto
a quelli di input e all'Event Begin Play.

% ---------------------------------------------------------------------
\section{Is Valid: esiste davvero?}
\label{ue4:sec:isvalid}

``Vogliamo essere sicuri che il nostro actor in questione sia proprio il player. Perché
un actor può essere anche un cubo, volendo.'' Prima, però, un controllo più semplice. Si
tira fuori da Other Actor un \textbf{Is Valid}: ``ce ne avete tanti, ma quello che ci
interessa è quello col \textbf{punto interrogativo}'', la \emph{macro} con due uscite di
esecuzione, \textbf{Is Valid} e \textbf{Is Not Valid}.

\paragraph{Che cosa controlla.} ``Is Valid, statemi a sentire, ci dà una risultante: è
valido o non è valido?'' Qui chiede: ``hai overlappato qualcosa? È vero, esiste questa
cosa?'' In generale, Is Valid serve a dire se esiste o no qualcosa, ``che può essere
anche una \textbf{variabile}, quindi una variabile vuota''.

\paragraph{L'esempio della variabile.} Per farlo vedere il docente costruisce un piccolo
esperimento:
\begin{enumerate}
  \item da Other Actor crea una variabile con \textbf{Promote to Variable};
  \item mette un Is Valid su quella variabile, all'inizio dell'evento di \emph{overlap};
  \item sul ramo \textbf{Is Not Valid} assegna alla variabile l'actor appena
        sovrapposto (il \emph{Third Person Character});
  \item sul ramo \textbf{Is Valid} stampa con un Print String ``valore registrato'' e il
        nome contenuto nella variabile.
\end{enumerate}
Prima entrata nella sfera: non succede nulla di visibile, ``è normale''. Il Blueprint ha
chiesto: ``questa variabile ha qualcosa dentro? No, non c'è niente. Quindi non è
valida'', e allora le ha assegnato l'actor. Si esce e si rientra: ``ah, me l'hai fatto
prima, quindi questa cosa c'ha qualcosa dentro'', e il ramo è Is Valid: compare
\emph{valore registrato: BP\_ThirdPersonCharacter\_0}. ``Qualcuno ha qualche dubbio
riguardo all'Is Valid? Abbastanza chiaro?'' L'esperimento si cancella: ``ci serve
solamente a capire cosa fa l'Is Valid''.

\paragraph{Approfondimento: la regola di Romero e Sewell.} Il quarto capitolo di
\emph{Blueprints Visual Scripting for Unreal Engine 5}, dedicato alla comunicazione tra
Blueprint, mette in evidenza una nota: ``usate sempre la macro Is Valid prima di
eseguire una funzione usando una variabile che contiene una referenza a un oggetto. La
variabile può non essere valida per diversi motivi, ed eseguire una funzione con una
variabile non valida causerà un errore durante l'esecuzione''. I motivi tipici sono
due: la variabile non è mai stata riempita (come nella prima entrata dell'esempio), o
l'oggetto a cui puntava è stato distrutto nel frattempo. In entrambi i casi un nodo che
prova a leggere la posizione di ``nessuno'' produce nel registro il classico errore
\emph{Accessed None}. Il libro usa lo stesso schema in molti punti: un Is Valid prima di
distruggere un actor, prima di applicare un danno a un bersaglio, prima di leggere un
salvataggio (se non è valido, lo si crea). La torretta lo userà due volte, nella
prossima sezione e nell'evento Move Turret.

% ---------------------------------------------------------------------
\section{I tag e Actor Has Tag}
\label{ue4:sec:tag}

Is Valid dice che \emph{qualcosa} è entrato. Per sapere \emph{chi}, si tira fuori da
Other Actor un \textbf{Actor Has Tag}, ``un blocco che va a esaminare una cosa chiamata
\textbf{tag}''. I tag ``sono dei \emph{label}, delle \textbf{etichette} che noi possiamo
mettere all'interno di vari \emph{asset} del gioco''.

\paragraph{Mettere il tag al player.} Si apre il \textbf{Third Person Character} e,
``non dobbiamo fare niente di particolare'', si scrive \emph{tag} nella ricerca dei
dettagli: compare \textbf{Actor} $\rightarrow$ \textbf{Tags}, un elenco vuoto con un
\textbf{+}. Si aggiunge un elemento e lo si chiama \textbf{Player}. ``Cosa abbiamo
fatto? Abbiamo letteralmente etichettato il player con un tag Player'': è come avere un
cartellino di riconoscimento sulla schiena, con scritto sopra ``tu sei il player''.

\paragraph{Il controllo.} Tornati nella torretta, Actor Has Tag chiede: il
\emph{target}, cioè l'actor sovrapposto (Other Actor), ha il tag \textbf{Player}? Il
nodo restituisce un valore \textbf{booleano}: ``sì, questo actor che ho overlappato ha
il tag Player'', e si va nel \textbf{True} di un \textbf{Branch}; ``se no, invece andrà
in False''. Il docente lo mostra con due Print String, ``Player trovato'' sul ramo True e
``Player non trovato'' sul ramo False: entrando nella sfera, ``Player trovato''.

\paragraph{Attenzione all'ortografia.} ``Non lo fate, ragazzi, faccio solo vedere per
esempio.'' Il docente modifica il tag, togliendo la R finale, e rientra: ``Player non
trovato''. ``Questo vi fa vedere anche che dovete essere precisi nel dire qual è il
tag.'' Il consiglio: ``fate Ctrl+C, Ctrl+V del tag che avete scritto'', copiandolo dal
\emph{character} nel nodo invece di riscriverlo a mano.

\paragraph{Approfondimento: tag, Cast o confronto?} Ci sono almeno tre modi per
rispondere alla domanda ``l'actor che è entrato è il player?''. Il primo è il
\textbf{tag}, quello della lezione: semplice, e indipendente dalla classe (lo stesso tag
\emph{Player} può stare su personaggi diversi, e un tag \emph{Enemy} su nemici di classi
diverse); esiste anche \emph{Component Has Tag}, per etichettare singoli componenti. Il
secondo è il \textbf{Cast To} (per esempio \emph{Cast To BP\_ThirdPersonCharacter}),
che verifica la classe e, se riesce, restituisce l'actor già ``convertito'', con
accesso alle sue variabili e funzioni. Il terzo è il confronto diretto tra oggetti: nel
capitolo sull'intelligenza artificiale della seconda edizione di Romero e Sewell, un
nemico che pattuglia usa un \emph{Event Actor Begin Overlap} seguito da un Branch su un
nodo \textbf{Equal (Object)}, che controlla se l'Other Actor è proprio il punto di
pattuglia salvato in una variabile. Una curiosità sui tag: sono di tipo \emph{Name}, e
in Unreal il confronto tra \emph{Name} \textbf{non distingue maiuscole e minuscole},
per cui \emph{player} e \emph{Player} risultano uguali; quello che conta, come ha
mostrato il docente, è che le lettere siano le stesse. Copiare e incollare il tag
resta comunque la via più sicura.

% ---------------------------------------------------------------------
\section{La variabile del bersaglio e i Custom Event}
\label{ue4:sec:custom}

Ora che il player è stato trovato, bisogna ricordarselo: ``andiamo a crearci una
variabile in cui mettiamo il nostro player che è stato trovato'', di tipo
\textbf{Actor}. Il modo più veloce: tirare fuori un cavo da Other Actor e scegliere
\textbf{Promote to Variable}. Si crea un nodo \emph{Set} di una variabile nuova, che si
rinomina \textbf{PlayerTarget}. Per rinominare, ``o vai qua e fai \textbf{F2}, che è più
facile, oppure fai tasto destro e Rinomina''.

\paragraph{Un Custom Event.} ``Manca mezz'ora'', e il docente si chiede se si riuscirà a
finire (``questo lo continueremo la prossima volta''), ma va avanti. ``Adesso qua dentro
andremo a fare un \textbf{Custom Event}.'' Si scrive \emph{custom event} nel grafo, si
sceglie \emph{Add Custom Event} e gli si dà un nome: \textbf{TrackTargetTimer}. ``Come
potete vedere, questo bellissimo evento è un Custom Event: è uguale agli altri'', come
il Begin Overlap, il Tick o il Begin Play, ``ma è un evento che abbiamo creato noi'',
``che possiamo richiamare quando vogliamo, con delle regole specifiche''. Gli eventi di
Unreal sono ``prefissati'': l'\emph{overlap} scatta perché l'editor di Unreal lo
controlla già per conto suo. I Custom Event, invece, partono quando li si chiama.

\paragraph{Chiamarlo.} Per mostrarlo, il docente collega al Custom Event un Print String
che stampa ``Hello'' e lo chiama dal ramo del Begin Overlap, subito dopo il Set di
PlayerTarget: si scrive il nome dell'evento (\emph{TrackTargetTimer}) e compare il nodo
che lo invoca. In Play, entrando nella sfera: ``Hello''. ``Ho creato un evento
customizzato che richiamo quando entro nella Sphere Collision. Compreso?'' ``Sappiate
che di Custom Event ne potete creare a milioni'': basta scrivere \emph{custom event} e
dare un nome.

\paragraph{Approfondimento: eventi, funzioni e macro.} La seconda edizione di Romero e
Sewell ricorda che nell'Event Graph ogni evento di Unreal può comparire una sola volta,
mentre i \textbf{Custom Event} si possono creare a volontà; e offre una tabella per
scegliere tra \emph{macro}, \emph{funzioni} ed \emph{eventi}. In sintesi: tutti e tre
accettano parametri in ingresso; solo macro e funzioni restituiscono valori; le funzioni
e gli eventi possono essere chiamati da altri Blueprint, le macro no; e solo macro ed
eventi possono contenere azioni \emph{latenti}, cioè che aspettano (come un
\emph{Delay}), mentre le \emph{Timeline} stanno solo negli eventi. Il capitolo sui
vincoli di gioco della stessa edizione usa proprio un Custom Event, \emph{Sprint Drain},
richiamato da un timer, per consumare la stamina del giocatore mentre corre: ``un Custom
Event permette di attivare i nodi collegati a lui ogni volta che un altro nodo lo
chiama; così gruppi di nodi dello stesso Blueprint possono comunicare tra loro anche
quando non sono collegati direttamente''. È lo schema esatto della torretta.

% ---------------------------------------------------------------------
\section{Il timer: Set Timer by Event}
\label{ue4:sec:timer}

``Cosa vogliamo fare con questo Custom Event?'' Durante tutto il tempo, ``come un Event
Tick'', la torretta deve cercare la posizione verso cui ruotare per seguire il player.
``Lo potremmo fare tranquillamente nell'Event Tick, volendo, ma possiamo fare una cosa
interessante: possiamo fare un \textbf{timer}. L'abbiamo chiamato timer apposta.''

\paragraph{Il nodo.} Dal Custom Event si tira fuori \textbf{Set Timer by Event}, un
blocco ``che ci permette di far partire un timer che ci chiede ogni quanto''. I suoi
parametri, così come li imposta il docente:
\begin{itemize}
  \item \textbf{Time}: ogni quanti secondi scatta. ``Vogliamo che lo faccia allo stesso
        tempo del Tick.'' Si potrebbe portare fin qui il cavo dei \emph{delta seconds}
        dall'Event Tick (``ma è un po' brutto''); c'è invece un nodo apposta,
        \textbf{Get World Delta Seconds}, ``un blocchetto semplicemente per comodità'',
        ``per non far sì di filare il Tick così, a cavolo'';
  \item \textbf{Looping}: spuntato, ``vogliamo che sia un timer che \emph{loopi}'';
  \item \textbf{Max Once Per Frame}: spuntato anche questo, ``massimo una volta per
        ogni \emph{frame}''. ``Ci permette almeno di dire: ok, per ogni \emph{frame} si
        calcola, tanto effettivamente si muoverà per ogni \emph{frame}. Questa è una
        cosa un po' più nel micro, però noi mettiamo \emph{max}, in modo che sia più
        ottimizzato.''
\end{itemize}

\paragraph{Quale evento?} ``Set Timer by Event: quale evento?'' Sul nodo c'è un
\textbf{quadratino rosso}, il pin \textbf{Event}, ``che guarda caso è uguale a quello del
Custom Event'': un pin di tipo \emph{delegato}. Vuol dire che lo si può tirare fuori e
creare da lì un \textbf{altro Custom Event}, quello che il timer chiamerà a ogni giro:
lo si chiama \textbf{MoveTurret}. È lì che si farà muovere la torretta.

\paragraph{Come ragiona.} Per farlo vedere, a MoveTurret si collega il solito Print
String, ``Hello''. La catena completa è: Begin Overlap, Is Valid, Branch su Actor Has
Tag, Set PlayerTarget, TrackTargetTimer; e TrackTargetTimer fa partire un timer, alla
velocità dei \emph{delta seconds}, in \emph{loop}, al massimo una volta per
\emph{frame}, che chiama MoveTurret, che stampa ``Hello''. ``Domanda: quando mi
stamperà Hello?'' ``Da quando entri in range.'' ``Ottimo, esatto.'' Appena il player
entra, Hello a ogni \emph{frame}. ``Esco fuori: si cancellano\ldots{} ah, no, scusate,
non si stanno cancellando, è vero. Perché non stoppiamo mai il timer.'' Dalla prima
entrata in poi, il timer continua all'infinito: ``ciò vuol dire che il timer a un certo
punto va disattivato; lo vedremo dopo''.

\paragraph{Approfondimento: timer o Event Tick?} Romero e Sewell presentano i timer, già
nel secondo capitolo, come un'alternativa all'Event Tick: tutto ciò che è collegato al
Tick viene eseguito a ogni \emph{frame}, e ciò che non ha bisogno di aggiornarsi così
spesso può girare a intervalli più lunghi. Nel quinto capitolo un Custom Event
\emph{ChangeDirection}, chiamato da un Set Timer by Event con Looping attivo, fa
invertire periodicamente la direzione di un bersaglio mobile; nella seconda edizione un
timer chiama ripetutamente \emph{Sprint Drain}, e un \emph{Clear Timer by Function Name}
lo ferma quando il giocatore smette di correre. Nella torretta della lezione l'intervallo
è pari ai \emph{delta seconds}, quindi il timer scatta (circa) una volta per
\emph{frame}, come un Tick: il vantaggio non è la frequenza, ma il fatto che il timer
\textbf{parte solo quando serve} (all'ingresso del player) e \textbf{si può fermare}.
Tre dettagli utili. Primo: il valore di Get World Delta Seconds viene letto una sola
volta, quando il timer viene impostato, e resta fisso anche se poi il \emph{frame rate}
cambia. Secondo: se l'intervallo di un timer è più breve di un \emph{frame}, Unreal può
eseguirlo più volte nello stesso \emph{frame} per ``recuperare''; \textbf{Max Once Per
Frame}, un'opzione delle versioni recenti di Unreal Engine 5, impedisce proprio questo.
Terzo: per fermarlo, Set Timer by Event restituisce un \textbf{Timer Handle}, una
``maniglia'' che conviene salvare in una variabile (Promote to Variable) e passare,
quando serve, a \textbf{Clear and Invalidate Timer by Handle}. Il posto naturale per
farlo è l'evento \textbf{On Component End Overlap} della sfera, insieme a
\emph{svuotare} la variabile PlayerTarget (un Set senza nulla collegato): altrimenti
l'Is Valid di MoveTurret continuerà a trovare un bersaglio valido anche dopo che il
player è uscito dal range.

% ---------------------------------------------------------------------
\section{Move Turret: la torretta che segue il player}
\label{ue4:sec:move}

``Cosa dobbiamo fare? L'abbiamo fatto prima, praticamente'': un \textbf{Find Look at
Rotation}, ``perché dobbiamo ruotare il nostro Pivot in direzione del player''.

\paragraph{Prima, un Is Valid.} ``Intanto facciamo un Is Valid. Perché? Perché diciamo:
ma ce l'ho un target?'' Si trascina nel grafo la variabile \textbf{PlayerTarget} e la si
collega a un Is Valid, in testa a MoveTurret. ``Ce l'ho un target effettivo, o non ce
l'ho più, l'ho perso, sono uscito fuori? Se ce l'ho, lo faccio e mi muovo. Se non ce
l'ho, non lo faccio'': senza questo controllo la torretta ``continuerebbe a ruotare a
cavolo''.

\paragraph{Start e Target.} ``Del Find Look at Rotation ci servono due punti nello
spazio. Quali sono?'' ``Lo Start è l'\emph{holder}, che tanto la \emph{gun} è collegata
all'\emph{holder}, e il Target è il player.'' ``Esatto.'' Per lo Start si usa
\textbf{Get Actor Location} della torretta stessa, ``la locazione generale
dell'actor''. Per il Target, il docente comincia con Get Player Character e Get Actor
Location, come per il mulino, ma uno studente fa notare che si può usare
\textbf{PlayerTarget}: ``bravo, mi piace che l'hai notato, effettivamente sarebbe stato
il secondo metodo''. Da PlayerTarget si tira fuori Get Actor Location, ``una sua
particolarità'', e la si collega al Target.

\paragraph{Lo yaw e RInterp To.} Il Pivot deve ruotare attorno alla \textbf{Z}:
``guardate un po' dove si modifica: sulla Z''. Si splitta l'uscita di Find Look at
Rotation e si fa ``quello che abbiamo fatto prima'', un \textbf{RInterp To}:
\begin{itemize}
  \item \textbf{Current}: la rotazione attuale del Pivot, con \textbf{Get Relative
        Rotation} sul componente Pivot;
  \item \textbf{Target}: splittato, con la sola Z collegata, ``perché non ci interessa
        la X e la Y'';
  \item \textbf{Delta Time}: 5;
  \item \textbf{Interp Speed}: \textbf{0,04}, ``perché, essendo che sta sui
        \emph{frame}, è molto veloce. Se voi mettete uno, questa \emph{snappa}
        diretto''.
\end{itemize}
Infine \textbf{Set Relative Rotation} sul Pivot (``prendete Pivot, Set Relative
Rotation''), con i pin splittati ``per sicurezza'', collegato al ramo Is Valid; e un
Print String per il debug (figura~\ref{ue4:fig:turretbp}).

\begin{figure}[htbp]
\centering
\begin{tikzpicture}[font=\footnotesize,
  ue4node/.style={draw, thick, rounded corners=4pt, align=center, fill=gray!8,
                  minimum height=0.85cm, text width=2.3cm},
  ue4event/.style={ue4node, fill=red!12},
  ue4exec/.style={-{Stealth[length=2.5mm]}, very thick, black!80},
  ue4data/.style={-{Stealth[length=2mm]}, thick, blue!60!black},
  ue4flt/.style={-{Stealth[length=2mm]}, thick, green!50!black},
  ue4dlg/.style={-{Stealth[length=2mm]}, thick, red!70!black, dashed},
  ue4comp/.style={draw, thick, rounded corners=3pt, minimum width=2.3cm,
                  minimum height=0.55cm, align=center, font=\scriptsize}]
  % --- riga 1: overlap
  \node[ue4event] (ov) at (0,0)
    {\textbf{On Component Begin Overlap}\\(Sphere)};
  \node[ue4node] (iv1) at (2.95,0) {\textbf{Is Valid}\\Other Actor};
  \node[ue4node] (br) at (5.9,0) {\textbf{Branch}\\Actor Has Tag ``Player''};
  \node[ue4node] (set) at (8.85,0) {\textbf{Set}\\\textbf{PlayerTarget}};
  \node[ue4node, text width=2.8cm] (call) at (12.05,0) {\textbf{TrackTargetTimer}\\(chiamata)};
  \draw[ue4exec] (ov) -- (iv1);
  \draw[ue4exec] (iv1) -- (br);
  \draw[ue4exec] (br) -- node[above, font=\scriptsize]{True} (set);
  \draw[ue4exec] (set) -- (call);
  \draw[ue4data] (ov.south) -- (0,-1.15) -| node[pos=0.2, below,
    font=\scriptsize]{Other Actor} (set.south);
  % --- riga 2: timer
  \node[ue4event, text width=2.8cm] (cev) at (0,-2.9)
    {\textbf{Custom Event}\\\textbf{TrackTargetTimer}};
  \node[ue4node, text width=3.0cm] (timer) at (3.65,-2.9)
    {\textbf{Set Timer by Event}\\Time: Get World Delta Seconds\\Looping \checkmark\\
    Max Once Per Frame \checkmark};
  \node[ue4event] (mv) at (7.1,-2.9) {\textbf{Custom Event}\\\textbf{MoveTurret}};
  \draw[ue4exec] (cev) -- (timer);
  \draw[ue4dlg] (timer.east) -- node[above, font=\scriptsize]{Event} (mv.west);
  % --- riga 3: move turret
  \node[ue4node] (iv2) at (0,-5.7) {\textbf{Is Valid}\\PlayerTarget};
  \node[ue4node] (look) at (3.1,-5.7) {\textbf{Find Look at Rotation}\\
    Start: self\\Target: PlayerTarget};
  \node[ue4node] (interp) at (6.3,-5.7) {\textbf{RInterp To}\\
    Current: Pivot\\Target: solo Z\\Speed 0,04--0,05};
  \node[ue4node] (srr) at (9.5,-5.7) {\textbf{Set Relative Rotation}\\Target: Pivot};
  \draw[ue4exec] (mv.south) -- (7.1,-4.35) -| (iv2.north);
  \draw[ue4exec] (iv2.south) -- (0,-7.1) -| node[pos=0.2, below,
    font=\scriptsize]{Is Valid} (srr.south);
  \draw[ue4flt] (look) -- node[above, font=\scriptsize]{Z} (interp);
  \draw[ue4flt] (interp) -- (srr);
  % --- gerarchia componenti
  \node[font=\scriptsize] at (11.6,-1.95) {i componenti};
  \node[ue4comp, fill=orange!25] (cb) at (11.4,-2.5) {Base (radice)};
  \node[ue4comp, fill=gray!20] (cp) at (11.9,-3.2) {Pivot};
  \node[ue4comp, fill=orange!12] (cc) at (12.4,-3.9) {Cannon};
  \node[ue4comp, fill=red!8] (cs) at (11.9,-4.6) {Sphere (r. 800)};
  \draw[thick, gray!60!black] ([xshift=-0.85cm]cb.south) |- (cp.west);
  \draw[thick, gray!60!black] ([xshift=-0.85cm]cp.south) |- (cc.west);
  \draw[thick, gray!60!black] ([xshift=-0.85cm]cb.south) |- (cs.west);
\end{tikzpicture}
\caption{La torretta della lezione. In alto: quando qualcosa entra nella sfera, se
esiste ed è etichettato \emph{Player}, lo si salva in PlayerTarget e si chiama
TrackTargetTimer. In mezzo: TrackTargetTimer fa partire un timer che, a ogni
\emph{frame}, chiama MoveTurret (il pin rosso tratteggiato è il delegato Event). In
basso: MoveTurret, se c'è ancora un bersaglio, calcola la rotazione verso di lui, ne
tiene lo \emph{yaw}, la interpola e la applica al Pivot (la Current dell'interpolazione
viene da Get Relative Rotation del Pivot; l'eventuale offset è discusso nel testo). A
destra la gerarchia dei componenti: il Cannon segue il Pivot, la sfera non ruota. Schema
ricostruito.}
\label{ue4:fig:turretbp}
\end{figure}

\paragraph{L'offset, e il pivot reimportato.} In Play la torretta ruota, ma ``ha
bisogno di un offset, come abbiamo visto con il nostro mulino''. Il docente prova con
$-90$, poi con $+90$: ``eccolo qua. Non è esattamente precisissimo''. Uno studente
osserva: ``noi questa cosa non la dobbiamo fare, perché l'altra volta abbiamo fatto un
\textbf{reimport} modificando i valori''. Ed è la chiave: senza alcun offset, la
torretta del docente ``guarda direttamente in faccia'', mentre quelle degli studenti no,
perché con Davide avevano reimportato il \textbf{Pivot} e il \textbf{Cannon} con una
\textbf{rotazione di 113} gradi sulla Z (la base no). Il docente ricostruisce i numeri
con la classe (si parla anche di $-67$ e di $-271$, ``proprio perfetti''), reimporta
prima il pivot, poi il cannone, prova con il solo pivot, e alla fine: ``ecco a voi.
Adesso la nostra torretta ci segue perfettamente''. Per gli studenti la regola pratica,
ripetuta più volte nel finale, è: se la torretta ha ``un \emph{offsettino}'',
\textbf{reimportare il Pivot mettendo $-113$} sulla rotazione Z\footnote{La trascrizione
di questo passaggio è tra le più confuse della lezione: i valori citati sono 113, $-113$,
$-67$, 118, $-270$ e $-271$, e non è chiaro in che ordine siano stati provati. Il senso
è comunque chiaro: la rotazione di 113 gradi aggiunta in import nella terza lezione
sposta la direzione ``in avanti'' della \emph{mesh}, e Find Look at Rotation, che
presuppone che l'oggetto guardi lungo la sua X, ne risulta sfasato. Le due correzioni
possibili sono equivalenti: annullare la rotazione in import, o sommare un offset fisso
allo \emph{yaw} nel grafo, come per il mulino.}.

\paragraph{Approfondimento: Find Look at Rotation e lo spazio relativo.} Il grafo della
lezione mescola due spazi: Find Look at Rotation lavora nello spazio del
\textbf{mondo} (le posizioni di torretta e player sono posizioni nel mondo), mentre
Set Relative Rotation imposta la rotazione del Pivot \textbf{relativa} alla base. Le
due cose coincidono finché la torretta nel livello non è ruotata: se la si piazza
girata di 30 gradi, il Pivot punterà 30 gradi a lato del player. Le soluzioni robuste
sono due: usare \textbf{Set World Rotation} sul Pivot (la rotazione calcolata è già
quella del mondo), oppure sottrarre allo \emph{yaw} calcolato quello dell'actor (Get
Actor Rotation). Anche la Current dell'interpolazione andrebbe presa nello stesso
spazio del Target: Get World Rotation con Set World Rotation, Get Relative Rotation con
Set Relative Rotation. Lo Start, infine, è la posizione dell'actor, cioè della base; per
una torretta alta conviene usare la posizione del Pivot o del cannone (\emph{Get World
Location} sul componente), specie quando si vorrà calcolare anche il \emph{pitch} per
alzare la canna.

% ---------------------------------------------------------------------
\section{Il tuning e i problemi degli studenti}
\label{ue4:sec:problemi}

\paragraph{La velocità di inseguimento.} ``Giratele intorno.'' Per aumentare o diminuire
la velocità con cui la torretta ``ci mette a cercare il target'' si cambia l'\textbf{Interp
Speed}. A 0,07 è ``molto più reattiva''; a 2 ``sarà proprio quasi uno \emph{snap}''
(``pam, \emph{snappa} proprio''). Il consiglio: ``mettetela tipo 0,05''. E, volendo, la
si può trasformare in variabile (Promote to Variable) ``e la potete settare
esternamente'', nei dettagli dell'istanza in scena, come le variabili \emph{Instance
Editable} della seconda lezione.

\paragraph{``A me snappa lo stesso.''} Uno studente ha Delta Time 5 e Interp Speed 0,05,
e la torretta scatta comunque. Si prova 0,02 (``vedo leggerissima differenza''), poi
Delta Time 7: ``stampa sempre la stessa velocità''. In condivisione il docente controlla
Get World Delta Seconds, Looping e Max Once Per Frame: tutto giusto. ``Prova ad andare di
\emph{tweak}, di \emph{tuning}: vedi un po' tu i valori.'' Abbassando ancora l'Interp
Speed, verso lo zero, la torretta rallenta: ``ora sì, già un po' più così''. (Nel
frattempo lo studente che lavora deve andare: ``ciao, ciao''.)

\paragraph{Approfondimento: perché ``snappa''.} La formula di RInterp To vista per il
mulino spiega tutto. Con Delta Time fisso a 5 e Interp Speed 0,04, la frazione di
strada percorsa a ogni chiamata è $5 \times 0{,}04 = 0{,}2$: il 20\% di quel che manca.
Dopo dieci chiamate manca $0{,}8^{10} \approx 11\%$ dell'angolo iniziale; a 100
\emph{frame} al secondo, sono un decimo di secondo. Con 0,05 la frazione è 0,25; con 7
e 0,02 è 0,14; con Interp Speed 2 il prodotto è 10, cioè più di 1, e la torretta arriva
in un solo \emph{frame}: lo \emph{snap}. In tutti questi casi, inoltre, la velocità
dipende dai \emph{frame}: con il doppio dei \emph{frame} la torretta è due volte più
rapida. L'uso previsto del nodo è il contrario di quello della lezione: \textbf{Delta
Time} collegato a \textbf{Get World Delta Seconds} (circa 0,01 a 100 \emph{fps}) e
\textbf{Interp Speed} nell'ordine delle unità (da 2 a 5 per un inseguimento morbido),
così che il prodotto resti piccolo e la velocità sia la stessa su ogni macchina. È lo
stesso principio dei \emph{delta seconds} della seconda lezione, applicato a
un'interpolazione. E un'avvertenza: un Interp Speed esattamente 0 non ferma la torretta,
ma la fa scattare subito sul bersaglio, perché Unreal lo interpreta come ``nessuna
interpolazione''.

\paragraph{La lista degli errori.} Il resto della lezione è una carrellata di
condivisioni dello schermo, molto istruttiva perché mostra gli errori più frequenti. La
tabella~\ref{ue4:tab:errori} li raccoglie. Il caso più ricco è quello di uno studente
a cui ``non va perché probabilmente ho mancato delle parti'':
\begin{itemize}
  \item per la Current di RInterp To aveva usato un \emph{Set} invece di \textbf{Get
        Relative Rotation} (``no, no, no, non \emph{set}: \emph{get}'');
  \item l'Is Valid di MoveTurret non era collegato alla variabile \textbf{PlayerTarget}
        (``a sinistra, proprio sulle variabili, non sul grafo'');
  \item MoveTurret non era un evento separato. Per riorganizzare il grafo il docente
        gli fa selezionare tutto tranne TrackTargetTimer, Set Timer by Event e Get
        World Delta Seconds (``tieni premuto Ctrl e clicca su quelli che vuoi
        escludere''), spostare in basso la selezione, tirare fuori dal pin Event del
        timer un nuovo Custom Event chiamato \textbf{MoveTurret} e collegarlo all'Is
        Valid;
  \item il timer non era in \textbf{Looping};
  \item il \textbf{Target} di Find Look at Rotation era scollegato: da PlayerTarget,
        Get Actor Location, al Target;
  \item il \emph{character} non aveva il \textbf{tag}: lo si aggiunge nel Third Person
        Character, ``\emph{Player} esattamente come l'hai scritto, quindi con la P
        maiuscola'';
  \item la sfera aveva ancora Hidden in Game attivo e raggio \textbf{32}: ``devi
        metterla a 800'';
  \item in scena c'era solo la \textbf{mesh} della torretta, non il
        \textbf{Blueprint}: ``questa è solo la \emph{mesh}, è solo il modello 3D che hai
        messo in scena''; si cancella e si trascina nel livello BP\_Turret;
  \item il player partiva già dentro il range, e quindi il Begin Overlap non scattava:
        ``il player parte da fuori, mettilo fuori'';
  \item in Play \textbf{girava la base}: nella gerarchia, la base era figlia del Pivot.
        ``Devi staccare la base'': trascinando il Pivot dentro la base, ``perfetto, ce
        l'hai perfettamente'';
  \item infine, l'\emph{offset} del pivot reimportato: ``devi reimportare il pivot a
        $-113$''.
\end{itemize}
Un altro studente ha la torretta ferma: ``Other Actor devi attaccarlo a PlayerTarget,
perché così non gli stai impostando il \emph{target} che è overlappato: non avrà un
\emph{target} da seguire''. Un terzo non vede la sfera in gioco, solo in modalità
\emph{spectator}: aveva saltato il passaggio di Hidden in Game.

\begin{table}[htbp]
\centering
\small
\begin{tabular}{@{}p{5.4cm}p{8.6cm}@{}}
\toprule
\textbf{Sintomo} & \textbf{Causa e rimedio} \\
\midrule
La pala gira attorno a un punto sbagliato & la \emph{mesh} ha il pivot fuori
centro: scegliere un cubo con il pivot al centro (o correggere il pivot) \\
La pala non gira & il risultato di Combine Rotators non è collegato a Set
Relative Rotation; oppure il componente non è Movable \\
Il mulino o la torretta restano fermi & mancano Delta Time e Interp Speed in
RInterp To; oppure non si è compilato \\
Non si riesce a collegare un valore a una sola componente & fare Split Struct Pin
(tasto destro sul pin) \\
Gira tutto l'actor, non solo una parte & il Rotating Movement fa girare la radice:
programmare la rotazione del componente \\
La torretta non si accorge del player & manca il tag \emph{Player} (o è scritto
diversamente); Other Actor non collegato al Set di PlayerTarget; il player parte già
dentro la sfera \\
La sfera non si vede in Play & Hidden in Game ancora attivo; raggio di default 32 \\
Non succede niente in Play & in scena c'è la \emph{mesh}, non il Blueprint \\
Gira la base & gerarchia sbagliata: la base deve essere la radice, il Pivot suo
figlio, il Cannon figlio del Pivot \\
La torretta guarda di lato & la \emph{mesh} non guarda lungo X (Pivot reimportato
a 113): reimportare a $-113$ o sommare un offset allo \emph{yaw} \\
La torretta scatta (\emph{snap}) & prodotto Delta Time $\times$ Interp Speed troppo
grande: usare Get World Delta Seconds e un Interp Speed di qualche unità \\
La torretta continua a seguire dopo l'uscita & il timer non viene mai fermato e
PlayerTarget resta pieno: End Overlap, Clear and Invalidate Timer, svuotare la
variabile \\
\bottomrule
\end{tabular}
\caption{I problemi incontrati durante la lezione, con la causa e il rimedio. Le ultime
due righe riguardano punti che il docente ha rimandato alla lezione successiva.}
\label{ue4:tab:errori}
\end{table}

% ---------------------------------------------------------------------
\section{La chiusura e la prossima lezione}
\label{ue4:sec:chiusura}

``Ok, perfetto, ragazzi.'' Per questa volta, ``nessun particolare compito da casa o
dimostrazione''. Chi vuole può fare una \textbf{torretta di colore diverso}, con un
materiale suo, verde o blu, ripassando quello che si è fatto con Davide.

\paragraph{La prossima volta.} ``La prossima volta faremo la nostra bellissima
\textbf{torretta che spara}'': \textbf{spawnerà} dei \textbf{proiettili}. Si vedrà
``come funziona il sistema di \emph{spawning} di Unreal'': come si basa una cosa
\emph{spawnata} ``con le sue regole'', come si fa apparire un oggetto in generale, e
come si fa il \textbf{line trace}. ``Altre domande? Se no, andiamo a cena tutti,
ragazzi.'' ``Vi auguro buona serata. E a martedì.''

\begin{table}[htbp]
\centering
\small
\begin{tabular}{@{}p{4.0cm}p{10.0cm}@{}}
\toprule
\textbf{Nodo o componente} & \textbf{Che cosa fa} \\
\midrule
Rotating Movement Component & fa ruotare la radice dell'actor a velocità costante
(Rotation Rate, default Z 180) \\
Event Tick, Delta Seconds & esegue a ogni \emph{frame}; i \emph{delta seconds} sono il
tempo tra due \emph{frame} \\
Get World Delta Seconds & gli stessi \emph{delta seconds}, ovunque nel grafo \\
Make Rotator & costruisce un \emph{rotator} da tre \emph{float} (Roll X, Pitch Y,
Yaw Z) \\
Combine Rotators & combina due rotazioni in una \\
Get / Set Relative Rotation & legge e scrive la rotazione di un componente rispetto al
genitore \\
Get Actor Location & la posizione dell'actor nel mondo \\
Get Player Character & una referenza al personaggio del giocatore \\
Find Look at Rotation & la rotazione per guardare da Start verso Target \\
RInterp To & avvicina gradualmente un \emph{rotator} a un altro \\
Clamp (float) & limita un valore tra un minimo e un massimo \\
Split Struct Pin & divide un pin di struttura (vettore, \emph{rotator}) nelle sue
componenti \\
Sphere Collision & un volume sferico per rilevare \emph{overlap} (Sphere Radius,
Hidden in Game) \\
On Component Begin Overlap & scatta quando qualcosa entra nel volume del componente;
Other Actor è chi è entrato \\
Is Valid (macro con il ``?'') & controlla che una referenza punti a qualcosa che
esiste \\
Actor Has Tag & controlla se un actor porta un'etichetta (Actor $\rightarrow$ Tags) \\
Promote to Variable & crea una variabile del tipo giusto da un pin \\
Custom Event & un evento creato da noi, che parte quando lo si chiama \\
Set Timer by Event & chiama un evento dopo un tempo, una volta o in \emph{loop} (Max
Once Per Frame) \\
Print String & stampa a schermo: lo strumento base del debug \\
\bottomrule
\end{tabular}
\caption{I nodi e i componenti della quarta lezione.}
\label{ue4:tab:nodi}
\end{table}

% ---------------------------------------------------------------------
\section*{Punti chiave del capitolo}
\addcontentsline{toc}{section}{Punti chiave del capitolo}

\begin{itemize}
  \item \textbf{Quixel Megascans} è un ``calderone'' gratuito di materiali, texture e
        \emph{mesh} scansionati, da usare con l'account Epic. Per la pallina che
        attraversa il piano: il \textbf{CCD} (\emph{Continuous Collision Detection}).
  \item Un \textbf{vettore} ha direzione, magnitudine e verso; in Unreal la struttura
        Vector (X, Y, Z) rappresenta sia punti sia spostamenti. Somma: si spostano i
        punti; \textbf{sottrazione}: B $-$ A va da A a B; \textbf{moltiplicazione}:
        allunga; \textbf{prodotto scalare}: un \emph{float}, 1 se concordi, 0 se
        perpendicolari, $-1$ se opposti (per vettori normalizzati).
  \item Un vettore \textbf{normalizzato} ha lunghezza 1: tiene solo la direzione.
  \item Un \textbf{rotator} ha tre angoli: \textbf{Roll} (X), \textbf{Pitch} (Y),
        \textbf{Yaw} (Z). Gli angoli di Eulero soffrono di \textbf{gimbal lock}; i
        \textbf{quaternioni}, usati internamente da Unreal, no.
  \item Il \textbf{Rotating Movement Component} fa girare un actor senza programmare,
        ma solo per intero. Ogni oggetto che si muove deve essere \textbf{Movable}.
  \item Per far girare un componente: \textbf{Event Tick} $\rightarrow$ velocità
        $\times$ \emph{delta seconds} $\rightarrow$ \textbf{Make Rotator}
        $\rightarrow$ \textbf{Combine Rotators} con \textbf{Get Relative Rotation}
        $\rightarrow$ \textbf{Set Relative Rotation}. Le rotazioni dei componenti si
        lavorano in spazio \textbf{relativo}.
  \item La \textbf{gerarchia} fa il resto: un componente figlio segue il genitore,
        mantenendo la propria posizione relativa (la seconda pala, il cubetto).
  \item \textbf{Find Look at Rotation} calcola la rotazione per guardare un punto;
        dello \emph{yaw} si tiene la Z (\textbf{Split Struct Pin}), con un
        \textbf{offset} se la \emph{mesh} non guarda lungo X. \textbf{RInterp To}
        rende il movimento graduale: Delta Time dai \emph{delta seconds}, Interp Speed
        di qualche unità, altrimenti ``snappa''.
  \item Il \textbf{Clamp} limita un valore; prima di sceglierne i limiti, si
        \textbf{stampano} i valori con Print String.
  \item La torretta è un Blueprint in \textbf{tre pezzi}: Base (radice), Pivot
        (figlio, ruota in \emph{yaw}), Cannon (figlio del Pivot). Browse to Asset aiuta
        a trovare le \emph{mesh}.
  \item Il range è una \textbf{Sphere Collision} (raggio 800, Hidden in Game tolto per
        vederla). \textbf{On Component Begin Overlap} scatta quando qualcosa ci entra;
        \textbf{Other Actor} dice chi.
  \item \textbf{Is Valid} prima di usare una referenza; \textbf{Actor Has Tag} con il tag
        \emph{Player} (scritto identico) per riconoscere il giocatore;
        \textbf{Promote to Variable} per ricordarlo in \emph{PlayerTarget}.
  \item I \textbf{Custom Event} sono eventi nostri, chiamati quando vogliamo;
        \textbf{Set Timer by Event} (Looping, Max Once Per Frame) chiama a ripetizione
        l'evento \textbf{MoveTurret}, che ruota il Pivot verso il bersaglio. Il timer
        va fermato (End Overlap, Clear and Invalidate Timer): sarà il prossimo passo.
  \item Prossima lezione: la torretta che \textbf{spara}, lo \textbf{spawn} dei
        proiettili e il \textbf{line trace}.
\end{itemize}
